\documentclass[10pt, openany]{book}
\usepackage[text={4.65in,7.45in}, centering, includefoot]{geometry}
\usepackage[table, x11names]{xcolor}
\usepackage{fontspec,realscripts}
\usepackage{polyglossia}
\setdefaultlanguage{sanskrit}
\setotherlanguage{english}
\setmainfont[
    Script        = Devanagari,
    Scale         = .9,
    Renderer      = HarfBuzz,
]{Shobhika-Bold}
\newfontfamily\la[Script=Devanagari, Scale=.9]{Shobhika-Bold}
\newfontfamily\regular{Linux Libertine O}
\newfontfamily\en[Language=English, Script=Latin]{Linux Libertine O}
\newfontfamily\vp[Script=Devanagari, Scale=1, Color=purple]{Shobhika-Bold}
\newfontfamily\vps[Script=Devanagari, Scale=.8, Color=purple]{Shobhika-Bold}
\newfontfamily\vpc[Script=Devanagari, Scale=1, Color=violet]{Shobhika-Bold}
\newfontfamily\qt[Script=Devanagari, Scale=.8, Color=violet]{Shobhika-Bold}
\XeTeXgenerateactualtext=1

%% Packages
\usepackage{enumerate}
\usepackage{fancyhdr}
\usepackage{afterpage}
\usepackage{multirow}
\usepackage{multicol}
\usepackage{wrapfig}
\usepackage{caption}
\usepackage{ragged2e}
\usepackage{microtype}
\usepackage{graphicx}
\usepackage{longtable}
\usepackage{setspace}
\usepackage{xspace}
\usepackage{array}
\usepackage{emptypage}

%% Sutra numbering in margin
\usepackage{marginnote}
\reversemarginpar

%% Section styling (keep from ashta.tex)
\usepackage{titlesec}
\titleformat{\chapter}[block]{\centering\large\bfseries}{}{0pt}{}
\titleformat{\section}[block]{\centering\normalsize\bfseries}{}{0pt}{}

%% Spacing & layout (keep from ashta.tex)
\usepackage{parskip}
\setlength{\parindent}{0pt}
\setlength{\parskip}{4pt}
\setstretch{1.4}

%% Hyperlinks
\usepackage[unicode,colorlinks=false,pdfusetitle]{hyperref}

%% ── Devanagari page numbers (set AFTER all packages) ──────────────────────
%% \la forces the Devanagari font; \number\value{page} gives a plain integer
%% that \devanagaridigits can safely convert.
\pagestyle{fancy}
\fancyhf{}
\renewcommand{\headrulewidth}{0pt}
\fancyfoot[C]{{\la\devanagaridigits{\number\value{page}}}}
%% book class uses 'plain' on chapter-opening pages — mirror the same style
\fancypagestyle{plain}{%
  \fancyhf{}%
  \renewcommand{\headrulewidth}{0pt}%
  \fancyfoot[C]{{\la\devanagaridigits{\number\value{page}}}}%
}

\title{अष्टाध्यायी अथवा सूत्रपाठ पाणिनीकृत}
\date{}
\begin{document}
\maketitle

\section*{॥ अध्याय १ ॥}

\marginnote{\scriptsize १.१.१}\noindent वृद्धिरादैच् ।

\marginnote{\scriptsize १.१.२}\noindent अदेङ् गुणः ।

\marginnote{\scriptsize १.१.३}\noindent इको गुणवृद्धी ।

\marginnote{\scriptsize १.१.४}\noindent न धातुलोप आर्धधातुके ।

\marginnote{\scriptsize १.१.५}\noindent क्ङिति च ।

\marginnote{\scriptsize १.१.६}\noindent दीधीवेवीटाम् ।

\marginnote{\scriptsize १.१.७}\noindent हलोऽनन्तराः संयोगः ।

\marginnote{\scriptsize १.१.८}\noindent मुखनासिकावचनोऽनुनासिकः ।

\marginnote{\scriptsize १.१.९}\noindent तुल्यास्यप्रयत्नं सवर्णम् ।

\marginnote{\scriptsize १.१.१०}\noindent नाज्झलौ ।

\marginnote{\scriptsize १.१.११}\noindent ईदूदेद्द्विवचनं प्रगृह्यम् ।

\marginnote{\scriptsize १.१.१२}\noindent अदसो मात् ।

\marginnote{\scriptsize १.१.१३}\noindent शे ।

\marginnote{\scriptsize १.१.१४}\noindent निपात एकाजनाङ् ।

\marginnote{\scriptsize १.१.१५}\noindent ओत् ।

\marginnote{\scriptsize १.१.१६}\noindent सम्बुद्धौ शाकल्यस्येतावनार्षे ।

\marginnote{\scriptsize १.१.१७}\noindent उञः ।

\marginnote{\scriptsize १.१.१८}\noindent ऊँ ।

\marginnote{\scriptsize १.१.१९}\noindent ईदूतौ च सप्तम्यर्थे ।

\marginnote{\scriptsize १.१.२०}\noindent दाधा घ्वदाप् ।

\marginnote{\scriptsize १.१.२१}\noindent आद्यन्तवदेकस्मिन् ।

\marginnote{\scriptsize १.१.२२}\noindent तरप्तमपौ घः ।

\marginnote{\scriptsize १.१.२३}\noindent बहुगणवतुडति संख्या ।

\marginnote{\scriptsize १.१.२४}\noindent ष्णान्ता षट् ।

\marginnote{\scriptsize १.१.२५}\noindent डति च ।

\marginnote{\scriptsize १.१.२६}\noindent क्तक्तवतू निष्ठा ।

\marginnote{\scriptsize १.१.२७}\noindent सर्वादीनि सर्वनामानि ।

\marginnote{\scriptsize १.१.२८}\noindent विभाषा दिक्समासे बहुव्रीहौ ।

\marginnote{\scriptsize १.१.२९}\noindent न बहुव्रीहौ ।

\marginnote{\scriptsize १.१.३०}\noindent तृतीयासमासे ।

\marginnote{\scriptsize १.१.३१}\noindent द्वन्द्वे च ।

\marginnote{\scriptsize १.१.३२}\noindent विभाषा जसि ।

\marginnote{\scriptsize १.१.३३}\noindent प्रथमचरमतयाल्पार्धकतिपयनेमाश्च ।

\marginnote{\scriptsize १.१.३४}\noindent पूर्वपरावरदक्षिणोत्तरापराधराणि व्यवस्थायामसंज्ञायाम् ।

\marginnote{\scriptsize १.१.३५}\noindent स्वमज्ञातिधनाख्यायाम् ।

\marginnote{\scriptsize १.१.३६}\noindent अन्तरं बहिर्योगोपसंव्यानयोः ।

\marginnote{\scriptsize १.१.३७}\noindent स्वरादिनिपातमव्ययम् ।

\marginnote{\scriptsize १.१.३८}\noindent तद्धितश्चासर्वविभक्तिः ।

\marginnote{\scriptsize १.१.३९}\noindent कृन्मेजन्तः ।

\marginnote{\scriptsize १.१.४०}\noindent क्त्वातोसुन्कसुनः ।

\marginnote{\scriptsize १.१.४१}\noindent अव्ययीभावश्च ।

\marginnote{\scriptsize १.१.४२}\noindent शि सर्वनामस्थानम् ।

\marginnote{\scriptsize १.१.४३}\noindent सुडनपुंसकस्य ।

\marginnote{\scriptsize १.१.४४}\noindent न वेति विभाषा ।

\marginnote{\scriptsize १.१.४५}\noindent इग्यणः सम्प्रसारणम् ।

\marginnote{\scriptsize १.१.४६}\noindent आद्यन्तौ टकितौ ।

\marginnote{\scriptsize १.१.४७}\noindent मिदचोऽन्त्यात्परः ।

\marginnote{\scriptsize १.१.४८}\noindent एच इग्घ्रस्वादेशे ।

\marginnote{\scriptsize १.१.४९}\noindent षष्ठी स्थानेयोगा ।

\marginnote{\scriptsize १.१.५०}\noindent स्थानेऽन्तरतमः ।

\marginnote{\scriptsize १.१.५१}\noindent उरण् रपरः ।

\marginnote{\scriptsize १.१.५२}\noindent अलोऽन्त्यस्य ।

\marginnote{\scriptsize १.१.५३}\noindent ङिच्च ।

\marginnote{\scriptsize १.१.५४}\noindent आदेः परस्य ।

\marginnote{\scriptsize १.१.५५}\noindent अनेकाल्शित्सर्वस्य ।

\marginnote{\scriptsize १.१.५६}\noindent स्थानिवदादेशोऽनल्विधौ ।

\marginnote{\scriptsize १.१.५७}\noindent अचः परस्मिन् पूर्वविधौ ।

\marginnote{\scriptsize १.१.५८}\noindent न पदान्तद्विर्वचनवरेयलोपस्वरसवर्णानुस्वारदीर्घ- जश्चर्विधिषु ।

\marginnote{\scriptsize १.१.५९}\noindent द्विर्वचनेऽचि ।

\marginnote{\scriptsize १.१.६०}\noindent अदर्शनं लोपः ।

\marginnote{\scriptsize १.१.६१}\noindent प्रत्ययस्य लुक्श्लुलुपः ।

\marginnote{\scriptsize १.१.६२}\noindent प्रत्ययलोपे प्रत्ययलक्षणम् ।

\marginnote{\scriptsize १.१.६३}\noindent न लुमताऽङ्गस्य ।

\marginnote{\scriptsize १.१.६४}\noindent अचोऽन्त्यादि टि ।

\marginnote{\scriptsize १.१.६५}\noindent अलोऽन्त्यात् पूर्व उपधा ।

\marginnote{\scriptsize १.१.६६}\noindent तस्मिन्निति निर्दिष्टे पूर्वस्य ।

\marginnote{\scriptsize १.१.६७}\noindent तस्मादित्युत्तरस्य ।

\marginnote{\scriptsize १.१.६८}\noindent स्वं रूपं शब्दस्याशब्दसंज्ञा ।

\marginnote{\scriptsize १.१.६९}\noindent अणुदित् सवर्णस्य चाप्रत्ययः ।

\marginnote{\scriptsize १.१.७०}\noindent तपरस्तत्कालस्य ।

\marginnote{\scriptsize १.१.७१}\noindent आदिरन्त्येन सहेता ।

\marginnote{\scriptsize १.१.७२}\noindent येन विधिस्तदन्तस्य ।

\marginnote{\scriptsize १.१.७३}\noindent वृद्धिर्यस्याचामादिस्तद् वृद्धम् ।

\marginnote{\scriptsize १.१.७४}\noindent त्यदादीनि च ।

\marginnote{\scriptsize १.१.७५}\noindent एङ् प्राचां देशे ।

\marginnote{\scriptsize १.२.१}\noindent गाङ्कुटादिभ्योऽञ्णिन्ङ् इत् ।

\marginnote{\scriptsize १.२.२}\noindent विज इट् ।

\marginnote{\scriptsize १.२.३}\noindent विभाषोर्णोः ।

\marginnote{\scriptsize १.२.४}\noindent सार्वधातुकमपित् ।

\marginnote{\scriptsize १.२.५}\noindent असंयोगाल्लिट् कित् ।

\marginnote{\scriptsize १.२.६}\noindent इन्धिभवतिभ्यां च ।

\marginnote{\scriptsize १.२.७}\noindent मृडमृदगुधकुषक्लिशवदवसः क्त्वा ।

\marginnote{\scriptsize १.२.८}\noindent रुदविदमुषग्रहिस्वपिप्रच्छः सँश्च ।

\marginnote{\scriptsize १.२.९}\noindent इको झल् ।

\marginnote{\scriptsize १.२.१०}\noindent हलन्ताच्च ।

\marginnote{\scriptsize १.२.११}\noindent लिङ्सिचावात्मनेपदेषु ।

\marginnote{\scriptsize १.२.१२}\noindent उश्च ।

\marginnote{\scriptsize १.२.१३}\noindent वा गमः ।

\marginnote{\scriptsize १.२.१४}\noindent हनः सिच् ।

\marginnote{\scriptsize १.२.१५}\noindent यमो गन्धने ।

\marginnote{\scriptsize १.२.१६}\noindent विभाषोपयमने ।

\marginnote{\scriptsize १.२.१७}\noindent स्थाघ्वोरिच्च ।

\marginnote{\scriptsize १.२.१८}\noindent न क्त्वा सेट् ।

\marginnote{\scriptsize १.२.१९}\noindent निष्ठा शीङ्स्विदिमिदिक्ष्विदिधृषः ।

\marginnote{\scriptsize १.२.२०}\noindent मृषस्तितिक्षायाम् ।

\marginnote{\scriptsize १.२.२१}\noindent उदुपधाद्भावादिकर्मणोरन्यतरस्याम् ।

\marginnote{\scriptsize १.२.२२}\noindent पूङः क्त्वा च ।

\marginnote{\scriptsize १.२.२३}\noindent नोपधात्थफान्ताद्वा ।

\marginnote{\scriptsize १.२.२४}\noindent वञ्चिलुञ्च्यृतश्च ।

\marginnote{\scriptsize १.२.२५}\noindent तृषिमृषिकृशेः काश्यपस्य ।

\marginnote{\scriptsize १.२.२६}\noindent रलो व्युपधाद्धलादेः संश्च ।

\marginnote{\scriptsize १.२.२७}\noindent ऊकालोऽज्झ्रस्वदीर्घप्लुतः ।

\marginnote{\scriptsize १.२.२८}\noindent अचश्च ।

\marginnote{\scriptsize १.२.२९}\noindent उच्चैरुदात्तः ।

\marginnote{\scriptsize १.२.३०}\noindent नीचैरनुदात्तः ।

\marginnote{\scriptsize १.२.३१}\noindent समाहारः स्वरितः ।

\marginnote{\scriptsize १.२.३२}\noindent तस्यादित उदात्तमर्धह्रस्वम् ।

\marginnote{\scriptsize १.२.३३}\noindent एकश्रुति दूरात् सम्बुद्धौ ।

\marginnote{\scriptsize १.२.३४}\noindent यज्ञकर्मण्यजपन्यूङ्खसामसु ।

\marginnote{\scriptsize १.२.३५}\noindent उच्चैस्तरां वा वषट्कारः ।

\marginnote{\scriptsize १.२.३६}\noindent विभाषा छन्दसि ।

\marginnote{\scriptsize १.२.३७}\noindent न सुब्रह्मण्यायां स्वरितस्य तूदात्तः ।

\marginnote{\scriptsize १.२.३८}\noindent देवब्रह्मणोरनुदात्तः ।

\marginnote{\scriptsize १.२.३९}\noindent स्वरितात् संहितायामनुदात्तानाम् ।

\marginnote{\scriptsize १.२.४०}\noindent उदात्तस्वरितपरस्य सन्नतरः ।

\marginnote{\scriptsize १.२.४१}\noindent अपृक्त एकाल् प्रत्ययः ।

\marginnote{\scriptsize १.२.४२}\noindent तत्पुरुषः समानाधिकरणः कर्मधारयः ।

\marginnote{\scriptsize १.२.४३}\noindent प्रथमानिर्दिष्टं समास उपसर्जनम् ।

\marginnote{\scriptsize १.२.४४}\noindent एकविभक्ति चापूर्वनिपाते ।

\marginnote{\scriptsize १.२.४५}\noindent अर्थवदधातुरप्रत्ययः प्रातिपदिकम् ।

\marginnote{\scriptsize १.२.४६}\noindent कृत्तद्धितसमासाश्च ।

\marginnote{\scriptsize १.२.४७}\noindent ह्रस्वो नपुंसके प्रातिपदिकस्य ।

\marginnote{\scriptsize १.२.४८}\noindent गोस्त्रियोरुपसर्जनस्य ।

\marginnote{\scriptsize १.२.४९}\noindent लुक् तद्धितलुकि ।

\marginnote{\scriptsize १.२.५०}\noindent इद्गोण्याः ।

\marginnote{\scriptsize १.२.५१}\noindent लुपि युक्तवद्व्यक्तिवचने ।

\marginnote{\scriptsize १.२.५२}\noindent विशेषणानां चाजातेः ।

\marginnote{\scriptsize १.२.५३}\noindent तदशिष्यं संज्ञाप्रमाणत्वात् ।

\marginnote{\scriptsize १.२.५४}\noindent लुब्योगाप्रख्यानात् ।

\marginnote{\scriptsize १.२.५५}\noindent योगप्रमाणे च तदभावेऽदर्शनं स्यात् ।

\marginnote{\scriptsize १.२.५६}\noindent प्रधानप्रत्ययार्थवचनमर्थस्यान्यप्रमाणत्वात् ।

\marginnote{\scriptsize १.२.५७}\noindent कालोपसर्जने च तुल्यम् ।

\marginnote{\scriptsize १.२.५८}\noindent जात्याख्यायामेकस्मिन् बहुवचनमन्यतरस्याम् ।

\marginnote{\scriptsize १.२.५९}\noindent अस्मदो द्वयोश्च ।

\marginnote{\scriptsize १.२.६०}\noindent फल्गुनीप्रोष्ठपदानां च नक्षत्रे ।

\marginnote{\scriptsize १.२.६१}\noindent छन्दसि पुनर्वस्वोरेकवचनम् ।

\marginnote{\scriptsize १.२.६२}\noindent विशाखयोश्च ।

\marginnote{\scriptsize १.२.६३}\noindent तिष्यपुनर्वस्वोर्नक्षत्रद्वंद्वे बहुवचनस्य द्विवचनं नित्यम् ।

\marginnote{\scriptsize १.२.६४}\noindent सरूपाणामेकशेष एकविभक्तौ ।

\marginnote{\scriptsize १.२.६५}\noindent वृद्धो यूना तल्लक्षणश्चेदेव विशेषः ।

\marginnote{\scriptsize १.२.६६}\noindent स्त्री पुंवच्च ।

\marginnote{\scriptsize १.२.६७}\noindent पुमान् स्त्रिया ।

\marginnote{\scriptsize १.२.६८}\noindent भ्रातृपुत्रौ स्वसृदुहितृभ्याम् ।

\marginnote{\scriptsize १.२.६९}\noindent नपुंसकमनपुंसकेनैकवच्चास्यान्यतरस्याम् ।

\marginnote{\scriptsize १.२.७०}\noindent पिता मात्रा ।

\marginnote{\scriptsize १.२.७१}\noindent श्वशुरः श्वश्र्वा ।

\marginnote{\scriptsize १.२.७२}\noindent त्यदादीनि सर्वैर्नित्यम् ।

\marginnote{\scriptsize १.२.७३}\noindent ग्राम्यपशुसंघेषु अतरुणेषु स्त्री ।

\marginnote{\scriptsize १.३.१}\noindent भूवादयो धातवः ।

\marginnote{\scriptsize १.३.२}\noindent उपदेशेऽजनुनासिक इत् ।

\marginnote{\scriptsize १.३.३}\noindent हलन्त्यम् ।

\marginnote{\scriptsize १.३.४}\noindent न विभक्तौ तुस्माः ।

\marginnote{\scriptsize १.३.५}\noindent आदिर्ञिटुडवः ।

\marginnote{\scriptsize १.३.६}\noindent षः प्रत्ययस्य ।

\marginnote{\scriptsize १.३.७}\noindent चुटू ।

\marginnote{\scriptsize १.३.८}\noindent लशक्वतद्धिते ।

\marginnote{\scriptsize १.३.९}\noindent तस्य लोपः ।

\marginnote{\scriptsize १.३.१०}\noindent यथासंख्यम् अनुदेशः समानाम् ।

\marginnote{\scriptsize १.३.११}\noindent स्वरितेनाधिकारः ।

\marginnote{\scriptsize १.३.१२}\noindent अनुदात्तङित आत्मनेपदम् ।

\marginnote{\scriptsize १.३.१३}\noindent भावकर्मणोः ।

\marginnote{\scriptsize १.३.१४}\noindent कर्तरि कर्मव्यतिहारे ।

\marginnote{\scriptsize १.३.१५}\noindent न गतिहिंसार्थेभ्यः ।

\marginnote{\scriptsize १.३.१६}\noindent इतरेतरान्योन्योपपदाच्च ।

\marginnote{\scriptsize १.३.१७}\noindent नेर्विशः ।

\marginnote{\scriptsize १.३.१८}\noindent परिव्यवेभ्यः क्रियः ।

\marginnote{\scriptsize १.३.१९}\noindent विपराभ्यां जेः ।

\marginnote{\scriptsize १.३.२०}\noindent आङो दोऽनास्यविहरणे ।

\marginnote{\scriptsize १.३.२१}\noindent क्रीडोऽनुसम्परिभ्यश्च ।

\marginnote{\scriptsize १.३.२२}\noindent समवप्रविभ्यः स्थः ।

\marginnote{\scriptsize १.३.२३}\noindent प्रकाशनस्थेयाख्ययोश्च ।

\marginnote{\scriptsize १.३.२४}\noindent उदोऽनूर्द्ध्वकर्मणि ।

\marginnote{\scriptsize १.३.२५}\noindent उपान्मन्त्रकरणे ।

\marginnote{\scriptsize १.३.२६}\noindent अकर्मकाच्च ।

\marginnote{\scriptsize १.३.२७}\noindent उद्विभ्यां तपः ।

\marginnote{\scriptsize १.३.२८}\noindent आङो यमहनः ।

\marginnote{\scriptsize १.३.२९}\noindent समो गम्यृच्छिप्रच्छिस्वरत्यर्तिश्रुविदिभ्यः ।

\marginnote{\scriptsize १.३.३०}\noindent निसमुपविभ्यो ह्वः ।

\marginnote{\scriptsize १.३.३१}\noindent स्पर्द्धायामाङः ।

\marginnote{\scriptsize १.३.३२}\noindent गन्धनावक्षेपणसेवनसाहसिक्य- प्रतियत्नप्रकथनोपयोगेषु कृञः ।

\marginnote{\scriptsize १.३.३३}\noindent अधेः प्रसहने ।

\marginnote{\scriptsize १.३.३४}\noindent वेः शब्दकर्मणः ।

\marginnote{\scriptsize १.३.३५}\noindent अकर्मकाच्च ।

\marginnote{\scriptsize १.३.३६}\noindent सम्माननोत्सञ्जनाचार्यकरणज्ञानभृतिविगणनव्ययेषु नियः ।

\marginnote{\scriptsize १.३.३७}\noindent कर्तृस्थे चाशरीरे कर्मणि ।

\marginnote{\scriptsize १.३.३८}\noindent वृत्तिसर्गतायनेषु क्रमः ।

\marginnote{\scriptsize १.३.३९}\noindent उपपराभ्याम् ।

\marginnote{\scriptsize १.३.४०}\noindent आङ उद्गमने ।

\marginnote{\scriptsize १.३.४१}\noindent वेः पादविहरणे ।

\marginnote{\scriptsize १.३.४२}\noindent प्रोपाभ्यां समर्थाभ्याम् ।

\marginnote{\scriptsize १.३.४३}\noindent अनुपसर्गाद्वा ।

\marginnote{\scriptsize १.३.४४}\noindent अपह्नवे ज्ञः ।

\marginnote{\scriptsize १.३.४५}\noindent अकर्मकाच्च ।

\marginnote{\scriptsize १.३.४६}\noindent सम्प्रतिभ्यामनाध्याने ।

\marginnote{\scriptsize १.३.४७}\noindent भासनोपसम्भाषाज्ञानयत्नविमत्युपमन्त्रणेषु वदः ।

\marginnote{\scriptsize १.३.४८}\noindent व्यक्तवाचां समुच्चारणे ।

\marginnote{\scriptsize १.३.४९}\noindent अनोरकर्मकात् ।

\marginnote{\scriptsize १.३.५०}\noindent विभाषा विप्रलापे ।

\marginnote{\scriptsize १.३.५१}\noindent अवाद्ग्रः ।

\marginnote{\scriptsize १.३.५२}\noindent समः प्रतिज्ञाने ।

\marginnote{\scriptsize १.३.५३}\noindent उदश्चरः सकर्मकात् ।

\marginnote{\scriptsize १.३.५४}\noindent समस्तृतीयायुक्तात् ।

\marginnote{\scriptsize १.३.५५}\noindent दाणश्च सा चेच्चतुर्थ्यर्थे ।

\marginnote{\scriptsize १.३.५६}\noindent उपाद्यमः स्वकरणे ।

\marginnote{\scriptsize १.३.५७}\noindent ज्ञाश्रुस्मृदृशां सनः ।

\marginnote{\scriptsize १.३.५८}\noindent नानोर्ज्ञः ।

\marginnote{\scriptsize १.३.५९}\noindent प्रत्याङ्भ्यां श्रुवः ।

\marginnote{\scriptsize १.३.६०}\noindent शदेः शितः ।

\marginnote{\scriptsize १.३.६१}\noindent म्रियतेर्लुङ्लिङोश्च ।

\marginnote{\scriptsize १.३.६२}\noindent पूर्ववत् सनः ।

\marginnote{\scriptsize १.३.६३}\noindent आम्प्रत्ययवत् कृञोऽनुप्रयोगस्य ।

\marginnote{\scriptsize १.३.६४}\noindent प्रोपाभ्यां युजेरयज्ञपात्रेषु ।

\marginnote{\scriptsize १.३.६५}\noindent समः क्ष्णुवः ।

\marginnote{\scriptsize १.३.६६}\noindent भुजोऽनवने ।

\marginnote{\scriptsize १.३.६७}\noindent णेरणौ यत् कर्म णौ चेत् स कर्ताऽनाध्याने ।

\marginnote{\scriptsize १.३.६८}\noindent भीस्म्योर्हेतुभये ।

\marginnote{\scriptsize १.३.६९}\noindent गृधिवञ्च्योः प्रलम्भने ।

\marginnote{\scriptsize १.३.७०}\noindent लियः सम्माननशालिनीकरणयोश्च ।

\marginnote{\scriptsize १.३.७१}\noindent मिथ्योपपदात् कृञोऽभ्यासे ।

\marginnote{\scriptsize १.३.७२}\noindent स्वरितञितः कर्त्रभिप्राये क्रियाफले ।

\marginnote{\scriptsize १.३.७३}\noindent अपाद्वदः ।

\marginnote{\scriptsize १.३.७४}\noindent णिचश्च ।

\marginnote{\scriptsize १.३.७५}\noindent समुदाङ्भ्यो यमोऽग्रन्थे ।

\marginnote{\scriptsize १.३.७६}\noindent अनुपसर्गाज्ज्ञः ।

\marginnote{\scriptsize १.३.७७}\noindent विभाषोपपदेन प्रतीयमाने ।

\marginnote{\scriptsize १.३.७८}\noindent शेषात् कर्तरि परस्मैपदम् ।

\marginnote{\scriptsize १.३.७९}\noindent अनुपराभ्यां कृञः ।

\marginnote{\scriptsize १.३.८०}\noindent अभिप्रत्यतिभ्यः क्षिपः ।

\marginnote{\scriptsize १.३.८१}\noindent प्राद्वहः ।

\marginnote{\scriptsize १.३.८२}\noindent परेर्मृषः ।

\marginnote{\scriptsize १.३.८३}\noindent व्याङ्परिभ्यो रमः ।

\marginnote{\scriptsize १.३.८४}\noindent उपाच्च ।

\marginnote{\scriptsize १.३.८५}\noindent विभाषाऽकर्मकात् ।

\marginnote{\scriptsize १.३.८६}\noindent बुधयुधनशजनेङ्प्रुद्रुस्रुभ्यो णेः ।

\marginnote{\scriptsize १.३.८७}\noindent निगरणचलनार्थेभ्यः ।

\marginnote{\scriptsize १.३.८८}\noindent अणावकर्मकाच्चित्तवत्कर्तृकात् ।

\marginnote{\scriptsize १.३.८९}\noindent न पादम्याङ्यमाङ्यसपरिमुहरुचिनृतिवदवसः ।

\marginnote{\scriptsize १.३.९०}\noindent वा क्यषः ।

\marginnote{\scriptsize १.३.९१}\noindent द्युद्भ्यो लुङि ।

\marginnote{\scriptsize १.३.९२}\noindent वृद्भ्यः स्यसनोः ।

\marginnote{\scriptsize १.३.९३}\noindent लुटि च कॢपः ।

\marginnote{\scriptsize १.४.१}\noindent आ कडारादेका संज्ञा ।

\marginnote{\scriptsize १.४.२}\noindent विप्रतिषेधे परं कार्यम् ।

\marginnote{\scriptsize १.४.३}\noindent यू स्त्र्याख्यौ नदी ।

\marginnote{\scriptsize १.४.४}\noindent नेयङुवङ्स्थानावस्त्री ।

\marginnote{\scriptsize १.४.५}\noindent वाऽऽमि ।

\marginnote{\scriptsize १.४.६}\noindent ङिति ह्रस्वश्च ।

\marginnote{\scriptsize १.४.७}\noindent शेषो घ्यसखि ।

\marginnote{\scriptsize १.४.८}\noindent पतिः समास एव ।

\marginnote{\scriptsize १.४.९}\noindent षष्ठीयुक्तश्छन्दसि वा ।

\marginnote{\scriptsize १.४.१०}\noindent ह्रस्वं लघु ।

\marginnote{\scriptsize १.४.११}\noindent संयोगे गुरु ।

\marginnote{\scriptsize १.४.१२}\noindent दीर्घं च ।

\marginnote{\scriptsize १.४.१३}\noindent यस्मात् प्रत्ययविधिस्तदादि प्रत्ययेऽङ्गम् ।

\marginnote{\scriptsize १.४.१४}\noindent सुप्तिङन्तं पदम् ।

\marginnote{\scriptsize १.४.१५}\noindent नः क्ये ।

\marginnote{\scriptsize १.४.१६}\noindent सिति च ।

\marginnote{\scriptsize १.४.१७}\noindent स्वादिष्वसर्वनमस्थाने ।

\marginnote{\scriptsize १.४.१८}\noindent यचि भम् ।

\marginnote{\scriptsize १.४.१९}\noindent तसौ मत्वर्थे ।

\marginnote{\scriptsize १.४.२०}\noindent अयस्मयादीनि च्छन्दसि ।

\marginnote{\scriptsize १.४.२१}\noindent बहुषु बहुवचनम् ।

\marginnote{\scriptsize १.४.२२}\noindent द्व्येकयोर्द्विवचनैकवचने ।

\marginnote{\scriptsize १.४.२३}\noindent कारके ।

\marginnote{\scriptsize १.४.२४}\noindent ध्रुवमपायेऽपादानम् ।

\marginnote{\scriptsize १.४.२५}\noindent भीत्रार्थानां भयहेतुः ।

\marginnote{\scriptsize १.४.२६}\noindent पराजेरसोढः ।

\marginnote{\scriptsize १.४.२७}\noindent वारणार्थानां ईप्सितः ।

\marginnote{\scriptsize १.४.२८}\noindent अन्तर्द्धौ येनादर्शनमिच्छति ।

\marginnote{\scriptsize १.४.२९}\noindent आख्यातोपयोगे ।

\marginnote{\scriptsize १.४.३०}\noindent जनिकर्तुः प्रकृतिः ।

\marginnote{\scriptsize १.४.३१}\noindent भुवः प्रभवः ।

\marginnote{\scriptsize १.४.३२}\noindent कर्मणा यमभिप्रैति स सम्प्रदानम् ।

\marginnote{\scriptsize १.४.३३}\noindent रुच्यर्थानां प्रीयमाणः ।

\marginnote{\scriptsize १.४.३४}\noindent श्लाघह्नुङ्स्थाशपां ज्ञीप्स्यमानः ।

\marginnote{\scriptsize १.४.३५}\noindent धारेरुत्तमर्णः ।

\marginnote{\scriptsize १.४.३६}\noindent स्पृहेरीप्सितः ।

\marginnote{\scriptsize १.४.३७}\noindent क्रुधद्रुहेर्ष्याऽसूयार्थानां यं प्रति कोपः ।

\marginnote{\scriptsize १.४.३८}\noindent क्रुधद्रुहोरुपसृष्टयोः कर्म ।

\marginnote{\scriptsize १.४.३९}\noindent राधीक्ष्योर्यस्य विप्रश्नः ।

\marginnote{\scriptsize १.४.४०}\noindent प्रत्याङ्भ्यां श्रुवः पूर्वस्य कर्ता ।

\marginnote{\scriptsize १.४.४१}\noindent अनुप्रतिगृणश्च ।

\marginnote{\scriptsize १.४.४२}\noindent साधकतमं करणम् ।

\marginnote{\scriptsize १.४.४३}\noindent दिवः कर्म च ।

\marginnote{\scriptsize १.४.४४}\noindent परिक्रयणे सम्प्रदानमन्यतरस्याम् ।

\marginnote{\scriptsize १.४.४५}\noindent आधारोऽधिकरणम् ।

\marginnote{\scriptsize १.४.४६}\noindent अधिशीङ्स्थाऽऽसां कर्म ।

\marginnote{\scriptsize १.४.४७}\noindent अभिनिविशश्च ।

\marginnote{\scriptsize १.४.४८}\noindent उपान्वध्याङ्वसः ।

\marginnote{\scriptsize १.४.४९}\noindent कर्तुरीप्सिततमं कर्म ।

\marginnote{\scriptsize १.४.५०}\noindent तथायुक्तं चानिप्सीतम् ।

\marginnote{\scriptsize १.४.५१}\noindent अकथितं च ।

\marginnote{\scriptsize १.४.५२}\noindent गतिबुद्धिप्रत्यवसानार्थशब्दकर्माकर्मकाणामणि कर्ता स णौ ।

\marginnote{\scriptsize १.४.५३}\noindent हृक्रोरन्यतरस्याम् ।

\marginnote{\scriptsize १.४.५४}\noindent स्वतन्त्रः कर्ता ।

\marginnote{\scriptsize १.४.५५}\noindent तत्प्रयोजको हेतुश्च ।

\marginnote{\scriptsize १.४.५६}\noindent प्राग्रीश्वरान्निपातः ।

\marginnote{\scriptsize १.४.५७}\noindent चादयोऽसत्त्वे ।

\marginnote{\scriptsize १.४.५८}\noindent प्रादयः ।

\marginnote{\scriptsize १.४.५९}\noindent उपसर्गाः क्रियायोगे ।

\marginnote{\scriptsize १.४.६०}\noindent गतिश्च ।

\marginnote{\scriptsize १.४.६१}\noindent ऊर्यादिच्विडाचश्च ।

\marginnote{\scriptsize १.४.६२}\noindent अनुकरणं चानितिपरम् ।

\marginnote{\scriptsize १.४.६३}\noindent आदरानादरयोः सदसती ।

\marginnote{\scriptsize १.४.६४}\noindent भूषणेऽलम् ।

\marginnote{\scriptsize १.४.६५}\noindent अन्तरपरिग्रहे ।

\marginnote{\scriptsize १.४.६६}\noindent कणेमनसी श्रद्धाप्रतीघाते ।

\marginnote{\scriptsize १.४.६७}\noindent पुरोऽव्ययम् ।

\marginnote{\scriptsize १.४.६८}\noindent अस्तं च ।

\marginnote{\scriptsize १.४.६९}\noindent अच्छ गत्यर्थवदेषु ।

\marginnote{\scriptsize १.४.७०}\noindent अदोऽनुपदेशे ।

\marginnote{\scriptsize १.४.७१}\noindent तिरोऽन्तर्धौ । variationर्द्धौ

\marginnote{\scriptsize १.४.७२}\noindent विभाषा कृञि ।

\marginnote{\scriptsize १.४.७३}\noindent उपाजेऽन्वाजे ।

\marginnote{\scriptsize १.४.७४}\noindent साक्षात्प्रभृतीनि च ।

\marginnote{\scriptsize १.४.७५}\noindent अनत्याधान उरसिमनसी ।

\marginnote{\scriptsize १.४.७६}\noindent मध्येपदेनिवचने च ।

\marginnote{\scriptsize १.४.७७}\noindent नित्यं हस्ते पाणावुपयमने ।

\marginnote{\scriptsize १.४.७८}\noindent प्राध्वं बन्धने ।

\marginnote{\scriptsize १.४.७९}\noindent जीविकोपनिषदावौपम्ये ।

\marginnote{\scriptsize १.४.८०}\noindent ते प्राग्धातोः ।

\marginnote{\scriptsize १.४.८१}\noindent छन्दसि परेऽपि ।

\marginnote{\scriptsize १.४.८२}\noindent व्यवहिताश्च ।

\marginnote{\scriptsize १.४.८३}\noindent कर्मप्रवचनीयाः ।

\marginnote{\scriptsize १.४.८४}\noindent अनुर्लक्षणे ।

\marginnote{\scriptsize १.४.८५}\noindent तृतीयाऽर्थे ।

\marginnote{\scriptsize १.४.८६}\noindent हीने ।

\marginnote{\scriptsize १.४.८७}\noindent उपोऽधिके च ।

\marginnote{\scriptsize १.४.८८}\noindent अपपरी वर्जने ।

\marginnote{\scriptsize १.४.८९}\noindent आङ् मर्यादावचने ।

\marginnote{\scriptsize १.४.९०}\noindent लक्षणेत्थम्भूताख्यानभागवीप्सासु प्रतिपर्यनवः ।

\marginnote{\scriptsize १.४.९१}\noindent अभिरभागे ।

\marginnote{\scriptsize १.४.९२}\noindent प्रतिः प्रतिनिधिप्रतिदानयोः ।

\marginnote{\scriptsize १.४.९३}\noindent अधिपरी अनर्थकौ ।

\marginnote{\scriptsize १.४.९४}\noindent सुः पूजायाम् ।

\marginnote{\scriptsize १.४.९५}\noindent अतिरतिक्रमणे च ।

\marginnote{\scriptsize १.४.९६}\noindent अपिः पदार्थसम्भावनान्ववसर्गगर्हासमुच्चयेषु ।

\marginnote{\scriptsize १.४.९७}\noindent अधिरीश्वरे ।

\marginnote{\scriptsize १.४.९८}\noindent विभाषा कृञि ।

\marginnote{\scriptsize १.४.९९}\noindent लः परस्मैपदम् ।

\marginnote{\scriptsize १.४.१००}\noindent तङानावात्मनेपदम् ।

\marginnote{\scriptsize १.४.१०१}\noindent तिङस्त्रीणि त्रीणि प्रथममध्यमोत्तमाः ।

\marginnote{\scriptsize १.४.१०२}\noindent तान्येकवचनद्विवचनबहुवचनान्येकशः ।

\marginnote{\scriptsize १.४.१०३}\noindent सुपः ।

\marginnote{\scriptsize १.४.१०४}\noindent विभक्तिश्च ।

\marginnote{\scriptsize १.४.१०५}\noindent युष्मद्युपपदे समानाधिकरणे स्थानिन्यपि मध्यमः ।

\marginnote{\scriptsize १.४.१०६}\noindent प्रहासे च मन्योपपदे मन्यतेरुत्तम एकवच्च ।

\marginnote{\scriptsize १.४.१०७}\noindent अस्मद्युत्तमः ।

\marginnote{\scriptsize १.४.१०८}\noindent शेषे प्रथमः ।

\marginnote{\scriptsize १.४.१०९}\noindent परः संनिकर्षः संहिता ।

\marginnote{\scriptsize १.४.११०}\noindent विरामोऽवसानम् ।

\section*{॥ अध्याय २ ॥}

\marginnote{\scriptsize २.१.१}\noindent समर्थः पदविधिः ।

\marginnote{\scriptsize २.१.२}\noindent सुबामन्त्रिते पराङ्गवत् स्वरे ।

\marginnote{\scriptsize २.१.३}\noindent प्राक् कडारात् समासः ।

\marginnote{\scriptsize २.१.४}\noindent सह सुपा ।

\marginnote{\scriptsize २.१.५}\noindent अव्ययीभावः ।

\marginnote{\scriptsize २.१.६}\noindent अव्ययं विभक्तिसमीपसमृद्धि- व्यृद्ध्यर्थाभावात्ययासम्प्रति- शब्दप्रादुर्भावपश्चाद्यथाऽऽनुपूर्व्ययौगपद्यसादृश्य- सम्पत्तिसाकल्यान्तवचनेषु ।

\marginnote{\scriptsize २.१.७}\noindent यथाऽसादृये ।

\marginnote{\scriptsize २.१.८}\noindent यावदवधारणे ।

\marginnote{\scriptsize २.१.९}\noindent सुप्प्रतिना मात्राऽर्थे ।

\marginnote{\scriptsize २.१.१०}\noindent अक्षशलाकासंख्याः परिणा ।

\marginnote{\scriptsize २.१.११}\noindent विभाषा ।

\marginnote{\scriptsize २.१.१२}\noindent अपपरिबहिरञ्चवः पञ्चम्या ।

\marginnote{\scriptsize २.१.१३}\noindent आङ् मर्यादाऽभिविध्योः ।

\marginnote{\scriptsize २.१.१४}\noindent लक्षणेनाभिप्रती आभिमुख्ये ।

\marginnote{\scriptsize २.१.१५}\noindent अनुर्यत्समया ।

\marginnote{\scriptsize २.१.१६}\noindent यस्य चायामः ।

\marginnote{\scriptsize २.१.१७}\noindent तिष्ठद्गुप्रभृतीनि च ।

\marginnote{\scriptsize २.१.१८}\noindent पारे मध्ये षष्ठ्या वा ।

\marginnote{\scriptsize २.१.१९}\noindent संख्या वंश्येन ।

\marginnote{\scriptsize २.१.२०}\noindent नदीभिश्च ।

\marginnote{\scriptsize २.१.२१}\noindent अन्यपदार्थे च संज्ञायाम् ।

\marginnote{\scriptsize २.१.२२}\noindent तत्पुरुषः ।

\marginnote{\scriptsize २.१.२३}\noindent द्विगुश्च ।

\marginnote{\scriptsize २.१.२४}\noindent द्वितीया श्रितातीतपतितगतात्यस्तप्राप्तापन्नैः ।

\marginnote{\scriptsize २.१.२५}\noindent स्वयं क्तेन ।

\marginnote{\scriptsize २.१.२६}\noindent खट्वा क्षेपे ।

\marginnote{\scriptsize २.१.२७}\noindent सामि ।

\marginnote{\scriptsize २.१.२८}\noindent कालाः ।

\marginnote{\scriptsize २.१.२९}\noindent अत्यन्तसंयोगे च ।

\marginnote{\scriptsize २.१.३०}\noindent तृतीया तत्कृतार्थेन गुणवचनेन ।

\marginnote{\scriptsize २.१.३१}\noindent पूर्वसदृशसमोनार्थकलहनिपुणमिश्रश्लक्ष्णैः ।

\marginnote{\scriptsize २.१.३२}\noindent कर्तृकरणे कृता बहुलम् ।

\marginnote{\scriptsize २.१.३३}\noindent कृत्यैरधिकार्थवचने ।

\marginnote{\scriptsize २.१.३४}\noindent अन्नेन व्यञ्जनम् ।

\marginnote{\scriptsize २.१.३५}\noindent भक्ष्येण मिश्रीकरणम् ।

\marginnote{\scriptsize २.१.३६}\noindent चतुर्थी तदर्थार्थबलिहितसुखरक्षितैः ।

\marginnote{\scriptsize २.१.३७}\noindent पञ्चमी भयेन ।

\marginnote{\scriptsize २.१.३८}\noindent अपेतापोढमुक्तपतितापत्रस्तैरल्पशः ।

\marginnote{\scriptsize २.१.३९}\noindent स्तोकान्तिकदूरार्थकृच्छ्राणि क्तेन ।

\marginnote{\scriptsize २.१.४०}\noindent सप्तमी शौण्डैः ।

\marginnote{\scriptsize २.१.४१}\noindent सिद्धशुष्कपक्वबन्धैश्च ।

\marginnote{\scriptsize २.१.४२}\noindent ध्वाङ्क्षेण क्षेपे ।

\marginnote{\scriptsize २.१.४३}\noindent कृत्यैरृणे ।

\marginnote{\scriptsize २.१.४४}\noindent संज्ञायाम् ।

\marginnote{\scriptsize २.१.४५}\noindent क्तेनाहोरात्रावयवाः ।

\marginnote{\scriptsize २.१.४६}\noindent तत्र ।

\marginnote{\scriptsize २.१.४७}\noindent क्षेपे ।

\marginnote{\scriptsize २.१.४८}\noindent पात्रेसमितादयश्च ।

\marginnote{\scriptsize २.१.४९}\noindent पूर्वकालैकसर्वजरत्पुराणनवकेवलाः समानाधिकरणेन ।

\marginnote{\scriptsize २.१.५०}\noindent दिक्संख्ये संज्ञायाम् ।

\marginnote{\scriptsize २.१.५१}\noindent तद्धितार्थोत्तरपदसमाहारे च ।

\marginnote{\scriptsize २.१.५२}\noindent संख्यापूर्वो द्विगुः ।

\marginnote{\scriptsize २.१.५३}\noindent कुत्सितानि कुत्सनैः ।

\marginnote{\scriptsize २.१.५४}\noindent पापाणके कुत्सितैः ।

\marginnote{\scriptsize २.१.५५}\noindent उपमानानि सामान्यवचनैः ।

\marginnote{\scriptsize २.१.५६}\noindent उपमितं व्याघ्रादिभिः सामान्याप्रयोगे ।

\marginnote{\scriptsize २.१.५७}\noindent विशेषणं विशेष्येण बहुलम् ।

\marginnote{\scriptsize २.१.५८}\noindent पूर्वापरप्रथमचरमजघन्यसमान- मध्यमध्यमवीराश्च ।

\marginnote{\scriptsize २.१.५९}\noindent श्रेण्यादयः कृतादिभिः ।

\marginnote{\scriptsize २.१.६०}\noindent क्तेन नञ्विशिष्टेनानञ् ।

\marginnote{\scriptsize २.१.६१}\noindent सन्महत्परमोत्तमोत्कृष्टाः पूज्यमानैः ।

\marginnote{\scriptsize २.१.६२}\noindent वृन्दारकनागकुञ्जरैः पूज्यमानम् ।

\marginnote{\scriptsize २.१.६३}\noindent कतरकतमौ जातिपरिप्रश्ने ।

\marginnote{\scriptsize २.१.६४}\noindent किं क्षेपे ।

\marginnote{\scriptsize २.१.६५}\noindent पोटायुवतिस्तोककतिपयगृष्टिधेनुवशा- वेहत्बष्कयणीप्रवक्तॄ- श्रोत्रियाध्यापकधूर्तैर्जातिः ।

\marginnote{\scriptsize २.१.६६}\noindent प्रशंसावचनैश्च ।

\marginnote{\scriptsize २.१.६७}\noindent युवा खलतिपलितवलिनजरतीभिः ।

\marginnote{\scriptsize २.१.६८}\noindent कृत्यतुल्याख्या अजात्या ।

\marginnote{\scriptsize २.१.६९}\noindent वर्णो वर्णेन ।

\marginnote{\scriptsize २.१.७०}\noindent कुमारः श्रमणाऽऽदिभिः ।

\marginnote{\scriptsize २.१.७१}\noindent चतुष्पादो गर्भिण्या ।

\marginnote{\scriptsize २.१.७२}\noindent मयूरव्यंसकादयश्च ।

\marginnote{\scriptsize २.२.१}\noindent पूर्वापराधरोत्तरमेकदेशिनैकाधिकरणे ।

\marginnote{\scriptsize २.२.२}\noindent अर्धं नपुंसकम् ।

\marginnote{\scriptsize २.२.३}\noindent द्वितीयतृतीयचतुर्थतुर्याण्यन्यतरस्याम् ।

\marginnote{\scriptsize २.२.४}\noindent प्राप्तापन्ने च द्वितीयया ।

\marginnote{\scriptsize २.२.५}\noindent कालाः परिमाणिना ।

\marginnote{\scriptsize २.२.६}\noindent नञ् ।

\marginnote{\scriptsize २.२.७}\noindent ईषदकृता ।

\marginnote{\scriptsize २.२.८}\noindent षष्ठी ।

\marginnote{\scriptsize २.२.९}\noindent याजकादिभिश्च ।

\marginnote{\scriptsize २.२.१०}\noindent न निर्धारणे ।

\marginnote{\scriptsize २.२.११}\noindent पूरणगुणसुहितार्थसदव्ययतव्यसमानाधिकरणेन ।

\marginnote{\scriptsize २.२.१२}\noindent क्तेन च पूजायाम् ।

\marginnote{\scriptsize २.२.१३}\noindent अधिकरणवाचिना च ।

\marginnote{\scriptsize २.२.१४}\noindent कर्मणि च ।

\marginnote{\scriptsize २.२.१५}\noindent तृजकाभ्यां कर्तरि ।

\marginnote{\scriptsize २.२.१६}\noindent कर्त्तरि च ।

\marginnote{\scriptsize २.२.१७}\noindent नित्यं क्रीडाजीविकयोः ।

\marginnote{\scriptsize २.२.१८}\noindent कुगतिप्रादयः ।

\marginnote{\scriptsize २.२.१९}\noindent उपपदमतिङ् ।

\marginnote{\scriptsize २.२.२०}\noindent अमैवाव्ययेन ।

\marginnote{\scriptsize २.२.२१}\noindent तृतीयाप्रभृतीन्यन्यतरस्याम् ।

\marginnote{\scriptsize २.२.२२}\noindent क्त्वा च ।

\marginnote{\scriptsize २.२.२३}\noindent शेषो बहुव्रीहिः ।

\marginnote{\scriptsize २.२.२४}\noindent अनेकमन्यपदार्थे ।

\marginnote{\scriptsize २.२.२५}\noindent संख्ययाऽव्ययासन्नादूराधिकसंख्याः संख्येये ।

\marginnote{\scriptsize २.२.२६}\noindent दिङ्नामान्यन्तराले ।

\marginnote{\scriptsize २.२.२७}\noindent तत्र तेनेदमिति सरूपे ।

\marginnote{\scriptsize २.२.२८}\noindent तेन सहेति तुल्ययोगे ।

\marginnote{\scriptsize २.२.२९}\noindent चार्थे द्वंद्वः ।

\marginnote{\scriptsize २.२.३०}\noindent उपसर्जनं पूर्वम् ।

\marginnote{\scriptsize २.२.३१}\noindent राजदन्तादिषु परम् ।

\marginnote{\scriptsize २.२.३२}\noindent द्वंद्वे घि ।

\marginnote{\scriptsize २.२.३३}\noindent अजाद्यदन्तम् ।

\marginnote{\scriptsize २.२.३४}\noindent अल्पाच्तरम् ।

\marginnote{\scriptsize २.२.३५}\noindent सप्तमीविशेषणे बहुव्रीहौ ।

\marginnote{\scriptsize २.२.३६}\noindent निष्ठा ।

\marginnote{\scriptsize २.२.३७}\noindent वाऽऽहिताग्न्यादिषु ।

\marginnote{\scriptsize २.२.३८}\noindent कडाराः कर्मधराये ।

\marginnote{\scriptsize २.३.१}\noindent अनभिहिते ।

\marginnote{\scriptsize २.३.२}\noindent कर्मणि द्वितीया ।

\marginnote{\scriptsize २.३.३}\noindent तृतीया च होश्छन्दसि ।

\marginnote{\scriptsize २.३.४}\noindent अन्तराऽन्तरेण युक्ते ।

\marginnote{\scriptsize २.३.५}\noindent कालाध्वनोरत्यन्तसंयोगे ।

\marginnote{\scriptsize २.३.६}\noindent अपवर्गे तृतीया ।

\marginnote{\scriptsize २.३.७}\noindent सप्तमीपञ्चम्यौ कारकमध्ये ।

\marginnote{\scriptsize २.३.८}\noindent कर्मप्रवचनीययुक्ते द्वितीया ।

\marginnote{\scriptsize २.३.९}\noindent यस्मादधिकं यस्य चेश्वरवचनं तत्र सप्तमी ।

\marginnote{\scriptsize २.३.१०}\noindent पञ्चमी अपाङ्परिभिः ।

\marginnote{\scriptsize २.३.११}\noindent प्रतिनिधिप्रतिदाने च यस्मात् ।

\marginnote{\scriptsize २.३.१२}\noindent गत्यर्थकर्मणि द्वितीयाचतुर्थ्यौ चेष्टायामनध्वनि ।

\marginnote{\scriptsize २.३.१३}\noindent चतुर्थी सम्प्रदाने ।

\marginnote{\scriptsize २.३.१४}\noindent क्रियार्थोपपदस्य च कर्मणि स्थानिनः ।

\marginnote{\scriptsize २.३.१५}\noindent तुमर्थाच्च भाववचनात् ।

\marginnote{\scriptsize २.३.१६}\noindent नमःस्वस्तिस्वाहास्वधालंवषड्योगाच्च ।

\marginnote{\scriptsize २.३.१७}\noindent मन्यकर्मण्यनादरे विभाषाऽप्राणिषु ।

\marginnote{\scriptsize २.३.१८}\noindent कर्तृकरणयोस्तृतीया ।

\marginnote{\scriptsize २.३.१९}\noindent सहयुक्तेऽप्रधाने ।

\marginnote{\scriptsize २.३.२०}\noindent येनाङ्गविकारः ।

\marginnote{\scriptsize २.३.२१}\noindent इत्थंभूतलक्षणे ।

\marginnote{\scriptsize २.३.२२}\noindent संज्ञोऽन्यतरस्यां कर्मणि ।

\marginnote{\scriptsize २.३.२३}\noindent हेतौ ।

\marginnote{\scriptsize २.३.२४}\noindent अकर्तर्यृणे पञ्चमी ।

\marginnote{\scriptsize २.३.२५}\noindent विभाषा गुणेऽस्त्रियाम् ।

\marginnote{\scriptsize २.३.२६}\noindent षष्ठी हेतुप्रयोगे ।

\marginnote{\scriptsize २.३.२७}\noindent सर्वनाम्नस्तृतीया च ।

\marginnote{\scriptsize २.३.२८}\noindent अपादाने पञ्चमी ।

\marginnote{\scriptsize २.३.२९}\noindent अन्यारादितरर्त्तेदिक्शब्दाञ्चूत्तरपदाजाहियुक्ते ।

\marginnote{\scriptsize २.३.३०}\noindent षष्ठ्यतसर्थप्रत्ययेन ।

\marginnote{\scriptsize २.३.३१}\noindent एनपा द्वितीया ।

\marginnote{\scriptsize २.३.३२}\noindent पृथग्विनानानाभिस्तृतीयाऽन्यतरस्याम् ।

\marginnote{\scriptsize २.३.३३}\noindent करणे च स्तोकाल्पकृच्छ्रकतिपयस्यासत्त्ववचनस्य ।

\marginnote{\scriptsize २.३.३४}\noindent दूरान्तिकार्थैः षष्ठ्यन्यतरस्याम् ।

\marginnote{\scriptsize २.३.३५}\noindent दूरान्तिकार्थेभ्यो द्वितीया च ।

\marginnote{\scriptsize २.३.३६}\noindent सप्तम्यधिकरणे च ।

\marginnote{\scriptsize २.३.३७}\noindent यस्य च भावेन भावलक्षणम् ।

\marginnote{\scriptsize २.३.३८}\noindent षष्ठी चानादरे ।

\marginnote{\scriptsize २.३.३९}\noindent स्वामीश्वराधिपतिदायादसाक्षिप्रतिभूप्रसूतैश्च ।

\marginnote{\scriptsize २.३.४०}\noindent आयुक्तकुशलाभ्यां चासेवायाम् ।

\marginnote{\scriptsize २.३.४१}\noindent यतश्च निर्धारणम् ।

\marginnote{\scriptsize २.३.४२}\noindent पञ्चमी विभक्ते ।

\marginnote{\scriptsize २.३.४३}\noindent साधुनिपुणाभ्याम् अर्चायां सप्तम्यप्रतेः ।

\marginnote{\scriptsize २.३.४४}\noindent प्रसितोत्सुकाभ्यां तृतीया च ।

\marginnote{\scriptsize २.३.४५}\noindent नक्षत्रे च लुपि ।

\marginnote{\scriptsize २.३.४६}\noindent प्रातिपदिकार्थलिङ्गपरिमाणवचनमात्रे प्रथमा ।

\marginnote{\scriptsize २.३.४७}\noindent सम्बोधने च ।

\marginnote{\scriptsize २.३.४८}\noindent साऽऽमन्त्रितम् ।

\marginnote{\scriptsize २.३.४९}\noindent एकवचनं संबुद्धिः ।

\marginnote{\scriptsize २.३.५०}\noindent षष्ठी शेषे ।

\marginnote{\scriptsize २.३.५१}\noindent ज्ञोऽविदर्थस्य करणे ।

\marginnote{\scriptsize २.३.५२}\noindent अधीगर्थदयेशां कर्मणि ।

\marginnote{\scriptsize २.३.५३}\noindent कृञः प्रतियत्ने ।

\marginnote{\scriptsize २.३.५४}\noindent रुजार्थानां भाववचनानामज्वरेः ।

\marginnote{\scriptsize २.३.५५}\noindent आशिषि नाथः ।

\marginnote{\scriptsize २.३.५६}\noindent जासिनिप्रहणनाटक्राथपिषां हिंसायाम् ।

\marginnote{\scriptsize २.३.५७}\noindent व्यवहृपणोः समर्थयोः ।

\marginnote{\scriptsize २.३.५८}\noindent दिवस्तदर्थस्य ।

\marginnote{\scriptsize २.३.५९}\noindent विभाषोपसर्गे ।

\marginnote{\scriptsize २.३.६०}\noindent द्वितीया ब्राह्मणे ।

\marginnote{\scriptsize २.३.६१}\noindent प्रेष्यब्रुवोर्हविषो देवतासम्प्रदाने ।

\marginnote{\scriptsize २.३.६२}\noindent चतुर्थ्यर्थे बहुलं छन्दसि ।

\marginnote{\scriptsize २.३.६३}\noindent यजेश्च करणे ।

\marginnote{\scriptsize २.३.६४}\noindent कृत्वोऽर्थप्रयोगे कालेऽधिकरणे ।

\marginnote{\scriptsize २.३.६५}\noindent कर्तृकर्मणोः कृति ।

\marginnote{\scriptsize २.३.६६}\noindent उभयप्राप्तौ कर्मणि ।

\marginnote{\scriptsize २.३.६७}\noindent क्तस्य च वर्तमाने ।

\marginnote{\scriptsize २.३.६८}\noindent अधिकरणवाचिनश्च ।

\marginnote{\scriptsize २.३.६९}\noindent न लोकाव्ययनिष्ठाखलर्थतृनाम् ।

\marginnote{\scriptsize २.३.७०}\noindent अकेनोर्भविष्यदाधमर्ण्ययोः ।

\marginnote{\scriptsize २.३.७१}\noindent कृत्यानां कर्तरि वा ।

\marginnote{\scriptsize २.३.७२}\noindent तुल्यार्थैरतुलोपमाभ्यां तृतीयाऽन्यतरस्याम् ।

\marginnote{\scriptsize २.३.७३}\noindent चतुर्थी चाशिष्यायुष्यमद्रभद्र- कुशलसुखार्थहितैः ।

\marginnote{\scriptsize २.४.१}\noindent द्विगुरेकवचनम् ।

\marginnote{\scriptsize २.४.२}\noindent द्वंद्वश्च प्राणितूर्यसेनाङ्गानाम् ।

\marginnote{\scriptsize २.४.३}\noindent अनुवादे चरणानाम् ।

\marginnote{\scriptsize २.४.४}\noindent अध्वर्युक्रतुरनपुंसकम्। ।

\marginnote{\scriptsize २.४.५}\noindent अध्ययनतोऽविप्रकृष्टाख्यानाम् ।

\marginnote{\scriptsize २.४.६}\noindent जातिरप्राणिनाम् ।

\marginnote{\scriptsize २.४.७}\noindent विशिष्टलिङ्गो नदी देशोऽग्रामाः ।

\marginnote{\scriptsize २.४.८}\noindent क्षुद्रजन्तवः ।

\marginnote{\scriptsize २.४.९}\noindent येषां च विरोधः शाश्वतिकः ।

\marginnote{\scriptsize २.४.१०}\noindent शूद्राणामनिरवसितानाम् ।

\marginnote{\scriptsize २.४.११}\noindent गवाश्वप्रभृतीनि च ।

\marginnote{\scriptsize २.४.१२}\noindent विभाषा वृक्षमृगतृणधान्यव्यञ्जन- पशुशकुन्यश्ववडवपूर्वापराधरोत्तराणाम् ।

\marginnote{\scriptsize २.४.१३}\noindent विप्रतिषिद्धं चानधिकरणवाचि ।

\marginnote{\scriptsize २.४.१४}\noindent न दधिपयादीनि ।

\marginnote{\scriptsize २.४.१५}\noindent अधिकरणैतावत्त्वे च ।

\marginnote{\scriptsize २.४.१६}\noindent विभाषा समीपे ।

\marginnote{\scriptsize २.४.१७}\noindent स नपुंसकम् ।

\marginnote{\scriptsize २.४.१८}\noindent अव्ययीभावश्च ।

\marginnote{\scriptsize २.४.१९}\noindent तत्पुरुषोऽनञ् कर्मधारयः ।

\marginnote{\scriptsize २.४.२०}\noindent संज्ञायां कन्थोशीनरेषु ।

\marginnote{\scriptsize २.४.२१}\noindent उपज्ञोपक्रमं तदाद्याचिख्यासायाम् ।

\marginnote{\scriptsize २.४.२२}\noindent छाया बाहुल्ये ।

\marginnote{\scriptsize २.४.२३}\noindent सभा राजाऽमनुष्यपूर्वा ।

\marginnote{\scriptsize २.४.२४}\noindent अशाला च ।

\marginnote{\scriptsize २.४.२५}\noindent विभाषा सेनासुराछायाशालानिशानाम् ।

\marginnote{\scriptsize २.४.२६}\noindent परवल्लिङ्गं द्वंद्वतत्पुरुषयोः ।

\marginnote{\scriptsize २.३.२७}\noindent पूर्ववदश्ववडवौ ।

\marginnote{\scriptsize २.४.२८}\noindent हेमन्तशिशिरावहोरात्रे च च्छन्दसि ।

\marginnote{\scriptsize २.४.२९}\noindent रात्राह्नाहाः पुंसि ।

\marginnote{\scriptsize २.४.३०}\noindent अपथं नपुंसकम् ।

\marginnote{\scriptsize २.४.३१}\noindent अर्धर्चाः पुंसि च ।

\marginnote{\scriptsize २.४.३२}\noindent इदमोऽन्वादेशेऽशनुदात्तस्तृतीयाऽऽदौ ।

\marginnote{\scriptsize २.४.३३}\noindent एतदस्त्रतसोस्त्रतसौ चानुदात्तौ ।

\marginnote{\scriptsize २.४.३४}\noindent द्वितीयाटौस्स्वेनः ।

\marginnote{\scriptsize २.४.३५}\noindent आर्द्धधातुके ।

\marginnote{\scriptsize २.४.३६}\noindent अदो जग्धिर्ल्यप्ति किति ।

\marginnote{\scriptsize २.४.३७}\noindent लुङ्सनोर्घसॢ ।

\marginnote{\scriptsize २.४.३८}\noindent घञपोश्च ।

\marginnote{\scriptsize २.४.३९}\noindent बहुलं छन्दसि ।

\marginnote{\scriptsize २.४.४०}\noindent लिट्यन्यतरस्याम् ।

\marginnote{\scriptsize २.४.४१}\noindent वेञो वयिः ।

\marginnote{\scriptsize २.४.४२}\noindent हनो वध लिङि ।

\marginnote{\scriptsize २.४.४३}\noindent लुङि च ।

\marginnote{\scriptsize २.४.४४}\noindent आत्मनेपदेष्वन्यतरस्याम् ।

\marginnote{\scriptsize २.४.४५}\noindent इणो गा लुङि ।

\marginnote{\scriptsize २.४.४६}\noindent णौ गमिरबोधने ।

\marginnote{\scriptsize २.४.४७}\noindent सनि च ।

\marginnote{\scriptsize २.४.४८}\noindent इङश्च ।

\marginnote{\scriptsize २.४.४९}\noindent गाङ् लिटि ।

\marginnote{\scriptsize २.४.५०}\noindent विभाषा लुङ्लृङोः ।

\marginnote{\scriptsize २.४.५१}\noindent णौ च सँश्चङोः ।

\marginnote{\scriptsize २.४.५२}\noindent अस्तेर्भूः ।

\marginnote{\scriptsize २.४.५३}\noindent ब्रुवो वचिः ।

\marginnote{\scriptsize २.४.५४}\noindent चक्षिङः ख्याञ् ।

\marginnote{\scriptsize २.४.५५}\noindent वा लिटि ।

\marginnote{\scriptsize २.४.५६}\noindent अजेर्व्यघञपोः ।

\marginnote{\scriptsize २.४.५७}\noindent वा यौ ।

\marginnote{\scriptsize २.४.५८}\noindent ण्यक्षत्रियार्षञितो यूनि लुगणिञोः ।

\marginnote{\scriptsize २.४.५९}\noindent पैलादिभ्यश्च ।

\marginnote{\scriptsize २.४.६०}\noindent इञः प्राचाम् ।

\marginnote{\scriptsize २.४.६१}\noindent न तौल्वलिभ्यः ।

\marginnote{\scriptsize २.४.६२}\noindent तद्राजस्य बहुषु तेनैवास्त्रियाम् ।

\marginnote{\scriptsize २.४.६३}\noindent यस्कादिभ्यो गोत्रे ।

\marginnote{\scriptsize २.४.६४}\noindent यञञोश्च ।

\marginnote{\scriptsize २.४.६५}\noindent अत्रिभृगुकुत्सवसिष्ठगोतमाङ्गिरोभ्यश्च ।

\marginnote{\scriptsize २.४.६६}\noindent बह्वचः इञः प्राच्यभरतेषु ।

\marginnote{\scriptsize २.४.६७}\noindent न गोपवनादिभ्यः ।

\marginnote{\scriptsize २.४.६८}\noindent तिककितवादिभ्यो द्वंद्वे ।

\marginnote{\scriptsize २.४.६९}\noindent उपकादिभ्योऽन्यतरस्यामद्वंद्वे ।

\marginnote{\scriptsize २.४.७०}\noindent आगस्त्यकौण्डिन्ययोरगस्तिकुण्डिनच् ।

\marginnote{\scriptsize २.४.७१}\noindent सुपो धातुप्रातिपदिकयोः ।

\marginnote{\scriptsize २.४.७२}\noindent अदिप्रभृतिभ्यः शपः ।

\marginnote{\scriptsize २.४.७३}\noindent बहुलं छन्दसि ।

\marginnote{\scriptsize २.४.७४}\noindent यङोऽचि च ।

\marginnote{\scriptsize २.४.७५}\noindent जुहोत्यादिभ्यः श्लुः ।

\marginnote{\scriptsize २.४.७६}\noindent बहुलं छन्दसि ।

\marginnote{\scriptsize २.४.७७}\noindent गातिस्थाघुपाभूभ्यः सिचः परस्मैपदेषु ।

\marginnote{\scriptsize २.४.७८}\noindent विभाषा घ्राधेट्शाच्छासः ।

\marginnote{\scriptsize २.४.७९}\noindent तनादिभ्यस्तथासोः ।

\marginnote{\scriptsize २.४.८०}\noindent मन्त्रे घसह्वरणशवृदहाद्वृच्कृगमिजनिभ्यो लेः ।

\marginnote{\scriptsize २.४.८१}\noindent आमः ।

\marginnote{\scriptsize २.४.८२}\noindent अव्ययादाप्सुपः ।

\marginnote{\scriptsize २.४.८३}\noindent नाव्ययीभावादतोऽम्त्वपञ्चम्याः ।

\marginnote{\scriptsize २.४.८४}\noindent तृतीयासप्तम्योर्बहुलम् ।

\marginnote{\scriptsize २.४.८५}\noindent लुटः प्रथमस्य डारौरसः ।

\section*{॥ अध्याय ३ ॥}

\marginnote{\scriptsize ३.१.१}\noindent प्रत्ययः ।

\marginnote{\scriptsize ३.१.२}\noindent परश्च ।

\marginnote{\scriptsize ३.१.३}\noindent आद्युदात्तश्च ।

\marginnote{\scriptsize ३.१.४}\noindent अनुदत्तौ सुप्पितौ ।

\marginnote{\scriptsize ३.१.५}\noindent गुप्तिज्किद्भ्यः सन् ।

\marginnote{\scriptsize ३.१.६}\noindent मान्बधदान्शान्भ्यो दीर्घश्चाभ्यासस्य ।

\marginnote{\scriptsize ३.१.७}\noindent धातोः कर्मणः समानकर्तृकादिच्छायां वा ।

\marginnote{\scriptsize ३.१.८}\noindent सुप आत्मनः क्यच् ।

\marginnote{\scriptsize ३.१.९}\noindent काम्यच्च ।

\marginnote{\scriptsize ३.१.१०}\noindent उपमानादाचारे ।

\marginnote{\scriptsize ३.१.११}\noindent कर्तुः क्यङ् सलोपश्च ।

\marginnote{\scriptsize ३.१.१२}\noindent भृशादिभ्यो भुव्यच्वेर्लोपश्च हलः ।

\marginnote{\scriptsize ३.१.१३}\noindent लोहितादिडाज्भ्यः क्यष्।

\marginnote{\scriptsize ३.१.१४}\noindent कष्टाय क्रमणे ।

\marginnote{\scriptsize ३.१.१५}\noindent कर्मणः रोमन्थतपोभ्यां वर्तिचरोः ।

\marginnote{\scriptsize ३.१.१६}\noindent बाष्पोष्माभ्यां उद्वमने ।

\marginnote{\scriptsize ३.१.१७}\noindent शब्दवैरकलहाभ्रकण्वमेघेभ्यः करणे ।

\marginnote{\scriptsize ३.१.१८}\noindent सुखादिभ्यः कर्तृवेदनायाम् ।

\marginnote{\scriptsize ३.१.१९}\noindent नमोवरिवश्चित्रङः क्यच् ।

\marginnote{\scriptsize ३.१.२०}\noindent पुच्छभाण्डचीवराण्णिङ् ।

\marginnote{\scriptsize ३.१.२१}\noindent मुण्डमिश्रश्लक्ष्णलवणव्रतवस्त्रहलकलकृततूस्तेभ्यो णिच् ।

\marginnote{\scriptsize ३.१.२२}\noindent धातोरेकाचो हलादेः क्रियासमभिहारे यङ् ।

\marginnote{\scriptsize ३.१.२३}\noindent नित्यं कौटिल्ये गतौ ।

\marginnote{\scriptsize ३.१.२४}\noindent लुपसदचरजपजभदहदशगॄभ्यो भावगर्हायाम् ।

\marginnote{\scriptsize ३.१.२५}\noindent सत्यापपाशरूपवीणातूलश्लोकसेनालोमत्वचवर्मवर्ण- चूर्णचुरादिभ्यो णिच् ।

\marginnote{\scriptsize ३.१.२६}\noindent हेतुमति च ।

\marginnote{\scriptsize ३.१.२७}\noindent कण्ड्वादिभ्यो यक् ।

\marginnote{\scriptsize ३.१.२८}\noindent गुपूधूपविच्छिपणिपनिभ्य आयः ।

\marginnote{\scriptsize ३.१.२९}\noindent ऋतेरीयङ् ।

\marginnote{\scriptsize ३.१.३०}\noindent कमेर्णिङ् ।

\marginnote{\scriptsize ३.१.३१}\noindent आयादय आर्धद्धातुके वा ।

\marginnote{\scriptsize ३.१.३२}\noindent सनाद्यन्ता धातवः ।

\marginnote{\scriptsize ३.१.३३}\noindent स्यतासी लृलुटोः ।

\marginnote{\scriptsize ३.१.३४}\noindent सिब्बहुलं लेटि ।

\marginnote{\scriptsize ३.१.३५}\noindent कास्प्रत्ययादाममन्त्रे लिटि ।

\marginnote{\scriptsize ३.१.३६}\noindent इजादेश्च गुरुमतोऽनृच्छः ।

\marginnote{\scriptsize ३.१.३७}\noindent दयायासश्च ।

\marginnote{\scriptsize ३.१.३८}\noindent उषविदजागृभ्योऽन्यतरस्याम् ।

\marginnote{\scriptsize ३.१.३९}\noindent भीह्रीभृहुवां श्लुवच्च ।

\marginnote{\scriptsize ३.१.४०}\noindent कृञ् चानुप्रयुज्यते लिटि ।

\marginnote{\scriptsize ३.१.४१}\noindent विदाङ्कुर्वन्त्वित्यन्यतरस्याम् ।

\marginnote{\scriptsize ३.१.४२}\noindent अभ्युत्सादयांप्रजनयांचिकयांरमयामकः पावयांक्रियाद्विदामक्रन्निति च्छन्दसि ।

\marginnote{\scriptsize ३.१.४३}\noindent च्लि लुङि ।

\marginnote{\scriptsize ३.१.४४}\noindent च्लेः सिच् ।

\marginnote{\scriptsize ३.१.४५}\noindent शल इगुपधादनिटः क्सः ।

\marginnote{\scriptsize ३.१.४६}\noindent श्लिष आलिङ्गने ।

\marginnote{\scriptsize ३.१.४७}\noindent न दृशः ।

\marginnote{\scriptsize ३.१.४८}\noindent णिश्रिद्रुस्रुभ्यः कर्तरि चङ् ।

\marginnote{\scriptsize ३.१.४९}\noindent विभाषा धेट्श्व्योः ।

\marginnote{\scriptsize ३.१.५०}\noindent गुपेश्छन्दसि ।

\marginnote{\scriptsize ३.१.५१}\noindent नोनयतिध्वनयत्येलयत्यर्दयतिभ्यः ।

\marginnote{\scriptsize ३.१.५२}\noindent अस्यतिवक्तिख्यातिभ्यः अङ् ।

\marginnote{\scriptsize ३.१.५३}\noindent लिपिसिचिह्वश्च ।

\marginnote{\scriptsize ३.१.५४}\noindent आत्मनेपदेष्वन्यतरस्याम् ।

\marginnote{\scriptsize ३.१.५५}\noindent पुषादिद्युताद्यॢदितः परस्मैपदेषु ।

\marginnote{\scriptsize ३.१.५६}\noindent सर्त्तिशास्त्यर्तिभ्यश्च ।

\marginnote{\scriptsize ३.१.५७}\noindent इरितो वा ।

\marginnote{\scriptsize ३.१.५८}\noindent जृस्तम्भुम्रुचुम्लुचुग्रुचुग्लुचुग्लुञ्चुश्विभ्यश्च ।

\marginnote{\scriptsize ३.१.५९}\noindent कृमृदृरुहिभ्यश्छन्दसि ।

\marginnote{\scriptsize ३.१.६०}\noindent चिण् ते पदः ।

\marginnote{\scriptsize ३.१.६१}\noindent दीपजनबुधपूरितायिप्यायिभ्योऽन्यतरस्याम् ।

\marginnote{\scriptsize ३.१.६२}\noindent अचः कर्मकर्तरि ।

\marginnote{\scriptsize ३.१.६३}\noindent दुहश्च ।

\marginnote{\scriptsize ३.१.६४}\noindent न रुधः ।

\marginnote{\scriptsize ३.१.६५}\noindent तपोऽनुतापे च ।

\marginnote{\scriptsize ३.१.६६}\noindent चिण् भावकर्मणोः ।

\marginnote{\scriptsize ३.१.६७}\noindent सार्वधातुके यक् ।

\marginnote{\scriptsize ३.१.६८}\noindent कर्तरि शप् ।

\marginnote{\scriptsize ३.१.६९}\noindent दिवादिभ्यः श्यन् ।

\marginnote{\scriptsize ३.१.७०}\noindent वा भ्राशभ्लाशभ्रमुक्रमुक्लमुत्रसित्रुटिलषः ।

\marginnote{\scriptsize ३.१.७१}\noindent यसोऽनुपसर्गात् ।

\marginnote{\scriptsize ३.१.७२}\noindent संयसश्च ।

\marginnote{\scriptsize ३.१.७३}\noindent स्वादिभ्यः श्नुः ।

\marginnote{\scriptsize ३.१.७४}\noindent श्रुवः श‍ृ च ।

\marginnote{\scriptsize ३.१.७५}\noindent अक्षोऽन्यतरस्याम् ।

\marginnote{\scriptsize ३.१.७६}\noindent तनूकरणे तक्षः ।

\marginnote{\scriptsize ३.१.७७}\noindent तुदादिभ्यः शः ।

\marginnote{\scriptsize ३.१.७८}\noindent रुधादिभ्यः श्नम् ।

\marginnote{\scriptsize ३.१.७९}\noindent तनादिकृञ्भ्य उः ।

\marginnote{\scriptsize ३.१.८०}\noindent धिन्विकृण्व्योर च ।

\marginnote{\scriptsize ३.१.८१}\noindent क्र्यादिभ्यः श्ना ।

\marginnote{\scriptsize ३.१.८२}\noindent स्तम्भुस्तुम्भुस्कम्भुस्कुम्भुस्कुञ्भ्यः श्नुश्च ।

\marginnote{\scriptsize ३.१.८३}\noindent हलः श्नः शानज्झौ ।

\marginnote{\scriptsize ३.१.८४}\noindent छन्दसि शायजपि ।

\marginnote{\scriptsize ३.१.८५}\noindent व्यत्ययो बहुलम् ।

\marginnote{\scriptsize ३.१.८६}\noindent लिङ्याशिष्यङ् ।

\marginnote{\scriptsize ३.१.८७}\noindent कर्मवत् कर्मणा तुल्यक्रियः ।

\marginnote{\scriptsize ३.१.८८}\noindent तपस्तपःकर्मकस्यैव ।

\marginnote{\scriptsize ३.१.८९}\noindent न दुहस्नुनमां यक्चिणौ ।

\marginnote{\scriptsize ३.१.९०}\noindent कुषिरजोः प्राचां श्यन् परस्मैपदं च ।

\marginnote{\scriptsize ३.१.९१}\noindent धातोः ।

\marginnote{\scriptsize ३.१.९२}\noindent तत्रोपपदं सप्तमीस्थम् ।

\marginnote{\scriptsize ३.१.९३}\noindent कृदतिङ् ।

\marginnote{\scriptsize ३.१.९४}\noindent वाऽसरूपोऽस्त्रियाम् ।

\marginnote{\scriptsize ३.१.९५}\noindent कृत्याः प्राङ् ण्वुलः ।

\marginnote{\scriptsize ३.१.९६}\noindent तव्यत्तव्यानीयरः ।

\marginnote{\scriptsize ३.१.९७}\noindent अचो यत् ।

\marginnote{\scriptsize ३.१.९८}\noindent पोरदुपधात् ।

\marginnote{\scriptsize ३.१.९९}\noindent शकिसहोश्च ।

\marginnote{\scriptsize ३.१.१००}\noindent गदमदचरयमश्चानुपसर्गे ।

\marginnote{\scriptsize ३.१.१०१}\noindent अवद्यपण्यवर्या गर्ह्यपणितव्यानिरोधेषु ।

\marginnote{\scriptsize ३.१.१०२}\noindent वह्यं करणम् ।

\marginnote{\scriptsize ३.१.१०३}\noindent अर्यः स्वामिवैश्ययोः ।

\marginnote{\scriptsize ३.१.१०४}\noindent उपसर्या काल्या प्रजने ।

\marginnote{\scriptsize ३.१.१०५}\noindent अजर्यं संगतम् ।

\marginnote{\scriptsize ३.१.१०६}\noindent वदः सुपि क्यप् च ।

\marginnote{\scriptsize ३.१.१०७}\noindent भुवो भावे ।

\marginnote{\scriptsize ३.१.१०८}\noindent हनस्त च ।

\marginnote{\scriptsize ३.१.१०९}\noindent एतिस्तुशस्वृदृजुषः क्यप् ।

\marginnote{\scriptsize ३.१.११०}\noindent ऋदुपधाच्चाकॢपिचृतेः ।

\marginnote{\scriptsize ३.१.१११}\noindent ई च खनः ।

\marginnote{\scriptsize ३.१.११२}\noindent भृञोऽसंज्ञायाम् ।

\marginnote{\scriptsize ३.१.११३}\noindent मृजेर्विभाषा ।

\marginnote{\scriptsize ३.१.११४}\noindent राजसूयसूर्यमृषोद्यरुच्यकुप्यकृष्टपच्याव्यथ्याः ।

\marginnote{\scriptsize ३.१.११५}\noindent भिद्योद्ध्यौ नदे ।

\marginnote{\scriptsize ३.१.११६}\noindent पुष्यसिद्ध्यौ नक्षत्रे ।

\marginnote{\scriptsize ३.१.११७}\noindent विपूयविनीयजित्या मुञ्जकल्कहलिषु ।

\marginnote{\scriptsize ३.१.११८}\noindent प्रत्यपिभ्यां ग्रहेश्छन्दसि ।

\marginnote{\scriptsize ३.१.११९}\noindent पदास्वैरिबाह्यापक्ष्येषु च ।

\marginnote{\scriptsize ३.१.१२०}\noindent विभाषा कृवृषोः ।

\marginnote{\scriptsize ३.१.१२१}\noindent युग्यं च पत्त्रे ।

\marginnote{\scriptsize ३.१.१२२}\noindent अमावस्यदन्यतरस्याम् ।

\marginnote{\scriptsize ३.१.१२३}\noindent छन्दसि निष्टर्क्यदेवहूयप्रणीयोन्नीयोच्छिष्य मर्यस्तर्याध्वर्यखन्यखान्यदेवयज्याऽऽपृच्छ्य प्रतिषीव्यब्रह्मवाद्यभाव्यस्ताव्योपचाय्यपृडानि ।

\marginnote{\scriptsize ३.१.१२४}\noindent ऋहलोर्ण्यत् ।

\marginnote{\scriptsize ३.१.१२५}\noindent ओरावश्यके ।

\marginnote{\scriptsize ३.१.१२६}\noindent आसुयुवपिरपिलपित्रपिचमश्च ।

\marginnote{\scriptsize ३.१.१२७}\noindent आनाय्योऽनित्ये ।

\marginnote{\scriptsize ३.१.१२८}\noindent प्रणाय्योऽसंमतौ ।

\marginnote{\scriptsize ३.१.१२९}\noindent पाय्यसान्नाय्यनिकाय्यधाय्या मानहविर्निवाससामिधेनीषु ।

\marginnote{\scriptsize ३.१.१३०}\noindent क्रतौ कुण्डपाय्यसंचाय्यौ ।

\marginnote{\scriptsize ३.१.१३१}\noindent अग्नौ परिचाय्योपचाय्यसमूह्याः ।

\marginnote{\scriptsize ३.१.१३२}\noindent चित्याग्निचित्ये च ।

\marginnote{\scriptsize ३.१.१३३}\noindent ण्वुल्तृचौ ।

\marginnote{\scriptsize ३.१.१३४}\noindent नन्दिग्रहिपचादिभ्यो ल्युणिन्यचः ।

\marginnote{\scriptsize ३.१.१३५}\noindent इगुपधज्ञाप्रीकिरः कः ।

\marginnote{\scriptsize ३.१.१३६}\noindent आतश्चोपसर्गे ।

\marginnote{\scriptsize ३.१.१३७}\noindent पाघ्राध्माधेट्दृशः शः ।

\marginnote{\scriptsize ३.१.१३८}\noindent अनुपसर्गाल्लिम्पविन्दधारिपारिवेद्युदेजिचेति- सातिसाहिभ्यश्च ।

\marginnote{\scriptsize ३.१.१३९}\noindent ददातिदधात्योर्विभाषा ।

\marginnote{\scriptsize ३.१.१४०}\noindent ज्वलितिकसन्तेभ्यो णः ।

\marginnote{\scriptsize ३.१.१४१}\noindent श्याऽऽद्व्यधास्रुसंस्र्वतीणवसाऽवहृलिह- श्लिषश्वसश्च ।

\marginnote{\scriptsize ३.१.१४२}\noindent दुन्योरनुपसर्गे ।

\marginnote{\scriptsize ३.१.१४३}\noindent विभाषा ग्रहेः ।

\marginnote{\scriptsize ३.१.१४४}\noindent गेहे कः ।

\marginnote{\scriptsize ३.१.१४५}\noindent शिल्पिनि ष्वुन् ।

\marginnote{\scriptsize ३.१.१४६}\noindent गस्थकन् ।

\marginnote{\scriptsize ३.१.१४७}\noindent ण्युट् च ।

\marginnote{\scriptsize ३.१.१४८}\noindent हश्च व्रीहिकालयोः ।

\marginnote{\scriptsize ३.१.१४९}\noindent प्रुसृल्वः समभिहारे वुन् ।

\marginnote{\scriptsize ३.१.१५०}\noindent आशिषि च ।

\marginnote{\scriptsize ३.२.१}\noindent कर्मण्यण् ।

\marginnote{\scriptsize ३.२.२}\noindent ह्वावामश्च ।

\marginnote{\scriptsize ३.२.३}\noindent आतोऽनुपसर्गे कः ।

\marginnote{\scriptsize ३.२.४}\noindent सुपि स्थः ।

\marginnote{\scriptsize ३.२.५}\noindent तुन्दशोकयोः परिमृजापनुदोः ।

\marginnote{\scriptsize ३.२.६}\noindent प्रे दाज्ञः ।

\marginnote{\scriptsize ३.२.७}\noindent समि ख्यः ।

\marginnote{\scriptsize ३.२.८}\noindent गापोष्टक् ।

\marginnote{\scriptsize ३.२.९}\noindent हरतेरनुद्यमनेऽच् ।

\marginnote{\scriptsize ३.२.१०}\noindent वयसि च ।

\marginnote{\scriptsize ३.२.११}\noindent आङि ताच्छील्ये ।

\marginnote{\scriptsize ३.२.१२}\noindent अर्हः ।

\marginnote{\scriptsize ३.२.१३}\noindent स्तम्बकर्णयोः रमिजपोः ।

\marginnote{\scriptsize ३.२.१४}\noindent शमि धातोः संज्ञायाम् ।

\marginnote{\scriptsize ३.२.१५}\noindent अधिकरणे शेतेः ।

\marginnote{\scriptsize ३.२.१६}\noindent चरेष्टः ।

\marginnote{\scriptsize ३.२.१७}\noindent भिक्षासेनाऽऽदायेषु च ।

\marginnote{\scriptsize ३.२.१८}\noindent पुरोऽग्रतोऽग्रेषु सर्तेः ।

\marginnote{\scriptsize ३.२.१९}\noindent पूर्वे कर्तरि ।

\marginnote{\scriptsize ३.२.२०}\noindent कृञो हेतुताच्छील्यानुलोम्येषु ।

\marginnote{\scriptsize ३.२.२१}\noindent दिवाविभानिशाप्रभाभास्करान्तानन्तादिबहुनान्दी- किम्लिपि लिबिबलिभक्तिकर्तृचित्रक्षेत्र- संख्याजङ्घाबाह्वहर्यत्तत्धनुररुष्षु ।

\marginnote{\scriptsize ३.२.२२}\noindent कर्मणि भृतौ ।

\marginnote{\scriptsize ३.२.२३}\noindent न शब्दश्लोककलहगाथावैरचाटुसूत्रमन्त्रपदेषु ।

\marginnote{\scriptsize ३.२.२४}\noindent स्तम्बशकृतोरिन् ।

\marginnote{\scriptsize ३.२.२५}\noindent हरतेर्दृतिनाथयोः पशौ ।

\marginnote{\scriptsize ३.२.२६}\noindent फलेग्रहिरात्मम्भरिश्च ।

\marginnote{\scriptsize ३.२.२७}\noindent छन्दसि वनसनरक्षिमथाम् ।

\marginnote{\scriptsize ३.२.२८}\noindent एजेः खश् ।

\marginnote{\scriptsize ३.२.२९}\noindent नासिकास्तनयोर्ध्माधेटोः ।

\marginnote{\scriptsize ३.२.३०}\noindent नाडीमुष्ट्योश्च ।

\marginnote{\scriptsize ३.२.३१}\noindent उदि कूले रुजिवहोः ।

\marginnote{\scriptsize ३.२.३२}\noindent वहाभ्रे लिहः ।

\marginnote{\scriptsize ३.२.३३}\noindent परिमाणे पचः ।

\marginnote{\scriptsize ३.२.३४}\noindent मितनखे च ।

\marginnote{\scriptsize ३.२.३५}\noindent विध्वरुषोः तुदः ।

\marginnote{\scriptsize ३.२.३६}\noindent असूर्यललाटयोर्दृशितपोः ।

\marginnote{\scriptsize ३.२.३७}\noindent उग्रम्पश्येरम्मदपाणिन्धमाश्च ।

\marginnote{\scriptsize ३.२.३८}\noindent प्रियवशे वदः खच् ।

\marginnote{\scriptsize ३.२.३९}\noindent द्विषत्परयोस्तापेः ।

\marginnote{\scriptsize ३.२.४०}\noindent वाचि यमो व्रते ।

\marginnote{\scriptsize ३.२.४१}\noindent पूःसर्वयोर्दारिसहोः ।

\marginnote{\scriptsize ३.२.४२}\noindent सर्वकूलाभ्रकरीषेषु कषः ।

\marginnote{\scriptsize ३.२.४३}\noindent मेघर्तिभयेषु कृञः ।

\marginnote{\scriptsize ३.२.४४}\noindent क्षेमप्रियमद्रेऽण् च ।

\marginnote{\scriptsize ३.२.४५}\noindent आशिते भुवः करणभावयोः ।

\marginnote{\scriptsize ३.२.४६}\noindent संज्ञायां भृतॄवृजिधारिसहितपिदमः ।

\marginnote{\scriptsize ३.२.४७}\noindent गमश्च ।

\marginnote{\scriptsize ३.२.४८}\noindent अन्तात्यन्ताध्वदूरपारसर्वानन्तेषु डः ।

\marginnote{\scriptsize ३.२.४९}\noindent आशिषि हनः ।

\marginnote{\scriptsize ३.२.५०}\noindent अपे क्लेशतमसोः ।

\marginnote{\scriptsize ३.२.५१}\noindent कुमारशीर्षयोर्णिनिः ।

\marginnote{\scriptsize ३.२.५२}\noindent लक्षणे जायापत्योष्टक् ।

\marginnote{\scriptsize ३.२.५३}\noindent अमनुष्यकर्तृके च ।

\marginnote{\scriptsize ३.२.५४}\noindent शक्तौ हस्तिकपाटयोः ।

\marginnote{\scriptsize ३.२.५५}\noindent पाणिघताडघौ शिल्पिनि ।

\marginnote{\scriptsize ३.२.५६}\noindent आढ्यसुभगस्थूलपलितनग्नान्धप्रियेषु च्व्य्र्थेष्वच्वौ कृञः करणे ख्युन् ।

\marginnote{\scriptsize ३.२.५७}\noindent कर्तरि भुवः खिष्णुच्खुकञौ ।

\marginnote{\scriptsize ३.२.५८}\noindent स्पृशोऽनुदके क्विन् ।

\marginnote{\scriptsize ३.२.५९}\noindent ऋत्विग्दधृक्स्रग्दिगुष्णिगञ्चुयुजिक्रुञ्चां च ।

\marginnote{\scriptsize ३.२.६०}\noindent त्यदादिषु दृशोऽनालोचने कञ् च । ३.२.६१ सत्सूद्विषद्रुहदुहयुजविदभिदच्छिद-जिनीराजामुपसर्गेऽपि क्विप् ।

\marginnote{\scriptsize ३.२.६२}\noindent भजो ण्विः ।

\marginnote{\scriptsize ३.२.६३}\noindent छन्दसि सहः ।

\marginnote{\scriptsize ३.२.६४}\noindent वहश्च ।

\marginnote{\scriptsize ३.२.६५}\noindent कव्यपुरीषपुरीष्येषु ञ्युट् ।

\marginnote{\scriptsize ३.२.६६}\noindent हव्येऽनन्तः पादम् ।

\marginnote{\scriptsize ३.२.६७}\noindent जनसनखनक्रमगमो विट् ।

\marginnote{\scriptsize ३.२.६८}\noindent अदोऽनन्ने ।

\marginnote{\scriptsize ३.२.६९}\noindent क्रव्ये च ।

\marginnote{\scriptsize ३.२.७०}\noindent दुहः कब् घश्च ।

\marginnote{\scriptsize ३.२.७१}\noindent मन्त्रे श्वेतवहौक्थशस्पुरोडाशो ण्विन् ।

\marginnote{\scriptsize ३.२.७२}\noindent अवे यजः ।

\marginnote{\scriptsize ३.२.७३}\noindent विजुपे छन्दसि ।

\marginnote{\scriptsize ३.२.७४}\noindent आतो मनिन्क्वनिप्वनिपश्च ।

\marginnote{\scriptsize ३.२.७५}\noindent अन्येभ्योऽपि दृश्यन्ते ।

\marginnote{\scriptsize ३.२.७६}\noindent क्विप् च ।

\marginnote{\scriptsize ३.२.७७}\noindent स्थः क च ।

\marginnote{\scriptsize ३.२.७८}\noindent सुप्यजातौ णिनिस्ताच्छिल्ये ।

\marginnote{\scriptsize ३.२.७९}\noindent कर्तर्युपमाने ।

\marginnote{\scriptsize ३.२.८०}\noindent व्रते ।

\marginnote{\scriptsize ३.२.८१}\noindent बहुलमाभीक्ष्ण्ये ।

\marginnote{\scriptsize ३.२.८२}\noindent मनः ।

\marginnote{\scriptsize ३.२.८३}\noindent आत्ममाने खश्च ।

\marginnote{\scriptsize ३.२.८४}\noindent भूते ।

\marginnote{\scriptsize ३.२.८५}\noindent करणे यजः ।

\marginnote{\scriptsize ३.२.८६}\noindent कर्मणि हनः ।

\marginnote{\scriptsize ३.२.८७}\noindent ब्रह्मभ्रूणवृत्रेषु क्विप् ।

\marginnote{\scriptsize ३.२.८८}\noindent बहुलं छन्दसि ।

\marginnote{\scriptsize ३.२.८९}\noindent सुकर्मपापमन्त्रपुण्येषु कृञः ।

\marginnote{\scriptsize ३.२.९०}\noindent सोमे सुञः ।

\marginnote{\scriptsize ३.२.९१}\noindent अग्नौ चेः ।

\marginnote{\scriptsize ३.२.९२}\noindent कर्मण्यग्न्याख्यायाम् ।

\marginnote{\scriptsize ३.२.९३}\noindent कर्मणीनिर्विक्रियः ।

\marginnote{\scriptsize ३.२.९४}\noindent दृशेः क्वनिप् ।

\marginnote{\scriptsize ३.२.९५}\noindent राजनि युधिकृञः ।

\marginnote{\scriptsize ३.२.९६}\noindent सहे च ।

\marginnote{\scriptsize ३.२.९७}\noindent सप्तम्यां जनेर्डः ।

\marginnote{\scriptsize ३.२.९८}\noindent पञ्चम्यामजातौ ।

\marginnote{\scriptsize ३.२.९९}\noindent उपसर्गे च संज्ञायाम् ।

\marginnote{\scriptsize ३.२.१००}\noindent अनौ कर्मणि ।

\marginnote{\scriptsize ३.२.१०१}\noindent अन्येष्वपि दृश्यते ।

\marginnote{\scriptsize ३.२.१०२}\noindent निष्ठा ।

\marginnote{\scriptsize ३.२.१०३}\noindent सुयजोर्ङ्वनिप् ।

\marginnote{\scriptsize ३.२.१०४}\noindent जीर्यतेरतृन् ।

\marginnote{\scriptsize ३.२.१०५}\noindent छन्दसि लिट् ।

\marginnote{\scriptsize ३.२.१०६}\noindent लिटः कानज्वा ।

\marginnote{\scriptsize ३.२.१०७}\noindent क्वसुश्च ।

\marginnote{\scriptsize ३.२.१०८}\noindent भाषायां सदवसश्रुवः ।

\marginnote{\scriptsize ३.२.१०९}\noindent उपेयिवाननाश्वाननूचानश्च ।

\marginnote{\scriptsize ३.२.११०}\noindent लुङ् ।

\marginnote{\scriptsize ३.२.१११}\noindent अनद्यतने लङ् ।

\marginnote{\scriptsize ३.२.११२}\noindent अभिज्ञावचने लृट् ।

\marginnote{\scriptsize ३.२.११३}\noindent न यदि ।

\marginnote{\scriptsize ३.२.११४}\noindent विभाषा साकाङ्क्षे ।

\marginnote{\scriptsize ३.२.११५}\noindent परोक्षे लिट् ।

\marginnote{\scriptsize ३.२.११६}\noindent हशश्वतोर्लङ् च ।

\marginnote{\scriptsize ३.२.११७}\noindent प्रश्ने चासन्नकाले ।

\marginnote{\scriptsize ३.२.११८}\noindent लट् स्मे ।

\marginnote{\scriptsize ३.२.११९}\noindent अपरोक्षे च ।

\marginnote{\scriptsize ३.२.१२०}\noindent ननौ पृष्टप्रतिवचने ।

\marginnote{\scriptsize ३.२.१२१}\noindent नन्वोर्विभाषा ।

\marginnote{\scriptsize ३.२.१२२}\noindent पुरि लुङ् चास्मे ।

\marginnote{\scriptsize ३.२.१२३}\noindent वर्तमाने लट् ।

\marginnote{\scriptsize ३.२.१२४}\noindent लटः शतृशनचावप्रथमासमानाधिकरणे ।

\marginnote{\scriptsize ३.२.१२५}\noindent सम्बोधने च ।

\marginnote{\scriptsize ३.२.१२६}\noindent लक्षणहेत्वोः क्रियायाः ।

\marginnote{\scriptsize ३.२.१२७}\noindent तौ सत् ।

\marginnote{\scriptsize ३.२.१२८}\noindent पूङ्यजोः शानन् ।

\marginnote{\scriptsize ३.२.१२९}\noindent ताच्छील्यवयोवचनशक्तिषु चानश् ।

\marginnote{\scriptsize ३.२.१३०}\noindent इङ्धार्योः शत्रकृच्छ्रिणि ।

\marginnote{\scriptsize ३.२.१३१}\noindent द्विषोऽमित्रे ।

\marginnote{\scriptsize ३.२.१३२}\noindent सुञो यज्ञसंयोगे ।

\marginnote{\scriptsize ३.२.१३३}\noindent अर्हः पूजायाम् ।

\marginnote{\scriptsize ३.२.१३४}\noindent आक्वेस्तच्छीलतद्धर्मतत्साधुकारिषु ।

\marginnote{\scriptsize ३.२.१३५}\noindent तृन् ।

\marginnote{\scriptsize ३.२.१३६}\noindent अलंकृञ्निराकृञ्प्रजनोत्पचोत्पतोन्मद- रुच्यपत्रपवृतुवृधुसहचर इष्णुच् ।

\marginnote{\scriptsize ३.२.१३७}\noindent णेश्छन्दसि ।

\marginnote{\scriptsize ३.२.१३८}\noindent भुवश्च ।

\marginnote{\scriptsize ३.२.१३९}\noindent ग्लाजिस्थश्च क्स्नुः ।

\marginnote{\scriptsize ३.२.१४०}\noindent त्रसिगृधिधृषिक्षिपेः क्नुः ।

\marginnote{\scriptsize ३.२.१४१}\noindent शमित्यष्टाभ्यो घिनुण् ।

\marginnote{\scriptsize ३.२.१४२}\noindent सम्पृचानुरुधाङ्यमाङ्यसपरिसृसंसृज- परिदेविसंज्वरपरिक्षिपपरिरटपरिवदपरिदहपरिमुह- दुषद्विषद्रुहदुहयुजाक्रीडविविचत्यजरज- भजातिचरापचरामुषाभ्याहनश्च ।

\marginnote{\scriptsize ३.२.१४३}\noindent वौ कषलसकत्थस्रम्भः ।

\marginnote{\scriptsize ३.२.१४४}\noindent अपे च लषः ।

\marginnote{\scriptsize ३.२.१४५}\noindent प्रे लपसृद्रुमथवदवसः ।

\marginnote{\scriptsize ३.२.१४६}\noindent निन्दहिंसक्लिशखादविनाशपरिक्षिपपरिरटपरिवादि- व्याभाषासूयो वुञ् ।

\marginnote{\scriptsize ३.२.१४७}\noindent देविक्रुशोश्चोपसर्गे ।

\marginnote{\scriptsize ३.२.१४८}\noindent चलनशब्दार्थादकर्मकाद्युच् ।

\marginnote{\scriptsize ३.२.१४९}\noindent अनुदात्तेतश्च हलादेः ।

\marginnote{\scriptsize ३.२.१५०}\noindent जुचङ्क्रम्यदन्द्रम्यसृगृधिज्वलशुचलषपतपदः ।

\marginnote{\scriptsize ३.२.१५१}\noindent क्रुधमण्डार्थेभ्यश्च ।

\marginnote{\scriptsize ३.२.१५२}\noindent न यः ।

\marginnote{\scriptsize ३.२.१५३}\noindent सूददीपदीक्षश्च ।

\marginnote{\scriptsize ३.२.१५४}\noindent लषपतपदस्थाभूवृषहनकमगमशॄभ्य उकञ् ।

\marginnote{\scriptsize ३.२.१५५}\noindent जल्पभिक्षकुट्टलुण्टवृङः षाकन् ।

\marginnote{\scriptsize ३.२.१५६}\noindent प्रजोरिनिः ।

\marginnote{\scriptsize ३.२.१५७}\noindent जिदृक्षिविश्रीण्वमाव्यथाभ्यमपरिभूप्रसूभ्यश्च ।

\marginnote{\scriptsize ३.२.१५८}\noindent स्पृहिगृहिपतिदयिनिद्रातन्द्राश्रद्धाभ्य आलुच् ।

\marginnote{\scriptsize ३.२.१५९}\noindent दाधेट्सिशदसदो रुः ।

\marginnote{\scriptsize ३.२.१६०}\noindent सृघस्यदः क्मरच् ।

\marginnote{\scriptsize ३.२.१६१}\noindent भञ्जभासमिदो घुरच् ।

\marginnote{\scriptsize ३.२.१६२}\noindent विदिभिदिच्छिदेः कुरच् ।

\marginnote{\scriptsize ३.२.१६३}\noindent इण्नश्जिसर्त्तिभ्यः क्वरप् ।

\marginnote{\scriptsize ३.२.१६४}\noindent गत्वरश्च ।

\marginnote{\scriptsize ३.२.१६५}\noindent जागुरूकः ।

\marginnote{\scriptsize ३.२.१६६}\noindent यजजपदशां यङः ।

\marginnote{\scriptsize ३.२.१६७}\noindent नमिकम्पिस्म्यजसकमहिंसदीपो रः ।

\marginnote{\scriptsize ३.२.१६८}\noindent सनाशंसभिक्ष उः ।

\marginnote{\scriptsize ३.२.१६९}\noindent विन्दुरिच्छुः ।

\marginnote{\scriptsize ३.२.१७०}\noindent क्याच्छन्दसि ।

\marginnote{\scriptsize ३.२.१७१}\noindent आदृगमहनजनः किकिनौ लिट् च ।

\marginnote{\scriptsize ३.२.१७२}\noindent स्वपितृषोर्नजिङ् ।

\marginnote{\scriptsize ३.२.१७३}\noindent शॄवन्द्योरारुः ।

\marginnote{\scriptsize ३.२.१७४}\noindent भियः क्रुक्लुकनौ ।

\marginnote{\scriptsize ३.२.१७५}\noindent स्थेशभासपिसकसो वरच् ।

\marginnote{\scriptsize ३.२.१७६}\noindent यश्च यङः ।

\marginnote{\scriptsize ३.२.१७७}\noindent भ्राजभासधुर्विद्युतोर्जिपॄजुग्रावस्तुवः क्विप् ।

\marginnote{\scriptsize ३.२.१७८}\noindent अन्येभ्योऽपि दृश्यते ।

\marginnote{\scriptsize ३.२.१७९}\noindent भुवः संज्ञाऽन्तरयोः ।

\marginnote{\scriptsize ३.२.१८०}\noindent विप्रसम्भ्यो ड्वसंज्ञायाम् ।

\marginnote{\scriptsize ३.२.१८१}\noindent धः कर्मणि ष्ट्रन् ।

\marginnote{\scriptsize ३.२.१८२}\noindent दाम्नीशसयुयुजस्तुतुदसिसिचमिहपतदशनहः करणे ।

\marginnote{\scriptsize ३.२.१८३}\noindent हलसूकरयोः पुवः ।

\marginnote{\scriptsize ३.२.१८४}\noindent अर्तिलूधूसूखनसहचर इत्रः ।

\marginnote{\scriptsize ३.२.१८५}\noindent पुवः संज्ञायाम् ।

\marginnote{\scriptsize ३.२.१८६}\noindent कर्तरि चर्षिदेवतयोः ।

\marginnote{\scriptsize ३.२.१८७}\noindent ञीतः क्तः ।

\marginnote{\scriptsize ३.२.१८८}\noindent मतिबुद्धिपूजार्थेभ्यश्च ।

\marginnote{\scriptsize ३.३.१}\noindent उणादयो बहुलम् ।

\marginnote{\scriptsize ३.३.२}\noindent भूतेऽपि दृश्यन्ते ।

\marginnote{\scriptsize ३.३.३}\noindent भविष्यति गम्यादयः ।

\marginnote{\scriptsize ३.३.४}\noindent यावत्पुरानिपातयोर्लट् ।

\marginnote{\scriptsize ३.३.५}\noindent विभाषा कदाकर्ह्योः ।

\marginnote{\scriptsize ३.३.६}\noindent किंवृत्ते लिप्सायाम् ।

\marginnote{\scriptsize ३.३.७}\noindent लिप्स्यमानसिद्धौ च ।

\marginnote{\scriptsize ३.३.८}\noindent लोडर्थलक्षणे च ।

\marginnote{\scriptsize ३.३.९}\noindent लिङ् चोर्ध्वमौहूर्तिके ।

\marginnote{\scriptsize ३.३.१०}\noindent तुमुन्ण्वुलौ क्रियायां क्रियार्थायाम् ।

\marginnote{\scriptsize ३.३.११}\noindent भाववचनाश्च ।

\marginnote{\scriptsize ३.३.१२}\noindent अण् कर्मणि च ।

\marginnote{\scriptsize ३.३.१३}\noindent लृट् शेषे च ।

\marginnote{\scriptsize ३.३.१४}\noindent लृटः सद् वा ।

\marginnote{\scriptsize ३.३.१५}\noindent अनद्यतने लुट् ।

\marginnote{\scriptsize ३.३.१६}\noindent पदरुजविशस्पृशो घञ् ।

\marginnote{\scriptsize ३.३.१७}\noindent सृ स्थिरे ।

\marginnote{\scriptsize ३.३.१८}\noindent भावे ।

\marginnote{\scriptsize ३.३.१९}\noindent अकर्तरि च कारके संज्ञायाम् ।

\marginnote{\scriptsize ३.३.२०}\noindent परिमणाख्यायां सर्वेभ्यः ।

\marginnote{\scriptsize ३.३.२१}\noindent इङश्च ।

\marginnote{\scriptsize ३.३.२२}\noindent उपसर्गे रुवः ।

\marginnote{\scriptsize ३.३.२३}\noindent समि युद्रुदुवः ।

\marginnote{\scriptsize ३.३.२४}\noindent श्रिणीभुवोऽनुपसर्गे ।

\marginnote{\scriptsize ३.३.२५}\noindent वौ क्षुश्रुवः ।

\marginnote{\scriptsize ३.३.२६}\noindent अवोदोर्नियः ।

\marginnote{\scriptsize ३.३.२७}\noindent प्रे द्रुस्तुस्रुवः ।

\marginnote{\scriptsize ३.३.२८}\noindent निरभ्योः पूल्वोः ।

\marginnote{\scriptsize ३.३.२९}\noindent उन्न्योर्ग्रः ।

\marginnote{\scriptsize ३.३.३०}\noindent कॄ धान्ये ।

\marginnote{\scriptsize ३.३.३१}\noindent यज्ञे समि स्तुवः ।

\marginnote{\scriptsize ३.३.३२}\noindent प्रे स्त्रोऽयज्ञे ।

\marginnote{\scriptsize ३.३.३३}\noindent प्रथने वावशब्दे ।

\marginnote{\scriptsize ३.३.३४}\noindent छन्दोनाम्नि च ।

\marginnote{\scriptsize ३.३.३५}\noindent उदि ग्रहः ।

\marginnote{\scriptsize ३.३.३६}\noindent समि मुष्टौ ।

\marginnote{\scriptsize ३.३.३७}\noindent परिन्योर्नीणोर्द्यूताभ्रेषयोः ।

\marginnote{\scriptsize ३.३.३८}\noindent परावनुपात्यय इणः ।

\marginnote{\scriptsize ३.३.३९}\noindent व्युपयोः शेतेः पर्याये ।

\marginnote{\scriptsize ३.३.४०}\noindent हस्तादाने चेरस्तेये ।

\marginnote{\scriptsize ३.३.४१}\noindent निवासचितिशरीरोपसमाधानेष्वादेश्च कः ।

\marginnote{\scriptsize ३.३.४२}\noindent संघे चानौत्तराधर्ये ।

\marginnote{\scriptsize ३.३.४३}\noindent कर्मव्यतिहारे णच् स्त्रियाम् ।

\marginnote{\scriptsize ३.३.४४}\noindent अभिविधौ भाव इनुण् ।

\marginnote{\scriptsize ३.३.४५}\noindent आक्रोशेऽवन्योर्ग्रहः ।

\marginnote{\scriptsize ३.३.४६}\noindent प्रे लिप्सायाम् ।

\marginnote{\scriptsize ३.३.४७}\noindent परौ यज्ञे ।

\marginnote{\scriptsize ३.३.४८}\noindent नौ वृ धान्ये ।

\marginnote{\scriptsize ३.३.४९}\noindent उदि श्रयतियौतिपूद्रुवः ।

\marginnote{\scriptsize ३.३.५०}\noindent विभाषाऽऽङि रुप्लुवोः ।

\marginnote{\scriptsize ३.३.५१}\noindent अवे ग्रहो वर्षप्रतिबन्धे ।

\marginnote{\scriptsize ३.३.५२}\noindent प्रे वणिजाम् ।

\marginnote{\scriptsize ३.३.५३}\noindent रश्मौ च ।

\marginnote{\scriptsize ३.३.५४}\noindent वृणोतेराच्छादने ।

\marginnote{\scriptsize ३.३.५५}\noindent परौ भुवोऽवज्ञाने ।

\marginnote{\scriptsize ३.३.५६}\noindent एरच् ।

\marginnote{\scriptsize ३.३.५७}\noindent ऋदोरप् ।

\marginnote{\scriptsize ३.३.५८}\noindent ग्रहवृदृनिश्चिगमश्च ।

\marginnote{\scriptsize ३.३.५९}\noindent उपसर्गेऽदः ।

\marginnote{\scriptsize ३.३.६०}\noindent नौ ण च ।

\marginnote{\scriptsize ३.३.६१}\noindent व्यधजपोरनुपसर्गे ।

\marginnote{\scriptsize ३.३.६२}\noindent स्वनहसोर्वा ।

\marginnote{\scriptsize ३.३.६३}\noindent यमः समुपनिविषु ।

\marginnote{\scriptsize ३.३.६४}\noindent नौ गदनदपठस्वनः ।

\marginnote{\scriptsize ३.३.६५}\noindent क्वणो वीणायां च ।

\marginnote{\scriptsize ३.३.६६}\noindent नित्यं पणः परिमाणे ।

\marginnote{\scriptsize ३.३.६७}\noindent मदोऽनुपसर्गे ।

\marginnote{\scriptsize ३.३.६८}\noindent प्रमदसम्मदौ हर्षे ।

\marginnote{\scriptsize ३.३.६९}\noindent समुदोरजः पशुषु ।

\marginnote{\scriptsize ३.३.७०}\noindent अक्षेषु ग्लहः ।

\marginnote{\scriptsize ३.३.७१}\noindent प्रजने सर्तेः ।

\marginnote{\scriptsize ३.३.७२}\noindent ह्वः सम्प्रसारणं च न्यभ्युपविषु ।

\marginnote{\scriptsize ३.३.७३}\noindent आङि युद्धे ।

\marginnote{\scriptsize ३.३.७४}\noindent निपानमाहावः ।

\marginnote{\scriptsize ३.३.७५}\noindent भावेऽनुपसर्गस्य ।

\marginnote{\scriptsize ३.३.७६}\noindent हनश्च वधः ।

\marginnote{\scriptsize ३.३.७७}\noindent मूर्तौ घनः ।

\marginnote{\scriptsize ३.३.७८}\noindent अन्तर्घनो देशे ।

\marginnote{\scriptsize ३.३.७९}\noindent अगारैकदेशे प्रघणः प्रघाणश्च ।

\marginnote{\scriptsize ३.३.८०}\noindent उद्घनोऽत्याधानम् ।

\marginnote{\scriptsize ३.३.८१}\noindent अपघनोऽङ्गम् ।

\marginnote{\scriptsize ३.३.८२}\noindent करणेऽयोविद्रुषु ।

\marginnote{\scriptsize ३.३.८३}\noindent स्तम्बे क च ।

\marginnote{\scriptsize ३.३.८४}\noindent परौ घः ।

\marginnote{\scriptsize ३.३.८५}\noindent उपघ्न आश्रये ।

\marginnote{\scriptsize ३.३.८६}\noindent संघोद्घौ गणप्रशंसयोः ।

\marginnote{\scriptsize ३.३.८७}\noindent निघो निमितम् ।

\marginnote{\scriptsize ३.३.८८}\noindent ड्वितः क्त्रिः ।

\marginnote{\scriptsize ३.३.८९}\noindent ट्वितोऽथुच् ।

\marginnote{\scriptsize ३.३.९०}\noindent यजयाचयतविच्छप्रच्छरक्षो नङ् ।

\marginnote{\scriptsize ३.३.९१}\noindent स्वपो नन् ।

\marginnote{\scriptsize ३.३.९२}\noindent उपसर्गे घोः किः ।

\marginnote{\scriptsize ३.३.९३}\noindent कर्मण्यधिकरणे च ।

\marginnote{\scriptsize ३.३.९४}\noindent स्त्रियां क्तिन् ।

\marginnote{\scriptsize ३.३.९५}\noindent स्थागापापचां भावे ।

\marginnote{\scriptsize ३.३.९६}\noindent मन्त्रे वृषेषपचमनविदभूवीरा उदात्तः ।

\marginnote{\scriptsize ३.३.९७}\noindent ऊतियूतिजूतिसातिहेतिकीर्तयश्च ।

\marginnote{\scriptsize ३.३.९८}\noindent व्रजयजोर्भावे क्यप् ।

\marginnote{\scriptsize ३.३.९९}\noindent संज्ञायां समजनिषदनिपतमनविदषुञ्शीङ्भृञिणः ।

\marginnote{\scriptsize ३.३.१००}\noindent कृञः श च ।

\marginnote{\scriptsize ३.३.१०१}\noindent इच्छा ।

\marginnote{\scriptsize ३.३.१०२}\noindent अ प्रत्ययात् ।

\marginnote{\scriptsize ३.३.१०३}\noindent गुरोश्च हलः ।

\marginnote{\scriptsize ३.३.१०४}\noindent षिद्भिदादिभ्योऽङ् ।

\marginnote{\scriptsize ३.३.१०५}\noindent चिन्तिपूजिकथिकुम्बिचर्चश्च ।

\marginnote{\scriptsize ३.३.१०६}\noindent आतश्चोपसर्गे ।

\marginnote{\scriptsize ३.३.१०७}\noindent ण्यासश्रन्थो युच् ।

\marginnote{\scriptsize ३.३.१०८}\noindent रोगाख्यायां ण्वुल् बहुलम् ।

\marginnote{\scriptsize ३.३.१०९}\noindent संज्ञायाम् ।

\marginnote{\scriptsize ३.३.११०}\noindent विभाषाऽऽख्यानपरिप्रश्नयोरिञ् च ।

\marginnote{\scriptsize ३.३.१११}\noindent पर्यायार्हर्णोत्पत्तिषु ण्वुच् ।

\marginnote{\scriptsize ३.३.११२}\noindent आक्रोशे नञ्यनिः ।

\marginnote{\scriptsize ३.३.११३}\noindent कृत्यल्युटो बहुलम् ।

\marginnote{\scriptsize ३.३.११४}\noindent नपुंसके भावे क्तः ।

\marginnote{\scriptsize ३.३.११५}\noindent ल्युट् च ।

\marginnote{\scriptsize ३.३.११६}\noindent कर्मणि च येन संस्पर्शात् कर्तुः शरीरसुखम् ।

\marginnote{\scriptsize ३.३.११७}\noindent करणाधिकरणयोश्च ।

\marginnote{\scriptsize ३.३.११८}\noindent पुंसि संज्ञायां घः प्रायेण ।

\marginnote{\scriptsize ३.३.११९}\noindent गोचरसंचरवहव्रजव्यजापणनिगमाश्च ।

\marginnote{\scriptsize ३.३.१२०}\noindent अवे तॄस्त्रोर्घञ् ।

\marginnote{\scriptsize ३.३.१२१}\noindent हलश्च ।

\marginnote{\scriptsize ३.३.१२२}\noindent अध्यायन्यायोद्यावसंहाराधारावयाश्च ।

\marginnote{\scriptsize ३.३.१२३}\noindent उदङ्कोऽनुदके ।

\marginnote{\scriptsize ३.३.१२४}\noindent जालमानायः ।

\marginnote{\scriptsize ३.३.१२५}\noindent खनो घ च ।

\marginnote{\scriptsize ३.३.१२६}\noindent ईषद्दुःसुषु कृच्छ्राकृच्छ्रार्थेषु खल् ।

\marginnote{\scriptsize ३.३.१२७}\noindent कर्तृकर्मणोश्च भूकृञोः ।

\marginnote{\scriptsize ३.३.१२८}\noindent आतो युच् ।

\marginnote{\scriptsize ३.३.१२९}\noindent छन्दसि गत्यर्थेभ्यः ।

\marginnote{\scriptsize ३.३.१३०}\noindent अन्येभ्योऽपि दृश्यते ।

\marginnote{\scriptsize ३.३.१३१}\noindent वर्तमानसामीप्ये वर्तमानवद्वा ।

\marginnote{\scriptsize ३.३.१३२}\noindent आशंसायां भूतवच्च ।

\marginnote{\scriptsize ३.३.१३३}\noindent क्षिप्रवचने लृट् ।

\marginnote{\scriptsize ३.३.१३४}\noindent आशंसावचने लिङ् ।

\marginnote{\scriptsize ३.३.१३५}\noindent नानद्यतनवत् क्रियाप्रबन्धसामीप्ययोः ।

\marginnote{\scriptsize ३.३.१३६}\noindent भविष्यति मर्यादावचनेऽवरस्मिन् ।

\marginnote{\scriptsize ३.३.१३७}\noindent कालविभागे चानहोरात्राणाम् ।

\marginnote{\scriptsize ३.३.१३८}\noindent परस्मिन् विभाषा ।

\marginnote{\scriptsize ३.३.१३९}\noindent लिङ्निमित्ते लृङ् क्रियाऽतिपत्तौ ।

\marginnote{\scriptsize ३.३.१४०}\noindent भूते च ।

\marginnote{\scriptsize ३.३.१४१}\noindent वोताप्योः ।

\marginnote{\scriptsize ३.३.१४२}\noindent गर्हायां लडपिजात्वोः ।

\marginnote{\scriptsize ३.३.१४३}\noindent विभाषा कथमि लिङ् च ।

\marginnote{\scriptsize ३.३.१४४}\noindent किंवृत्ते लिङ्लृटौ ।

\marginnote{\scriptsize ३.३.१४५}\noindent अनवकॢप्त्यमर्षयोरकिंवृत्ते अपि ।

\marginnote{\scriptsize ३.३.१४६}\noindent किंकिलास्त्यर्थेषु लृट् ।

\marginnote{\scriptsize ३.३.१४७}\noindent जातुयदोर्लिङ् ।

\marginnote{\scriptsize ३.३.१४८}\noindent यच्चयत्रयोः ।

\marginnote{\scriptsize ३.३.१४९}\noindent गर्हायां च ।

\marginnote{\scriptsize ३.३.१५०}\noindent चित्रीकरणे च ।

\marginnote{\scriptsize ३.३.१५१}\noindent शेषे लृडयदौ ।

\marginnote{\scriptsize ३.३.१५२}\noindent उताप्योः समर्थयोर्लिङ् ।

\marginnote{\scriptsize ३.३.१५३}\noindent कामप्रवेदनेऽकच्चिति ।

\marginnote{\scriptsize ३.३.१५४}\noindent सम्भवानेऽलमिति चेत् सिद्धाप्रयोगे ।

\marginnote{\scriptsize ३.३.१५५}\noindent विभाषा धातौ सम्भावनवचनेऽयदि ।

\marginnote{\scriptsize ३.३.१५६}\noindent हेतुहेतुमतोर्लिङ् ।

\marginnote{\scriptsize ३.३.१५७}\noindent इच्छार्थेषु लिङ्लोटौ ।

\marginnote{\scriptsize ३.३.१५८}\noindent समानकर्तृकेषु तुमुन् ।

\marginnote{\scriptsize ३.३.१५९}\noindent लिङ् च ।

\marginnote{\scriptsize ३.३.१६०}\noindent इच्छार्थेभ्यो विभाषा वर्तमाने ।

\marginnote{\scriptsize ३.३.१६१}\noindent विधिनिमन्त्रणामन्त्रण अधीष्टसम्प्रश्नप्रार्थनेषु लिङ् ।

\marginnote{\scriptsize ३.३.१६२}\noindent लोट् च ।

\marginnote{\scriptsize ३.३.१६३}\noindent प्रैषातिसर्गप्राप्तकालेषु कृत्याश्च ।

\marginnote{\scriptsize ३.३.१६४}\noindent लिङ् चोर्ध्वमौहूर्तिके ।

\marginnote{\scriptsize ३.३.१६५}\noindent स्मे लोट् ।

\marginnote{\scriptsize ३.३.१६६}\noindent अधीष्टे च ।

\marginnote{\scriptsize ३.३.१६७}\noindent कालसमयवेलासु तुमुन् ।

\marginnote{\scriptsize ३.३.१६८}\noindent लिङ् यदि ।

\marginnote{\scriptsize ३.३.१६९}\noindent अर्हे कृत्यतृचश्च ।

\marginnote{\scriptsize ३.३.१७०}\noindent आवश्यकाधमर्ण्ययोर्णिनिः ।

\marginnote{\scriptsize ३.३.१७१}\noindent कृत्याश्च ।

\marginnote{\scriptsize ३.३.१७२}\noindent शकि लिङ् च ।

\marginnote{\scriptsize ३.३.१७३}\noindent आशिषि लिङ्लोटौ ।

\marginnote{\scriptsize ३.३.१७४}\noindent क्तिच्क्तौ च संज्ञायाम् ।

\marginnote{\scriptsize ३.३.१७५}\noindent माङि लुङ् ।

\marginnote{\scriptsize ३.३.१७६}\noindent स्मोत्तरे लङ् च ।

\marginnote{\scriptsize ३.४.१}\noindent धातुसम्बन्धे प्रत्ययाः ।

\marginnote{\scriptsize ३.४.२}\noindent क्रियासमभिहारे लोट्, लोटो हिस्वौ, वा च तध्वमोः ।

\marginnote{\scriptsize ३.४.३}\noindent समुच्चयेऽन्यतरस्याम् ।

\marginnote{\scriptsize ३.४.४}\noindent यथाविध्यनुप्रयोगः पूर्वस्मिन् ।

\marginnote{\scriptsize ३.४.५}\noindent समुच्चये सामान्यवचनस्य ।

\marginnote{\scriptsize ३.४.६}\noindent छन्दसि लुङ्लङ्लिटः ।

\marginnote{\scriptsize ३.४.७}\noindent लिङर्थे लेट् ।

\marginnote{\scriptsize ३.४.८}\noindent उपसंवादाशङ्कयोश्च ।

\marginnote{\scriptsize ३.४.९}\noindent तुमर्थे सेसेनसेअसेन्क्सेकसेनध्यैअध्यैन्कध्यैकध्यैन्- शध्यैशध्यैन्तवैतवेङ्तवेनः ।

\marginnote{\scriptsize ३.४.१०}\noindent प्रयै रोहिष्यै अव्यथिष्यै ।

\marginnote{\scriptsize ३.४.११}\noindent दृशे विख्ये च ।

\marginnote{\scriptsize ३.४.१२}\noindent शकि णमुल्कमुलौ ।

\marginnote{\scriptsize ३.४.१३}\noindent ईश्वरे तोसुन्कसुनौ ।

\marginnote{\scriptsize ३.४.१४}\noindent कृत्यार्थे तवैकेन्केन्यत्वनः ।

\marginnote{\scriptsize ३.४.१५}\noindent अवचक्षे च ।

\marginnote{\scriptsize ३.४.१६}\noindent भावलक्षणे स्थेङ्कृञ्वदिचरिहुतमिजनिभ्यस्तोसुन् ।

\marginnote{\scriptsize ३.४.१७}\noindent सृपितृदोः कसुन् ।

\marginnote{\scriptsize ३.४.१८}\noindent अलङ्खल्वोः प्रतिषेधयोः प्राचां क्त्वा ।

\marginnote{\scriptsize ३.४.१९}\noindent उदीचां माङो व्यतीहारे ।

\marginnote{\scriptsize ३.४.२०}\noindent परावरयोगे च ।

\marginnote{\scriptsize ३.४.२१}\noindent समानकर्तृकयोः पूर्वकाले ।

\marginnote{\scriptsize ३.४.२२}\noindent आभीक्ष्ण्ये णमुल् च ।

\marginnote{\scriptsize ३.४.२३}\noindent न यद्यनाकाङ्क्षे ।

\marginnote{\scriptsize ३.४.२४}\noindent विभाषाऽग्रेप्रथमपूर्वेषु ।

\marginnote{\scriptsize ३.४.२५}\noindent कर्मण्याक्रोशे कृञः खमुञ् ।

\marginnote{\scriptsize ३.४.२६}\noindent स्वादुमि णमुल् ।

\marginnote{\scriptsize ३.४.२७}\noindent अन्यथैवंकथमित्थंसु सिद्धाप्रयोगश्चेत् ।

\marginnote{\scriptsize ३.४.२८}\noindent यथातथयोरसूयाप्रतिवचने ।

\marginnote{\scriptsize ३.४.२९}\noindent कर्मणि दृशिविदोः साकल्ये ।

\marginnote{\scriptsize ३.४.३०}\noindent यावति विन्दजीवोः ।

\marginnote{\scriptsize ३.४.३१}\noindent चर्मोदरयोः पूरेः ।

\marginnote{\scriptsize ३.४.३२}\noindent वर्षप्रमाण ऊलोपश्चास्यान्यतरास्यम् ।

\marginnote{\scriptsize ३.४.३३}\noindent चेले क्नोपेः ।

\marginnote{\scriptsize ३.४.३४}\noindent निमूलसमूलयोः कषः ।

\marginnote{\scriptsize ३.४.३५}\noindent शुष्कचूर्णरूक्षेषु पिषः ।

\marginnote{\scriptsize ३.४.३६}\noindent समूलाकृतजीवेषु हन्कृङ्ग्रहः ।

\marginnote{\scriptsize ३.४.३७}\noindent करणे हनः ।

\marginnote{\scriptsize ३.४.३८}\noindent स्नेहने पिषः ।

\marginnote{\scriptsize ३.४.३९}\noindent हस्ते वर्त्तिग्रहोः ।

\marginnote{\scriptsize ३.४.४०}\noindent स्वे पुषः ।

\marginnote{\scriptsize ३.४.४१}\noindent अधिकरणे बन्धः ।

\marginnote{\scriptsize ३.४.४२}\noindent संज्ञायाम् ।

\marginnote{\scriptsize ३.४.४३}\noindent कर्त्रोर्जीवपुरुषयोर्नशिवहोः ।

\marginnote{\scriptsize ३.४.४४}\noindent ऊर्ध्वे शुषिपूरोः ।

\marginnote{\scriptsize ३.४.४५}\noindent उपमाने कर्मणि च ।

\marginnote{\scriptsize ३.४.४६}\noindent कषादिषु यथाविध्यनुप्रयोगः ।

\marginnote{\scriptsize ३.४.४७}\noindent उपदंशस्तृतीयायाम् ।

\marginnote{\scriptsize ३.४.४८}\noindent हिंसार्थानां च समानकर्मकाणाम् ।

\marginnote{\scriptsize ३.४.४९}\noindent सप्तम्यां चोपपीडरुधकर्षः ।

\marginnote{\scriptsize ३.४.५०}\noindent समासत्तौ ।

\marginnote{\scriptsize ३.४.५१}\noindent प्रमाणे च ।

\marginnote{\scriptsize ३.४.५२}\noindent अपादाने परीप्सायाम् ।

\marginnote{\scriptsize ३.४.५३}\noindent द्वितीयायां च ।

\marginnote{\scriptsize ३.४.५४}\noindent स्वाङ्गेऽध्रुवे ।

\marginnote{\scriptsize ३.४.५५}\noindent परिक्लिश्यमाने च ।

\marginnote{\scriptsize ३.४.५६}\noindent विशिपतिपदिस्कन्दां व्याप्यमानासेव्यमानयोः ।

\marginnote{\scriptsize ३.४.५७}\noindent अस्यतितृषोः क्रियाऽन्तरे कालेषु ।

\marginnote{\scriptsize ३.४.५८}\noindent नाम्न्यादिशिग्रहोः ।

\marginnote{\scriptsize ३.४.५९}\noindent अव्ययेऽयथाभिप्रेताख्याने कृञः क्त्वाणमुलौ ।

\marginnote{\scriptsize ३.४.६०}\noindent तिर्यच्यपवर्गे ।

\marginnote{\scriptsize ३.४.६१}\noindent स्वाङ्गे तस्प्रत्यये कृभ्वोः ।

\marginnote{\scriptsize ३.४.६२}\noindent नाधाऽर्थप्रत्यये च्व्यर्थे ।

\marginnote{\scriptsize ३.४.६३}\noindent तूष्णीमि भुवः ।

\marginnote{\scriptsize ३.४.६४}\noindent अन्वच्यानुलोम्ये । ३.४.६५ शकधृषज्ञाग्लाघटरभलभक्रमसहार्हास्त्यर्थेषु तुमुन् ।

\marginnote{\scriptsize ३.४.६६}\noindent पर्याप्तिवचनेष्वलमर्थेषु ।

\marginnote{\scriptsize ३.४.६७}\noindent कर्तरि कृत् ।

\marginnote{\scriptsize ३.४.६८}\noindent भव्यगेयप्रवचनीयोपस्थानीयजन्याप्लाव्यापात्या वा ।

\marginnote{\scriptsize ३.४.६९}\noindent लः कर्मणि च भावे चाकर्मकेभ्यः। ।

\marginnote{\scriptsize ३.४.७०}\noindent तयोरेव कृत्यक्तखलर्थाः ।

\marginnote{\scriptsize ३.४.७१}\noindent अदिकर्मणि क्तः कर्तरि च । ३.४.७२ गत्यर्थाकर्मकश्लिषशीङ्स्थाऽऽसवसजनरुहजीर्यतिभ्यश्च ।

\marginnote{\scriptsize ३.४.७३}\noindent दाशगोघ्नौ सम्प्रदाने ।

\marginnote{\scriptsize ३.४.७४}\noindent भीमादयोऽपादाने ।

\marginnote{\scriptsize ३.४.७५}\noindent ताभ्यामन्यत्रोणादयः ।

\marginnote{\scriptsize ३.४.७६}\noindent क्तोऽधिकरणे च ध्रौव्यगतिप्रत्यवसानार्थेभ्यः ।

\marginnote{\scriptsize ३.४.७७}\noindent लस्य ।

\marginnote{\scriptsize ३.४.७८}\noindent तिप्तस्झिसिप्थस्थमिब्वस्मस्- तातांझथासाथांध्वमिड्वहिमहिङ् ।

\marginnote{\scriptsize ३.४.७९}\noindent टित आत्मनेपदानां टेरे ।

\marginnote{\scriptsize ३.४.८०}\noindent थासस्से ।

\marginnote{\scriptsize ३.४.८१}\noindent लिटस्तझयोरेशिरेच् ।

\marginnote{\scriptsize ३.४.८२}\noindent परस्मैपदानां णलतुसुस्थलथुसणल्वमाः ।

\marginnote{\scriptsize ३.४.८३}\noindent विदो लटो वा ।

\marginnote{\scriptsize ३.४.८४}\noindent ब्रुवः पञ्चानामादित आहो ब्रुवः ।

\marginnote{\scriptsize ३.४.८५}\noindent लोटो लङ्वत् ।

\marginnote{\scriptsize ३.४.८६}\noindent एरुः ।

\marginnote{\scriptsize ३.४.८७}\noindent सेर्ह्यपिच्च ।

\marginnote{\scriptsize ३.४.८८}\noindent वा छन्दसि ।

\marginnote{\scriptsize ३.४.८९}\noindent मेर्निः ।

\marginnote{\scriptsize ३.४.९०}\noindent आमेतः ।

\marginnote{\scriptsize ३.४.९१}\noindent सवाभ्यां वामौ ।

\marginnote{\scriptsize ३.४.९२}\noindent आडुत्तमस्य पिच्च ।

\marginnote{\scriptsize ३.४.९३}\noindent एत ऐ ।

\marginnote{\scriptsize ३.४.९४}\noindent लेटोऽडाटौ ।

\marginnote{\scriptsize ३.४.९५}\noindent आत ऐ ।

\marginnote{\scriptsize ३.४.९६}\noindent वैतोऽन्यत्र ।

\marginnote{\scriptsize ३.४.९७}\noindent इतश्च लोपः परस्मैपदेषु ।

\marginnote{\scriptsize ३.४.९८}\noindent स उत्तमस्य ।

\marginnote{\scriptsize ३.४.९९}\noindent नित्यं ङितः ।

\marginnote{\scriptsize ३.४.१००}\noindent इतश्च ।

\marginnote{\scriptsize ३.४.१०१}\noindent तस्थस्थमिपां तांतंतामः ।

\marginnote{\scriptsize ३.४.१०२}\noindent लिङस्सीयुट् ।

\marginnote{\scriptsize ३.४.१०३}\noindent यासुट् परस्मैपदेषूदात्तो ङिच्च ।

\marginnote{\scriptsize ३.४.१०४}\noindent किदाशिषि ।

\marginnote{\scriptsize ३.४.१०५}\noindent झस्य रन् ।

\marginnote{\scriptsize ३.४.१०६}\noindent इटोऽत् ।

\marginnote{\scriptsize ३.४.१०७}\noindent सुट् तिथोः ।

\marginnote{\scriptsize ३.४.१०८}\noindent झेर्जुस् ।

\marginnote{\scriptsize ३.४.१०९}\noindent सिजभ्यस्तविदिभ्यः च ।

\marginnote{\scriptsize ३.४.११०}\noindent आतः ।

\marginnote{\scriptsize ३.४.१११}\noindent लङः शाकटायनस्यैव ।

\marginnote{\scriptsize ३.४.११२}\noindent द्विषश्च ।

\marginnote{\scriptsize ३.४.११३}\noindent तिङ्शित्सार्वधातुकम् ।

\marginnote{\scriptsize ३.४.११४}\noindent आर्द्धधातुकं शेषः ।

\marginnote{\scriptsize ३.४.११५}\noindent लिट् च ।

\marginnote{\scriptsize ३.४.११६}\noindent लिङाशिषि ।

\marginnote{\scriptsize ३.४.११७}\noindent छन्दस्युभयथा ।

\section*{॥ अध्याय ४ ॥}

\marginnote{\scriptsize ४.१.१}\noindent ङ्याप्प्रातिपदिकात् ।

\marginnote{\scriptsize ४.१.२}\noindent स्वौजसमौट्छष्टाभ्याम्भिस्ङेभ्याम्भ्यस्ङसिभ्याम्भ्यस्- ङसोसाम्ङ्योस्सुप् ।

\marginnote{\scriptsize ४.१.३}\noindent स्त्रियाम् ।

\marginnote{\scriptsize ४.१.४}\noindent अजाद्यतष्टाप् ।

\marginnote{\scriptsize ४.१.५}\noindent ऋन्नेभ्यो ङीप् ।

\marginnote{\scriptsize ४.१.६}\noindent उगितश्च ।

\marginnote{\scriptsize ४.१.७}\noindent वनो र च ।

\marginnote{\scriptsize ४.१.८}\noindent पादोऽन्यतरस्याम् ।

\marginnote{\scriptsize ४.१.९}\noindent टाबृचि ।

\marginnote{\scriptsize ४.१.१०}\noindent न षट्स्वस्रादिभ्यः ।

\marginnote{\scriptsize ४.१.११}\noindent मनः ।

\marginnote{\scriptsize ४.१.१२}\noindent अनो बहुव्रीहेः ।

\marginnote{\scriptsize ४.१.१३}\noindent डाबुभाभ्यामन्यतरस्याम् ।

\marginnote{\scriptsize ४.१.१४}\noindent अनुपसर्जनात् ।

\marginnote{\scriptsize ४.१.१५}\noindent टिड्ढाणञ्द्वयसज्दघ्नञ्मात्रच्तयप्ठक्ठङ्कङ्क्वरपः ।

\marginnote{\scriptsize ४.१.१६}\noindent यञश्च ।

\marginnote{\scriptsize ४.१.१७}\noindent प्राचां ष्फ तद्धितः ।

\marginnote{\scriptsize ४.१.१८}\noindent सर्वत्र लोहितादिकतान्तेभ्यः ।

\marginnote{\scriptsize ४.१.१९}\noindent कौरव्यमाण्डूकाभ्यां च ।

\marginnote{\scriptsize ४.१.२०}\noindent वयसि प्रथमे ।

\marginnote{\scriptsize ४.१.२१}\noindent द्विगोः ।

\marginnote{\scriptsize ४.१.२२}\noindent अपरिमाणबिस्ताचितकम्बल्येभ्यो न तद्धितलुकि ।

\marginnote{\scriptsize ४.१.२३}\noindent काण्डान्तात् क्षेत्रे ।

\marginnote{\scriptsize ४.१.२४}\noindent पुरुषात् प्रमाणेऽन्यतरस्याम् ।

\marginnote{\scriptsize ४.१.२५}\noindent बहुव्रीहेरूधसो ङीष्।

\marginnote{\scriptsize ४.१.२६}\noindent संख्याऽव्ययादेर्ङीप् ।

\marginnote{\scriptsize ४.१.२७}\noindent दामहायनान्ताच्च ।

\marginnote{\scriptsize ४.१.२८}\noindent अन उपधालोपिनोन्यतरस्याम् ।

\marginnote{\scriptsize ४.१.२९}\noindent नित्यं संज्ञाछन्दसोः ।

\marginnote{\scriptsize ४.१.३०}\noindent केवलमामकभागधेयपापापरसमानार्यकृत- सुमङ्गलभेषजाच्च ।

\marginnote{\scriptsize ४.१.३१}\noindent रात्रेश्चाजसौ ।

\marginnote{\scriptsize ४.१.३२}\noindent अन्तर्वत्पतिवतोर्नुक् ।

\marginnote{\scriptsize ४.१.३३}\noindent पत्युर्नो यज्ञसंयोगे ।

\marginnote{\scriptsize ४.१.३४}\noindent विभाषा सपूर्वस्य ।

\marginnote{\scriptsize ४.१.३५}\noindent नित्यं सपत्न्यादिषु ।

\marginnote{\scriptsize ४.१.३६}\noindent पूतक्रतोरै च ।

\marginnote{\scriptsize ४.१.३७}\noindent वृषाकप्यग्निकुसितकुसीदानामुदात्तः ।

\marginnote{\scriptsize ४.१.३८}\noindent मनोरौ वा ।

\marginnote{\scriptsize ४.१.३९}\noindent वर्णादनुदात्तात्तोपधात्तो नः ।

\marginnote{\scriptsize ४.१.४०}\noindent अन्यतो ङीष्।

\marginnote{\scriptsize ४.१.४१}\noindent षिद्गौरादिभ्यश्च ।

\marginnote{\scriptsize ४.१.४२}\noindent जानपदकुण्डगोणस्थलभाजनागकालनीलकुशकामुक- कबराद्वृत्त्यमत्रावपनाकृत्रिमाश्राणास्थौल्य- वर्णानाच्छादनायोविकारमैथुनेच्छाकेशवेशेषु ।

\marginnote{\scriptsize ४.१.४३}\noindent शोणात् प्राचाम् ।

\marginnote{\scriptsize ४.१.४४}\noindent वोतो गुणवचनात् ।

\marginnote{\scriptsize ४.१.४५}\noindent बह्वादिभ्यश्च ।

\marginnote{\scriptsize ४.१.४६}\noindent नित्यं छन्दसि ।

\marginnote{\scriptsize ४.१.४७}\noindent भुवश्च ।

\marginnote{\scriptsize ४.१.४८}\noindent पुंयोगादाख्यायाम् ।

\marginnote{\scriptsize ४.१.४९}\noindent इन्द्रवरुणभवशर्वरुद्रमृडहिमारण्ययवयवन- मातुलाचार्याणामानुक् ।

\marginnote{\scriptsize ४.१.५०}\noindent क्रीतात् करणपूर्वात् ।

\marginnote{\scriptsize ४.१.५१}\noindent क्तादल्पाख्यायाम् ।

\marginnote{\scriptsize ४.१.५२}\noindent बहुव्रीहेश्चान्तोदात्तात् ।

\marginnote{\scriptsize ४.१.५३}\noindent अस्वाङ्गपूर्वपदाद्वा ।

\marginnote{\scriptsize ४.१.५४}\noindent स्वाङ्गाच्चोपसर्जनादसंयोगोपधात् ।

\marginnote{\scriptsize ४.१.५५}\noindent नासिकोदरौष्ठजङ्घादन्तकर्णश‍ृङ्गाच्च ।

\marginnote{\scriptsize ४.१.५६}\noindent न क्रोडादिबह्वचः ।

\marginnote{\scriptsize ४.१.५७}\noindent सहनञ्विद्यमानपूर्वाच्च ।

\marginnote{\scriptsize ४.१.५८}\noindent नखमुखात् संज्ञायाम् ।

\marginnote{\scriptsize ४.१.५९}\noindent दीर्घजिह्वी च च्छन्दसि ।

\marginnote{\scriptsize ४.१.६०}\noindent दिक्पूर्वपदान्ङीप् ।

\marginnote{\scriptsize ४.१.६१}\noindent वाहः ।

\marginnote{\scriptsize ४.१.६२}\noindent सख्यशिश्वीति भाषायाम् ।

\marginnote{\scriptsize ४.१.६३}\noindent जातेरस्त्रीविषयादयोपधात् ।

\marginnote{\scriptsize ४.१.६४}\noindent पाककर्णपर्णपुष्पफलमूलबालोत्तरपदाच्च ।

\marginnote{\scriptsize ४.१.६५}\noindent इतो मनुष्यजातेः ।

\marginnote{\scriptsize ४.१.६६}\noindent ऊङुतः ।

\marginnote{\scriptsize ४.१.६७}\noindent बाह्वन्तात् संज्ञायाम् ।

\marginnote{\scriptsize ४.१.६८}\noindent पङ्गोश्च ।

\marginnote{\scriptsize ४.१.६९}\noindent ऊरूत्तरपदादौपम्ये ।

\marginnote{\scriptsize ४.१.७०}\noindent संहितशफलक्षणवामादेश्च ।

\marginnote{\scriptsize ४.१.७१}\noindent कद्रुकमण्डल्वोश्छन्दसि ।

\marginnote{\scriptsize ४.१.७२}\noindent संज्ञायाम् ।

\marginnote{\scriptsize ४.१.७३}\noindent शार्ङ्गरवाद्यञो ङीन् ।

\marginnote{\scriptsize ४.१.७४}\noindent यङश्चाप् ।

\marginnote{\scriptsize ४.१.७५}\noindent आवट्याच्च ।

\marginnote{\scriptsize ४.१.७६}\noindent तद्धिताः ।

\marginnote{\scriptsize ४.१.७७}\noindent यूनस्तिः ।

\marginnote{\scriptsize ४.१.७८}\noindent अणिञोरनार्षयोर्गुरूपोत्तमयोः ष्यङ् गोत्रे ।

\marginnote{\scriptsize ४.१.७९}\noindent गोत्रावयवात् ।

\marginnote{\scriptsize ४.१.८०}\noindent क्रौड्यादिभ्यश्च ।

\marginnote{\scriptsize ४.१.८१}\noindent दैवयज्ञिशौचिवृक्षिसात्यमुग्रि- काण्ठेविद्धिभ्योऽन्यतरस्याम् ।

\marginnote{\scriptsize ४.१.८२}\noindent समर्थानां प्रथमाद्वा ।

\marginnote{\scriptsize ४.१.८३}\noindent प्राग्दीव्यतोऽण् ।

\marginnote{\scriptsize ४.१.८४}\noindent अश्वपत्यादिभ्यश्च ।

\marginnote{\scriptsize ४.१.८५}\noindent दित्यदित्यादित्यपत्युत्तरपदाण्ण्यः ।

\marginnote{\scriptsize ४.१.८६}\noindent उत्सादिभ्योऽञ् ।

\marginnote{\scriptsize ४.१.८७}\noindent स्त्रीपुंसाभ्यां नञ्स्नञौ भवनात् ।

\marginnote{\scriptsize ४.१.८८}\noindent द्विगोर्लुगनपत्ये ।

\marginnote{\scriptsize ४.१.८९}\noindent गोत्रेऽलुगचि ।

\marginnote{\scriptsize ४.१.९०}\noindent यूनि लुक् ।

\marginnote{\scriptsize ४.१.९१}\noindent फक्फिञोरन्यतरस्याम् ।

\marginnote{\scriptsize ४.१.९२}\noindent तस्यापत्यम् ।

\marginnote{\scriptsize ४.१.९३}\noindent एको गोत्रे ।

\marginnote{\scriptsize ४.१.९४}\noindent गोत्राद्यून्यस्त्रियाम् ।

\marginnote{\scriptsize ४.१.९५}\noindent अत इञ् ।

\marginnote{\scriptsize ४.१.९६}\noindent बाह्वादिभ्यश्च ।

\marginnote{\scriptsize ४.१.९७}\noindent सुधातुरकङ् च ।

\marginnote{\scriptsize ४.१.९८}\noindent गोत्रे कुञ्जादिभ्यश्च्फञ् ।

\marginnote{\scriptsize ४.१.९९}\noindent नडादिभ्यः फक् ।

\marginnote{\scriptsize ४.१.१००}\noindent हरितादिभ्योऽञः ।

\marginnote{\scriptsize ४.१.१०१}\noindent यञिञोश्च ।

\marginnote{\scriptsize ४.१.१०२}\noindent शरद्वच्छुनकदर्भाद्भृगुवत्साग्रायणेषु ।

\marginnote{\scriptsize ४.१.१०३}\noindent द्रोणपर्वतजीवन्तादन्यतरयाम् ।

\marginnote{\scriptsize ४.१.१०४}\noindent अनृष्यानन्तर्ये बिदादिभ्योऽञ् ।

\marginnote{\scriptsize ४.१.१०५}\noindent गर्गादिभ्यो यञ् ।

\marginnote{\scriptsize ४.१.१०६}\noindent मधुबभ्र्वोर्ब्राह्मणकौशिकयोः ।

\marginnote{\scriptsize ४.१.१०७}\noindent कपिबोधादाङ्गिरसे ।

\marginnote{\scriptsize ४.१.१०८}\noindent वतण्डाच्च ।

\marginnote{\scriptsize ४.१.१०९}\noindent लुक् स्त्रियाम् ।

\marginnote{\scriptsize ४.१.११०}\noindent अश्वादिभ्यः फञ् ।

\marginnote{\scriptsize ४.१.१११}\noindent भर्गात् त्रैगर्ते ।

\marginnote{\scriptsize ४.१.११२}\noindent शिवादिभ्योऽण् ।

\marginnote{\scriptsize ४.१.११३}\noindent अवृद्धाभ्यो नदीमानुषीभ्यस्तन्नामिकाभ्यः ।

\marginnote{\scriptsize ४.१.११४}\noindent ऋष्यन्धकवृष्णिकुरुभ्यश्च ।

\marginnote{\scriptsize ४.१.११५}\noindent मातुरुत् संख्यासम्भद्रपूर्वायाः ।

\marginnote{\scriptsize ४.१.११६}\noindent कन्यायाः कनीन च ।

\marginnote{\scriptsize ४.१.११७}\noindent विकर्णशुङ्गच्छगलाद्वत्सभरद्वाजात्रिषु ।

\marginnote{\scriptsize ४.१.११८}\noindent पीलाया वा ।

\marginnote{\scriptsize ४.१.११९}\noindent ढक् च मण्डूकात् ।

\marginnote{\scriptsize ४.१.१२०}\noindent स्त्रीभ्यो ढक् ।

\marginnote{\scriptsize ४.१.१२१}\noindent द्व्यचः ।

\marginnote{\scriptsize ४.१.१२२}\noindent इतश्चानिञः ।

\marginnote{\scriptsize ४.१.१२३}\noindent शुभ्रादिभ्यश्च ।

\marginnote{\scriptsize ४.१.१२४}\noindent विकर्णकुषीतकात् काश्यपे ।

\marginnote{\scriptsize ४.१.१२५}\noindent भ्रुवो वुक् च ।

\marginnote{\scriptsize ४.१.१२६}\noindent कल्याण्यादीनामिनङ् ।

\marginnote{\scriptsize ४.१.१२७}\noindent कुलटाया वा ।

\marginnote{\scriptsize ४.१.१२८}\noindent चटकाया ऐरक् ।

\marginnote{\scriptsize ४.१.१२९}\noindent गोधाया ढ्रक् ।

\marginnote{\scriptsize ४.१.१३०}\noindent आरगुदीचाम् ।

\marginnote{\scriptsize ४.१.१३१}\noindent क्षुद्राभ्यो वा ।

\marginnote{\scriptsize ४.१.१३२}\noindent पितृष्वसुश्छण् ।

\marginnote{\scriptsize ४.१.१३३}\noindent ढकि लोपः ।

\marginnote{\scriptsize ४.१.१३४}\noindent मातृष्वसुश्च ।

\marginnote{\scriptsize ४.१.१३५}\noindent चतुष्पाद्भ्यो ढञ् ।

\marginnote{\scriptsize ४.१.१३६}\noindent गृष्ट्यादिभ्यश्च ।

\marginnote{\scriptsize ४.१.१३७}\noindent राजश्वशुराद्यत् ।

\marginnote{\scriptsize ४.१.१३८}\noindent क्षत्राद्घः ।

\marginnote{\scriptsize ४.१.१३९}\noindent कुलात् खः ।

\marginnote{\scriptsize ४.१.१४०}\noindent अपूर्वपदादन्यतरस्यां यड्ढकञौ ।

\marginnote{\scriptsize ४.१.१४१}\noindent महाकुलादङ्खञौ ।

\marginnote{\scriptsize ४.१.१४२}\noindent दुष्कुलाड्ढक् ।

\marginnote{\scriptsize ४.१.१४३}\noindent स्वसुश्छः ।

\marginnote{\scriptsize ४.१.१४४}\noindent भ्रातुर्व्यच्च ।

\marginnote{\scriptsize ४.१.१४५}\noindent व्यन् सपत्ने ।

\marginnote{\scriptsize ४.१.१४६}\noindent रेवत्यादिभ्यष्ठक् ।

\marginnote{\scriptsize ४.१.१४७}\noindent गोत्रस्त्रियाः कुत्सने ण च ।

\marginnote{\scriptsize ४.१.१४८}\noindent वृद्धाट्ठक् सौवीरेषु बहुलम् ।

\marginnote{\scriptsize ४.१.१४९}\noindent फेश्छ च ।

\marginnote{\scriptsize ४.१.१५०}\noindent फाण्टाहृतिमिमताभ्यां णफिञौ ।

\marginnote{\scriptsize ४.१.१५१}\noindent कुर्वादिभ्यो ण्यः ।

\marginnote{\scriptsize ४.१.१५२}\noindent सेनान्तलक्षणकारिभ्यश्च ।

\marginnote{\scriptsize ४.१.१५३}\noindent उदीचामिञ् ।

\marginnote{\scriptsize ४.१.१५४}\noindent तिकादिभ्यः फिञ् ।

\marginnote{\scriptsize ४.१.१५५}\noindent कौसल्यकार्मार्याभ्यां च ।

\marginnote{\scriptsize ४.१.१५६}\noindent अणो द्व्यचः ।

\marginnote{\scriptsize ४.१.१५७}\noindent उदीचां वृद्धादगोत्रात् ।

\marginnote{\scriptsize ४.१.१५८}\noindent वाकिनादीनां कुक् च ।

\marginnote{\scriptsize ४.१.१५९}\noindent पुत्रान्तादन्यतरस्याम् ।

\marginnote{\scriptsize ४.१.१६०}\noindent प्राचामवृद्धात् फिन् बहुलम् ।

\marginnote{\scriptsize ४.१.१६१}\noindent मनोर्जातावञ्यतौ षुक् च ।

\marginnote{\scriptsize ४.१.१६२}\noindent अपत्यं पौत्रप्रभृति गोत्रम् ।

\marginnote{\scriptsize ४.१.१६३}\noindent जीवति तु वंश्ये युवा ।

\marginnote{\scriptsize ४.१.१६४}\noindent भ्रातरि च ज्यायसि ।

\marginnote{\scriptsize ४.१.१६५}\noindent वाऽन्यस्मिन् सपिण्डे स्थविरतरे जीवति ।

\marginnote{\scriptsize ४.१.१६६}\noindent वृद्धस्य च पूजायाम् ।

\marginnote{\scriptsize ४.१.१६७}\noindent यूनश्च कुत्सायाम् ।

\marginnote{\scriptsize ४.१.१६८}\noindent जनपदशब्दात् क्षत्रियादञ् ।

\marginnote{\scriptsize ४.१.१६९}\noindent साल्वेयगान्धारिभ्यां च ।

\marginnote{\scriptsize ४.१.१७०}\noindent द्व्यञ्मगधकलिङ्गसूरमसादण् ।

\marginnote{\scriptsize ४.१.१७१}\noindent वृद्धेत्कोसलाजादाञ्ञ्यङ् ।

\marginnote{\scriptsize ४.१.१७२}\noindent कुरुणादिभ्यो ण्यः ।

\marginnote{\scriptsize ४.१.१७३}\noindent साल्वावयवप्रत्यग्रथकलकूटाश्मकादिञ् ।

\marginnote{\scriptsize ४.१.१७४}\noindent ते तद्राजाः ।

\marginnote{\scriptsize ४.१.१७५}\noindent कम्बोजाल्लुक् ।

\marginnote{\scriptsize ४.१.१७६}\noindent स्त्रियामवन्तिकुन्तिकुरुभ्यश्च ।

\marginnote{\scriptsize ४.१.१७७}\noindent अतश्च ।

\marginnote{\scriptsize ४.१.१७८}\noindent न प्राच्यभर्गादियौधेयादिभ्यः ।

\marginnote{\scriptsize ४.२.१}\noindent तेन रक्तं रागात् ।

\marginnote{\scriptsize ४.२.२}\noindent लाक्षारोचना(शकलकर्दमा)ट्ठक् ।

\marginnote{\scriptsize ४.२.३}\noindent नक्षत्रेण युक्तः कालः ।

\marginnote{\scriptsize ४.२.४}\noindent लुबविशेषे ।

\marginnote{\scriptsize ४.२.५}\noindent संज्ञायां श्रवणाश्वत्थाभ्याम् ।

\marginnote{\scriptsize ४.२.६}\noindent द्वंद्वाच्छः ।

\marginnote{\scriptsize ४.२.७}\noindent दृष्टं साम ।

\marginnote{\scriptsize ४.२.८}\noindent कलेर्ढक् ।

\marginnote{\scriptsize ४.२.९}\noindent वामदेवाड्ड्यड्ड्यौ ।

\marginnote{\scriptsize ४.२.१०}\noindent परिवृतो रथः ।

\marginnote{\scriptsize ४.२.११}\noindent पाण्डुकम्बलादिनिः ।

\marginnote{\scriptsize ४.२.१२}\noindent द्वैपवैयाघ्रादञ् ।

\marginnote{\scriptsize ४.२.१३}\noindent कौमारापूर्ववचने ।

\marginnote{\scriptsize ४.२.१४}\noindent तत्रोद्धृतममत्रेभ्यः ।

\marginnote{\scriptsize ४.२.१५}\noindent स्थण्डिलाच्छयितरि व्रते ।

\marginnote{\scriptsize ४.२.१६}\noindent संस्कृतं भक्षाः ।

\marginnote{\scriptsize ४.२.१७}\noindent शूलोखाद्यत् ।

\marginnote{\scriptsize ४.२.१८}\noindent दध्नष्ठक् ।

\marginnote{\scriptsize ४.२.१९}\noindent उदश्वितोऽन्यतरस्याम् ।

\marginnote{\scriptsize ४.२.२०}\noindent क्षीराड्ढञ् ।

\marginnote{\scriptsize ४.२.२१}\noindent साऽस्मिन् पौर्णमासीति (संज्ञायाम्) ।

\marginnote{\scriptsize ४.२.२२}\noindent आग्रहायण्यश्वत्थाट्ठक् ।

\marginnote{\scriptsize ४.२.२३}\noindent विभाषा फाल्गुनीश्रवणाकार्त्तिकीचैत्रीभ्यः ।

\marginnote{\scriptsize ४.२.२४}\noindent साऽस्य देवता ।

\marginnote{\scriptsize ४.२.२५}\noindent कस्येत् ।

\marginnote{\scriptsize ४.२.२६}\noindent शुक्राद्घन् ।

\marginnote{\scriptsize ४.२.२७}\noindent अपोनप्त्रपान्नप्तृभ्यां घः ।

\marginnote{\scriptsize ४.२.२८}\noindent छ च ।

\marginnote{\scriptsize ४.२.२९}\noindent महेन्द्राद्घाणौ च ।

\marginnote{\scriptsize ४.२.३०}\noindent सोमाट्ट्यण् ।

\marginnote{\scriptsize ४.२.३१}\noindent वाय्वृतुपित्रुषसो यत् ।

\marginnote{\scriptsize ४.२.३२}\noindent द्यावापृथिवीशुनासीरमरुत्वदग्नीषोमवास्तोष्पति- गृहमेधाच्छ च ।

\marginnote{\scriptsize ४.२.३३}\noindent अग्नेर्ढक् ।

\marginnote{\scriptsize ४.२.३४}\noindent कालेभ्यो भववत् ।

\marginnote{\scriptsize ४.२.३५}\noindent महाराजप्रोष्ठपदाट्ठञ् ।

\marginnote{\scriptsize ४.२.३६}\noindent पितृव्यमातुलमातामहपितामहाः ।

\marginnote{\scriptsize ४.२.३७}\noindent तस्य समूहः ।

\marginnote{\scriptsize ४.२.३८}\noindent भिक्षाऽऽदिभ्योऽण् ।

\marginnote{\scriptsize ४.२.३९}\noindent गोत्रोक्षोष्ट्रोरभ्रराजराजन्यराजपुत्रवत्स- मनुष्याजाद्वुञ् ।

\marginnote{\scriptsize ४.२.४०}\noindent केदाराद्यञ् च ।

\marginnote{\scriptsize ४.२.४१}\noindent ठञ् कवचिनश्च ।

\marginnote{\scriptsize ४.२.४२}\noindent ब्राह्मणमाणववाडवाद्यन् ।

\marginnote{\scriptsize ४.२.४३}\noindent ग्रामजनबन्धुसहायेभ्यः तल् ।

\marginnote{\scriptsize ४.२.४४}\noindent अनुदात्तादेरञ् ।

\marginnote{\scriptsize ४.२.४५}\noindent खण्डिकादिभ्यश्च ।

\marginnote{\scriptsize ४.२.४६}\noindent चरणेभ्यो धर्मवत् ।

\marginnote{\scriptsize ४.२.४७}\noindent अचित्तहस्तिधेनोष्ठक् ।

\marginnote{\scriptsize ४.२.४८}\noindent केशाश्वाभ्यां यञ्छावन्यतरस्याम् ।

\marginnote{\scriptsize ४.२.४९}\noindent पाशादिभ्यो यः ।

\marginnote{\scriptsize ४.२.५०}\noindent खलगोरथात् ।

\marginnote{\scriptsize ४.२.५१}\noindent इनित्रकट्यचश्च ।

\marginnote{\scriptsize ४.२.५२}\noindent विषयो देशे ।

\marginnote{\scriptsize ४.२.५३}\noindent राजन्यादिभ्यो वुञ् ।

\marginnote{\scriptsize ४.२.५४}\noindent भौरिक्याद्यैषुकार्यादिभ्यो विधल्भक्तलौ ।

\marginnote{\scriptsize ४.२.५५}\noindent सोऽस्यादिरिति च्छन्दसः प्रगाथेषु ।

\marginnote{\scriptsize ४.२.५६}\noindent संग्रामे प्रयोजनयोद्धृभ्यः ।

\marginnote{\scriptsize ४.२.५७}\noindent तदस्यां प्रहरणमिति क्रीडायाम् णः ।

\marginnote{\scriptsize ४.२.५८}\noindent घञः साऽस्यां क्रियेति ञः ।

\marginnote{\scriptsize ४.२.५९}\noindent तदधीते तद्वेद ।

\marginnote{\scriptsize ४.२.६०}\noindent क्रतूक्थादिसूत्रान्ताट्ठक् ।

\marginnote{\scriptsize ४.२.६१}\noindent क्रमादिभ्यो वुन् ।

\marginnote{\scriptsize ४.२.६२}\noindent अनुब्राह्मणादिनिः ।

\marginnote{\scriptsize ४.२.६३}\noindent वसन्तादिभ्यष्ठक् ।

\marginnote{\scriptsize ४.२.६४}\noindent प्रोक्ताल्लुक् ।

\marginnote{\scriptsize ४.२.६५}\noindent सूत्राच्च कोपधात् ।

\marginnote{\scriptsize ४.२.६६}\noindent छन्दोब्राह्मणानि च तद्विषयाणि ।

\marginnote{\scriptsize ४.२.६७}\noindent तदस्मिन्नस्तीति देशे तन्नाम्नि ।

\marginnote{\scriptsize ४.२.६८}\noindent तेन निर्वृत्तम् ।

\marginnote{\scriptsize ४.२.६९}\noindent तस्य निवासः ।

\marginnote{\scriptsize ४.२.७०}\noindent अदूरभवश्च ।

\marginnote{\scriptsize ४.२.७१}\noindent ओरञ् ।

\marginnote{\scriptsize ४.२.७२}\noindent मतोश्च बह्वजङ्गात् ।

\marginnote{\scriptsize ४.२.७३}\noindent बह्वचः कूपेषु ।

\marginnote{\scriptsize ४.२.७४}\noindent उदक् च विपाशः ।

\marginnote{\scriptsize ४.२.७५}\noindent संकलादिभ्यश्च ।

\marginnote{\scriptsize ४.२.७६}\noindent स्त्रीषु सौवीरसाल्वप्राक्षु ।

\marginnote{\scriptsize ४.२.७७}\noindent सुवास्त्वादिभ्योऽण् ।

\marginnote{\scriptsize ४.२.७८}\noindent रोणी ।

\marginnote{\scriptsize ४.२.७९}\noindent कोपधाच्च ।

\marginnote{\scriptsize ४.२.८०}\noindent वुञ्छङ्कठजिलशेनिरढञ्ण्ययफक्फिञिञ्ञ्य- कक्ठकोऽरीहणकृशाश्वर्श्यकुमुदकाशतृणप्रेक्षाऽश्मसखि- संकाशबलपक्षकर्णसुतंगमप्रगदिन्वराहकुमुदादिभ्यः ।

\marginnote{\scriptsize ४.२.८१}\noindent जनपदे लुप् ।

\marginnote{\scriptsize ४.२.८२}\noindent वरणादिभ्यश्च ।

\marginnote{\scriptsize ४.२.८३}\noindent शर्कराया वा ।

\marginnote{\scriptsize ४.२.८४}\noindent ठक्छौ च ।

\marginnote{\scriptsize ४.२.८५}\noindent नद्यां मतुप् ।

\marginnote{\scriptsize ४.२.८६}\noindent मध्वादिभ्यश्च ।

\marginnote{\scriptsize ४.२.८७}\noindent कुमुदनडवेतसेभ्यो ड्मतुप् ।

\marginnote{\scriptsize ४.२.८८}\noindent नडशादाड्ड्वलच् ।

\marginnote{\scriptsize ४.२.८९}\noindent शिखाया वलच् ।

\marginnote{\scriptsize ४.२.९०}\noindent उत्करादिभ्यश्छः ।

\marginnote{\scriptsize ४.२.९१}\noindent नडादीनां कुक् च ।

\marginnote{\scriptsize ४.२.९२}\noindent शेषे ।

\marginnote{\scriptsize ४.२.९३}\noindent राष्ट्रावारपाराद्घखौ ।

\marginnote{\scriptsize ४.२.९४}\noindent ग्रामाद्यखञौ ।

\marginnote{\scriptsize ४.२.९५}\noindent कत्त्र्यादिभ्यो ढकञ् ।

\marginnote{\scriptsize ४.२.९६}\noindent कुलकुक्षिग्रीवाभ्यः श्वास्यलंकारेषु ।

\marginnote{\scriptsize ४.२.९७}\noindent नद्यादिभ्यो ढक् ।

\marginnote{\scriptsize ४.२.९८}\noindent दक्षिणापश्चात्पुरसस्त्यक् ।

\marginnote{\scriptsize ४.२.९९}\noindent कापिश्याः ष्फक् ।

\marginnote{\scriptsize ४.२.१००}\noindent रंकोरमनुष्येऽण् च ।

\marginnote{\scriptsize ४.२.१०१}\noindent द्युप्रागपागुदक्प्रतीचो यत् ।

\marginnote{\scriptsize ४.२.१०२}\noindent कन्थायाष्ठक् ।

\marginnote{\scriptsize ४.२.१०३}\noindent वर्णौ वुक् ।

\marginnote{\scriptsize ४.२.१०४}\noindent अव्ययात्त्यप् ।

\marginnote{\scriptsize ४.२.१०५}\noindent ऐषमोह्यःश्वसोऽन्यतरस्याम् ।

\marginnote{\scriptsize ४.२.१०६}\noindent तीररूप्योत्तरपदादञ्ञौ ।

\marginnote{\scriptsize ४.२.१०७}\noindent दिक्पूर्वपदादसंज्ञायां ञः ।

\marginnote{\scriptsize ४.२.१०८}\noindent मद्रेभ्योऽञ् ।

\marginnote{\scriptsize ४.२.१०९}\noindent उदीच्यग्रामाच्च बह्वचोऽन्तोदात्तात् ।

\marginnote{\scriptsize ४.२.११०}\noindent प्रस्थोत्तरपदपलद्यादिकोपधादण् ।

\marginnote{\scriptsize ४.२.१११}\noindent कण्वादिभ्यो गोत्रे ।

\marginnote{\scriptsize ४.२.११२}\noindent इञश्च ।

\marginnote{\scriptsize ४.२.११३}\noindent न द्व्यचः प्राच्यभरतेषु ।

\marginnote{\scriptsize ४.२.११४}\noindent वृद्धाच्छः ।

\marginnote{\scriptsize ४.२.११५}\noindent भवतष्ठक्छसौ ।

\marginnote{\scriptsize ४.२.११६}\noindent काश्यादिभ्यष्ठञ्ञिठौ ।

\marginnote{\scriptsize ४.२.११७}\noindent वाहीकग्रामेभ्यश्च ।

\marginnote{\scriptsize ४.२.११८}\noindent विभाषोशीनरेषु ।

\marginnote{\scriptsize ४.२.११९}\noindent ओर्देशे ठञ् ।

\marginnote{\scriptsize ४.२.१२०}\noindent वृद्धात् प्राचाम् ।

\marginnote{\scriptsize ४.२.१२१}\noindent धन्वयोपधाद्वुञ् ।

\marginnote{\scriptsize ४.२.१२२}\noindent प्रस्थपुरवहान्ताच्च ।

\marginnote{\scriptsize ४.२.१२३}\noindent रोपधेतोः प्राचाम् ।

\marginnote{\scriptsize ४.२.१२४}\noindent जनपदतदवध्योश्च ।

\marginnote{\scriptsize ४.२.१२५}\noindent अवृद्धादपि बहुवचनविषयात् ।

\marginnote{\scriptsize ४.२.१२६}\noindent कच्छाग्निवक्त्रगर्त्तोत्तरपदात् ।

\marginnote{\scriptsize ४.२.१२७}\noindent धूमादिभ्यश्च ।

\marginnote{\scriptsize ४.२.१२८}\noindent नगरात् कुत्सनप्रावीण्ययोः ।

\marginnote{\scriptsize ४.२.१२९}\noindent अरण्यान्मनुष्ये ।

\marginnote{\scriptsize ४.२.१३०}\noindent विभाषा कुरुयुगन्धराभ्याम् ।

\marginnote{\scriptsize ४.२.१३१}\noindent मद्रवृज्योः कन् ।

\marginnote{\scriptsize ४.२.१३२}\noindent कोपधादण् ।

\marginnote{\scriptsize ४.२.१३३}\noindent कच्छादिभ्यश्च ।

\marginnote{\scriptsize ४.२.१३४}\noindent मनुष्यतत्स्थयोर्वुञ् ।

\marginnote{\scriptsize ४.२.१३५}\noindent अपदातौ साल्वात् ।

\marginnote{\scriptsize ४.२.१३६}\noindent गोयवाग्वोश्च ।

\marginnote{\scriptsize ४.२.१३७}\noindent गर्तोत्तरपदाच्छः ।

\marginnote{\scriptsize ४.२.१३८}\noindent गहादिभ्यश्च ।

\marginnote{\scriptsize ४.२.१३९}\noindent प्राचां कटादेः ।

\marginnote{\scriptsize ४.२.१४०}\noindent राज्ञः क च ।

\marginnote{\scriptsize ४.२.१४१}\noindent वृद्धादकेकान्तखोपधात् ।

\marginnote{\scriptsize ४.२.१४२}\noindent कन्थापलदनगरग्रामह्रदोत्तरपदात् ।

\marginnote{\scriptsize ४.२.१४३}\noindent पर्वताच्च ।

\marginnote{\scriptsize ४.२.१४४}\noindent विभाषाऽमनुष्ये ।

\marginnote{\scriptsize ४.२.१४५}\noindent कृकणपर्णाद्भारद्वाजे ।

\marginnote{\scriptsize ४.३.१}\noindent युष्मदस्मदोरन्यतरस्यां खञ् च ।

\marginnote{\scriptsize ४.३.२}\noindent तस्मिन् नणि च युष्माकास्माकौ ।

\marginnote{\scriptsize ४.३.३}\noindent तवकममकावेकवचने ।

\marginnote{\scriptsize ४.३.४}\noindent अर्धाद्यत् ।

\marginnote{\scriptsize ४.३.५}\noindent परावराधमोत्तमपूर्वाच्च ।

\marginnote{\scriptsize ४.३.६}\noindent दिक्पूर्वपदाट्ठञ् च ।

\marginnote{\scriptsize ४.३.७}\noindent ग्रामजनपदैकदेशादञ्ठञौ ।

\marginnote{\scriptsize ४.३.८}\noindent मध्यान्मः ।

\marginnote{\scriptsize ४.३.९}\noindent अ साम्प्रतिके ।

\marginnote{\scriptsize ४.३.१०}\noindent द्वीपादनुसमुद्रं यञ् ।

\marginnote{\scriptsize ४.३.११}\noindent कालाट्ठञ् ।

\marginnote{\scriptsize ४.३.१२}\noindent श्राद्धे शरदः ।

\marginnote{\scriptsize ४.३.१३}\noindent विभाषा रोगातपयोः ।

\marginnote{\scriptsize ४.३.१४}\noindent निशाप्रदोषाभ्यां च ।

\marginnote{\scriptsize ४.३.१५}\noindent श्वसस्तुट् च ।

\marginnote{\scriptsize ४.३.१६}\noindent संधिवेलाऽऽद्यृतुनक्षत्रेभ्योऽण् ।

\marginnote{\scriptsize ४.३.१७}\noindent प्रावृष एण्यः ।

\marginnote{\scriptsize ४.३.१८}\noindent वर्षाभ्यष्ठक् ।

\marginnote{\scriptsize ४.३.१९}\noindent छन्दसि ठञ् ।

\marginnote{\scriptsize ४.३.२०}\noindent वसन्ताच्च ।

\marginnote{\scriptsize ४.३.२१}\noindent हेमन्ताच्च ।

\marginnote{\scriptsize ४.३.२२}\noindent सर्वत्राण् च तलोपश्च ।

\marginnote{\scriptsize ४.३.२३}\noindent सायंचिरम्प्राह्णेप्रगेऽव्ययेभ्यष्ट्युट्युलौ तुट् च ।

\marginnote{\scriptsize ४.३.२४}\noindent विभाषा पूर्वाह्णापराह्णाभ्याम् ।

\marginnote{\scriptsize ४.३.२५}\noindent तत्र जातः ।

\marginnote{\scriptsize ४.३.२६}\noindent प्रावृषष्ठप् ।

\marginnote{\scriptsize ४.३.२७}\noindent संज्ञायां शरदो वुञ् ।

\marginnote{\scriptsize ४.३.२८}\noindent पूर्वाह्णापराह्णार्द्रामूलप्रदोषावस्कराद्वुन् ।

\marginnote{\scriptsize ४.३.२९}\noindent पथः पन्थ च ।

\marginnote{\scriptsize ४.३.३०}\noindent अमावास्याया वा ।

\marginnote{\scriptsize ४.३.३१}\noindent अ च ।

\marginnote{\scriptsize ४.३.३२}\noindent सिन्ध्वपकराभ्यां कन् ।

\marginnote{\scriptsize ४.३.३३}\noindent अणञौ च ।

\marginnote{\scriptsize ४.३.३४}\noindent श्रविष्ठाफल्गुन्यनुराधास्वातितिष्यपुनर्वसुहस्त- विशाखाऽषाढाबहुलाल्लुक् ।

\marginnote{\scriptsize ४.३.३५}\noindent स्थानान्तगोशालखरशालाच्च ।

\marginnote{\scriptsize ४.३.३६}\noindent वत्सशालाऽभिजिदश्वयुक्छतभिषजो वा ।

\marginnote{\scriptsize ४.३.३७}\noindent नक्षत्रेभ्यो बहुलम् ।

\marginnote{\scriptsize ४.३.३८}\noindent कृतलब्धक्रीतकुशलाः ।

\marginnote{\scriptsize ४.३.३९}\noindent प्रायभवः ।

\marginnote{\scriptsize ४.३.४०}\noindent उपजानूपकर्णोपनीवेष्ठक् ।

\marginnote{\scriptsize ४.३.४१}\noindent संभूते ।

\marginnote{\scriptsize ४.३.४२}\noindent कोशाड्ढञ् ।

\marginnote{\scriptsize ४.३.४३}\noindent कालात् साधुपुष्प्यत्पच्यमानेषु ।

\marginnote{\scriptsize ४.३.४४}\noindent उप्ते च ।

\marginnote{\scriptsize ४.३.४५}\noindent आश्वयुज्या वुञ् ।

\marginnote{\scriptsize ४.३.४६}\noindent ग्रीष्मवसन्तादन्यतरस्याम् ।

\marginnote{\scriptsize ४.३.४७}\noindent देयमृणे ।

\marginnote{\scriptsize ४.३.४८}\noindent कलाप्यश्वत्थयवबुसाद्वुन् ।

\marginnote{\scriptsize ४.३.४९}\noindent ग्रीष्मावरसमाद्वुञ् ।

\marginnote{\scriptsize ४.३.५०}\noindent संवत्सराग्रहायणीभ्यां ठञ् च ।

\marginnote{\scriptsize ४.३.५१}\noindent व्याहरति मृगः ।

\marginnote{\scriptsize ४.३.५२}\noindent तदस्य सोढम् ।

\marginnote{\scriptsize ४.३.५३}\noindent तत्र भवः ।

\marginnote{\scriptsize ४.३.५४}\noindent दिगादिभ्यो यत् ।

\marginnote{\scriptsize ४.३.५५}\noindent शरीरावयवाच्च ।

\marginnote{\scriptsize ४.३.५६}\noindent दृतिकुक्षिकलशिवस्त्यस्त्यहेर्ढञ् ।

\marginnote{\scriptsize ४.३.५७}\noindent ग्रीवाभ्योऽण् च ।

\marginnote{\scriptsize ४.३.५८}\noindent गम्भीराञ्ञ्यः ।

\marginnote{\scriptsize ४.३.५९}\noindent अव्ययीभावाच्च ।

\marginnote{\scriptsize ४.३.६०}\noindent अन्तःपूर्वपदाट्ठञ् ।

\marginnote{\scriptsize ४.३.६१}\noindent ग्रामात् पर्यनुपूर्वात् ।

\marginnote{\scriptsize ४.३.६२}\noindent जिह्वामूलाङ्गुलेश्छः ।

\marginnote{\scriptsize ४.३.६३}\noindent वर्गान्ताच्च ।

\marginnote{\scriptsize ४.३.६४}\noindent अशब्दे यत्खावन्यतरस्याम् ।

\marginnote{\scriptsize ४.३.६५}\noindent कर्णललाटात् कनलंकारे ।

\marginnote{\scriptsize ४.३.६६}\noindent तस्य व्याख्यान इति च व्याख्यातव्यनाम्नः ।

\marginnote{\scriptsize ४.३.६७}\noindent बह्वचोऽन्तोदात्ताट्ठञ् ।

\marginnote{\scriptsize ४.३.६८}\noindent क्रतुयज्ञेभ्यश्च ।

\marginnote{\scriptsize ४.३.६९}\noindent अध्यायेष्वेवर्षेः ।

\marginnote{\scriptsize ४.३.७०}\noindent पौरोडाशपुरोडाशात् ष्ठन् ।

\marginnote{\scriptsize ४.३.७१}\noindent छन्दसो यदणौ ।

\marginnote{\scriptsize ४.३.७२}\noindent द्व्यजृद्ब्राह्मणर्क्प्रथमाध्वरपुरश्चरण- नामाख्याताट्ठक् ।

\marginnote{\scriptsize ४.३.७३}\noindent अणृगयनादिभ्यः ।

\marginnote{\scriptsize ४.३.७४}\noindent तत आगतः ।

\marginnote{\scriptsize ४.३.७५}\noindent ठगायस्थानेभ्यः ।

\marginnote{\scriptsize ४.३.७६}\noindent शुण्डिकादिभ्योऽण् ।

\marginnote{\scriptsize ४.३.७७}\noindent विद्यायोनिसंबन्धेभ्यो वुञ् ।

\marginnote{\scriptsize ४.३.७८}\noindent ऋतश्ठञ् ।

\marginnote{\scriptsize ४.३.७९}\noindent पितुर्यच्च ।

\marginnote{\scriptsize ४.३.८०}\noindent गोत्रादङ्कवत् ।

\marginnote{\scriptsize ४.३.८१}\noindent हेतुमनुष्येभ्योऽन्यतरस्यां रूप्यः ।

\marginnote{\scriptsize ४.३.८२}\noindent मयट् च ।

\marginnote{\scriptsize ४.३.८३}\noindent प्रभवति ।

\marginnote{\scriptsize ४.३.८४}\noindent विदूराञ्ञ्यः ।

\marginnote{\scriptsize ४.३.८५}\noindent तद्गच्छति पथिदूतयोः ।

\marginnote{\scriptsize ४.३.८६}\noindent अभिनिष्क्रामति द्वारम् ।

\marginnote{\scriptsize ४.३.८७}\noindent अधिकृत्य कृते ग्रन्थे ।

\marginnote{\scriptsize ४.३.८८}\noindent शिशुक्रन्दयमसभद्वंद्वेन्द्रजननादिभ्यश्छः ।

\marginnote{\scriptsize ४.३.८९}\noindent सोऽस्य निवासः ।

\marginnote{\scriptsize ४.३.९०}\noindent अभिजनश्च ।

\marginnote{\scriptsize ४.३.९१}\noindent आयुधजीविभ्यश्छः पर्वते ।

\marginnote{\scriptsize ४.३.९२}\noindent शण्डिकादिभ्यो ञ्यः ।

\marginnote{\scriptsize ४.३.९३}\noindent सिन्धुतक्षशिलाऽऽदिभ्योऽणञौ ।

\marginnote{\scriptsize ४.३.९४}\noindent तूदीशलातुरवर्मतीकूचवाराड्ढक्छण्ढञ्यकः ।

\marginnote{\scriptsize ४.३.९५}\noindent भक्तिः ।

\marginnote{\scriptsize ४.३.९६}\noindent अचित्ताददेशकालाट्ठक् ।

\marginnote{\scriptsize ४.३.९७}\noindent महाराजाट्ठञ् ।

\marginnote{\scriptsize ४.३.९८}\noindent वासुदेवार्जुनाभ्यां वुन् ।

\marginnote{\scriptsize ४.३.९९}\noindent गोत्रक्षत्रियाख्येभ्यो बहुलं वुञ् ।

\marginnote{\scriptsize ४.३.१००}\noindent जनपदिनां जनपदवत् सर्वं जनपदेन समानशब्दानां बहुवचने ।

\marginnote{\scriptsize ४.३.१०१}\noindent तेन प्रोक्तम् ।

\marginnote{\scriptsize ४.३.१०२}\noindent तित्तिरिवरतन्तुखण्डिकोखाच्छण् ।

\marginnote{\scriptsize ४.३.१०३}\noindent काश्यपकौशिकाभ्यामृषिभ्यां णिनिः ।

\marginnote{\scriptsize ४.३.१०४}\noindent कलापिवैशम्पायनान्तेवासिभ्यश्च ।

\marginnote{\scriptsize ४.३.१०५}\noindent पुराणप्रोक्तेषु ब्राह्मणकल्पेषु ।

\marginnote{\scriptsize ४.३.१०६}\noindent शौनकादिभ्यश्छन्दसि ।

\marginnote{\scriptsize ४.३.१०७}\noindent कठचरकाल्लुक् ।

\marginnote{\scriptsize ४.३.१०८}\noindent कलापिनोऽण् ।

\marginnote{\scriptsize ४.३.१०९}\noindent छगलिनो ढिनुक् ।

\marginnote{\scriptsize ४.३.११०}\noindent पाराशर्यशिलालिभ्यां भिक्षुनटसूत्रयोः ।

\marginnote{\scriptsize ४.३.१११}\noindent कर्मन्दकृशाश्वादिनिः ।

\marginnote{\scriptsize ४.३.११२}\noindent तेनैकदिक् ।

\marginnote{\scriptsize ४.३.११३}\noindent तसिश्च ।

\marginnote{\scriptsize ४.३.११४}\noindent उरसो यच्च ।

\marginnote{\scriptsize ४.३.११५}\noindent उपज्ञाते ।

\marginnote{\scriptsize ४.३.११६}\noindent कृते ग्रन्थे ।

\marginnote{\scriptsize ४.३.११७}\noindent संज्ञायाम् ।

\marginnote{\scriptsize ४.३.११८}\noindent कुलालादिभ्यो वुञ् ।

\marginnote{\scriptsize ४.३.११९}\noindent क्षुद्राभ्रमरवटरपादपादञ् ।

\marginnote{\scriptsize ४.३.१२०}\noindent तस्येदम् ।

\marginnote{\scriptsize ४.३.१२१}\noindent रथाद्यत् ।

\marginnote{\scriptsize ४.३.१२२}\noindent पत्त्रपूर्वादञ् ।

\marginnote{\scriptsize ४.३.१२३}\noindent पत्त्राध्वर्युपरिषदश्च ।

\marginnote{\scriptsize ४.३.१२४}\noindent हलसीराट्ठक् ।

\marginnote{\scriptsize ४.३.१२५}\noindent द्वंद्वाद्वुन् वैरमैथुनिकयोः ।

\marginnote{\scriptsize ४.३.१२६}\noindent गोत्रचरणाद्वुञ् ।

\marginnote{\scriptsize ४.३.१२७}\noindent संघाङ्कलक्षणेष्वञ्यञिञामण् ।

\marginnote{\scriptsize ४.३.१२८}\noindent शाकलाद्वा ।

\marginnote{\scriptsize ४.३.१२९}\noindent छन्दोगौक्थिकयाज्ञिकबह्वृचनटाञ्ञ्यः ।

\marginnote{\scriptsize ४.३.१३०}\noindent न दण्डमाणवान्तेवासिषु ।

\marginnote{\scriptsize ४.३.१३१}\noindent रैवतिकादिभ्यश्छः ।

\marginnote{\scriptsize ४.३.१३२}\noindent कौपिञ्जलहास्तिपदादण् ।

\marginnote{\scriptsize ४.३.१३३}\noindent आथर्वणिकस्येकलोपश्च ।

\marginnote{\scriptsize ४.३.१३४}\noindent तस्य विकारः ।

\marginnote{\scriptsize ४.३.१३५}\noindent अवयवे च प्राण्योषधिवृक्षेभ्यः ।

\marginnote{\scriptsize ४.३.१३६}\noindent बिल्वादिभ्योऽण् ।

\marginnote{\scriptsize ४.३.१३७}\noindent कोपधाच्च ।

\marginnote{\scriptsize ४.३.१३८}\noindent त्रपुजतुनोः षुक् ।

\marginnote{\scriptsize ४.३.१३९}\noindent ओरञ् ।

\marginnote{\scriptsize ४.३.१४०}\noindent अनुदात्तादेश्च ।

\marginnote{\scriptsize ४.३.१४१}\noindent पलाशादिभ्यो वा ।

\marginnote{\scriptsize ४.३.१४२}\noindent शम्याष्ट्लञ् ।

\marginnote{\scriptsize ४.३.१४३}\noindent मयड्वैतयोर्भाषायामभक्ष्याच्छादनयोः ।

\marginnote{\scriptsize ४.३.१४४}\noindent नित्यं वृद्धशरादिभ्यः ।

\marginnote{\scriptsize ४.३.१४५}\noindent गोश्च पुरीषे ।

\marginnote{\scriptsize ४.३.१४६}\noindent पिष्टाच्च ।

\marginnote{\scriptsize ४.३.१४७}\noindent संज्ञायां कन् ।

\marginnote{\scriptsize ४.३.१४८}\noindent व्रीहेः पुरोडाशे ।

\marginnote{\scriptsize ४.३.१४९}\noindent असंज्ञायां तिलयवाभ्याम् ।

\marginnote{\scriptsize ४.३.१५०}\noindent द्व्यचश्छन्दसि ।

\marginnote{\scriptsize ४.३.१५१}\noindent नोत्वद्वर्ध्रबिल्वात् ।

\marginnote{\scriptsize ४.३.१५२}\noindent तालादिभ्योऽण् ।

\marginnote{\scriptsize ४.३.१५३}\noindent जातरूपेभ्यः परिमाणे ।

\marginnote{\scriptsize ४.३.१५४}\noindent प्राणिरजतादिभ्योऽञ् ।

\marginnote{\scriptsize ४.३.१५५}\noindent ञितश्च तत्प्रत्ययात् ।

\marginnote{\scriptsize ४.३.१५६}\noindent क्रीतवत् परिमाणात् ।

\marginnote{\scriptsize ४.३.१५७}\noindent उष्ट्राद्वुञ् ।

\marginnote{\scriptsize ४.३.१५८}\noindent उमोर्णयोर्वा ।

\marginnote{\scriptsize ४.३.१५९}\noindent एण्या ढञ् ।

\marginnote{\scriptsize ४.३.१६०}\noindent गोपयसोर्यत् ।

\marginnote{\scriptsize ४.३.१६१}\noindent द्रोश्च ।

\marginnote{\scriptsize ४.३.१६२}\noindent माने वयः ।

\marginnote{\scriptsize ४.३.१६३}\noindent फले लुक् ।

\marginnote{\scriptsize ४.३.१६४}\noindent प्लक्षादिभ्योऽण् ।

\marginnote{\scriptsize ४.३.१६५}\noindent जम्ब्वा वा ।

\marginnote{\scriptsize ४.३.१६६}\noindent लुप् च ।

\marginnote{\scriptsize ४.३.१६७}\noindent हरीतक्यादिभ्यश्च ।

\marginnote{\scriptsize ४.३.१६८}\noindent कंसीयपरशव्ययोर्यञञौ लुक् च ।

\marginnote{\scriptsize ४.४.१}\noindent प्राग्वहतेष्ठक् ।

\marginnote{\scriptsize ४.४.२}\noindent तेन दीव्यति खनति जयति जितम् ।

\marginnote{\scriptsize ४.४.३}\noindent संस्कृतम् ।

\marginnote{\scriptsize ४.४.४}\noindent कुलत्थकोपधादण् ।

\marginnote{\scriptsize ४.४.५}\noindent तरति ।

\marginnote{\scriptsize ४.४.६}\noindent गोपुच्छाट्ठञ् ।

\marginnote{\scriptsize ४.४.७}\noindent नौद्व्यचष्ठन् ।

\marginnote{\scriptsize ४.४.८}\noindent चरति ।

\marginnote{\scriptsize ४.४.९}\noindent आकर्षात् ष्ठल् ।

\marginnote{\scriptsize ४.४.१०}\noindent पर्पादिभ्यः ष्ठन् ।

\marginnote{\scriptsize ४.४.११}\noindent श्वगणाट्ठञ्च ।

\marginnote{\scriptsize ४.४.१२}\noindent वेतनादिभ्यो जीवति ।

\marginnote{\scriptsize ४.४.१३}\noindent वस्नक्रयविक्रयाट्ठन् ।

\marginnote{\scriptsize ४.४.१४}\noindent आयुधाच्छ च ।

\marginnote{\scriptsize ४.४.१५}\noindent हरत्युत्सङ्गादिभ्यः ।

\marginnote{\scriptsize ४.४.१६}\noindent भस्त्राऽऽदिभ्यः ष्ठन् ।

\marginnote{\scriptsize ४.४.१७}\noindent विभाषा विवधवीवधात् ।

\marginnote{\scriptsize ४.४.१८}\noindent अण् कुटिलिकायाः ।

\marginnote{\scriptsize ४.४.१९}\noindent निर्वृत्तेऽक्षद्यूतादिभ्यः ।

\marginnote{\scriptsize ४.४.२०}\noindent क्त्रेर्मम् नित्यं ।

\marginnote{\scriptsize ४.४.२१}\noindent अपमित्ययाचिताभ्यां कक्कनौ ।

\marginnote{\scriptsize ४.४.२२}\noindent संसृष्टे ।

\marginnote{\scriptsize ४.४.२३}\noindent चूर्णादिनिः ।

\marginnote{\scriptsize ४.४.२४}\noindent लवणाल्लुक् ।

\marginnote{\scriptsize ४.४.२५}\noindent मुद्गादण् ।

\marginnote{\scriptsize ४.४.२६}\noindent व्यञ्जनैरुपसिक्ते ।

\marginnote{\scriptsize ४.४.२७}\noindent ओजस्सहोऽम्भसा वर्तते ।

\marginnote{\scriptsize ४.४.२८}\noindent तत् प्रत्यनुपूर्वमीपलोमकूलम् ।

\marginnote{\scriptsize ४.४.२९}\noindent परिमुखं च ।

\marginnote{\scriptsize ४.४.३०}\noindent प्रयच्छति गर्ह्यम् ।

\marginnote{\scriptsize ४.४.३१}\noindent कुसीददशैकादशात् ष्ठन्ष्ठचौ ।

\marginnote{\scriptsize ४.४.३२}\noindent उञ्छति ।

\marginnote{\scriptsize ४.४.३३}\noindent रक्षति ।

\marginnote{\scriptsize ४.४.३४}\noindent शब्ददर्दुरं करोति ।

\marginnote{\scriptsize ४.४.३५}\noindent पक्षिमत्स्यमृगान् हन्ति ।

\marginnote{\scriptsize ४.४.३६}\noindent परिपन्थं च तिष्ठति ।

\marginnote{\scriptsize ४.४.३७}\noindent माथोत्तरपदपदव्यनुपदं धावति ।

\marginnote{\scriptsize ४.४.३८}\noindent आक्रन्दाट्ठञ्च ।

\marginnote{\scriptsize ४.४.३९}\noindent पदोत्तरपदं गृह्णाति ।

\marginnote{\scriptsize ४.४.४०}\noindent प्रतिकण्ठार्थललामं च ।

\marginnote{\scriptsize ४.४.४१}\noindent धर्मं चरति ।

\marginnote{\scriptsize ४.४.४२}\noindent प्रतिपथमेति ठंश्च ।

\marginnote{\scriptsize ४.४.४३}\noindent समवायान् समवैति ।

\marginnote{\scriptsize ४.४.४४}\noindent परिषदो ण्यः ।

\marginnote{\scriptsize ४.४.४५}\noindent सेनाया वा ।

\marginnote{\scriptsize ४.४.४६}\noindent संज्ञायां ललाटकुक्कुट्यौ पश्यति ।

\marginnote{\scriptsize ४.४.४७}\noindent तस्य धर्म्यम् ।

\marginnote{\scriptsize ४.४.४८}\noindent अण् महिष्यादिभ्यः ।

\marginnote{\scriptsize ४.४.४९}\noindent ऋतोऽञ् ।

\marginnote{\scriptsize ४.४.५०}\noindent अवक्रयः ।

\marginnote{\scriptsize ४.४.५१}\noindent तदस्य पण्यम् ।

\marginnote{\scriptsize ४.४.५२}\noindent लवणाट्ठञ् ।

\marginnote{\scriptsize ४.४.५३}\noindent किशरादिभ्यः ष्ठन् ।

\marginnote{\scriptsize ४.४.५४}\noindent शलालुनोऽन्यतरस्याम् ।

\marginnote{\scriptsize ४.४.५५}\noindent शिल्पम् ।

\marginnote{\scriptsize ४.४.५६}\noindent मड्डुकझर्झरादणन्यतरस्याम् ।

\marginnote{\scriptsize ४.४.५७}\noindent प्रहरणम् ।

\marginnote{\scriptsize ४.४.५८}\noindent परश्वधाट्ठञ्च ।

\marginnote{\scriptsize ४.४.५९}\noindent शक्तियष्ट्योरीकक् ।

\marginnote{\scriptsize ४.४.६०}\noindent अस्तिनास्तिदिष्टं मतिः ।

\marginnote{\scriptsize ४.४.६१}\noindent शीलम् ।

\marginnote{\scriptsize ४.४.६२}\noindent छत्रादिभ्यो णः ।

\marginnote{\scriptsize ४.४.६३}\noindent कर्माध्ययने वृत्तम् ।

\marginnote{\scriptsize ४.४.६४}\noindent बह्वच्पूर्वपदाट्ठच् ।

\marginnote{\scriptsize ४.४.६५}\noindent हितं भक्षाः ।

\marginnote{\scriptsize ४.४.६६}\noindent तदस्मै दीयते नियुक्तम् ।

\marginnote{\scriptsize ४.४.६७}\noindent श्राणामांसौदनाट्टिठन् ।

\marginnote{\scriptsize ४.४.६८}\noindent भक्तादणन्यतरस्याम् ।

\marginnote{\scriptsize ४.४.६९}\noindent तत्र नियुक्तः ।

\marginnote{\scriptsize ४.४.७०}\noindent अगारान्ताट्ठन् ।

\marginnote{\scriptsize ४.४.७१}\noindent अध्यायिन्यदेशकालात् ।

\marginnote{\scriptsize ४.४.७२}\noindent कठिनान्तप्रस्तारसंस्थानेषु व्यवहरति ।

\marginnote{\scriptsize ४.४.७३}\noindent निकटे वसति ।

\marginnote{\scriptsize ४.४.७४}\noindent आवसथात् ष्ठल् ।

\marginnote{\scriptsize ४.४.७५}\noindent प्राग्घिताद्यत् ।

\marginnote{\scriptsize ४.४.७६}\noindent तद्वहति रथयुगप्रासङ्गम् ।

\marginnote{\scriptsize ४.४.७७}\noindent धुरो यड्ढकौ ।

\marginnote{\scriptsize ४.४.७८}\noindent खः सर्वधुरात् ।

\marginnote{\scriptsize ४.४.७९}\noindent एकधुराल्लुक् च ।

\marginnote{\scriptsize ४.४.८०}\noindent शकटादण् ।

\marginnote{\scriptsize ४.४.८१}\noindent हलसीराट्ठक् ।

\marginnote{\scriptsize ४.४.८२}\noindent संज्ञायां जन्याः ।

\marginnote{\scriptsize ४.४.८३}\noindent विध्यत्यधनुषा ।

\marginnote{\scriptsize ४.४.८४}\noindent धनगणं लब्धा ।

\marginnote{\scriptsize ४.४.८५}\noindent अन्नाण्णः ।

\marginnote{\scriptsize ४.४.८६}\noindent वशं गतः ।

\marginnote{\scriptsize ४.४.८७}\noindent पदमस्मिन् दृश्यम् ।

\marginnote{\scriptsize ४.४.८८}\noindent मूलमस्याबर्हि ।

\marginnote{\scriptsize ४.४.८९}\noindent संज्ञायां धेनुष्या ।

\marginnote{\scriptsize ४.४.९०}\noindent गृहपतिना संयुक्ते ञ्यः ।

\marginnote{\scriptsize ४.४.९१}\noindent नौवयोधर्मविषमूलमूलसीतातुलाभ्यस्तार्यतुल्यप्राप्य- वध्यानाम्यसमसमितसम्मितेषु ।

\marginnote{\scriptsize ४.४.९२}\noindent धर्मपथ्यर्थन्यायादनपेते ।

\marginnote{\scriptsize ४.४.९३}\noindent छन्दसो निर्मिते ।

\marginnote{\scriptsize ४.४.९४}\noindent उरसोऽण् च ।

\marginnote{\scriptsize ४.४.९५}\noindent हृदयस्य प्रियः ।

\marginnote{\scriptsize ४.४.९६}\noindent बन्धने चर्षौ ।

\marginnote{\scriptsize ४.४.९७}\noindent मतजनहलात् करणजल्पकर्षेषु ।

\marginnote{\scriptsize ४.४.९८}\noindent तत्र साधुः ।

\marginnote{\scriptsize ४.४.९९}\noindent प्रतिजनादिभ्यः खञ् ।

\marginnote{\scriptsize ४.४.१००}\noindent भक्ताण्णः ।

\marginnote{\scriptsize ४.४.१०१}\noindent परिषदो ण्यः ।

\marginnote{\scriptsize ४.४.१०२}\noindent कथाऽऽदिभ्यष्ठक् ।

\marginnote{\scriptsize ४.४.१०३}\noindent गुडादिभ्यष्ठञ् ।

\marginnote{\scriptsize ४.४.१०४}\noindent पथ्यतिथिवसतिस्वपतेर्ढञ् ।

\marginnote{\scriptsize ४.४.१०५}\noindent सभाया यः ।

\marginnote{\scriptsize ४.४.१०६}\noindent ढश्छन्दसि ।

\marginnote{\scriptsize ४.४.१०७}\noindent समानतीर्थे वासी ।

\marginnote{\scriptsize ४.४.१०८}\noindent समानोदरे शयित ओ चोदात्तः ।

\marginnote{\scriptsize ४.४.१०९}\noindent सोदराद्यः ।

\marginnote{\scriptsize ४.४.११०}\noindent भवे छन्दसि ।

\marginnote{\scriptsize ४.४.१११}\noindent पाथोनदीभ्यां ड्यण् ।

\marginnote{\scriptsize ४.४.११२}\noindent वेशन्तहिमवद्भ्यामण् ।

\marginnote{\scriptsize ४.४.११३}\noindent स्रोतसो विभाषा ड्यड्ड्यौ ।

\marginnote{\scriptsize ४.४.११४}\noindent सगर्भसयूथसनुताद्यन् ।

\marginnote{\scriptsize ४.४.११५}\noindent तुग्राद्घन् ।

\marginnote{\scriptsize ४.४.११६}\noindent अग्राद्यत् ।

\marginnote{\scriptsize ४.४.११७}\noindent घच्छौ च ।

\marginnote{\scriptsize ४.४.११८}\noindent समुद्राभ्राद्घः ।

\marginnote{\scriptsize ४.४.११९}\noindent बर्हिषि दत्तम् ।

\marginnote{\scriptsize ४.४.१२०}\noindent दूतस्य भागकर्मणी ।

\marginnote{\scriptsize ४.४.१२१}\noindent रक्षोयातूनां हननी ।

\marginnote{\scriptsize ४.४.१२२}\noindent रेवतीजगतीहविष्याभ्यः प्रशस्ये ।

\marginnote{\scriptsize ४.४.१२३}\noindent असुरस्य स्वम् ।

\marginnote{\scriptsize ४.४.१२४}\noindent मायायामण् ।

\marginnote{\scriptsize ४.४.१२५}\noindent तद्वानासामुपधानो मन्त्र इतीष्टकासु लुक् च मतोः ।

\marginnote{\scriptsize ४.४.१२६}\noindent अश्विमानण् ।

\marginnote{\scriptsize ४.४.१२७}\noindent वयस्यासु मूर्ध्नो मतुप् ।

\marginnote{\scriptsize ४.४.१२८}\noindent मत्वर्थे मासतन्वोः । ४ ४.१२९ मधोर्ञ च ।

\marginnote{\scriptsize ४.४.१३०}\noindent ओजसोऽहनि यत्खौ ।

\marginnote{\scriptsize ४.४.१३१}\noindent वेशोयशादेर्भगाद्यल् ।

\marginnote{\scriptsize ४.४.१३२}\noindent ख च ।

\marginnote{\scriptsize ४.४.१३३}\noindent पूर्वैः कृतमिनियौ च ।

\marginnote{\scriptsize ४.४.१३४}\noindent अद्भिः संस्कृतम् ।

\marginnote{\scriptsize ४.४.१३५}\noindent सहस्रेण संमितौ घः ।

\marginnote{\scriptsize ४.४.१३६}\noindent मतौ च ।

\marginnote{\scriptsize ४.४.१३७}\noindent सोममर्हति यः ।

\marginnote{\scriptsize ४.४.१३८}\noindent मये च ।

\marginnote{\scriptsize ४.४.१३९}\noindent मधोः ।

\marginnote{\scriptsize ४.४.१४०}\noindent वसोः समूहे च ।

\marginnote{\scriptsize ४.४.१४१}\noindent नक्षत्राद्घः ।

\marginnote{\scriptsize ४.४.१४२}\noindent सर्वदेवात् तातिल् ।

\marginnote{\scriptsize ४.४.१४३}\noindent शिवशमरिष्टस्य करे ।

\marginnote{\scriptsize ४.४.१४४}\noindent भावे च ।

\section*{॥ अध्याय ५ ॥}

\marginnote{\scriptsize ५.१.१}\noindent प्राक् क्रीताच्छः ।

\marginnote{\scriptsize ५.१.२}\noindent उगवादिभ्योऽत् ।

\marginnote{\scriptsize ५.१.३}\noindent कम्बलाच्च संज्ञायाम् ।

\marginnote{\scriptsize ५.१.४}\noindent विभाषा हविरपूपादिभ्यः ।

\marginnote{\scriptsize ५.१.५}\noindent तस्मै हितम् ।

\marginnote{\scriptsize ५.१.६}\noindent शरीरावयवाद्यत् ।

\marginnote{\scriptsize ५.१.७}\noindent खलयवमाषतिलवृषब्रह्मणश्च ।

\marginnote{\scriptsize ५.१.८}\noindent अजाविभ्यां थ्यन् ।

\marginnote{\scriptsize ५.१.९}\noindent आत्मन्विश्वजनभोगोत्तरपदात् खः ।

\marginnote{\scriptsize ५.१.१०}\noindent सर्वपुरुषाभ्यां णढञौ ।

\marginnote{\scriptsize ५.१.११}\noindent माणवचरकाभ्यां खञ् ।

\marginnote{\scriptsize ५.१.१२}\noindent तदर्थं विकृतेः प्रकृतौ ।

\marginnote{\scriptsize ५.१.१३}\noindent छदिरुपधिबलेः ढञ् ।

\marginnote{\scriptsize ५.१.१४}\noindent ऋषभोपानहोर्ञ्यः ।

\marginnote{\scriptsize ५.१.१५}\noindent चर्म्मणोऽञ् ।

\marginnote{\scriptsize ५.१.१६}\noindent तदस्य तदस्मिन् स्यादिति ।

\marginnote{\scriptsize ५.१.१७}\noindent परिखाया ढञ् ।

\marginnote{\scriptsize ५.१.१८}\noindent प्राग्वतेष्ठञ् ।

\marginnote{\scriptsize ५.१.१९}\noindent आर्हादगोपुच्छसंख्यापरिमाणाट्ठक् ।

\marginnote{\scriptsize ५.१.२०}\noindent असमासे निष्कादिभ्यः ।

\marginnote{\scriptsize ५.१.२१}\noindent शताच्च ठन्यतावशते ।

\marginnote{\scriptsize ५.१.२२}\noindent संख्याया अतिशदन्तायाः कन् ।

\marginnote{\scriptsize ५.१.२३}\noindent वतोरिड्वा ।

\marginnote{\scriptsize ५.१.२४}\noindent विंशतित्रिंशद्भ्यां ड्वुन्नसंज्ञायाम् ।

\marginnote{\scriptsize ५.१.२५}\noindent कंसाट्टिठन् ।

\marginnote{\scriptsize ५.१.२६}\noindent शूर्पादञन्यतरस्याम् ।

\marginnote{\scriptsize ५.१.२७}\noindent शतमानविंशतिकसहस्रवसनादण् ।

\marginnote{\scriptsize ५.१.२८}\noindent अध्यर्धपूर्वद्विगोर्लुगसंज्ञायाम् ।

\marginnote{\scriptsize ५.१.२९}\noindent विभाषा कार्षापणसहस्राभ्याम् ।

\marginnote{\scriptsize ५.१.३०}\noindent द्वित्रिपूर्वान्निष्कात् ।

\marginnote{\scriptsize ५.१.३१}\noindent बिस्ताच्च ।

\marginnote{\scriptsize ५.१.३२}\noindent विंशतिकात् खः ।

\marginnote{\scriptsize ५.१.३३}\noindent खार्या ईकन् ।

\marginnote{\scriptsize ५.१.३४}\noindent पणपादमाषशतादत् ।

\marginnote{\scriptsize ५.१.३५}\noindent शाणाद्वा ।

\marginnote{\scriptsize ५.१.३६}\noindent द्वित्रिपूर्वादण् च ।

\marginnote{\scriptsize ५.१.३७}\noindent तेन क्रीतम् ।

\marginnote{\scriptsize ५.१.३८}\noindent तस्य निमित्तं संयोगोत्पातौ ।

\marginnote{\scriptsize ५.१.३९}\noindent गोद्व्यचोरसंख्यापरिमाणाश्वादेर्यत् ।

\marginnote{\scriptsize ५.१.४०}\noindent पुत्राच्छ च ।

\marginnote{\scriptsize ५.१.४१}\noindent सर्वभूमिपृथिवीभ्यामणञौ ।

\marginnote{\scriptsize ५.१.४२}\noindent तस्येश्वरः ।

\marginnote{\scriptsize ५.१.४३}\noindent तत्र विदित इति च ।

\marginnote{\scriptsize ५.१.४४}\noindent लोकसर्वलोकाट्ठञ् ।

\marginnote{\scriptsize ५.१.४५}\noindent तस्य वापः ।

\marginnote{\scriptsize ५.१.४६}\noindent पात्रात् ष्ठन् ।

\marginnote{\scriptsize ५.१.४७}\noindent तदस्मिन् वृद्ध्यायलाभशुल्कोपदा दीयते ।

\marginnote{\scriptsize ५.१.४८}\noindent पूरणार्धाट्ठन् ।

\marginnote{\scriptsize ५.१.४९}\noindent भागाद्यच्च ।

\marginnote{\scriptsize ५.१.५०}\noindent तद्धरति वहत्यावहति भाराद्वंशादिभ्यः ।

\marginnote{\scriptsize ५.१.५१}\noindent वस्नद्रव्याभ्यां ठन्कनौ ।

\marginnote{\scriptsize ५.१.५२}\noindent सम्भवत्यवहरति पचति ।

\marginnote{\scriptsize ५.१.५३}\noindent आढकाचितपात्रात् खोऽन्यतरयाम् ।

\marginnote{\scriptsize ५.१.५४}\noindent द्विगोष्ठंश्च ।

\marginnote{\scriptsize ५.१.५५}\noindent कुलिजाल्लुक्खौ च ।

\marginnote{\scriptsize ५.१.५६}\noindent सोऽस्यांशवस्नभृतयः ।

\marginnote{\scriptsize ५.१.५७}\noindent तदस्य परिमाणम् ।

\marginnote{\scriptsize ५.१.५८}\noindent संख्यायाः संज्ञासंघसूत्राध्ययनेषु ।

\marginnote{\scriptsize ५.१.५९}\noindent पङ्क्तिविंशतित्रिंशत्चत्वारिंशत्पञ्चाशत्षष्टि- सप्तत्यशीतिनवतिशतम् ।

\marginnote{\scriptsize ५.१.६०}\noindent पञ्चद्दशतौ वर्गे वा ।

\marginnote{\scriptsize ५.१.६१}\noindent सप्तनोऽञ् छन्दसि ।

\marginnote{\scriptsize ५.१.६२}\noindent त्रिंशच्चत्वारिंशतोर्ब्राह्मणे संज्ञायां डण् ।

\marginnote{\scriptsize ५.१.६३}\noindent तद् अर्हति ।

\marginnote{\scriptsize ५.१.६४}\noindent छेदादिभ्यो नित्यम् ।

\marginnote{\scriptsize ५.१.६५}\noindent शीर्षच्छेदाद्यच्च ।

\marginnote{\scriptsize ५.१.६६}\noindent दण्डादिभ्यः ।

\marginnote{\scriptsize ५.१.६७}\noindent छन्दसि च ।

\marginnote{\scriptsize ५.१.६८}\noindent पात्राद्घंश्च ।

\marginnote{\scriptsize ५.१.६९}\noindent कडङ्गरदक्षिणाच्छ च ।

\marginnote{\scriptsize ५.१.७०}\noindent स्थालीबिलात् ।

\marginnote{\scriptsize ५.१.७१}\noindent यज्ञर्त्विग्भ्यां घखञौ ।

\marginnote{\scriptsize ५.१.७२}\noindent पारायणतुरायणचान्द्रायणं वर्तयति ।

\marginnote{\scriptsize ५.१.७३}\noindent संशयमापन्नः ।

\marginnote{\scriptsize ५.१.७४}\noindent योजनं गच्छति ।

\marginnote{\scriptsize ५.१.७५}\noindent पथः ष्कन् ।

\marginnote{\scriptsize ५.१.७६}\noindent पन्थो ण नित्यम् ।

\marginnote{\scriptsize ५.१.७७}\noindent उत्तरपथेनाहृतं च ।

\marginnote{\scriptsize ५.१.७८}\noindent कालात् ।

\marginnote{\scriptsize ५.१.७९}\noindent तेन निर्वृत्तम् ।

\marginnote{\scriptsize ५.१.८०}\noindent तमधीष्टो भृतो भूतो भावी ।

\marginnote{\scriptsize ५.१.८१}\noindent मासाद्वयसि यत्खञौ ।

\marginnote{\scriptsize ५.१.८२}\noindent द्विगोर्यप् ।

\marginnote{\scriptsize ५.१.८३}\noindent षण्मासाण्ण्यच्च ।

\marginnote{\scriptsize ५.१.८४}\noindent अवयसि ठंश्च ।

\marginnote{\scriptsize ५.१.८५}\noindent समायाः खः ।

\marginnote{\scriptsize ५.१.८६}\noindent द्विगोर्वा ।

\marginnote{\scriptsize ५.१.८७}\noindent रात्र्यहस्संवत्सराच्च ।

\marginnote{\scriptsize ५.१.८८}\noindent वर्षाल्लुक् च ।

\marginnote{\scriptsize ५.१.८९}\noindent चित्तवति नित्यम् ।

\marginnote{\scriptsize ५.१.९०}\noindent षष्टिकाः षष्टिरात्रेण पच्यन्ते ।

\marginnote{\scriptsize ५.१.९१}\noindent वत्सरान्ताच्छश्छन्दसि ।

\marginnote{\scriptsize ५.१.९२}\noindent सम्परिपूर्वात् ख च ।

\marginnote{\scriptsize ५.१.९३}\noindent तेन परिजय्यलभ्यकार्यसुकरम् ।

\marginnote{\scriptsize ५.१.९४}\noindent तदस्य ब्रह्मचर्यम् ।

\marginnote{\scriptsize ५.१.९५}\noindent तस्य च दक्षिणा यज्ञाख्येभ्यः ।

\marginnote{\scriptsize ५.१.९६}\noindent तत्र च दीयते कार्यं भववत् ।

\marginnote{\scriptsize ५.१.९७}\noindent व्युष्टादिभ्योऽण् ।

\marginnote{\scriptsize ५.१.९८}\noindent तेन यथाकथाचहस्ताभ्यां णयतौ ।

\marginnote{\scriptsize ५.१.९९}\noindent सम्पादिनि ।

\marginnote{\scriptsize ५.१.१००}\noindent कर्मवेषाद्यत् ।

\marginnote{\scriptsize ५.१.१०१}\noindent तस्मै प्रभवति संतापादिभ्यः ।

\marginnote{\scriptsize ५.१.१०२}\noindent योगाद्यच्च ।

\marginnote{\scriptsize ५.१.१०३}\noindent कर्मण उकञ् ।

\marginnote{\scriptsize ५.१.१०४}\noindent समयस्तदस्य प्राप्तम् ।

\marginnote{\scriptsize ५.१.१०५}\noindent ऋतोरण् ।

\marginnote{\scriptsize ५.१.१०६}\noindent छन्दसि घस् ।

\marginnote{\scriptsize ५.१.१०७}\noindent कालाद्यत् ।

\marginnote{\scriptsize ५.१.१०८}\noindent प्रकृष्टे ठञ् ।

\marginnote{\scriptsize ५.१.१०९}\noindent प्रयोजनम् ।

\marginnote{\scriptsize ५.१.११०}\noindent विशाखाऽऽषाढादण् मन्थदण्डयोः ।

\marginnote{\scriptsize ५.१.१११}\noindent अनुप्रवचनादिभ्यश्छः ।

\marginnote{\scriptsize ५.१.११२}\noindent समापनात् सपूर्वपदात् ।

\marginnote{\scriptsize ५.१.११३}\noindent ऐकागारिकट् चौरे ।

\marginnote{\scriptsize ५.१.११४}\noindent आकालिकडाद्यन्तवचने ।

\marginnote{\scriptsize ५.१.११५}\noindent तेन तुल्यं क्रिया चेद्वतिः ।

\marginnote{\scriptsize ५.१.११६}\noindent तत्र तस्येव ।

\marginnote{\scriptsize ५.१.११७}\noindent तदर्हम् ।

\marginnote{\scriptsize ५.१.११८}\noindent उपसर्गाच्छन्दसि धात्वर्थे ।

\marginnote{\scriptsize ५.१.११९}\noindent तस्य भावस्त्वतलौ ।

\marginnote{\scriptsize ५.१.१२०}\noindent आ च त्वात् ।

\marginnote{\scriptsize ५.१.१२१}\noindent न नञ्पूर्वात्तत्पुरुषादचतुरसंगतलवणवटयुधकतर- सलसेभ्यः ।

\marginnote{\scriptsize ५.१.१२२}\noindent पृथ्वादिभ्य इमनिज्वा ।

\marginnote{\scriptsize ५.१.१२३}\noindent वर्णदृढादिभ्यः ष्यञ् च ।

\marginnote{\scriptsize ५.१.१२४}\noindent गुणवचनब्राह्मणादिभ्यः कर्मणि च ।

\marginnote{\scriptsize ५.१.१२५}\noindent स्तेनाद्यन्नलोपश्च ।

\marginnote{\scriptsize ५.१.१२६}\noindent सख्युर्यः ।

\marginnote{\scriptsize ५.१.१२७}\noindent कपिज्ञात्योर्ढक् ।

\marginnote{\scriptsize ५.१.१२८}\noindent पत्यन्तपुरोहितादिभ्यो यक् ।

\marginnote{\scriptsize ५.१.१२९}\noindent प्राणभृज्जातिवयोवचनोद्गात्रादिभ्योऽञ् ।

\marginnote{\scriptsize ५.१.१३०}\noindent हायनान्तयुवादिभ्योऽण् ।

\marginnote{\scriptsize ५.१.१३१}\noindent इगन्ताच्च लघुपूर्वात् ।

\marginnote{\scriptsize ५.१.१३२}\noindent योपधाद्गुरूपोत्तमाद्वुञ् ।

\marginnote{\scriptsize ५.१.१३३}\noindent द्वंद्वमनोज्ञादिभ्यश्च ।

\marginnote{\scriptsize ५.१.१३४}\noindent गोत्रचरणाच्श्लाघाऽत्याकारतदवेतेषु ।

\marginnote{\scriptsize ५.१.१३५}\noindent होत्राभ्यश्छः ।

\marginnote{\scriptsize ५.१.१३६}\noindent ब्रह्मणस्त्वः ।

\marginnote{\scriptsize ५.२.१}\noindent धान्यानां भवने क्षेत्रे खञ् ।

\marginnote{\scriptsize ५.२.२}\noindent व्रीहिशाल्योर्ढक् ।

\marginnote{\scriptsize ५.२.३}\noindent यवयवकषष्टिकादत् ।

\marginnote{\scriptsize ५.२.४}\noindent विभाषा तिलमाषोमाभङ्गाऽणुभ्यः ।

\marginnote{\scriptsize ५.२.५}\noindent सर्वचर्मणः कृतः खखञौ ।

\marginnote{\scriptsize ५.२.६}\noindent यथामुखसंमुखस्य दर्शनः खः ।

\marginnote{\scriptsize ५.२.७}\noindent तत्सर्वादेः पथ्यङ्गकर्मपत्रपात्रं व्याप्नोति ।

\marginnote{\scriptsize ५.२.८}\noindent आप्रपदं प्राप्नोति ।

\marginnote{\scriptsize ५.२.९}\noindent अनुपदसर्वान्नायानयं बद्धाभक्षयतिनेयेषु ।

\marginnote{\scriptsize ५.२.१०}\noindent परोवरपरम्परपुत्रपौत्रमनुभवति ।

\marginnote{\scriptsize ५.२.११}\noindent अवारपारात्यन्तानुकामं गामी ।

\marginnote{\scriptsize ५.२.१२}\noindent समांसमां विजायते ।

\marginnote{\scriptsize ५.२.१३}\noindent अद्यश्वीनाऽवष्टब्धे ।

\marginnote{\scriptsize ५.२.१४}\noindent आगवीनः ।

\marginnote{\scriptsize ५.२.१५}\noindent अनुग्वलंगामी ।

\marginnote{\scriptsize ५.२.१६}\noindent अध्वनो यत्खौ ।

\marginnote{\scriptsize ५.२.१७}\noindent अभ्यमित्राच्छ च ।

\marginnote{\scriptsize ५.२.१८}\noindent गोष्ठात् खञ् भूतपूर्वे ।

\marginnote{\scriptsize ५.२.१९}\noindent अश्वस्यैकाहगमः ।

\marginnote{\scriptsize ५.२.२०}\noindent शालीनकौपीने अधृष्टाकार्ययोः ।

\marginnote{\scriptsize ५.२.२१}\noindent व्रातेन जीवति ।

\marginnote{\scriptsize ५.२.२२}\noindent साप्तपदीनं सख्यम् ।

\marginnote{\scriptsize ५.२.२३}\noindent हैयंगवीनं संज्ञायाम् ।

\marginnote{\scriptsize ५.२.२४}\noindent तस्य पाकमूले पील्वदिकर्णादिभ्यः कुणब्जाहचौ ।

\marginnote{\scriptsize ५.२.२५}\noindent पक्षात्तिः ।

\marginnote{\scriptsize ५.२.२६}\noindent तेन वित्तश्चुञ्चुप्चणपौ ।

\marginnote{\scriptsize ५.२.२७}\noindent विनञ्भ्यां नानाञौ नसह ।

\marginnote{\scriptsize ५.२.२८}\noindent वेः शालच्छङ्कटचौ ।

\marginnote{\scriptsize ५.२.२९}\noindent सम्प्रोदश्च कटच् ।

\marginnote{\scriptsize ५.२.३०}\noindent अवात् कुटारच्च ।

\marginnote{\scriptsize ५.२.३१}\noindent नते नासिकायाः संज्ञायां टीटञ्नाटज्भ्राटचः ।

\marginnote{\scriptsize ५.२.३२}\noindent नेर्बिडज्बिरीसचौ ।

\marginnote{\scriptsize ५.२.३३}\noindent इनच्पिटच्चिकचि च ।

\marginnote{\scriptsize ५.२.३४}\noindent उपाधिभ्यां त्यकन्नासन्नारूढयोः ।

\marginnote{\scriptsize ५.२.३५}\noindent कर्मणि घटोऽठच् ।

\marginnote{\scriptsize ५.२.३६}\noindent तदस्य संजातं तारकाऽऽदिभ्य इतच् ।

\marginnote{\scriptsize ५.२.३७}\noindent प्रमाणे द्वयसज्दघ्नञ्मात्रचः ।

\marginnote{\scriptsize ५.२.३८}\noindent पुरुषहस्तिभ्यामण् च ।

\marginnote{\scriptsize ५.२.३९}\noindent यद्तदेतेभ्यः परिमाणे वतुप् ।

\marginnote{\scriptsize ५.२.४०}\noindent किमिदंभ्यां वो घः ।

\marginnote{\scriptsize ५.२.४१}\noindent किमः संख्यापरिमाणे डति च ।

\marginnote{\scriptsize ५.२.४२}\noindent संख्याया अवयवे तयप् ।

\marginnote{\scriptsize ५.२.४३}\noindent द्वित्रिभ्यां तयस्यायज्वा ।

\marginnote{\scriptsize ५.२.४४}\noindent उभादुदात्तो नित्यम् ।

\marginnote{\scriptsize ५.२.४५}\noindent तदस्मिन्नधिकमिति दशान्ताड्डः ।

\marginnote{\scriptsize ५.२.४६}\noindent शदन्तविंशतेश्च ।

\marginnote{\scriptsize ५.२.४७}\noindent संख्याया गुणस्य निमाने मयट् ।

\marginnote{\scriptsize ५.२.४८}\noindent तस्य पूरणे डट् ।

\marginnote{\scriptsize ५.२.४९}\noindent नान्तादसंख्याऽऽदेर्मट् ।

\marginnote{\scriptsize ५.२.५०}\noindent थट् च च्छन्दसि ।

\marginnote{\scriptsize ५.२.५१}\noindent षट्कतिकतिपयचतुरां थुक् ।

\marginnote{\scriptsize ५.२.५२}\noindent बहुपूगगणसंघस्य तिथुक् ।

\marginnote{\scriptsize ५.२.५३}\noindent वतोरिथुक् ।

\marginnote{\scriptsize ५.२.५४}\noindent द्वेस्तीयः ।

\marginnote{\scriptsize ५.२.५५}\noindent त्रेः सम्प्रसारणम् च ।

\marginnote{\scriptsize ५.२.५६}\noindent विंशत्यादिभ्यस्तमडन्यतरस्याम् ।

\marginnote{\scriptsize ५.२.५७}\noindent नित्यं शतादिमासार्धमाससंवत्सराच्च ।

\marginnote{\scriptsize ५.२.५८}\noindent षष्ट्यादेश्चासंख्याऽऽदेः ।

\marginnote{\scriptsize ५.२.५९}\noindent मतौ च्छः सूक्तसाम्नोः ।

\marginnote{\scriptsize ५.२.६०}\noindent अध्यायानुवाकयोर्लुक् ।

\marginnote{\scriptsize ५.२.६१}\noindent विमुक्तादिभ्योऽण् ।

\marginnote{\scriptsize ५.२.६२}\noindent गोषदादिभ्यो वुन् ।

\marginnote{\scriptsize ५.२.६३}\noindent तत्र कुशलः पथः ।

\marginnote{\scriptsize ५.२.६४}\noindent आकर्षादिभ्यः कन् ।

\marginnote{\scriptsize ५.२.६५}\noindent धनहिरण्यात् कामे ।

\marginnote{\scriptsize ५.२.६६}\noindent स्वाङ्गेभ्यः प्रसिते ।

\marginnote{\scriptsize ५.२.६७}\noindent उदराट्ठगाद्यूने ।

\marginnote{\scriptsize ५.२.६८}\noindent सस्येन परिजातः ।

\marginnote{\scriptsize ५.२.६९}\noindent अंशं हारी ।

\marginnote{\scriptsize ५.२.७०}\noindent तन्त्रादचिरापहृते ।

\marginnote{\scriptsize ५.२.७१}\noindent ब्राह्मणकोष्णिके संज्ञायाम् ।

\marginnote{\scriptsize ५.२.७२}\noindent शीतोष्णाभ्यां कारिणि ।

\marginnote{\scriptsize ५.२.७३}\noindent अधिकम् ।

\marginnote{\scriptsize ५.२.७४}\noindent अनुकाभिकाभीकः कमिता ।

\marginnote{\scriptsize ५.२.७५}\noindent पार्श्वेनान्विच्छति ।

\marginnote{\scriptsize ५.२.७६}\noindent अयःशूलदण्डाजिनाभ्यां ठक्ठञौ ।

\marginnote{\scriptsize ५.२.७७}\noindent तावतिथं ग्रहणमिति लुग्वा ।

\marginnote{\scriptsize ५.२.७८}\noindent स एषां ग्रामणीः ।

\marginnote{\scriptsize ५.२.७९}\noindent श‍ृङ्खलमस्य बन्धनं करभे ।

\marginnote{\scriptsize ५.२.८०}\noindent उत्क उन्मनाः ।

\marginnote{\scriptsize ५.२.८१}\noindent कालप्रयोजनाद्रोगे ।

\marginnote{\scriptsize ५.२.८२}\noindent तदस्मिन्नन्नं प्राये संज्ञायाम् ।

\marginnote{\scriptsize ५.२.८३}\noindent कुल्माषादञ् ।

\marginnote{\scriptsize ५.२.८४}\noindent श्रोत्रियंश्छन्दोऽधीते ।

\marginnote{\scriptsize ५.२.८५}\noindent श्राद्धमनेन भुक्तमिनिठनौ ।

\marginnote{\scriptsize ५.२.८६}\noindent पूर्वादिनिः ।

\marginnote{\scriptsize ५.२.८७}\noindent सपूर्वाच्च ।

\marginnote{\scriptsize ५.२.८८}\noindent इष्टादिभ्यश्च ।

\marginnote{\scriptsize ५.२.८९}\noindent छन्दसि परिपन्थिपरिपरिणौ पर्यवस्थातरि ।

\marginnote{\scriptsize ५.२.९०}\noindent अनुपद्यन्वेष्टा ।

\marginnote{\scriptsize ५.२.९१}\noindent साक्षाद्द्रष्टरि संज्ञायाम् ।

\marginnote{\scriptsize ५.२.९२}\noindent क्षेत्रियच् परक्षेत्रे चिकित्स्यः । ५.२.९३ इन्द्रियमिन्द्रलिंगमिन्द्रदृष्टमिन्द्रसृष्टमिन्द्रजुष्टम्- इन्द्रदत्तमिति वा ।

\marginnote{\scriptsize ५.२.९४}\noindent तदस्यास्त्यस्मिन्निति मतुप् ।

\marginnote{\scriptsize ५.२.९५}\noindent रसादिभ्यश्च ।

\marginnote{\scriptsize ५.२.९६}\noindent प्राणिस्थादातो लजन्यतरस्याम् ।

\marginnote{\scriptsize ५.२.९७}\noindent सिध्मादिभ्यश्च ।

\marginnote{\scriptsize ५.२.९८}\noindent वत्सांसाभ्यां कामबले ।

\marginnote{\scriptsize ५.२.९९}\noindent फेनादिलच् च ।

\marginnote{\scriptsize ५.२.१००}\noindent लोमादिपामादिपिच्छादिभ्यः शनेलचः ।

\marginnote{\scriptsize ५.२.१०१}\noindent प्रज्ञाश्रद्धाऽर्चावृत्तिभ्यो णः ।

\marginnote{\scriptsize ५.२.१०२}\noindent तपःसहस्राभ्यां विनीनी ।

\marginnote{\scriptsize ५.२.१०३}\noindent अण् च ।

\marginnote{\scriptsize ५.२.१०४}\noindent सिकताशर्कराभ्यां च ।

\marginnote{\scriptsize ५.२.१०५}\noindent देशे लुबिलचौ च ।

\marginnote{\scriptsize ५.२.१०६}\noindent दन्त उन्नत उरच् ।

\marginnote{\scriptsize ५.२.१०७}\noindent ऊषसुषिमुष्कमधो रः ।

\marginnote{\scriptsize ५.२.१०८}\noindent द्युद्रुभ्यां मः ।

\marginnote{\scriptsize ५.२.१०९}\noindent केशाद्वोऽन्यतरस्याम् ।

\marginnote{\scriptsize ५.२.११०}\noindent गाण्ड्यजगात् संज्ञायाम् ।

\marginnote{\scriptsize ५.२.१११}\noindent काण्डाण्डादीरन्नीरचौ ।

\marginnote{\scriptsize ५.२.११२}\noindent रजःकृष्यासुतिपरिषदो वलच् ।

\marginnote{\scriptsize ५.२.११३}\noindent दन्तशिखात् संज्ञायाम् ।

\marginnote{\scriptsize ५.२.११४}\noindent ज्योत्स्नातमिस्राश‍ृङ्गिणोजस्विन्नूर्जस्वलगोमिन्- मलिनमलीमसाः ।

\marginnote{\scriptsize ५.२.११५}\noindent अत इनिठनौ ।

\marginnote{\scriptsize ५.२.११६}\noindent व्रीह्यादिभ्यश्च ।

\marginnote{\scriptsize ५.२.११७}\noindent तुन्दादिभ्य इलच् च ।

\marginnote{\scriptsize ५.२.११८}\noindent एकगोपूर्वाट्ठञ् नित्यम् ।

\marginnote{\scriptsize ५.२.११९}\noindent शतसहस्रान्ताच्च निष्कात् ।

\marginnote{\scriptsize ५.२.१२०}\noindent रूपादाहतप्रशंसयोरप् ।

\marginnote{\scriptsize ५.२.१२१}\noindent अस्मायामेधास्रजो विनिः ।

\marginnote{\scriptsize ५.२.१२२}\noindent बहुलं छन्दसि ।

\marginnote{\scriptsize ५.२.१२३}\noindent ऊर्णाया युस् ।

\marginnote{\scriptsize ५.२.१२४}\noindent वाचो ग्मिनिः ।

\marginnote{\scriptsize ५.२.१२५}\noindent आलजाटचौ बहुभाषिणि ।

\marginnote{\scriptsize ५.२.१२६}\noindent स्वामिन्नैश्वर्ये ।

\marginnote{\scriptsize ५.२.१२७}\noindent अर्शादिभ्योऽच् ।

\marginnote{\scriptsize ५.२.१२८}\noindent द्वंद्वोपतापगर्ह्यात् प्राणिस्थादिनिः ।

\marginnote{\scriptsize ५.२.१२९}\noindent वातातिसाराभ्यां कुक् च ।

\marginnote{\scriptsize ५.२.१३०}\noindent वयसि पूरणात् ।

\marginnote{\scriptsize ५.२.१३१}\noindent सुखादिभ्यश्च ।

\marginnote{\scriptsize ५.२.१३२}\noindent धर्मशीलवर्णान्ताच्च ।

\marginnote{\scriptsize ५.२.१३३}\noindent हस्ताज्जातौ ।

\marginnote{\scriptsize ५.२.१३४}\noindent वर्णाद्ब्रह्मचारिणि ।

\marginnote{\scriptsize ५.२.१३५}\noindent पुष्करादिभ्यो देशे ।

\marginnote{\scriptsize ५.२.१३६}\noindent बलादिभ्यो मतुबन्यतरस्याम् ।

\marginnote{\scriptsize ५.२.१३७}\noindent संज्ञायां मन्माभ्याम् ।

\marginnote{\scriptsize ५.२.१३८}\noindent कंशंभ्यां बभयुस्तितुतयसः ।

\marginnote{\scriptsize ५.२.१३९}\noindent तुन्दिवलिवटेर्भः ।

\marginnote{\scriptsize ५.२.१४०}\noindent अहंशुभमोर्युस् ।

\marginnote{\scriptsize ५.३.१}\noindent प्राग्दिशो विभक्तिः ।

\marginnote{\scriptsize ५.३.२}\noindent किंसर्वनामबहुभ्योऽद्व्यादिभ्यः ।

\marginnote{\scriptsize ५.३.३}\noindent इदम इश् ।

\marginnote{\scriptsize ५.३.४}\noindent एतेतौ रथोः ।

\marginnote{\scriptsize ५.३.५}\noindent एतदोऽश् ।

\marginnote{\scriptsize ५.३.६}\noindent सर्वस्य सोऽन्यतरस्यां दि ।

\marginnote{\scriptsize ५.३.७}\noindent पञ्चम्यास्तसिल् ।

\marginnote{\scriptsize ५.३.८}\noindent तसेश्च ।

\marginnote{\scriptsize ५.३.९}\noindent पर्यभिभ्यां च ।

\marginnote{\scriptsize ५.३.१०}\noindent सप्तम्यास्त्रल् ।

\marginnote{\scriptsize ५.३.११}\noindent इदमो हः ।

\marginnote{\scriptsize ५.३.१२}\noindent किमोऽत् ।

\marginnote{\scriptsize ५.३.१३}\noindent वा ह च च्छन्दसि ।

\marginnote{\scriptsize ५.३.१४}\noindent इतराभ्योऽपि दृश्यन्ते ।

\marginnote{\scriptsize ५.३.१५}\noindent सर्वैकान्यकिंयत्तदः काले दा ।

\marginnote{\scriptsize ५.३.१६}\noindent इदमो र्हिल् ।

\marginnote{\scriptsize ५.३.१७}\noindent अधुना ।

\marginnote{\scriptsize ५.३.१८}\noindent दानीं च ।

\marginnote{\scriptsize ५.३.१९}\noindent तदो दा च ।

\marginnote{\scriptsize ५.३.२०}\noindent तयोर्दार्हिलौ च च्छन्दसि ।

\marginnote{\scriptsize ५.३.२१}\noindent अनद्यतने र्हिलन्यतरस्याम् ।

\marginnote{\scriptsize ५.३.२२}\noindent सद्यःपरुत्परार्यैषमःपरेद्यव्यद्यपूर्वेद्युरन्येद्युर्- अन्यतरेद्युरितरेद्युरपरेद्युरधरेद्युरुभयेद्युरुत्तरेद्युः ।

\marginnote{\scriptsize ५.३.२३}\noindent प्रकारवचने थाल् ।

\marginnote{\scriptsize ५.३.२४}\noindent इदमस्थमुः ।

\marginnote{\scriptsize ५.३.२५}\noindent किमश्च ।

\marginnote{\scriptsize ५.३.२६}\noindent था हेतौ च च्छन्दसि ।

\marginnote{\scriptsize ५.३.२७}\noindent दिक्शब्देभ्यः सप्तमीपञ्चमीप्रथमाभ्यो दिग्देशकालेष्वस्तातिः ।

\marginnote{\scriptsize ५.३.२८}\noindent दक्षिणोत्तराभ्यामतसुच् ।

\marginnote{\scriptsize ५.३.२९}\noindent विभाषा परावराभ्याम् ।

\marginnote{\scriptsize ५.३.३०}\noindent अञ्चेर्लुक् ।

\marginnote{\scriptsize ५.३.३१}\noindent उपर्युपरिष्टात् ।

\marginnote{\scriptsize ५.३.३२}\noindent पश्चात् ।

\marginnote{\scriptsize ५.३.३३}\noindent पश्च पश्चा च च्छन्दसि ।

\marginnote{\scriptsize ५.३.३४}\noindent उत्तराधरदक्षिणादातिः ।

\marginnote{\scriptsize ५.३.३५}\noindent एनबन्यतरस्यामदूरेऽपञ्चम्याः ।

\marginnote{\scriptsize ५.३.३६}\noindent दक्षिणादाच् ।

\marginnote{\scriptsize ५.३.३७}\noindent आहि च दूरे ।

\marginnote{\scriptsize ५.३.३८}\noindent उत्तराच्च ।

\marginnote{\scriptsize ५.३.३९}\noindent पूर्वाधरावराणामसि पुरधवश्चैषाम् ।

\marginnote{\scriptsize ५.३.४०}\noindent अस्ताति च ।

\marginnote{\scriptsize ५.३.४१}\noindent विभाषाऽवरस्य ।

\marginnote{\scriptsize ५.३.४२}\noindent संख्याया विधाऽर्थे धा ।

\marginnote{\scriptsize ५.३.४३}\noindent अधिकरणविचाले च ।

\marginnote{\scriptsize ५.३.४४}\noindent एकाद्धो ध्यमुञन्यारयाम् ।

\marginnote{\scriptsize ५.३.४५}\noindent द्वित्र्योश्च धमुञ् ।

\marginnote{\scriptsize ५.३.४६}\noindent एधाच्च ।

\marginnote{\scriptsize ५.३.४७}\noindent याप्ये पाशप् ।

\marginnote{\scriptsize ५.३.४८}\noindent पूरणाद्भागे तीयादन् ।

\marginnote{\scriptsize ५.३.४९}\noindent प्रागेकादशभ्योऽच्छन्दसि ।

\marginnote{\scriptsize ५.३.५०}\noindent षष्ठाष्टमाभ्यां ञ च ।

\marginnote{\scriptsize ५.३.५१}\noindent मानपश्वङ्गयोः कन्लुकौ च ।

\marginnote{\scriptsize ५.३.५२}\noindent एकादाकिनिच्चासहाये ।

\marginnote{\scriptsize ५.३.५३}\noindent भूतपूर्वे चरट् ।

\marginnote{\scriptsize ५.३.५४}\noindent षष्ठ्या रूप्य च ।

\marginnote{\scriptsize ५.३.५५}\noindent अतिशायने तमबिष्ठनौ ।

\marginnote{\scriptsize ५.३.५६}\noindent तिङश्च ।

\marginnote{\scriptsize ५.३.५७}\noindent द्विवचनविभज्योपपदे तरबीयसुनौ ।

\marginnote{\scriptsize ५.३.५८}\noindent अजादी गुणवचनादेव ।

\marginnote{\scriptsize ५.३.५९}\noindent तुश्छन्दसि ।

\marginnote{\scriptsize ५.३.६०}\noindent प्रशस्यस्य श्रः ।

\marginnote{\scriptsize ५.३.६१}\noindent ज्य च ।

\marginnote{\scriptsize ५.३.६२}\noindent वृद्धस्य च ।

\marginnote{\scriptsize ५.३.६३}\noindent अन्तिकबाढयोर्नेदसाधौ ।

\marginnote{\scriptsize ५.३.६४}\noindent युवाल्पयोः कनन्यतरस्याम् ।

\marginnote{\scriptsize ५.३.६५}\noindent विन्मतोर्लुक् ।

\marginnote{\scriptsize ५.३.६६}\noindent प्रशंसायां रूपप् ।

\marginnote{\scriptsize ५.३.६७}\noindent ईषदसमाप्तौ कल्पब्देश्यदेशीयरः ।

\marginnote{\scriptsize ५.३.६८}\noindent विभाषा सुपो बहुच् पुरस्तात्तु ।

\marginnote{\scriptsize ५.३.६९}\noindent प्रकारवचने जातीयर्।

\marginnote{\scriptsize ५.३.७०}\noindent प्रागिवात्कः ।

\marginnote{\scriptsize ५.३.७१}\noindent अव्ययसर्वनाम्नामकच् प्राक् टेः ।

\marginnote{\scriptsize ५.३.७२}\noindent कस्य च दः ।

\marginnote{\scriptsize ५.३.७३}\noindent अज्ञाते ।

\marginnote{\scriptsize ५.३.७४}\noindent कुत्सिते ।

\marginnote{\scriptsize ५.३.७५}\noindent संज्ञायां कन् ।

\marginnote{\scriptsize ५.३.७६}\noindent अनुकम्पायाम् ।

\marginnote{\scriptsize ५.३.७७}\noindent नीतौ च तद्युक्तात् ।

\marginnote{\scriptsize ५.३.७८}\noindent बह्वचो मनुष्यनाम्नष्ठज्वा ।

\marginnote{\scriptsize ५.३.७९}\noindent घनिलचौ च ।

\marginnote{\scriptsize ५.३.८०}\noindent प्राचामुपादेरडज्वुचौ च ।

\marginnote{\scriptsize ५.३.८१}\noindent जातिनाम्नः कन् ।

\marginnote{\scriptsize ५.३.८२}\noindent अजिनान्तस्योत्तरपदलोपश्च ।

\marginnote{\scriptsize ५.३.८३}\noindent ठाजादावूर्ध्वं द्वितीयादचः ।

\marginnote{\scriptsize ५.३.८४}\noindent शेवलसुपरिविशालवरुणार्यमादीनां तृतीयात् ।

\marginnote{\scriptsize ५.३.८५}\noindent अल्पे ।

\marginnote{\scriptsize ५.३.८६}\noindent ह्रस्वे ।

\marginnote{\scriptsize ५.३.८७}\noindent संज्ञायां कन् ।

\marginnote{\scriptsize ५.३.८८}\noindent कुटीशमीशुण्डाभ्यो रः ।

\marginnote{\scriptsize ५.३.८९}\noindent कुत्वा डुपच् ।

\marginnote{\scriptsize ५.३.९०}\noindent कासूगोणीभ्यां ष्टरच् ।

\marginnote{\scriptsize ५.३.९१}\noindent वत्सोक्षाश्वर्षभेभ्यश्च तनुत्वे ।

\marginnote{\scriptsize ५.३.९२}\noindent किंयत्तदो निर्द्धारणे द्वयोरेकस्य डतरच् ।

\marginnote{\scriptsize ५.३.९३}\noindent वा बहूनां जातिपरिप्रश्ने डतमच् ।

\marginnote{\scriptsize ५.३.९४}\noindent एकाच्च प्राचाम् ।

\marginnote{\scriptsize ५.३.९५}\noindent अवक्षेपणे कन् ।

\marginnote{\scriptsize ५.३.९६}\noindent इवे प्रतिकृतौ ।

\marginnote{\scriptsize ५.३.९७}\noindent संज्ञायां च ।

\marginnote{\scriptsize ५.३.९८}\noindent लुम्मनुष्ये ।

\marginnote{\scriptsize ५.३.९९}\noindent जीविकाऽर्थे चापण्ये ।

\marginnote{\scriptsize ५.३.१००}\noindent देवपथादिभ्यश्च ।

\marginnote{\scriptsize ५.३.१०१}\noindent वस्तेर्ढञ् ।

\marginnote{\scriptsize ५.३.१०२}\noindent शिलाया ढः ।

\marginnote{\scriptsize ५.३.१०३}\noindent शाखाऽऽदिभ्यो यत् ।

\marginnote{\scriptsize ५.३.१०४}\noindent द्रव्यं च भव्ये ।

\marginnote{\scriptsize ५.३.१०५}\noindent कुशाग्राच्छः ।

\marginnote{\scriptsize ५.३.१०६}\noindent समासाच्च तद्विषयात् ।

\marginnote{\scriptsize ५.३.१०७}\noindent शर्कराऽऽदिभ्योऽण् ।

\marginnote{\scriptsize ५.३.१०८}\noindent अङ्गुल्यादिभ्यष्ठक् ।

\marginnote{\scriptsize ५.३.१०९}\noindent एकशालायाष्ठजन्यतरस्याम् ।

\marginnote{\scriptsize ५.३.११०}\noindent कर्कलोहितादीकक् ।

\marginnote{\scriptsize ५.३.१११}\noindent प्रत्नपूर्वविश्वेमात्थाल् छन्दसि ।

\marginnote{\scriptsize ५.३.११२}\noindent पूगाञ्ञ्योऽग्रामणीपूर्वात् ।

\marginnote{\scriptsize ५.३.११३}\noindent व्रातच्फञोरस्त्रियाम् ।

\marginnote{\scriptsize ५.३.११४}\noindent आयुधजीविसंघाञ्ञ्यड्वाहीकेष्वब्राह्मणराजन्यात् ।

\marginnote{\scriptsize ५.३.११५}\noindent वृकाट्टेण्यण् ।

\marginnote{\scriptsize ५.३.११६}\noindent दामन्यादित्रिगर्तषष्ठाच्छः ।

\marginnote{\scriptsize ५.३.११७}\noindent पर्श्वादियौधेयादिभ्यामणञौ । ५.३.११८ अभिजिद्विदभृच्छालावच्छिखावच्छमीवदूर्णावच्छ्रुमदणो यञ् ।

\marginnote{\scriptsize ५.३.११९}\noindent ञ्यादयस्तद्राजाः ।

\marginnote{\scriptsize ५.४.१}\noindent पादशतस्य संख्याऽऽदेर्वीप्सायां वुन् लोपश्च ।

\marginnote{\scriptsize ५.४.२}\noindent दण्डव्यवसर्गयोश्च ।

\marginnote{\scriptsize ५.४.३}\noindent स्थूलादिभ्यः प्रकारवचने कन् ।

\marginnote{\scriptsize ५.४.४}\noindent अनत्यन्तगतौ क्तात् ।

\marginnote{\scriptsize ५.४.५}\noindent न सामिवचने ।

\marginnote{\scriptsize ५.४.६}\noindent बृहत्या आच्छादने ।

\marginnote{\scriptsize ५.४.७}\noindent अषडक्षाशितङ्ग्वलंकर्मालम्पुरुषाध्युत्तरपदात् खः ।

\marginnote{\scriptsize ५.४.८}\noindent विभाषा अञ्चेरदिक्स्त्रियाम् ।

\marginnote{\scriptsize ५.४.९}\noindent जात्यन्ताच्छ बन्धुनि ।

\marginnote{\scriptsize ५.४.१०}\noindent स्थानान्ताद्विभाषा सस्थानेनेति चेत् ।

\marginnote{\scriptsize ५.४.११}\noindent किमेत्तिङव्ययघादाम्वद्रव्यप्रकर्षे ।

\marginnote{\scriptsize ५.४.१२}\noindent अमु च च्छन्दसि ।

\marginnote{\scriptsize ५.४.१३}\noindent अनुगादिनष्ठक् ।

\marginnote{\scriptsize ५.४.१४}\noindent णचः स्त्रियामञ् ।

\marginnote{\scriptsize ५.४.१५}\noindent अणिनुणः ।

\marginnote{\scriptsize ५.४.१६}\noindent विसारिणो मत्स्ये ।

\marginnote{\scriptsize ५.४.१७}\noindent संख्यायाः क्रियाऽभ्यावृत्तिगणने कृत्वसुच् ।

\marginnote{\scriptsize ५.४.१८}\noindent द्वित्रिचतुर्भ्यः सुच् ।

\marginnote{\scriptsize ५.४.१९}\noindent एकस्य सकृच्च ।

\marginnote{\scriptsize ५.४.२०}\noindent विभाषा बहोर्धाऽविप्रकृष्टकाले ।

\marginnote{\scriptsize ५.४.२१}\noindent तत्प्रकृतवचने मयट् ।

\marginnote{\scriptsize ५.४.२२}\noindent समूहवच्च बहुषु ।

\marginnote{\scriptsize ५.४.२३}\noindent अनन्तावसथेतिहभेषजाञ्ञ्यः ।

\marginnote{\scriptsize ५.४.२४}\noindent देवतान्तात्तादर्थ्ये यत् ।

\marginnote{\scriptsize ५.४.२५}\noindent पादार्घाभ्यां च ।

\marginnote{\scriptsize ५.४.२६}\noindent अतिथेर्ञ्यः ।

\marginnote{\scriptsize ५.४.२७}\noindent देवात्तल् ।

\marginnote{\scriptsize ५.४.२८}\noindent अवेः कः ।

\marginnote{\scriptsize ५.४.२९}\noindent यावादिभ्यः कन् ।

\marginnote{\scriptsize ५.४.३०}\noindent लोहितान्मणौ ।

\marginnote{\scriptsize ५.४.३१}\noindent वर्णे चानित्ये ।

\marginnote{\scriptsize ५.४.३२}\noindent रक्ते ।

\marginnote{\scriptsize ५.४.३३}\noindent कालाच्च ।

\marginnote{\scriptsize ५.४.३४}\noindent विनयादिभ्यष्ठक् ।

\marginnote{\scriptsize ५.४.३५}\noindent वाचो व्याहृतार्थायाम् ।

\marginnote{\scriptsize ५.४.३६}\noindent तद्युक्तात् कर्मणोऽण् ।

\marginnote{\scriptsize ५.४.३७}\noindent ओषधेरजातौ ।

\marginnote{\scriptsize ५.४.३८}\noindent प्रज्ञादिभ्यश्च ।

\marginnote{\scriptsize ५.४.३९}\noindent मृदस्तिकन् ।

\marginnote{\scriptsize ५.४.४०}\noindent सस्नौ प्रशंसायाम् ।

\marginnote{\scriptsize ५.४.४१}\noindent वृकज्येष्ठाभ्यां तिल्तातिलौ च च्छन्दसि ।

\marginnote{\scriptsize ५.४.४२}\noindent बह्वल्पार्थाच्छस् कारकादन्यतरस्याम् ।

\marginnote{\scriptsize ५.४.४३}\noindent संख्यैकवचनाच्च वीप्सायाम् ।

\marginnote{\scriptsize ५.४.४४}\noindent प्रतियोगे पञ्चम्यास्तसिः ।

\marginnote{\scriptsize ५.४.४५}\noindent अपादाने चाहीयरुहोः ।

\marginnote{\scriptsize ५.४.४६}\noindent अतिग्रहाव्यथनक्षेपेष्वकर्तरि तृतीयायाः ।

\marginnote{\scriptsize ५.४.४७}\noindent हीयमानपापयोगाच्च ।

\marginnote{\scriptsize ५.४.४८}\noindent षष्ठ्या व्याश्रये ।

\marginnote{\scriptsize ५.४.४९}\noindent रोगाच्चापनयने ।

\marginnote{\scriptsize ५.४.५०}\noindent अभूततद्भावे कृभ्वस्तियोगे सम्पद्यकर्तरि च्विः ।

\marginnote{\scriptsize ५.४.५१}\noindent अरुर्मनश्चक्षुश्चेतोरहोरजसां लोपश्च ।

\marginnote{\scriptsize ५.४.५२}\noindent विभाषा साति कार्त्स्न्ये ।

\marginnote{\scriptsize ५.४.५३}\noindent अभिविधौ सम्पदा च ।

\marginnote{\scriptsize ५.४.५४}\noindent तदधीनवचने ।

\marginnote{\scriptsize ५.४.५५}\noindent देये त्रा च ।

\marginnote{\scriptsize ५.४.५६}\noindent देवमनुष्यपुरुषमर्त्येभ्यो द्वितीयासप्तम्योर्बहुलम् ।

\marginnote{\scriptsize ५.४.५७}\noindent अव्यक्तानुकरणाद्द्व्यजवरार्धादनितौ डाच् ।

\marginnote{\scriptsize ५.४.५८}\noindent कृञो द्वितीयतृतीयशम्बबीजात् कृषौ ।

\marginnote{\scriptsize ५.४.५९}\noindent संख्यायाश्च गुणान्तायाः ।

\marginnote{\scriptsize ५.४.६०}\noindent समयाच्च यापनायाम् ।

\marginnote{\scriptsize ५.४.६१}\noindent सपत्त्रनिष्पत्रादतिव्यथने ।

\marginnote{\scriptsize ५.४.६२}\noindent निष्कुलान्निष्कोषणे ।

\marginnote{\scriptsize ५.४.६३}\noindent सुखप्रियादानुलोम्ये ।

\marginnote{\scriptsize ५.४.६४}\noindent दुःखात् प्रातिलोम्ये ।

\marginnote{\scriptsize ५.४.६५}\noindent शूलात् पाके ।

\marginnote{\scriptsize ५.४.६६}\noindent सत्यादशपथे ।

\marginnote{\scriptsize ५.४.६७}\noindent मद्रात् परिवापणे ।

\marginnote{\scriptsize ५.४.६८}\noindent समासान्ताः ।

\marginnote{\scriptsize ५.४.६९}\noindent न पूजनात् ।

\marginnote{\scriptsize ५.४.७०}\noindent किमः क्षेपे ।

\marginnote{\scriptsize ५.४.७१}\noindent नञस्तत्पुरुषात् ।

\marginnote{\scriptsize ५.४.७२}\noindent पथो विभाषा ।

\marginnote{\scriptsize ५.४.७३}\noindent बहुव्रीहौ संख्येये डजबहुगणात् ।

\marginnote{\scriptsize ५.४.७४}\noindent ऋक्पूरप्धूःपथामानक्षे ।

\marginnote{\scriptsize ५.४.७५}\noindent अच् प्रत्यन्ववपूर्वात् सामलोम्नः ।

\marginnote{\scriptsize ५.४.७६}\noindent अक्ष्णोऽदर्शनात् ।

\marginnote{\scriptsize ५.४.७७}\noindent अचतुरविचतुरसुचतुरस्त्रीपुंसधेन्वनडुहर्क्साम- वाङ्मनसाक्षिभ्रुवदारगवोर्वष्ठीवपदष्ठीवनक्तंदिव- रत्रिंदिवाहर्दिवसरजसनिःश्रेयसपुरुषायुषद्व्यायुष- त्र्यायुषर्ग्यजुषजातोक्षमहोक्षवृद्धोक्षोपशुनगोष्ठश्वाः ।

\marginnote{\scriptsize ५.४.७८}\noindent ब्रह्महस्तिभ्याम् वर्च्चसः ।

\marginnote{\scriptsize ५.४.७९}\noindent अवसमन्धेभ्यस्तमसः ।

\marginnote{\scriptsize ५.४.८०}\noindent श्वसो वसीयःश्रेयसः ।

\marginnote{\scriptsize ५.४.८१}\noindent अन्ववतप्ताद्रहसः ।

\marginnote{\scriptsize ५.४.८२}\noindent प्रतेरुरसः सप्तमीस्थात् ।

\marginnote{\scriptsize ५.४.८३}\noindent अनुगवमायामे ।

\marginnote{\scriptsize ५.४.८४}\noindent द्विस्तावा त्रिस्तावा वेदिः ।

\marginnote{\scriptsize ५.४.८५}\noindent उपसर्गादध्वनः ।

\marginnote{\scriptsize ५.४.८६}\noindent तत्पुरुषस्याङ्गुलेः संख्याऽव्ययादेः ।

\marginnote{\scriptsize ५.४.८७}\noindent अहस्सर्वैकदेशसंख्यातपुण्याच्च रात्रेः ।

\marginnote{\scriptsize ५.४.८८}\noindent अह्नोऽह्न एतेभ्यः ।

\marginnote{\scriptsize ५.४.८९}\noindent न संख्याऽऽदेः समाहारे ।

\marginnote{\scriptsize ५.४.९०}\noindent उत्तमैकाभ्यां च ।

\marginnote{\scriptsize ५.४.९१}\noindent राजाऽहस्सखिभ्यष्टच् ।

\marginnote{\scriptsize ५.४.९२}\noindent गोरतद्धितलुकि ।

\marginnote{\scriptsize ५.४.९३}\noindent अग्राख्यायामुरसः ।

\marginnote{\scriptsize ५.४.९४}\noindent अनोऽश्मायस्सरसाम् जातिसंज्ञयोः ।

\marginnote{\scriptsize ५.४.९५}\noindent ग्रामकौटाभ्यां च तक्ष्णः ।

\marginnote{\scriptsize ५.४.९६}\noindent अतेः शुनः ।

\marginnote{\scriptsize ५.४.९७}\noindent उपमानादप्राणिषु ।

\marginnote{\scriptsize ५.४.९८}\noindent उत्तरमृगपूर्वाच्च सक्थ्नः ।

\marginnote{\scriptsize ५.४.९९}\noindent नावो द्विगोः ।

\marginnote{\scriptsize ५.४.१००}\noindent अर्धाच्च ।

\marginnote{\scriptsize ५.४.१०१}\noindent खार्याः प्राचाम् ।

\marginnote{\scriptsize ५.४.१०२}\noindent द्वित्रिभ्यामञ्जलेः ।

\marginnote{\scriptsize ५.४.१०३}\noindent अनसन्तान्नपुंसकाच्छन्दसि ।

\marginnote{\scriptsize ५.४.१०४}\noindent ब्रह्मणो जानपदाख्यायाम् ।

\marginnote{\scriptsize ५.४.१०५}\noindent कुमहद्भ्यामन्यतरस्याम् ।

\marginnote{\scriptsize ५.४.१०६}\noindent द्वंद्वाच्चुदषहान्तात् समाहारे ।

\marginnote{\scriptsize ५.४.१०७}\noindent अव्ययीभावे शरत्प्रभृतिभ्यः ।

\marginnote{\scriptsize ५.४.१०८}\noindent अनश्च ।

\marginnote{\scriptsize ५.४.१०९}\noindent नपुंसकादन्यतरस्याम् ।

\marginnote{\scriptsize ५.४.११०}\noindent नदीपौर्णमास्याग्रहायणीभ्यः ।

\marginnote{\scriptsize ५.४.१११}\noindent झयः ।

\marginnote{\scriptsize ५.४.११२}\noindent गिरेश्च सेनकस्य ।

\marginnote{\scriptsize ५.४.११३}\noindent बहुव्रीहौ सक्थ्यक्ष्णोः स्वाङ्गात् षच् ।

\marginnote{\scriptsize ५.४.११४}\noindent अङ्गुलेर्दारुणि ।

\marginnote{\scriptsize ५.४.११५}\noindent द्वित्रिभ्यां ष मूर्ध्नः ।

\marginnote{\scriptsize ५.४.११६}\noindent अप् पूरणीप्रमाण्योः ।

\marginnote{\scriptsize ५.४.११७}\noindent अन्तर्बहिर्भ्यां च लोम्नः ।

\marginnote{\scriptsize ५.४.११८}\noindent अञ्नासिकायाः संज्ञायां नसं चास्थूलात् ।

\marginnote{\scriptsize ५.४.११९}\noindent उपसर्गाच्च ।

\marginnote{\scriptsize ५.४.१२०}\noindent सुप्रातसुश्वसुदिवशारिकुक्षचतुरश्रैणीपदाजपद- प्रोष्ठपदाः ।

\marginnote{\scriptsize ५.४.१२१}\noindent नञ्दुःसुभ्यो हलिसक्थ्योरन्यतरस्याम् ।

\marginnote{\scriptsize ५.४.१२२}\noindent नित्यमसिच् प्रजामेधयोः ।

\marginnote{\scriptsize ५.४.१२३}\noindent बहुप्रजाश्छन्दसि ।

\marginnote{\scriptsize ५.४.१२४}\noindent धर्मादनिच् केवलात् ।

\marginnote{\scriptsize ५.४.१२५}\noindent जम्भा सुहरिततृणसोमेभ्यः ।

\marginnote{\scriptsize ५.४.१२६}\noindent दक्षिणेर्मा लुब्धयोगे ।

\marginnote{\scriptsize ५.४.१२७}\noindent इच् कर्मव्यतिहारे ।

\marginnote{\scriptsize ५.४.१२८}\noindent द्विदण्ड्यादिभ्यश्च ।

\marginnote{\scriptsize ५.४.१२९}\noindent प्रसम्भ्यां जानुनोर्ज्ञुः ।

\marginnote{\scriptsize ५.४.१३०}\noindent ऊर्ध्वाद्विभाषा ।

\marginnote{\scriptsize ५.४.१३१}\noindent ऊधसोऽनङ् ।

\marginnote{\scriptsize ५.४.१३२}\noindent धनुषश्च ।

\marginnote{\scriptsize ५.४.१३३}\noindent वा संज्ञायाम् ।

\marginnote{\scriptsize ५.४.१३४}\noindent जायाया निङ् ।

\marginnote{\scriptsize ५.४.१३५}\noindent गन्धस्येदुत्पूतिसुसुरभिभ्यः ।

\marginnote{\scriptsize ५.४.१३६}\noindent अल्पाख्यायाम् ।

\marginnote{\scriptsize ५.४.१३७}\noindent उपमानाच्च ।

\marginnote{\scriptsize ५.४.१३८}\noindent पादस्य लोपोऽहस्त्यादिभ्यः ।

\marginnote{\scriptsize ५.४.१३९}\noindent कुम्भपदीषु च ।

\marginnote{\scriptsize ५.४.१४०}\noindent संख्यासुपूर्वस्य ।

\marginnote{\scriptsize ५.४.१४१}\noindent वयसि दन्तस्य दतृ ।

\marginnote{\scriptsize ५.४.१४२}\noindent छन्दसि च ।

\marginnote{\scriptsize ५.४.१४३}\noindent स्त्रियां संज्ञायाम् ।

\marginnote{\scriptsize ५.४.१४४}\noindent विभाषा श्यावारोकाभ्याम् ।

\marginnote{\scriptsize ५.४.१४५}\noindent अग्रान्तशुद्धशुभ्रवृषवराहेभ्यश्च ।

\marginnote{\scriptsize ५.४.१४६}\noindent ककुदस्यावस्थायां लोपः ।

\marginnote{\scriptsize ५.४.१४७}\noindent त्रिककुत् पर्वते ।

\marginnote{\scriptsize ५.४.१४८}\noindent उद्विभ्यां काकुदस्य ।

\marginnote{\scriptsize ५.४.१४९}\noindent पूर्णाद्विभाषा ।

\marginnote{\scriptsize ५.४.१५०}\noindent सुहृद्दुर्हृदौ मित्रामित्रयोः ।

\marginnote{\scriptsize ५.४.१५१}\noindent उरःप्रभृतिभ्यः कप् ।

\marginnote{\scriptsize ५.४.१५२}\noindent इनः स्त्रियाम् ।

\marginnote{\scriptsize ५.४.१५३}\noindent नद्यृतश्च ।

\marginnote{\scriptsize ५.४.१५४}\noindent शेषाद्विभाषा ।

\marginnote{\scriptsize ५.४.१५५}\noindent न संज्ञायाम् ।

\marginnote{\scriptsize ५.४.१५६}\noindent ईयसश्च ।

\marginnote{\scriptsize ५.४.१५७}\noindent वन्दिते भ्रातुः ।

\marginnote{\scriptsize ५.४.१५८}\noindent ऋतश्छन्दसि ।

\marginnote{\scriptsize ५.४.१५९}\noindent नाडीतन्त्र्योः स्वाङ्गे ।

\marginnote{\scriptsize ५.४.१६०}\noindent निष्प्रवाणिश्च ।

\section*{॥ अध्याय ६ ॥}

\marginnote{\scriptsize ६.१.१}\noindent एकाचो द्वे प्रथमस्य ।

\marginnote{\scriptsize ६.१.२}\noindent अजादेर्द्वितीयस्य ।

\marginnote{\scriptsize ६.१.३}\noindent न न्द्राः संयोगादयः ।

\marginnote{\scriptsize ६.१.४}\noindent पूर्वोऽभ्यासः ।

\marginnote{\scriptsize ६.१.५}\noindent उभे अभ्यस्तम् ।

\marginnote{\scriptsize ६.१.६}\noindent जक्षित्यादयः षट् ।

\marginnote{\scriptsize ६.१.७}\noindent तुजादीनां दीर्घोऽभ्यासस्य ।

\marginnote{\scriptsize ६.१.८}\noindent लिटि धातोरनभ्यासस्य ।

\marginnote{\scriptsize ६.१.९}\noindent सन्यङोः ।

\marginnote{\scriptsize ६.१.१०}\noindent श्लौ ।

\marginnote{\scriptsize ६.१.११}\noindent चङि ।

\marginnote{\scriptsize ६.१.१२}\noindent दाश्वान् साह्वान् मीढ्वांश्च ।

\marginnote{\scriptsize ६.१.१३}\noindent ष्यङः सम्प्रसारणं पुत्रपत्योस्तत्पुरुषे ।

\marginnote{\scriptsize ६.१.१४}\noindent बन्धुनि बहुव्रीहौ ।

\marginnote{\scriptsize ६.१.१५}\noindent वचिस्वपियजादिनां किति ।

\marginnote{\scriptsize ६.१.१६}\noindent ग्रहिज्यावयिव्यधिवष्टिविचतिवृश्चतिपृच्छति- भृज्जतीनां ङिति च ।

\marginnote{\scriptsize ६.१.१७}\noindent लिट्यभ्यासस्योभयेषाम् ।

\marginnote{\scriptsize ६.१.१८}\noindent स्वापेश्चङि ।

\marginnote{\scriptsize ६.१.१९}\noindent स्वपिस्यमिव्येञां यङि ।

\marginnote{\scriptsize ६.१.२०}\noindent न वशः ।

\marginnote{\scriptsize ६.१.२१}\noindent चायः की ।

\marginnote{\scriptsize ६.१.२२}\noindent स्फायः स्फी निष्ठायाम् ।

\marginnote{\scriptsize ६.१.२३}\noindent स्त्यः प्रपूर्वस्य ।

\marginnote{\scriptsize ६.१.२४}\noindent द्रवमूर्तिस्पर्शयोः श्यः ।

\marginnote{\scriptsize ६.१.२५}\noindent प्रतेश्च ।

\marginnote{\scriptsize ६.१.२६}\noindent विभाषाऽभ्यवपूर्वस्य ।

\marginnote{\scriptsize ६.१.२७}\noindent श‍ृतं पाके ।

\marginnote{\scriptsize ६.१.२८}\noindent प्यायः पी ।

\marginnote{\scriptsize ६.१.२९}\noindent लिड्यङोश्च ।

\marginnote{\scriptsize ६.१.३०}\noindent विभाषा श्वेः ।

\marginnote{\scriptsize ६.१.३१}\noindent णौ च संश्चङोः ।

\marginnote{\scriptsize ६.१.३२}\noindent ह्वः सम्प्रसारणम् ।

\marginnote{\scriptsize ६.१.३३}\noindent अभ्यस्तस्य च ।

\marginnote{\scriptsize ६.१.३४}\noindent बहुलं छन्दसि ।

\marginnote{\scriptsize ६.१.३५}\noindent चायः की ।

\marginnote{\scriptsize ६.१.३६}\noindent अपस्पृधेथामानृचुरानृहुश्चिच्युषेतित्याज- श्राताःश्रितमाशीराशीर्त्तः ।

\marginnote{\scriptsize ६.१.३७}\noindent न सम्प्रसारणे सम्प्रसारणम् ।

\marginnote{\scriptsize ६.१.३८}\noindent लिटि वयो यः ।

\marginnote{\scriptsize ६.१.३९}\noindent वश्चास्यान्यतरस्याम् किति ।

\marginnote{\scriptsize ६.१.४०}\noindent वेञः ।

\marginnote{\scriptsize ६.१.४१}\noindent ल्यपि च ।

\marginnote{\scriptsize ६.१.४२}\noindent ज्यश्च ।

\marginnote{\scriptsize ६.१.४३}\noindent व्यश्च ।

\marginnote{\scriptsize ६.१.४४}\noindent विभाषा परेः ।

\marginnote{\scriptsize ६.१.४५}\noindent आदेच उपदेशेऽशिति ।

\marginnote{\scriptsize ६.१.४६}\noindent न व्यो लिटि ।

\marginnote{\scriptsize ६.१.४७}\noindent स्फुरतिस्फुलत्योर्घञि ।

\marginnote{\scriptsize ६.१.४८}\noindent क्रीङ्जीनां णौ ।

\marginnote{\scriptsize ६.१.४९}\noindent सिध्यतेरपारलौकिके ।

\marginnote{\scriptsize ६.१.५०}\noindent मीनातिमिनोतिदीङां ल्यपि च ।

\marginnote{\scriptsize ६.१.५१}\noindent विभाषा लीयतेः ।

\marginnote{\scriptsize ६.१.५२}\noindent खिदेश्छन्दसि ।

\marginnote{\scriptsize ६.१.५३}\noindent अपगुरो णमुलि ।

\marginnote{\scriptsize ६.१.५४}\noindent चिस्फुरोर्णौ ।

\marginnote{\scriptsize ६.१.५५}\noindent प्रजने वीयतेः ।

\marginnote{\scriptsize ६.१.५६}\noindent बिभेतेर्हेतुभये ।

\marginnote{\scriptsize ६.१.५७}\noindent नित्यं स्मयतेः ।

\marginnote{\scriptsize ६.१.५८}\noindent सृजिदृशोर्झल्यमकिति ।

\marginnote{\scriptsize ६.१.५९}\noindent अनुदात्तस्य चर्दुपधस्यान्यतरस्याम् ।

\marginnote{\scriptsize ६.१.६०}\noindent शीर्षंश्छन्दसि ।

\marginnote{\scriptsize ६.१.६१}\noindent ये च तद्धिते ।

\marginnote{\scriptsize ६.१.६२}\noindent अचि शीर्षः । ६.१.६३ पद्दन्नोमाश‍ृन्निशसन्यूषन्दोषन्यकञ्छकनुदन्नासञ्छस्- प्रभृतिषु ।

\marginnote{\scriptsize ६.१.६४}\noindent धात्वादेः षः सः ।

\marginnote{\scriptsize ६.१.६५}\noindent णो नः ।

\marginnote{\scriptsize ६.१.६६}\noindent लोपो व्योर्वलि ।

\marginnote{\scriptsize ६.१.६७}\noindent वेरपृक्तस्य ।

\marginnote{\scriptsize ६.१.६८}\noindent हल्ङ्याब्भ्यो दीर्घात् सुतिस्यपृक्तं हल् ।

\marginnote{\scriptsize ६.१.६९}\noindent एङ्ह्रस्वात् सम्बुद्धेः ।

\marginnote{\scriptsize ६.१.७०}\noindent शेश्छन्दसि बहुलम् ।

\marginnote{\scriptsize ६.१.७१}\noindent ह्रस्वस्य पिति कृति तुक् ।

\marginnote{\scriptsize ६.१.७२}\noindent संहितायाम् ।

\marginnote{\scriptsize ६.१.७३}\noindent छे च ।

\marginnote{\scriptsize ६.१.७४}\noindent आङ्माङोश्च ।

\marginnote{\scriptsize ६.१.७५}\noindent दीर्घात् ।

\marginnote{\scriptsize ६.१.७६}\noindent पदान्ताद्वा ।

\marginnote{\scriptsize ६.१.७७}\noindent इको यणचि ।

\marginnote{\scriptsize ६.१.७८}\noindent एचोऽयवायावः ।

\marginnote{\scriptsize ६.१.७९}\noindent वान्तो यि प्रत्यये ।

\marginnote{\scriptsize ६.१.८०}\noindent धातोस्तन्निमित्तस्यैव ।

\marginnote{\scriptsize ६.१.८१}\noindent क्षय्यजय्यौ शक्यार्थे ।

\marginnote{\scriptsize ६.१.८२}\noindent क्रय्यस्तदर्थे ।

\marginnote{\scriptsize ६.१.८३}\noindent भय्यप्रवय्ये च च्छन्दसि ।

\marginnote{\scriptsize ६.१.८४}\noindent एकः पूर्वपरयोः ।

\marginnote{\scriptsize ६.१.८५}\noindent अन्तादिवच्च ।

\marginnote{\scriptsize ६.१.८६}\noindent षत्वतुकोरसिद्धः ।

\marginnote{\scriptsize ६.१.८७}\noindent आद्गुणः ।

\marginnote{\scriptsize ६.१.८८}\noindent वृद्धिरेचि ।

\marginnote{\scriptsize ६.१.८९}\noindent एत्येधत्यूठ्सु ।

\marginnote{\scriptsize ६.१.९०}\noindent आटश्च ।

\marginnote{\scriptsize ६.१.९१}\noindent उपसर्गादृति धातौ ।

\marginnote{\scriptsize ६.१.९२}\noindent वा सुप्यापिशलेः ।

\marginnote{\scriptsize ६.१.९३}\noindent ओतोऽम्शसोः ।

\marginnote{\scriptsize ६.१.९४}\noindent एङि पररूपम् ।

\marginnote{\scriptsize ६.१.९५}\noindent ओमाङोश्च ।

\marginnote{\scriptsize ६.१.९६}\noindent उस्यपदान्तात् ।

\marginnote{\scriptsize ६.१.९७}\noindent अतो गुणे ।

\marginnote{\scriptsize ६.१.९८}\noindent अव्यक्तानुकरणस्यात इतौ ।

\marginnote{\scriptsize ६.१.९९}\noindent नाम्रेडितस्यान्त्यस्य तु वा ।

\marginnote{\scriptsize ६.१.१००}\noindent नित्यमाम्रेडिते डाचि ।

\marginnote{\scriptsize ६.१.१०१}\noindent अकः सवर्णे दीर्घः ।

\marginnote{\scriptsize ६.१.१०२}\noindent प्रथमयोः पूर्वसवर्णः ।

\marginnote{\scriptsize ६.१.१०३}\noindent तस्माच्छसो नः पुंसि ।

\marginnote{\scriptsize ६.१.१०४}\noindent नादिचि ।

\marginnote{\scriptsize ६.१.१०५}\noindent दीर्घाज्जसि च ।

\marginnote{\scriptsize ६.१.१०६}\noindent वा छन्दसि ।

\marginnote{\scriptsize ६.१.१०७}\noindent अमि पूर्वः ।

\marginnote{\scriptsize ६.१.१०८}\noindent सम्प्रसारणाच्च ।

\marginnote{\scriptsize ६.१.१०९}\noindent एङः पदान्तादति ।

\marginnote{\scriptsize ६.१.११०}\noindent ङसिङसोश्च ।

\marginnote{\scriptsize ६.१.१११}\noindent ऋत उत् ।

\marginnote{\scriptsize ६.१.११२}\noindent ख्यत्यात् परस्य ।

\marginnote{\scriptsize ६.१.११३}\noindent अतो रोरप्लुतादप्लुते ।

\marginnote{\scriptsize ६.१.११४}\noindent हशि च ।

\marginnote{\scriptsize ६.१.११५}\noindent प्रकृत्याऽन्तःपादमव्यपरे ।

\marginnote{\scriptsize ६.१.११६}\noindent अव्यादवद्यादवक्रमुरव्रतायमवन्त्ववस्युषु च ।

\marginnote{\scriptsize ६.१.११७}\noindent यजुष्युरः ।

\marginnote{\scriptsize ६.१.११८}\noindent आपोजुषाणोवृष्णोवर्षिष्ठेऽम्बेऽम्बालेऽम्बिकेपूर्वे ।

\marginnote{\scriptsize ६.१.११९}\noindent अङ्ग इत्यादौ च ।

\marginnote{\scriptsize ६.१.१२०}\noindent अनुदात्ते च कुधपरे ।

\marginnote{\scriptsize ६.१.१२१}\noindent अवपथासि च ।

\marginnote{\scriptsize ६.१.१२२}\noindent सर्वत्र विभाषा गोः ।

\marginnote{\scriptsize ६.१.१२३}\noindent अवङ् स्फोटायनस्य ।

\marginnote{\scriptsize ६.१.१२४}\noindent इन्द्रे च (नित्यम्) ।

\marginnote{\scriptsize ६.१.१२५}\noindent प्लुतप्रगृह्या अचि नित्यम् ।

\marginnote{\scriptsize ६.१.१२६}\noindent आङोऽनुनासिकश्छन्दसि ।

\marginnote{\scriptsize ६.१.१२७}\noindent इकोऽसवर्णे शाकल्यस्य ह्रस्वश्च ।

\marginnote{\scriptsize ६.१.१२८}\noindent ऋत्यकः ।

\marginnote{\scriptsize ६.१.१२९}\noindent अप्लुतवदुपस्थिते ।

\marginnote{\scriptsize ६.१.१३०}\noindent ई३ चाक्रवर्मणस्य ।

\marginnote{\scriptsize ६.१.१३१}\noindent दिव उत् ।

\marginnote{\scriptsize ६.१.१३२}\noindent एतत्तदोः सुलोपोऽकोरनञ्समासे हलि ।

\marginnote{\scriptsize ६.१.१३३}\noindent स्यश्छन्दसि बहुलम् ।

\marginnote{\scriptsize ६.१.१३४}\noindent सोऽचि लोपे चेत् पादपूरणम् ।

\marginnote{\scriptsize ६.१.१३५}\noindent सुट् कात् पूर्वः ।

\marginnote{\scriptsize ६.१.१३६}\noindent अडभ्यासव्यवायेऽपि ।

\marginnote{\scriptsize ६.१.१३७}\noindent सम्पर्युपेभ्यः करोतौ भूषणे ।

\marginnote{\scriptsize ६.१.१३८}\noindent समवाये च ।

\marginnote{\scriptsize ६.१.१३९}\noindent उपात् प्रतियत्नवैकृतवाक्याध्याहारेषु ।

\marginnote{\scriptsize ६.१.१४०}\noindent किरतौ लवने ।

\marginnote{\scriptsize ६.१.१४१}\noindent हिंसायां प्रतेश्च ।

\marginnote{\scriptsize ६.१.१४२}\noindent अपाच्चतुष्पाच्छकुनिष्वालेखने ।

\marginnote{\scriptsize ६.१.१४३}\noindent कुस्तुम्बुरूणि जातिः ।

\marginnote{\scriptsize ६.१.१४४}\noindent अपरस्पराः क्रियासातत्ये ।

\marginnote{\scriptsize ६.१.१४५}\noindent गोष्पदं सेवितासेवितप्रमाणेषु ।

\marginnote{\scriptsize ६.१.१४६}\noindent आस्पदं प्रतिष्ठायाम् ।

\marginnote{\scriptsize ६.१.१४७}\noindent आश्चर्यमनित्ये ।

\marginnote{\scriptsize ६.१.१४८}\noindent वर्चस्केऽवस्करः ।

\marginnote{\scriptsize ६.१.१४९}\noindent अपस्करो रथाङ्गम् ।

\marginnote{\scriptsize ६.१.१५०}\noindent विष्किरः शकुनिर्विकरो वा ।

\marginnote{\scriptsize ६.१.१५१}\noindent ह्रस्वाच्चन्द्रोत्तरपदे मन्त्रे ।

\marginnote{\scriptsize ६.१.१५२}\noindent प्रतिष्कशश्च कशेः ।

\marginnote{\scriptsize ६.१.१५३}\noindent प्रस्कण्वहरिश्चन्द्रावृषी ।

\marginnote{\scriptsize ६.१.१५४}\noindent मस्करमस्करिणौ वेणुपरिव्राजकयोः ।

\marginnote{\scriptsize ६.१.१५५}\noindent कास्तीराजस्तुन्दे नगरे ।

\marginnote{\scriptsize ६.१.१५६}\noindent कारस्करो वृक्षः ।

\marginnote{\scriptsize ६.१.१५७}\noindent पारस्करप्रभृतीनि च संज्ञायाम् ।

\marginnote{\scriptsize ६.१.१५८}\noindent अनुदात्तं पदमेकवर्जम् ।

\marginnote{\scriptsize ६.१.१५९}\noindent कर्षात्वतो घञोऽन्त उदात्तः ।

\marginnote{\scriptsize ६.१.१६०}\noindent उञ्छादीनां च ।

\marginnote{\scriptsize ६.१.१६१}\noindent अनुदात्तस्य च यत्रोदात्तलोपः ।

\marginnote{\scriptsize ६.१.१६२}\noindent धातोः ।

\marginnote{\scriptsize ६.१.१६३}\noindent चितः ।

\marginnote{\scriptsize ६.१.१६४}\noindent तद्धितस्य ।

\marginnote{\scriptsize ६.१.१६५}\noindent कितः ।

\marginnote{\scriptsize ६.१.१६६}\noindent तिसृभ्यो जसः ।

\marginnote{\scriptsize ६.१.१६७}\noindent चतुरः शसि ।

\marginnote{\scriptsize ६.१.१६८}\noindent सावेकाचस्तृतीयाऽऽदिविभक्तिः ।

\marginnote{\scriptsize ६.१.१६९}\noindent अन्तोदत्तादुत्तरपदादन्यतरस्यामनित्यसमासे ।

\marginnote{\scriptsize ६.१.१७०}\noindent अञ्चेश्छन्दस्यसर्वनामस्थानम् ।

\marginnote{\scriptsize ६.१.१७१}\noindent ऊडिदम्पदाद्यप्पुम्रैद्युभ्यः ।

\marginnote{\scriptsize ६.१.१७२}\noindent अष्टनो दीर्घात् ।

\marginnote{\scriptsize ६.१.१७३}\noindent शतुरनुमो नद्यजादी ।

\marginnote{\scriptsize ६.१.१७४}\noindent उदात्तयणो हल्पूर्वात् ।

\marginnote{\scriptsize ६.१.१७५}\noindent नोङ्धात्वोः ।

\marginnote{\scriptsize ६.१.१७६}\noindent ह्रस्वनुड्भ्यां मतुप् ।

\marginnote{\scriptsize ६.१.१७७}\noindent नामन्यतरस्याम् ।

\marginnote{\scriptsize ६.१.१७८}\noindent ङ्याश्छन्दसि बहुलम् ।

\marginnote{\scriptsize ६.१.१७९}\noindent षट्त्रिचतुर्भ्यो हलादिः ।

\marginnote{\scriptsize ६.१.१८०}\noindent झल्युपोत्तमम् ।

\marginnote{\scriptsize ६.१.१८१}\noindent विभाषा भाषायाम् ।

\marginnote{\scriptsize ६.१.१८२}\noindent न गोश्वन्त्साववर्णराडङ्क्रुङ्कृद्भ्यः ।

\marginnote{\scriptsize ६.१.१८३}\noindent दिवो झल् ।

\marginnote{\scriptsize ६.१.१८४}\noindent नृ चान्यतरस्याम् ।

\marginnote{\scriptsize ६.१.१८५}\noindent तित्स्वरितम् ।

\marginnote{\scriptsize ६.१.१८६}\noindent तास्यनुदात्तेन्ङिददुपदेशाल्लसार्वधातुकम्- अनुदात्तमहन्विङोः ।

\marginnote{\scriptsize ६.१.१८७}\noindent आदिः सिचोऽन्यतरस्याम् ।

\marginnote{\scriptsize ६.१.१८८}\noindent स्वपादिर्हिंसामच्यनिटि ।

\marginnote{\scriptsize ६.१.१८९}\noindent अभ्यस्तानामादिः ।

\marginnote{\scriptsize ६.१.१९०}\noindent अनुदात्ते च ।

\marginnote{\scriptsize ६.१.१९१}\noindent सर्वस्य सुपि ।

\marginnote{\scriptsize ६.१.१९२}\noindent भीह्रीभृहुमदजनधनदरिद्राजागरां प्रत्ययात् पूर्वम् पिति ।

\marginnote{\scriptsize ६.१.१९३}\noindent लिति ।

\marginnote{\scriptsize ६.१.१९४}\noindent आदिर्णमुल्यन्यतरस्याम् ।

\marginnote{\scriptsize ६.१.१९५}\noindent अचः कर्तृयकि ।

\marginnote{\scriptsize ६.१.१९६}\noindent थलि च सेटीडन्तो वा ।

\marginnote{\scriptsize ६.१.१९७}\noindent ञ्णित्यादिर्नित्यम् ।

\marginnote{\scriptsize ६.१.१९८}\noindent आमन्त्रितस्य च ।

\marginnote{\scriptsize ६.१.१९९}\noindent पथिमथोः सर्वनामस्थाने ।

\marginnote{\scriptsize ६.१.२००}\noindent अन्तश्च तवै युगपत् ।

\marginnote{\scriptsize ६.१.२०१}\noindent क्षयो निवासे ।

\marginnote{\scriptsize ६.१.२०२}\noindent जयः करणम् ।

\marginnote{\scriptsize ६.१.२०३}\noindent वृषादीनां च ।

\marginnote{\scriptsize ६.१.२०४}\noindent संज्ञायामुपमानम् ।

\marginnote{\scriptsize ६.१.२०५}\noindent निष्ठा च द्व्यजनात् ।

\marginnote{\scriptsize ६.१.२०६}\noindent शुष्कधृष्टौ ।

\marginnote{\scriptsize ६.१.२०७}\noindent आशितः कर्ता ।

\marginnote{\scriptsize ६.१.२०८}\noindent रिक्ते विभाषा ।

\marginnote{\scriptsize ६.१.२०९}\noindent जुष्टार्पिते च छन्दसि ।

\marginnote{\scriptsize ६.१.२१०}\noindent नित्यं मन्त्रे ।

\marginnote{\scriptsize ६.१.२११}\noindent युष्मदस्मदोर्ङसि ।

\marginnote{\scriptsize ६.१.२१२}\noindent ङयि च ।

\marginnote{\scriptsize ६.१.२१३}\noindent यतोऽनावः ।

\marginnote{\scriptsize ६.१.२१४}\noindent ईडवन्दवृशंसदुहां ण्यतः ।

\marginnote{\scriptsize ६.१.२१५}\noindent विभाषा वेण्विन्धानयोः ।

\marginnote{\scriptsize ६.१.२१६}\noindent त्यागरागहासकुहश्वठक्रथानाम् ।

\marginnote{\scriptsize ६.१.२१७}\noindent उपोत्तमं रिति ।

\marginnote{\scriptsize ६.१.२१८}\noindent चङ्यन्यतरस्याम् ।

\marginnote{\scriptsize ६.१.२१९}\noindent मतोः पूर्वमात् संज्ञायां स्त्रियाम् ।

\marginnote{\scriptsize ६.१.२२०}\noindent अन्तोऽवत्याः ।

\marginnote{\scriptsize ६.१.२२१}\noindent ईवत्याः ।

\marginnote{\scriptsize ६.१.२२२}\noindent चौ ।

\marginnote{\scriptsize ६.१.२२३}\noindent समासस्य ।

\marginnote{\scriptsize ६.२.१}\noindent बहुव्रीहौ प्रकृत्या पूर्वपदम् ।

\marginnote{\scriptsize ६.२.२}\noindent तत्पुरुषे तुल्यार्थतृतीयासप्तम्युपमानाव्ययद्वितीयाकृत्याः ।

\marginnote{\scriptsize ६.२.३}\noindent वर्णः वर्णेष्वनेते ।

\marginnote{\scriptsize ६.२.४}\noindent गाधलवणयोः प्रमाणे ।

\marginnote{\scriptsize ६.२.५}\noindent दायाद्यं दायादे ।

\marginnote{\scriptsize ६.२.६}\noindent प्रतिबन्धि चिरकृच्छ्रयोः ।

\marginnote{\scriptsize ६.२.७}\noindent पदेऽपदेशे ।

\marginnote{\scriptsize ६.२.८}\noindent निवाते वातत्राणे ।

\marginnote{\scriptsize ६.२.९}\noindent शारदेअनार्तवे ।

\marginnote{\scriptsize ६.२.१०}\noindent अध्वर्युकषाययोर्जातौ ।

\marginnote{\scriptsize ६.२.११}\noindent सदृशप्रतिरूपयोः सादृश्ये ।

\marginnote{\scriptsize ६.२.१२}\noindent द्विगौ प्रमाणे ।

\marginnote{\scriptsize ६.२.१३}\noindent गन्तव्यपण्यं वाणिजे ।

\marginnote{\scriptsize ६.२.१४}\noindent मात्रोपज्ञोपक्रमच्छाये नपुंसके ।

\marginnote{\scriptsize ६.२.१५}\noindent सुखप्रिययोर्हिते ।

\marginnote{\scriptsize ६.२.१६}\noindent प्रीतौ च ।

\marginnote{\scriptsize ६.२.१७}\noindent स्वं स्वामिनि ।

\marginnote{\scriptsize ६.२.१८}\noindent पत्यावैश्वर्ये ।

\marginnote{\scriptsize ६.२.१९}\noindent न भूवाक्चिद्दिधिषु ।

\marginnote{\scriptsize ६.२.२०}\noindent वा भुवनम् ।

\marginnote{\scriptsize ६.२.२१}\noindent आशङ्काबाधनेदीयस्सु संभावने ।

\marginnote{\scriptsize ६.२.२२}\noindent पूर्वे भूतपूर्वे ।

\marginnote{\scriptsize ६.२.२३}\noindent सविधसनीडसमर्यादसवेशसदेशेषु सामीप्ये ।

\marginnote{\scriptsize ६.२.२४}\noindent विस्पष्टादीनि गुणवचनेषु ।

\marginnote{\scriptsize ६.२.२५}\noindent श्रज्याऽवमकन्पापवत्सु भावे कर्मधारये ।

\marginnote{\scriptsize ६.२.२६}\noindent कुमारश्च ।

\marginnote{\scriptsize ६.२.२७}\noindent आदिः प्रत्येनसि ।

\marginnote{\scriptsize ६.२.२८}\noindent पूगेष्वन्यतरस्याम् ।

\marginnote{\scriptsize ६.२.२९}\noindent इगन्तकालकपालभगालशरावेषु द्विगौ ।

\marginnote{\scriptsize ६.२.३०}\noindent बह्वन्यतरस्याम् ।

\marginnote{\scriptsize ६.२.३१}\noindent दिष्टिवितस्त्योश्च ।

\marginnote{\scriptsize ६.२.३२}\noindent सप्तमी सिद्धशुष्कपक्वबन्धेष्वकालात् ।

\marginnote{\scriptsize ६.२.३३}\noindent परिप्रत्युपापा वर्ज्यमानाहोरात्रावयवेषु ।

\marginnote{\scriptsize ६.२.३४}\noindent राजन्यबहुवचनद्वंद्वेऽन्धकवृष्णिषु ।

\marginnote{\scriptsize ६.२.३५}\noindent संख्या ।

\marginnote{\scriptsize ६.२.३६}\noindent आचार्योपसर्जनश्चान्तेवासी ।

\marginnote{\scriptsize ६.२.३७}\noindent कार्तकौजपादयश्च ।

\marginnote{\scriptsize ६.२.३८}\noindent महान् व्रीह्यपराह्णगृष्टीष्वासजाबाल- भारभारतहैलिहिलरौरवप्रवृद्धेषु ।

\marginnote{\scriptsize ६.२.३९}\noindent क्षुल्लकश्च वैश्वदेवे ।

\marginnote{\scriptsize ६.२.४०}\noindent उष्ट्रः सादिवाम्योः ।

\marginnote{\scriptsize ६.२.४१}\noindent गौः सादसादिसारथिषु ।

\marginnote{\scriptsize ६.२.४२}\noindent कुरुगार्हपतरिक्तगुर्वसूतजरत्यश्लीलदृढरूपा- पारेवडवातैतिलकद्रूःपण्यकम्बलो दासीभाराणां च ।

\marginnote{\scriptsize ६.२.४३}\noindent चतुर्थी तदर्थे ।

\marginnote{\scriptsize ६.२.४४}\noindent अर्थे ।

\marginnote{\scriptsize ६.२.४५}\noindent क्ते च ।

\marginnote{\scriptsize ६.२.४६}\noindent कर्मधारयेऽनिष्ठा ।

\marginnote{\scriptsize ६.२.४७}\noindent अहीने द्वितीया ।

\marginnote{\scriptsize ६.२.४८}\noindent तृतीया कर्मणि ।

\marginnote{\scriptsize ६.२.४९}\noindent गतिरनन्तरः ।

\marginnote{\scriptsize ६.२.५०}\noindent तादौ च निति कृत्यतौ ।

\marginnote{\scriptsize ६.२.५१}\noindent तवै चान्तश्च युगपत् ।

\marginnote{\scriptsize ६.२.५२}\noindent अनिगन्तोऽञ्चतौ वप्रत्यये ।

\marginnote{\scriptsize ६.२.५३}\noindent न्यधी च ।

\marginnote{\scriptsize ६.२.५४}\noindent ईषदन्यतरस्याम् ।

\marginnote{\scriptsize ६.२.५५}\noindent हिरण्यपरिमाणं धने ।

\marginnote{\scriptsize ६.२.५६}\noindent प्रथमोऽचिरोपसम्पत्तौ ।

\marginnote{\scriptsize ६.२.५७}\noindent कतरकतमौ कर्मधारये ।

\marginnote{\scriptsize ६.२.५८}\noindent आर्यो ब्राह्मणकुमारयोः ।

\marginnote{\scriptsize ६.२.५९}\noindent राजा च ।

\marginnote{\scriptsize ६.२.६०}\noindent षष्ठी प्रत्येनसि ।

\marginnote{\scriptsize ६.२.६१}\noindent क्ते नित्यार्थे ।

\marginnote{\scriptsize ६.२.६२}\noindent ग्रामः शिल्पिनि ।

\marginnote{\scriptsize ६.२.६३}\noindent राजा च प्रशंसायाम् ।

\marginnote{\scriptsize ६.२.६४}\noindent आदिरुदात्तः ।

\marginnote{\scriptsize ६.२.६५}\noindent सप्तमीहारिणौ धर्म्येऽहरणे ।

\marginnote{\scriptsize ६.२.६६}\noindent युक्ते च ।

\marginnote{\scriptsize ६.२.६७}\noindent विभाषाऽध्यक्षे ।

\marginnote{\scriptsize ६.२.६८}\noindent पापं च शिल्पिनि ।

\marginnote{\scriptsize ६.२.६९}\noindent गोत्रान्तेवासिमाणवब्राह्मणेषु क्षेपे ।

\marginnote{\scriptsize ६.२.७०}\noindent अङ्गानि मैरेये ।

\marginnote{\scriptsize ६.२.७१}\noindent भक्ताख्यास्तदर्थेषु ।

\marginnote{\scriptsize ६.२.७२}\noindent गोबिडालसिंहसैन्धवेषूपमाने ।

\marginnote{\scriptsize ६.२.७३}\noindent अके जीविकाऽर्थे ।

\marginnote{\scriptsize ६.२.७४}\noindent प्राचां क्रीडायाम् ।

\marginnote{\scriptsize ६.२.७५}\noindent अणि नियुक्ते ।

\marginnote{\scriptsize ६.२.७६}\noindent शिल्पिनि चाकृञः ।

\marginnote{\scriptsize ६.२.७७}\noindent संज्ञायां च ।

\marginnote{\scriptsize ६.२.७८}\noindent गोतन्तियवं पाले ।

\marginnote{\scriptsize ६.२.७९}\noindent णिनि ।

\marginnote{\scriptsize ६.२.८०}\noindent उपमानं शब्दार्थप्रकृतावेव ।

\marginnote{\scriptsize ६.२.८१}\noindent युक्तारोह्यादयश्च ।

\marginnote{\scriptsize ६.२.८२}\noindent दीर्घकाशतुषभ्राष्ट्रवटं जे ।

\marginnote{\scriptsize ६.२.८३}\noindent अन्त्यात् पूर्वं बह्वचः ।

\marginnote{\scriptsize ६.२.८४}\noindent ग्रामेऽनिवसन्तः ।

\marginnote{\scriptsize ६.२.८५}\noindent घोषादिषु ।

\marginnote{\scriptsize ६.२.८६}\noindent छात्र्यादयः शालायाम् ।

\marginnote{\scriptsize ६.२.८७}\noindent प्रस्थेऽवृद्धमकर्क्यादीनाम् ।

\marginnote{\scriptsize ६.२.८८}\noindent मालाऽऽदीनां च ।

\marginnote{\scriptsize ६.२.८९}\noindent अमहन्नवं नगरेऽनुदीचाम् ।

\marginnote{\scriptsize ६.२.९०}\noindent अर्मे चावर्णं द्व्यच्त्र्यच् ।

\marginnote{\scriptsize ६.२.९१}\noindent न भूताधिकसंजीवमद्राश्मकज्जलम् ।

\marginnote{\scriptsize ६.२.९२}\noindent अन्तः ।

\marginnote{\scriptsize ६.२.९३}\noindent सर्वं गुणकार्त्स्न्ये ।

\marginnote{\scriptsize ६.२.९४}\noindent संज्ञायां गिरिनिकाययोः ।

\marginnote{\scriptsize ६.२.९५}\noindent कुमार्यां वयसि ।

\marginnote{\scriptsize ६.२.९६}\noindent उदकेऽकेवले ।

\marginnote{\scriptsize ६.२.९७}\noindent द्विगौ क्रतौ ।

\marginnote{\scriptsize ६.२.९८}\noindent सभायां नपुंसके ।

\marginnote{\scriptsize ६.२.९९}\noindent पुरे प्राचाम् ।

\marginnote{\scriptsize ६.२.१००}\noindent अरिष्टगौडपूर्वे च ।

\marginnote{\scriptsize ६.२.१०१}\noindent न हास्तिनफलकमार्देयाः ।

\marginnote{\scriptsize ६.२.१०२}\noindent कुसूलकूपकुम्भशालं बिले ।

\marginnote{\scriptsize ६.२.१०३}\noindent दिक्शब्दा ग्रामजनपदाख्यानचानराटेषु ।

\marginnote{\scriptsize ६.२.१०४}\noindent आचार्योपसर्जनश्चान्तेवासिनि ।

\marginnote{\scriptsize ६.२.१०५}\noindent उत्तरपदवृद्धौ सर्वं च ।

\marginnote{\scriptsize ६.२.१०६}\noindent बहुव्रीहौ विश्वं संज्ञयाम् ।

\marginnote{\scriptsize ६.२.१०७}\noindent उदराश्वेषुषु ।

\marginnote{\scriptsize ६.२.१०८}\noindent क्षेपे ।

\marginnote{\scriptsize ६.२.१०९}\noindent नदी बन्धुनि ।

\marginnote{\scriptsize ६.२.११०}\noindent निष्ठोपसर्गपूर्वमन्यतरस्याम् ।

\marginnote{\scriptsize ६.२.१११}\noindent उत्तरपदादिः ।

\marginnote{\scriptsize ६.२.११२}\noindent कर्णो वर्णलक्षणात् ।

\marginnote{\scriptsize ६.२.११३}\noindent संज्ञौपम्ययोश्च ।

\marginnote{\scriptsize ६.२.११४}\noindent कण्ठपृष्ठग्रीवाजंघं च ।

\marginnote{\scriptsize ६.२.११५}\noindent श‍ृङ्गमवस्थायां च ।

\marginnote{\scriptsize ६.२.११६}\noindent नञो जरमरमित्रमृताः ।

\marginnote{\scriptsize ६.२.११७}\noindent सोर्मनसी अलोमोषसी ।

\marginnote{\scriptsize ६.२.११८}\noindent क्रत्वादयश्च ।

\marginnote{\scriptsize ६.२.११९}\noindent आद्युदात्तं द्व्यच् छन्दसि ।

\marginnote{\scriptsize ६.२.१२०}\noindent वीरवीर्यौ च ।

\marginnote{\scriptsize ६.२.१२१}\noindent कूलतीरतूलमूलशालाऽक्षसममव्ययीभावे ।

\marginnote{\scriptsize ६.२.१२२}\noindent कंसमन्थशूर्पपाय्यकाण्डं द्विगौ ।

\marginnote{\scriptsize ६.२.१२३}\noindent तत्पुरुषे शालायां नपुंसके ।

\marginnote{\scriptsize ६.२.१२४}\noindent कन्था च ।

\marginnote{\scriptsize ६.२.१२५}\noindent आदिश्चिहणादीनाम् ।

\marginnote{\scriptsize ६.२.१२६}\noindent चेलखेटकटुककाण्डं गर्हायाम् ।

\marginnote{\scriptsize ६.२.१२७}\noindent चीरमुपमानम् ।

\marginnote{\scriptsize ६.२.१२८}\noindent पललसूपशाकं मिश्रे ।

\marginnote{\scriptsize ६.२.१२९}\noindent कूलसूदस्थलकर्षाः संज्ञायाम् ।

\marginnote{\scriptsize ६.२.१३०}\noindent अकर्मधारये राज्यम् ।

\marginnote{\scriptsize ६.२.१३१}\noindent वर्ग्यादयश्च ।

\marginnote{\scriptsize ६.२.१३२}\noindent पुत्रः पुंभ्यः ।

\marginnote{\scriptsize ६.२.१३३}\noindent नाचार्यराजर्त्विक्संयुक्तज्ञात्याख्येभ्यः ।

\marginnote{\scriptsize ६.२.१३४}\noindent चूर्णादीन्यप्राणिषष्ठ्याः ।

\marginnote{\scriptsize ६.२.१३५}\noindent षट् च काण्डादीनि ।

\marginnote{\scriptsize ६.२.१३६}\noindent कुण्डं वनम् ।

\marginnote{\scriptsize ६.२.१३७}\noindent प्रकृत्या भगालम् ।

\marginnote{\scriptsize ६.२.१३८}\noindent शितेर्नित्याबह्वज्बहुव्रीहावभसत् ।

\marginnote{\scriptsize ६.२.१३९}\noindent गतिकारकोपपदात् कृत् ।

\marginnote{\scriptsize ६.२.१४०}\noindent उभे वनस्पत्यादिषु युगपत् ।

\marginnote{\scriptsize ६.२.१४१}\noindent देवताद्वंद्वे च ।

\marginnote{\scriptsize ६.२.१४२}\noindent नोत्तरपदेऽनुदात्तादावपृथिवीरुद्रपूषमन्थिषु ।

\marginnote{\scriptsize ६.२.१४३}\noindent अन्तः ।

\marginnote{\scriptsize ६.२.१४४}\noindent थाथघङ्क्ताजबित्रकाणाम् ।

\marginnote{\scriptsize ६.२.१४५}\noindent सूपमानात् क्तः ।

\marginnote{\scriptsize ६.२.१४६}\noindent संज्ञायामनाचितादीनाम् ।

\marginnote{\scriptsize ६.२.१४७}\noindent प्रवृद्धादीनां च ।

\marginnote{\scriptsize ६.२.१४८}\noindent कारकाद्दत्तश्रुतयोरेवाशिषि ।

\marginnote{\scriptsize ६.२.१४९}\noindent इत्थम्भूतेन कृतमिति च ।

\marginnote{\scriptsize ६.२.१५०}\noindent अनो भावकर्मवचनः ।

\marginnote{\scriptsize ६.२.१५१}\noindent मन्क्तिन्व्याख्यानशयनासनस्थानयाजकादिक्रीताः ।

\marginnote{\scriptsize ६.२.१५२}\noindent सप्तम्याः पुण्यम् ।

\marginnote{\scriptsize ६.२.१५३}\noindent ऊनार्थकलहं तृतीयायाः ।

\marginnote{\scriptsize ६.२.१५४}\noindent मिश्रं चानुपसर्गमसंधौ ।

\marginnote{\scriptsize ६.२.१५५}\noindent नञो गुणप्रतिषेधे सम्पाद्यर्हहितालमर्थास्तद्धिताः ।

\marginnote{\scriptsize ६.२.१५६}\noindent ययतोश्चातदर्थे ।

\marginnote{\scriptsize ६.२.१५७}\noindent अच्कावशक्तौ ।

\marginnote{\scriptsize ६.२.१५८}\noindent आक्रोशे च ।

\marginnote{\scriptsize ६.२.१५९}\noindent संज्ञायाम् ।

\marginnote{\scriptsize ६.२.१६०}\noindent कृत्योकेष्णुच्चार्वादयश्च ।

\marginnote{\scriptsize ६.२.१६१}\noindent विभाषा तृन्नन्नतीक्ष्णशुचिषु ।

\marginnote{\scriptsize ६.२.१६२}\noindent बहुव्रीहाविदमेतत्तद्भ्यः प्रथमपूरणयोः क्रियागणने ।

\marginnote{\scriptsize ६.२.१६३}\noindent संख्यायाः स्तनः ।

\marginnote{\scriptsize ६.२.१६४}\noindent विभाषा छन्दसि ।

\marginnote{\scriptsize ६.२.१६५}\noindent संज्ञायां मित्राजिनयोः ।

\marginnote{\scriptsize ६.२.१६६}\noindent व्यवायिनोऽन्तरम् ।

\marginnote{\scriptsize ६.२.१६७}\noindent मुखं स्वाङ्गम् ।

\marginnote{\scriptsize ६.२.१६८}\noindent नाव्ययदिक्शब्दगोमहत्स्थूलमुष्टिपृथुवत्सेभ्यः ।

\marginnote{\scriptsize ६.२.१६९}\noindent निष्ठोपमानादन्यतरस्याम् ।

\marginnote{\scriptsize ६.२.१७०}\noindent जातिकालसुखादिभ्योऽनाच्छादनात् क्तोऽकृतमितप्रतिपन्नाः ।

\marginnote{\scriptsize ६.२.१७१}\noindent वा जाते ।

\marginnote{\scriptsize ६.२.१७२}\noindent नञ्सुभ्याम् ।

\marginnote{\scriptsize ६.२.१७३}\noindent कपि पूर्वम् ।

\marginnote{\scriptsize ६.२.१७४}\noindent ह्रस्वान्तेऽन्त्यात् पूर्वम् ।

\marginnote{\scriptsize ६.२.१७५}\noindent बहोर्नञ्वदुत्तरपदभूम्नि ।

\marginnote{\scriptsize ६.२.१७६}\noindent न गुणादयोऽवयवाः ।

\marginnote{\scriptsize ६.२.१७७}\noindent उपसर्गात् स्वाङ्गं ध्रुवमपर्शु ।

\marginnote{\scriptsize ६.२.१७८}\noindent वनं समासे ।

\marginnote{\scriptsize ६.२.१७९}\noindent अन्तः ।

\marginnote{\scriptsize ६.२.१८०}\noindent अन्तश्च ।

\marginnote{\scriptsize ६.२.१८१}\noindent न निविभ्याम् ।

\marginnote{\scriptsize ६.२.१८२}\noindent परेरभितोभाविमण्डलम् ।

\marginnote{\scriptsize ६.२.१८३}\noindent प्रादस्वाङ्गं संज्ञायाम् ।

\marginnote{\scriptsize ६.२.१८४}\noindent निरुदकादीनि च ।

\marginnote{\scriptsize ६.२.१८५}\noindent अभेर्मुखम् ।

\marginnote{\scriptsize ६.२.१८६}\noindent अपाच्च ।

\marginnote{\scriptsize ६.२.१८७}\noindent स्फिगपूतवीणाऽञ्जोऽध्वकुक्षिसीरनामनाम च ।

\marginnote{\scriptsize ६.२.१८८}\noindent अधेरुपरिस्थम् ।

\marginnote{\scriptsize ६.२.१८९}\noindent अनोरप्रधानकनीयसी ।

\marginnote{\scriptsize ६.२.१९०}\noindent पुरुषश्चान्वादिष्टः ।

\marginnote{\scriptsize ६.२.१९१}\noindent अतेरकृत्पदे ।

\marginnote{\scriptsize ६.२.१९२}\noindent नेरनिधाने ।

\marginnote{\scriptsize ६.२.१९३}\noindent प्रतेरंश्वादयस्तत्पुरुषे ।

\marginnote{\scriptsize ६.२.१९४}\noindent उपाद् द्व्यजजिनमगौरादयः ।

\marginnote{\scriptsize ६.२.१९५}\noindent सोरवक्षेपणे ।

\marginnote{\scriptsize ६.२.१९६}\noindent विभाषोत्पुच्छे ।

\marginnote{\scriptsize ६.२.१९७}\noindent द्वित्रिभ्यां पाद्दन्मूर्धसु बहुव्रीहौ ।

\marginnote{\scriptsize ६.२.१९८}\noindent सक्थं चाक्रान्तात् ।

\marginnote{\scriptsize ६.२.१९९}\noindent परादिश्छन्दसि बहुलम् ।

\marginnote{\scriptsize ६.३.१}\noindent अलुगुत्तरपदे ।

\marginnote{\scriptsize ६.३.२}\noindent पञ्चम्याः स्तोकादिभ्यः ।

\marginnote{\scriptsize ६.३.३}\noindent ओजःसहोऽम्भस्तमसः तृतीयायाः ।

\marginnote{\scriptsize ६.३.४}\noindent मनसः संज्ञायाम् ।

\marginnote{\scriptsize ६.३.५}\noindent आज्ञायिनि च ।

\marginnote{\scriptsize ६.३.६}\noindent आत्मनश्च पूरणे ।

\marginnote{\scriptsize ६.३.७}\noindent वैयाकरणाख्यायां चतुर्थ्याः ।

\marginnote{\scriptsize ६.३.८}\noindent परस्य च ।

\marginnote{\scriptsize ६.३.९}\noindent हलदन्तात् सप्तम्याः संज्ञायाम् ।

\marginnote{\scriptsize ६.३.१०}\noindent कारनाम्नि च प्राचां हलादौ ।

\marginnote{\scriptsize ६.३.११}\noindent मध्याद्गुरौ ।

\marginnote{\scriptsize ६.३.१२}\noindent अमूर्धमस्तकात् स्वाङ्गादकामे ।

\marginnote{\scriptsize ६.३.१३}\noindent बन्धे च विभाषा ।

\marginnote{\scriptsize ६.३.१४}\noindent तत्पुरुषे कृति बहुलम् ।

\marginnote{\scriptsize ६.३.१५}\noindent प्रावृट्शरत्कालदिवां जे ।

\marginnote{\scriptsize ६.३.१६}\noindent विभाषा वर्षक्षरशरवरात् ।

\marginnote{\scriptsize ६.३.१७}\noindent घकालतनेषु कालनाम्नः ।

\marginnote{\scriptsize ६.३.१८}\noindent शयवासवासिषु अकालात् ।

\marginnote{\scriptsize ६.३.१९}\noindent नेन्सिद्धबध्नातिषु ।

\marginnote{\scriptsize ६.३.२०}\noindent स्थे च भाषायाम् ।

\marginnote{\scriptsize ६.३.२१}\noindent षष्ठ्या आक्रोशे ।

\marginnote{\scriptsize ६.३.२२}\noindent पुत्रेऽन्यतरस्याम् ।

\marginnote{\scriptsize ६.३.२३}\noindent ऋतो विद्यायोनिसम्बन्धेभ्यः ।

\marginnote{\scriptsize ६.३.२४}\noindent विभाषा स्वसृपत्योः ।

\marginnote{\scriptsize ६.३.२५}\noindent आनङ् ऋतो द्वंद्वे ।

\marginnote{\scriptsize ६.३.२६}\noindent देवताद्वंद्वे च ।

\marginnote{\scriptsize ६.३.२७}\noindent ईदग्नेः सोमवरुणयोः ।

\marginnote{\scriptsize ६.३.२८}\noindent इद्वृद्धौ ।

\marginnote{\scriptsize ६.३.२९}\noindent दिवो द्यावा ।

\marginnote{\scriptsize ६.३.३०}\noindent दिवसश्च पृथिव्याम् ।

\marginnote{\scriptsize ६.३.३१}\noindent उषासोषसः ।

\marginnote{\scriptsize ६.३.३२}\noindent मातरपितरावुदीचाम् ।

\marginnote{\scriptsize ६.३.३३}\noindent पितरामातरा च च्छन्दसि ।

\marginnote{\scriptsize ६.३.३४}\noindent स्त्रियाः पुंवद्भाषितपुंस्कादनूङ् समानाधिकरणे स्त्रियामपूरणीप्रियाऽऽदिषु ।

\marginnote{\scriptsize ६.३.३५}\noindent तसिलादिषु आकृत्वसुचः ।

\marginnote{\scriptsize ६.३.३६}\noindent क्यङ्मानिनोश्च ।

\marginnote{\scriptsize ६.३.३७}\noindent न कोपधायाः ।

\marginnote{\scriptsize ६.३.३८}\noindent संज्ञापूरण्योश्च ।

\marginnote{\scriptsize ६.३.३९}\noindent वृद्धिनिमित्तस्य च तद्धितस्यारक्तविकारे ।

\marginnote{\scriptsize ६.३.४०}\noindent स्वाङ्गाच्चेतोऽमानिनि ।

\marginnote{\scriptsize ६.३.४१}\noindent जातेश्च ।

\marginnote{\scriptsize ६.३.४२}\noindent पुंवत् कर्मधारयजातीयदेशीयेषु ।

\marginnote{\scriptsize ६.३.४३}\noindent घरूपकल्पचेलड्ब्रुवगोत्रमतहतेषु ङ्योऽनेकाचो ह्रस्वः ।

\marginnote{\scriptsize ६.३.४४}\noindent नद्याः शेषस्यान्यतरस्याम् ।

\marginnote{\scriptsize ६.३.४५}\noindent उगितश्च ।

\marginnote{\scriptsize ६.३.४६}\noindent आन्महतः समानाधिकरणजातीययोः ।

\marginnote{\scriptsize ६.३.४७}\noindent द्व्यष्टनः संख्यायामबहुव्रीह्यशीत्योः ।

\marginnote{\scriptsize ६.३.४८}\noindent त्रेस्त्रयः ।

\marginnote{\scriptsize ६.३.४९}\noindent विभाषा चत्वारिंशत्प्रभृतौ सर्वेषाम् ।

\marginnote{\scriptsize ६.३.५०}\noindent हृदयस्य हृल्लेखयदण्लासेषु ।

\marginnote{\scriptsize ६.३.५१}\noindent वा शोकष्यञ्रोगेषु ।

\marginnote{\scriptsize ६.३.५२}\noindent पादस्य पदाज्यातिगोपहतेषु ।

\marginnote{\scriptsize ६.३.५३}\noindent पद् यत्यतदर्थे ।

\marginnote{\scriptsize ६.३.५४}\noindent हिमकाषिहतिषु च ।

\marginnote{\scriptsize ६.३.५५}\noindent ऋचः शे ।

\marginnote{\scriptsize ६.३.५६}\noindent वा घोषमिश्रशब्देषु ।

\marginnote{\scriptsize ६.३.५७}\noindent उदकस्योदः संज्ञायाम् ।

\marginnote{\scriptsize ६.३.५८}\noindent पेषंवासवाहनधिषु च ।

\marginnote{\scriptsize ६.३.५९}\noindent एकहलादौ पूरयितव्येऽन्यतरस्याम् ।

\marginnote{\scriptsize ६.३.६०}\noindent मन्थौदनसक्तुबिन्दुवज्रभारहारवीवधगाहेषु च ।

\marginnote{\scriptsize ६.३.६१}\noindent इको ह्रस्वोऽङ्यो गालवस्य ।

\marginnote{\scriptsize ६.३.६२}\noindent एक तद्धिते च ।

\marginnote{\scriptsize ६.३.६३}\noindent ङ्यापोः संज्ञाछन्दसोर्बहुलम् ।

\marginnote{\scriptsize ६.३.६४}\noindent त्वे च ।

\marginnote{\scriptsize ६.३.६५}\noindent इष्टकेषीकामालानां चिततूलभारिषु ।

\marginnote{\scriptsize ६.३.६६}\noindent खित्यनव्ययस्य ।

\marginnote{\scriptsize ६.३.६७}\noindent अरुर्द्विषदजन्तस्य मुम् ।

\marginnote{\scriptsize ६.३.६८}\noindent इच एकाचोऽम्प्रत्ययवच्च ।

\marginnote{\scriptsize ६.३.६९}\noindent वाचंयमपुरंदरौ च ।

\marginnote{\scriptsize ६.३.७०}\noindent कारे सत्यागदस्य ।

\marginnote{\scriptsize ६.३.७१}\noindent श्येनतिलस्य पाते ञे ।

\marginnote{\scriptsize ६.३.७२}\noindent रात्रेः कृति विभाषा ।

\marginnote{\scriptsize ६.३.७३}\noindent नलोपः नञः ।

\marginnote{\scriptsize ६.३.७४}\noindent तस्मान्नुडचि ।

\marginnote{\scriptsize ६.३.७५}\noindent नभ्राण्नपान्नवेदानासत्यानमुचिनकुलनख- नपुंसकनक्षत्रनक्रनाकेषु प्रकृत्या ।

\marginnote{\scriptsize ६.३.७६}\noindent एकादिश्चैकस्य चादुक् ।

\marginnote{\scriptsize ६.३.७७}\noindent नगोऽप्राणिष्वन्यतरस्याम् ।

\marginnote{\scriptsize ६.३.७८}\noindent सहस्य सः संज्ञायाम् ।

\marginnote{\scriptsize ६.३.७९}\noindent ग्रन्थान्ताधिके च ।

\marginnote{\scriptsize ६.३.८०}\noindent द्वितीये चानुपाख्ये ।

\marginnote{\scriptsize ६.३.८१}\noindent अव्ययीभावे चाकाले ।

\marginnote{\scriptsize ६.३.८२}\noindent वोपसर्जनस्य ।

\marginnote{\scriptsize ६.३.८३}\noindent प्रकृत्याऽऽशिष्यगोवत्सहलेषु ।

\marginnote{\scriptsize ६.३.८४}\noindent समानस्य छन्दस्यमूर्धप्रभृत्युदर्केषु ।

\marginnote{\scriptsize ६.३.८५}\noindent ज्योतिर्जनपदरात्रिनाभिनामगोत्ररूपस्थानवर्ण- वयोवचनबन्धुषु ।

\marginnote{\scriptsize ६.३.८६}\noindent चरणे ब्रह्मचारिणि ।

\marginnote{\scriptsize ६.३.८७}\noindent तीर्थे ये ।

\marginnote{\scriptsize ६.३.८८}\noindent विभाषोदरे ।

\marginnote{\scriptsize ६.३.८९}\noindent दृग्दृशवतुषु ।

\marginnote{\scriptsize ६.३.९०}\noindent इदङ्किमोरीश्की ।

\marginnote{\scriptsize ६.३.९१}\noindent आ सर्वनाम्नः ।

\marginnote{\scriptsize ६.३.९२}\noindent विष्वग्देवयोश्च टेरद्र्यञ्चतौ वप्रत्यये ।

\marginnote{\scriptsize ६.३.९३}\noindent समः समि ।

\marginnote{\scriptsize ६.३.९४}\noindent तिरसस्तिर्यलोपे ।

\marginnote{\scriptsize ६.३.९५}\noindent सहस्य सध्रिः ।

\marginnote{\scriptsize ६.३.९६}\noindent सध मादस्थयोश्छन्दसि ।

\marginnote{\scriptsize ६.३.९७}\noindent द्व्यन्तरुपसर्गेभ्योऽप ईत् ।

\marginnote{\scriptsize ६.३.९८}\noindent ऊदनोर्देशे ।

\marginnote{\scriptsize ६.३.९९}\noindent अषष्ठ्यतृतीयास्थस्यान्यस्य दुगाशिराशाऽऽस्थाऽऽस्थितोत्सुकोतिकारकरागच्छेषु ।

\marginnote{\scriptsize ६.३.१००}\noindent अर्थे विभाषा ।

\marginnote{\scriptsize ६.३.१०१}\noindent कोः कत् तत्पुरुषेऽचि ।

\marginnote{\scriptsize ६.३.१०२}\noindent रथवदयोश्च ।

\marginnote{\scriptsize ६.३.१०३}\noindent तृणे च जातौ ।

\marginnote{\scriptsize ६.३.१०४}\noindent का पथ्यक्षयोः ।

\marginnote{\scriptsize ६.३.१०५}\noindent ईषदर्थे ।

\marginnote{\scriptsize ६.३.१०६}\noindent विभाषा पुरुषे ।

\marginnote{\scriptsize ६.३.१०७}\noindent कवं चोष्णे ।

\marginnote{\scriptsize ६.३.१०८}\noindent पथि च च्छन्दसि ।

\marginnote{\scriptsize ६.३.१०९}\noindent पृषोदरादीनि यथोपदिष्टम् ।

\marginnote{\scriptsize ६.३.११०}\noindent संख्याविसायपूर्वस्याह्नस्याहन्नन्यतरस्यां ङौ ।

\marginnote{\scriptsize ६.३.१११}\noindent ढ्रलोपे पूर्वस्य दीर्घोऽणः ।

\marginnote{\scriptsize ६.३.११२}\noindent सहिवहोरोदवर्णस्य ।

\marginnote{\scriptsize ६.३.११३}\noindent साढ्यै साढ्वा साढेति निगमे ।

\marginnote{\scriptsize ६.३.११४}\noindent संहितायाम् ।

\marginnote{\scriptsize ६.३.११५}\noindent कर्णे लक्षणस्याविष्टाष्टपञ्चमणिभिन्न- छिन्नछिद्रस्रुवस्वस्तिकस्य ।

\marginnote{\scriptsize ६.३.११६}\noindent नहिवृतिवृषिव्यधिरुचिसहितनिषु क्वौ ।

\marginnote{\scriptsize ६.३.११७}\noindent वनगिर्योः संज्ञायां कोटरकिंशुलकादीनाम् ।

\marginnote{\scriptsize ६.३.११८}\noindent वले ।

\marginnote{\scriptsize ६.३.११९}\noindent मतौ बह्वचोऽनजिरादीनाम् ।

\marginnote{\scriptsize ६.३.१२०}\noindent शरादीनां च ।

\marginnote{\scriptsize ६.३.१२१}\noindent इकः वहे अपीलोः ।

\marginnote{\scriptsize ६.३.१२२}\noindent उपसर्गस्य घञ्यमनुष्ये बहुलम् ।

\marginnote{\scriptsize ६.३.१२३}\noindent इकः काशे ।

\marginnote{\scriptsize ६.३.१२४}\noindent दस्ति ।

\marginnote{\scriptsize ६.३.१२५}\noindent अष्टनः संज्ञायाम् ।

\marginnote{\scriptsize ६.३.१२६}\noindent छन्दसि च ।

\marginnote{\scriptsize ६.३.१२७}\noindent चितेः कपि ।

\marginnote{\scriptsize ६.३.१२८}\noindent विश्वस्य वसुराटोः ।

\marginnote{\scriptsize ६.३.१२९}\noindent नरे संज्ञायाम् ।

\marginnote{\scriptsize ६.३.१३०}\noindent मित्रे चर्षौ ।

\marginnote{\scriptsize ६.३.१३१}\noindent मन्त्रे सोमाश्वेन्द्रियविश्वदेव्यस्य मतौ ।

\marginnote{\scriptsize ६.३.१३२}\noindent ओषधेश्च विभक्तावप्रथमायाम् ।

\marginnote{\scriptsize ६.३.१३३}\noindent ऋचि तुनुघमक्षुतङ्कुत्रोरुष्याणाम् ।

\marginnote{\scriptsize ६.३.१३४}\noindent इकः सुञि ।

\marginnote{\scriptsize ६.३.१३५}\noindent द्व्यचोऽतस्तिङः ।

\marginnote{\scriptsize ६.३.१३६}\noindent निपातस्य च ।

\marginnote{\scriptsize ६.३.१३७}\noindent अन्येषामपि दृश्यते ।

\marginnote{\scriptsize ६.३.१३८}\noindent चौ ।

\marginnote{\scriptsize ६.३.१३९}\noindent सम्प्रसारणस्य ।

\marginnote{\scriptsize ६.४.१}\noindent अङ्गस्य ।

\marginnote{\scriptsize ६.४.२}\noindent हलः ।

\marginnote{\scriptsize ६.४.३}\noindent नामि ।

\marginnote{\scriptsize ६.४.४}\noindent न तिसृचतसृ ।

\marginnote{\scriptsize ६.४.५}\noindent छन्दस्युभयथा ।

\marginnote{\scriptsize ६.४.६}\noindent नृ च ।

\marginnote{\scriptsize ६.४.७}\noindent नोपधायाः ।

\marginnote{\scriptsize ६.४.८}\noindent सर्वनामस्थाने चासम्बुद्धौ ।

\marginnote{\scriptsize ६.४.९}\noindent वा षपूर्वस्य निगमे ।

\marginnote{\scriptsize ६.४.१०}\noindent सान्तमहतः संयोगस्य । ६.४.११ अप्तृन्तृच्स्वसृनप्तृनेष्टृत्वष्टृक्षत्तृहोतृपोतॄ- प्रशास्तॄणाम् ।

\marginnote{\scriptsize ६.४.१२}\noindent इन्हन्पूषार्यम्णां शौ ।

\marginnote{\scriptsize ६.४.१३}\noindent सौ च ।

\marginnote{\scriptsize ६.४.१४}\noindent अत्वसन्तस्य चाधातोः ।

\marginnote{\scriptsize ६.४.१५}\noindent अनुनासिकस्य क्विझलोः क्ङिति ।

\marginnote{\scriptsize ६.४.१६}\noindent अज्झनगमां सनि ।

\marginnote{\scriptsize ६.४.१७}\noindent तनोतेर्विभाषा ।

\marginnote{\scriptsize ६.४.१८}\noindent क्रमश्च क्त्वि ।

\marginnote{\scriptsize ६.४.१९}\noindent च्छ्वोः शूडनुनासिके च ।

\marginnote{\scriptsize ६.४.२०}\noindent ज्वरत्वरश्रिव्यविमवामुपधायाश्च ।

\marginnote{\scriptsize ६.४.२१}\noindent राल्लोपः ।

\marginnote{\scriptsize ६.४.२२}\noindent असिद्धवदत्राभात् ।

\marginnote{\scriptsize ६.४.२३}\noindent श्नान्नलोपः ।

\marginnote{\scriptsize ६.४.२४}\noindent अनिदितां हल उपधायाः क्ङिति ।

\marginnote{\scriptsize ६.४.२५}\noindent दन्शसञ्जस्वञ्जां शपि ।

\marginnote{\scriptsize ६.४.२६}\noindent रञ्जेश्च ।

\marginnote{\scriptsize ६.४.२७}\noindent घञि च भावकरणयोः ।

\marginnote{\scriptsize ६.४.२८}\noindent स्यदो जवे ।

\marginnote{\scriptsize ६.४.२९}\noindent अवोदैधौद्मप्रश्रथहिमश्रथाः ।

\marginnote{\scriptsize ६.४.३०}\noindent नाञ्चेः पूजायाम् ।

\marginnote{\scriptsize ६.४.३१}\noindent क्त्वि स्कन्दिस्यन्दोः ।

\marginnote{\scriptsize ६.४.३२}\noindent जान्तनशां विभाषा ।

\marginnote{\scriptsize ६.४.३३}\noindent भञ्जेश्च चिणि ।

\marginnote{\scriptsize ६.४.३४}\noindent शास इदङ्हलोः ।

\marginnote{\scriptsize ६.४.३५}\noindent शा हौ ।

\marginnote{\scriptsize ६.४.३६}\noindent हन्तेर्जः ।

\marginnote{\scriptsize ६.४.३७}\noindent अनुदात्तोपदेशवनतितनोत्यादीनामनुनासिकलोपो झलि क्ङिति ।

\marginnote{\scriptsize ६.४.३८}\noindent वा ल्यपि ।

\marginnote{\scriptsize ६.४.३९}\noindent न क्तिचि दीर्घश्च ।

\marginnote{\scriptsize ६.४.४०}\noindent गमः क्वौ ।

\marginnote{\scriptsize ६.४.४१}\noindent विड्वनोरनुनासिकस्यात् ।

\marginnote{\scriptsize ६.४.४२}\noindent जनसनखनां सञ्झलोः ।

\marginnote{\scriptsize ६.४.४३}\noindent ये विभाषा ।

\marginnote{\scriptsize ६.४.४४}\noindent तनोतेर्यकि ।

\marginnote{\scriptsize ६.४.४५}\noindent सनः क्तिचि लोपश्चास्यान्यतरस्याम् ।

\marginnote{\scriptsize ६.४.४६}\noindent आर्धधातुके ।

\marginnote{\scriptsize ६.४.४७}\noindent भ्रस्जो रोपधयोः रमन्यतरस्याम् ।

\marginnote{\scriptsize ६.४.४८}\noindent अतो लोपः ।

\marginnote{\scriptsize ६.४.४९}\noindent यस्य हलः ।

\marginnote{\scriptsize ६.४.५०}\noindent क्यस्य विभाषा ।

\marginnote{\scriptsize ६.४.५१}\noindent णेरनिटि ।

\marginnote{\scriptsize ६.४.५२}\noindent निष्ठायां सेटि ।

\marginnote{\scriptsize ६.४.५३}\noindent जनिता मन्त्रे ।

\marginnote{\scriptsize ६.४.५४}\noindent शमिता यज्ञे ।

\marginnote{\scriptsize ६.४.५५}\noindent अयामन्ताल्वाय्येत्न्विष्णुषु ।

\marginnote{\scriptsize ६.४.५६}\noindent ल्यपि लघुपूर्वात् ।

\marginnote{\scriptsize ६.४.५७}\noindent विभाषाऽऽपः ।

\marginnote{\scriptsize ६.४.५८}\noindent युप्लुवोर्दीर्घश्छन्दसि ।

\marginnote{\scriptsize ६.४.५९}\noindent क्षियः ।

\marginnote{\scriptsize ६.४.६०}\noindent निष्ठायां अण्यदर्थे ।

\marginnote{\scriptsize ६.४.६१}\noindent वाऽऽक्रोशदैन्ययोः ।

\marginnote{\scriptsize ६.४.६२}\noindent स्यसिच्सीयुट्तासिषु भावकर्मणोर्- उपदेशेऽज्झनग्रहदृशां वा चिण्वदिट् च ।

\marginnote{\scriptsize ६.४.६३}\noindent दीङो युडचि क्ङिति ।

\marginnote{\scriptsize ६.४.६४}\noindent आतो लोप इटि च ।

\marginnote{\scriptsize ६.४.६५}\noindent ईद्यति ।

\marginnote{\scriptsize ६.४.६६}\noindent घुमास्थागापाजहातिसां हलि ।

\marginnote{\scriptsize ६.४.६७}\noindent एर्लिङि ।

\marginnote{\scriptsize ६.४.६८}\noindent वाऽन्यस्य संयोगादेः ।

\marginnote{\scriptsize ६.४.६९}\noindent न ल्यपि ।

\marginnote{\scriptsize ६.४.७०}\noindent मयतेरिदन्यतरस्याम् ।

\marginnote{\scriptsize ६.४.७१}\noindent लुङ्लङ्लृङ्क्ष्वडुदात्तः ।

\marginnote{\scriptsize ६.४.७२}\noindent आडजादीनाम् ।

\marginnote{\scriptsize ६.४.७३}\noindent छन्दस्यपि दृश्यते ।

\marginnote{\scriptsize ६.४.७४}\noindent न माङ्योगे ।

\marginnote{\scriptsize ६.४.७५}\noindent बहुलं छन्दस्यमाङ्योगेऽपि ।

\marginnote{\scriptsize ६.४.७६}\noindent इरयो रे ।

\marginnote{\scriptsize ६.४.७७}\noindent अचि श्नुधातुभ्रुवां य्वोरियङुवङौ ।

\marginnote{\scriptsize ६.४.७८}\noindent अभ्यासस्यासवर्णे ।

\marginnote{\scriptsize ६.४.७९}\noindent स्त्रियाः ।

\marginnote{\scriptsize ६.४.८०}\noindent वाऽम्शसोः ।

\marginnote{\scriptsize ६.४.८१}\noindent इणो यण् ।

\marginnote{\scriptsize ६.४.८२}\noindent एरनेकाचोऽसंयोगपूर्वस्य ।

\marginnote{\scriptsize ६.४.८३}\noindent ओः सुपि ।

\marginnote{\scriptsize ६.४.८४}\noindent वर्षाभ्वश्च ।

\marginnote{\scriptsize ६.४.८५}\noindent न भूसुधियोः ।

\marginnote{\scriptsize ६.४.८६}\noindent छन्दस्युभयथा ।

\marginnote{\scriptsize ६.४.८७}\noindent हुश्नुवोः सार्वधातुके ।

\marginnote{\scriptsize ६.४.८८}\noindent भुवो वुग्लुङ्लिटोः ।

\marginnote{\scriptsize ६.४.८९}\noindent ऊदुपधाया गोहः ।

\marginnote{\scriptsize ६.४.९०}\noindent दोषो णौ ।

\marginnote{\scriptsize ६.४.९१}\noindent वा चित्तविरागे ।

\marginnote{\scriptsize ६.४.९२}\noindent मितां ह्रस्वः ।

\marginnote{\scriptsize ६.४.९३}\noindent चिण्णमुलोर्दीर्घोऽन्यतरस्याम् ।

\marginnote{\scriptsize ६.४.९४}\noindent खचि ह्रस्वः ।

\marginnote{\scriptsize ६.४.९५}\noindent ह्लादो निष्ठायाम् ।

\marginnote{\scriptsize ६.४.९६}\noindent छादेर्घेऽद्व्युपसर्गस्य ।

\marginnote{\scriptsize ६.४.९७}\noindent इस्मन्त्रन्क्विषु च ।

\marginnote{\scriptsize ६.४.९८}\noindent गमहनजनखनघसां लोपः क्ङित्यनङि ।

\marginnote{\scriptsize ६.४.९९}\noindent तनिपत्योश्छन्दसि ।

\marginnote{\scriptsize ६.४.१००}\noindent घसिभसोर्हलि च ।

\marginnote{\scriptsize ६.४.१०१}\noindent हुझल्भ्यो हेर्धिः ।

\marginnote{\scriptsize ६.४.१०२}\noindent श्रुश‍ृणुपॄकृवृभ्यश्छन्दसि ।

\marginnote{\scriptsize ६.४.१०३}\noindent अङितश्च ।

\marginnote{\scriptsize ६.४.१०४}\noindent चिणो लुक् ।

\marginnote{\scriptsize ६.४.१०५}\noindent अतो हेः ।

\marginnote{\scriptsize ६.४.१०६}\noindent उतश्च प्रत्ययादसंयोगपूर्वात् ।

\marginnote{\scriptsize ६.४.१०७}\noindent लोपश्चास्यान्यतरस्यां म्वोः ।

\marginnote{\scriptsize ६.४.१०८}\noindent नित्यं करोतेः ।

\marginnote{\scriptsize ६.४.१०९}\noindent ये च ।

\marginnote{\scriptsize ६.४.११०}\noindent अत उत् सार्वधातुके ।

\marginnote{\scriptsize ६.४.१११}\noindent श्नसोरल्लोपः ।

\marginnote{\scriptsize ६.४.११२}\noindent श्नाऽभ्यस्तयोरातः ।

\marginnote{\scriptsize ६.४.११३}\noindent ई हल्यघोः ।

\marginnote{\scriptsize ६.४.११४}\noindent इद्दरिद्रस्य ।

\marginnote{\scriptsize ६.४.११५}\noindent भियोऽन्यतरस्याम् ।

\marginnote{\scriptsize ६.४.११६}\noindent जहातेश्च ।

\marginnote{\scriptsize ६.४.११७}\noindent आ च हौ ।

\marginnote{\scriptsize ६.४.११८}\noindent लोपो यि ।

\marginnote{\scriptsize ६.४.११९}\noindent घ्वसोरेद्धावभ्यासलोपश्च ।

\marginnote{\scriptsize ६.४.१२०}\noindent अत एकहल्मध्येऽनादेशादेर्लिटि ।

\marginnote{\scriptsize ६.४.१२१}\noindent थलि च सेटि ।

\marginnote{\scriptsize ६.४.१२२}\noindent तॄफलभजत्रपश्च ।

\marginnote{\scriptsize ६.४.१२३}\noindent राधो हिंसायाम् ।

\marginnote{\scriptsize ६.४.१२४}\noindent वा जॄभ्रमुत्रसाम् ।

\marginnote{\scriptsize ६.४.१२५}\noindent फणां च सप्तानाम् ।

\marginnote{\scriptsize ६.४.१२६}\noindent न शसददवादिगुणानाम् ।

\marginnote{\scriptsize ६.४.१२७}\noindent अर्वणस्त्रसावनञः ।

\marginnote{\scriptsize ६.४.१२८}\noindent मघवा बहुलम् ।

\marginnote{\scriptsize ६.४.१२९}\noindent भस्य ।

\marginnote{\scriptsize ६.४.१३०}\noindent पादः पत् ।

\marginnote{\scriptsize ६.४.१३१}\noindent वसोः सम्प्रसारणम् ।

\marginnote{\scriptsize ६.४.१३२}\noindent वाह ऊठ् ।

\marginnote{\scriptsize ६.४.१३३}\noindent श्वयुवमघोनामतद्धिते ।

\marginnote{\scriptsize ६.४.१३४}\noindent अल्लोपोऽनः ।

\marginnote{\scriptsize ६.४.१३५}\noindent षपूर्वहन्धृतराज्ञामणि ।

\marginnote{\scriptsize ६.४.१३६}\noindent विभाषा ङिश्योः ।

\marginnote{\scriptsize ६.४.१३७}\noindent न संयोगाद्वमान्तात् ।

\marginnote{\scriptsize ६.४.१३८}\noindent अचः ।

\marginnote{\scriptsize ६.४.१३९}\noindent उद ईत् ।

\marginnote{\scriptsize ६.४.१४०}\noindent आतो धातोः ।

\marginnote{\scriptsize ६.४.१४१}\noindent मन्त्रेष्वाङ्यादेरात्मनः ।

\marginnote{\scriptsize ६.४.१४२}\noindent ति विंशतेर्डिति ।

\marginnote{\scriptsize ६.४.१४३}\noindent टेः ।

\marginnote{\scriptsize ६.४.१४४}\noindent नस्तद्धिते ।

\marginnote{\scriptsize ६.४.१४५}\noindent अह्नष्टखोरेव ।

\marginnote{\scriptsize ६.४.१४६}\noindent ओर्गुणः ।

\marginnote{\scriptsize ६.४.१४७}\noindent ढे लोपोऽकद्र्वाः ।

\marginnote{\scriptsize ६.४.१४८}\noindent यस्येति च ।

\marginnote{\scriptsize ६.४.१४९}\noindent सूर्यतिष्यागस्त्यमत्स्यानां य उपधायाः ।

\marginnote{\scriptsize ६.४.१५०}\noindent हलस्तद्धितस्य ।

\marginnote{\scriptsize ६.४.१५१}\noindent आपत्यस्य च तद्धितेऽनाति ।

\marginnote{\scriptsize ६.४.१५२}\noindent क्यच्व्योश्च ।

\marginnote{\scriptsize ६.४.१५३}\noindent बिल्वकादिभ्यश्छस्य लुक् ।

\marginnote{\scriptsize ६.४.१५४}\noindent तुरिष्ठेमेयस्सु ।

\marginnote{\scriptsize ६.४.१५५}\noindent टेः ।

\marginnote{\scriptsize ६.४.१५६}\noindent स्थूलदूरयुवह्रस्वक्षिप्रक्षुद्राणां यणादिपरं पूर्वस्य च गुणः ।

\marginnote{\scriptsize ६.४.१५७}\noindent प्रियस्थिरस्फिरोरुबहुलगुरुवृद्धतृप्रदीर्घ- वृन्दारकाणां प्रस्थस्फवर्बंहिगर्वर्षित्रब्द्राघिवृन्दाः ।

\marginnote{\scriptsize ६.४.१५८}\noindent बहोर्लोपो भू च बहोः ।

\marginnote{\scriptsize ६.४.१५९}\noindent इष्ठस्य यिट् च ।

\marginnote{\scriptsize ६.४.१६०}\noindent ज्यादादीयसः ।

\marginnote{\scriptsize ६.४.१६१}\noindent र ऋतो हलादेर्लघोः ।

\marginnote{\scriptsize ६.४.१६२}\noindent विभाषर्जोश्छन्दसि ।

\marginnote{\scriptsize ६.४.१६३}\noindent प्रकृत्यैकाच् ।

\marginnote{\scriptsize ६.४.१६४}\noindent इनण्यनपत्ये ।

\marginnote{\scriptsize ६.४.१६५}\noindent गाथिविदथिकेशिगणिपणिनश्च ।

\marginnote{\scriptsize ६.४.१६६}\noindent संयोगादिश्च ।

\marginnote{\scriptsize ६.४.१६७}\noindent अन् ।

\marginnote{\scriptsize ६.४.१६८}\noindent ये चाभावकर्मणोः ।

\marginnote{\scriptsize ६.४.१६९}\noindent आत्माध्वानौ खे ।

\marginnote{\scriptsize ६.४.१७०}\noindent न मपूर्वोऽपत्येऽवर्मणः ।

\marginnote{\scriptsize ६.४.१७१}\noindent ब्राह्मोअजातौ ।

\marginnote{\scriptsize ६.४.१७२}\noindent कार्मस्ताच्छील्ये ।

\marginnote{\scriptsize ६.४.१७३}\noindent औक्षमनपत्ये ।

\marginnote{\scriptsize ६.४.१७४}\noindent दाण्डिनायनहास्तिनायनाथर्वणिकजैह्माशिनेय- वाशिनायनिभ्रौणहत्यधैवत्यसारवैक्ष्वाकमैत्रेयहिरण्मयानि ।

\marginnote{\scriptsize ६.४.१७५}\noindent ऋत्व्यवास्त्व्यवास्त्वमाध्वीहिरण्ययानि च्छन्दसि ।

\section*{॥ अध्याय ७ ॥}

\marginnote{\scriptsize ७.१.१}\noindent युवोरनाकौ ।

\marginnote{\scriptsize ७.१.२}\noindent आयनेयीनीयियः फढखच्छघां प्रत्ययादीनाम् ।

\marginnote{\scriptsize ७.१.३}\noindent झोऽन्तः ।

\marginnote{\scriptsize ७.१.४}\noindent अदभ्यस्तात् ।

\marginnote{\scriptsize ७.१.५}\noindent आत्मनेपदेष्वनतः ।

\marginnote{\scriptsize ७.१.६}\noindent शीङो रुट् ।

\marginnote{\scriptsize ७.१.७}\noindent वेत्तेर्विभाषा ।

\marginnote{\scriptsize ७.१.८}\noindent बहुलं छन्दसि ।

\marginnote{\scriptsize ७.१.९}\noindent अतो भिस ऐस् ।

\marginnote{\scriptsize ७.१.१०}\noindent बहुलं छन्दसि ।

\marginnote{\scriptsize ७.१.११}\noindent नेदमदसोरकोः ।

\marginnote{\scriptsize ७.१.१२}\noindent टाङसिङसामिनात्स्याः ।

\marginnote{\scriptsize ७.१.१३}\noindent ङेर्यः ।

\marginnote{\scriptsize ७.१.१४}\noindent सर्वनाम्नः स्मै ।

\marginnote{\scriptsize ७.१.१५}\noindent ङसिङ्योः स्मात्स्मिनौ ।

\marginnote{\scriptsize ७.१.१६}\noindent पूर्वादिभ्यो नवभ्यो वा ।

\marginnote{\scriptsize ७.१.१७}\noindent जसः शी ।

\marginnote{\scriptsize ७.१.१८}\noindent औङ आपः ।

\marginnote{\scriptsize ७.१.१९}\noindent नपुंसकाच्च ।

\marginnote{\scriptsize ७.१.२०}\noindent जस्शसोः शिः ।

\marginnote{\scriptsize ७.१.२१}\noindent अष्टाभ्य औश् ।

\marginnote{\scriptsize ७.१.२२}\noindent षड्भ्यो लुक् ।

\marginnote{\scriptsize ७.१.२३}\noindent स्वमोर्नपुंसकात् ।

\marginnote{\scriptsize ७.१.२४}\noindent अतोऽम् ।

\marginnote{\scriptsize ७.१.२५}\noindent अद्ड् डतरादिभ्यः पञ्चभ्यः ।

\marginnote{\scriptsize ७.१.२६}\noindent नेतराच्छन्दसि ।

\marginnote{\scriptsize ७.१.२७}\noindent युष्मदस्मद्भ्यां ङसोऽश् ।

\marginnote{\scriptsize ७.१.२८}\noindent ङे प्रथमयोरम् ।

\marginnote{\scriptsize ७.१.२९}\noindent शसो न ।

\marginnote{\scriptsize ७.१.३०}\noindent भ्यसः भ्यम् ।

\marginnote{\scriptsize ७.१.३१}\noindent पञ्चम्या अत् ।

\marginnote{\scriptsize ७.१.३२}\noindent एकवचनस्य च ।

\marginnote{\scriptsize ७.१.३३}\noindent साम आकम् ।

\marginnote{\scriptsize ७.१.३४}\noindent आत औ णलः ।

\marginnote{\scriptsize ७.१.३५}\noindent तुह्योस्तातङाशिष्यन्यतरस्याम् ।

\marginnote{\scriptsize ७.१.३६}\noindent विदेः शतुर्वसुः ।

\marginnote{\scriptsize ७.१.३७}\noindent समासेऽनञ्पूर्वे क्त्वो ल्यप् ।

\marginnote{\scriptsize ७.१.३८}\noindent क्त्वाऽपि छन्दसि ।

\marginnote{\scriptsize ७.१.३९}\noindent सुपां सुलुक्पूर्वसवर्णाऽऽच्छेयाडाड्यायाजालः ।

\marginnote{\scriptsize ७.१.४०}\noindent अमो मश् ।

\marginnote{\scriptsize ७.१.४१}\noindent लोपस्त आत्मनेपदेषु ।

\marginnote{\scriptsize ७.१.४२}\noindent ध्वमो ध्वात् ।

\marginnote{\scriptsize ७.१.४३}\noindent यजध्वैनमिति च ।

\marginnote{\scriptsize ७.१.४४}\noindent तस्य तात् ।

\marginnote{\scriptsize ७.१.४५}\noindent तप्तनप्तनथनाश्च ।

\marginnote{\scriptsize ७.१.४६}\noindent इदन्तो मसि ।

\marginnote{\scriptsize ७.१.४७}\noindent क्त्वो यक् ।

\marginnote{\scriptsize ७.१.४८}\noindent इष्ट्वीनमिति च ।

\marginnote{\scriptsize ७.१.४९}\noindent स्नात्व्यादयश्च ।

\marginnote{\scriptsize ७.१.५०}\noindent आज्जसेरसुक् ।

\marginnote{\scriptsize ७.१.५१}\noindent अश्वक्षीरवृषलवणानामात्मप्रीतौ क्यचि ।

\marginnote{\scriptsize ७.१.५२}\noindent आमि सर्वनाम्नः सुट् ।

\marginnote{\scriptsize ७.१.५३}\noindent त्रेस्त्रयः ।

\marginnote{\scriptsize ७.१.५४}\noindent ह्रस्वनद्यापो नुट् ।

\marginnote{\scriptsize ७.१.५५}\noindent षट्चतुर्भ्यश्च ।

\marginnote{\scriptsize ७.१.५६}\noindent श्रीग्रामण्योश्छन्दसि ।

\marginnote{\scriptsize ७.१.५७}\noindent गोः पादान्ते ।

\marginnote{\scriptsize ७.१.५८}\noindent इदितो नुम् धातोः ।

\marginnote{\scriptsize ७.१.५९}\noindent शे मुचादीनाम् ।

\marginnote{\scriptsize ७.१.६०}\noindent मस्जिनशोर्झलि ।

\marginnote{\scriptsize ७.१.६१}\noindent रधिजभोरचि ।

\marginnote{\scriptsize ७.१.६२}\noindent नेट्यलिटि रधेः ।

\marginnote{\scriptsize ७.१.६३}\noindent रभेरशब्लिटोः ।

\marginnote{\scriptsize ७.१.६४}\noindent लभेश्च ।

\marginnote{\scriptsize ७.१.६५}\noindent आङो यि ।

\marginnote{\scriptsize ७.१.६६}\noindent उपात् प्रशंसायाम् ।

\marginnote{\scriptsize ७.१.६७}\noindent उपसर्गात् खल्घञोः ।

\marginnote{\scriptsize ७.१.६८}\noindent न सुदुर्भ्यां केवलाभ्याम् ।

\marginnote{\scriptsize ७.१.६९}\noindent विभाषा चिण्णमुलोः ।

\marginnote{\scriptsize ७.१.७०}\noindent उगिदचां सर्वनामस्थानेऽधातोः ।

\marginnote{\scriptsize ७.१.७१}\noindent युजेरसमासे ।

\marginnote{\scriptsize ७.१.७२}\noindent नपुंसकस्य झलचः ।

\marginnote{\scriptsize ७.१.७३}\noindent इकोऽचि विभक्तौ ।

\marginnote{\scriptsize ७.१.७४}\noindent तृतीयाऽऽदिषु भाषितपुंस्कं पुंवद्गालवस्य ।

\marginnote{\scriptsize ७.१.७५}\noindent अस्थिदधिसक्थ्यक्ष्णामनङुदात्तः ।

\marginnote{\scriptsize ७.१.७६}\noindent छन्दस्यपि दृश्यते ।

\marginnote{\scriptsize ७.१.७७}\noindent ई च द्विवचने ।

\marginnote{\scriptsize ७.१.७८}\noindent नाभ्यस्ताच्छतुः ।

\marginnote{\scriptsize ७.१.७९}\noindent वा नपुंसकस्य ।

\marginnote{\scriptsize ७.१.८०}\noindent आच्छीनद्योर्नुम् ।

\marginnote{\scriptsize ७.१.८१}\noindent शप्श्यनोर्नित्यम् ।

\marginnote{\scriptsize ७.१.८२}\noindent सावनडुहः ।

\marginnote{\scriptsize ७.१.८३}\noindent दृक्स्ववस्स्वतवसां छन्दसि ।

\marginnote{\scriptsize ७.१.८४}\noindent दिव औत् ।

\marginnote{\scriptsize ७.१.८५}\noindent पथिमथ्यृभुक्षामात् ।

\marginnote{\scriptsize ७.१.८६}\noindent इतोऽत् सर्वनामस्थाने ।

\marginnote{\scriptsize ७.१.८७}\noindent थो न्थः ।

\marginnote{\scriptsize ७.१.८८}\noindent भस्य टेर्लोपः ।

\marginnote{\scriptsize ७.१.८९}\noindent पुंसोऽसुङ् ।

\marginnote{\scriptsize ७.१.९०}\noindent गोतो णित् ।

\marginnote{\scriptsize ७.१.९१}\noindent णलुत्तमो वा ।

\marginnote{\scriptsize ७.१.९२}\noindent सख्युरसम्बुद्धौ ।

\marginnote{\scriptsize ७.१.९३}\noindent अनङ् सौ ।

\marginnote{\scriptsize ७.१.९४}\noindent ऋदुशनस्पुरुदंशोऽनेहसां च ।

\marginnote{\scriptsize ७.१.९५}\noindent तृज्वत् क्रोष्टुः ।

\marginnote{\scriptsize ७.१.९६}\noindent स्त्रियां च ।

\marginnote{\scriptsize ७.१.९७}\noindent विभाषा तृतीयाऽऽदिष्वचि ।

\marginnote{\scriptsize ७.१.९८}\noindent चतुरनडुहोरामुदात्तः ।

\marginnote{\scriptsize ७.१.९९}\noindent अम् सम्बुद्धौ ।

\marginnote{\scriptsize ७.१.१००}\noindent ॠत इद्धातोः ।

\marginnote{\scriptsize ७.१.१०१}\noindent उपधायाश्च ।

\marginnote{\scriptsize ७.१.१०२}\noindent उदोष्ठ्यपूर्वस्य ।

\marginnote{\scriptsize ७.१.१०३}\noindent बहुलं छन्दसि ।

\marginnote{\scriptsize ७.२.१}\noindent सिचि वृद्धिः परस्मैपदेषु ।

\marginnote{\scriptsize ७.२.२}\noindent अतो र्लान्तस्य ।

\marginnote{\scriptsize ७.२.३}\noindent वदव्रजहलन्तस्याचः ।

\marginnote{\scriptsize ७.२.४}\noindent नेटि ।

\marginnote{\scriptsize ७.२.५}\noindent ह्म्यन्तक्षणश्वसजागृणिश्व्येदिताम् ।

\marginnote{\scriptsize ७.२.६}\noindent ऊर्णोतेर्विभाषा ।

\marginnote{\scriptsize ७.२.७}\noindent अतो हलादेर्लघोः ।

\marginnote{\scriptsize ७.२.८}\noindent नेड् वशि कृति ।

\marginnote{\scriptsize ७.२.९}\noindent तितुत्रतथसिसुसरकसेषु च ।

\marginnote{\scriptsize ७.२.१०}\noindent एकाच उपदेशेऽनुदात्तात् ।

\marginnote{\scriptsize ७.२.११}\noindent श्र्युकः किति ।

\marginnote{\scriptsize ७.२.१२}\noindent सनि ग्रहगुहोश्च ।

\marginnote{\scriptsize ७.२.१३}\noindent कृसृभृवृस्तुद्रुस्रुश्रुवो लिटि ।

\marginnote{\scriptsize ७.२.१४}\noindent श्वीदितो निष्ठायाम् ।

\marginnote{\scriptsize ७.२.१५}\noindent यस्य विभाषा ।

\marginnote{\scriptsize ७.२.१६}\noindent आदितश्च ।

\marginnote{\scriptsize ७.२.१७}\noindent विभाषा भावादिकर्मणोः ।

\marginnote{\scriptsize ७.२.१८}\noindent क्षुब्धस्वान्तध्वान्तलग्नम्लिष्टविरिब्धफाण्टबाढानि मन्थमनस्तमःसक्ताविस्पष्टस्वरानायासभृशेषु ।

\marginnote{\scriptsize ७.२.१९}\noindent धृषिशसी वैयात्ये ।

\marginnote{\scriptsize ७.२.२०}\noindent दृढः स्थूलबलयोः ।

\marginnote{\scriptsize ७.२.२१}\noindent प्रभौ परिवृढः ।

\marginnote{\scriptsize ७.२.२२}\noindent कृच्छ्रगहनयोः कषः ।

\marginnote{\scriptsize ७.२.२३}\noindent घुषिरविशब्दने ।

\marginnote{\scriptsize ७.२.२४}\noindent अर्देः संनिविभ्यः ।

\marginnote{\scriptsize ७.२.२५}\noindent अभेश्चाविदूर्ये ।

\marginnote{\scriptsize ७.२.२६}\noindent णेरध्ययने वृत्तम् ।

\marginnote{\scriptsize ७.२.२७}\noindent वा दान्तशान्तपूर्णदस्तस्पष्टच्छन्नज्ञप्ताः ।

\marginnote{\scriptsize ७.२.२८}\noindent रुष्यमत्वरसंघुषास्वनाम् ।

\marginnote{\scriptsize ७.२.२९}\noindent हृषेर्लोमसु ।

\marginnote{\scriptsize ७.२.३०}\noindent अपचितश्च ।

\marginnote{\scriptsize ७.२.३१}\noindent ह्रु ह्वरेश्छन्दसि ।

\marginnote{\scriptsize ७.२.३२}\noindent अपरिह्वृताश्च ।

\marginnote{\scriptsize ७.२.३३}\noindent सोमे ह्वरितः ।

\marginnote{\scriptsize ७.२.३४}\noindent ग्रसितस्कभितस्तभितोत्तभितचत्तविकस्तविशस्तॄ- शंस्तृशास्तृतरुतृतरूतृवरुतृवरूतृवरुत्रीरुज्ज्वलितिक्षरिति- क्षमितिवमित्यमितीति च ।

\marginnote{\scriptsize ७.२.३५}\noindent आर्धधातुकस्येड् वलादेः ।

\marginnote{\scriptsize ७.२.३६}\noindent स्नुक्रमोरनात्मनेपदनिमित्ते ।

\marginnote{\scriptsize ७.२.३७}\noindent ग्रहोऽलिटि दीर्घः ।

\marginnote{\scriptsize ७.२.३८}\noindent वॄतो वा ।

\marginnote{\scriptsize ७.२.३९}\noindent न लिङि ।

\marginnote{\scriptsize ७.२.४०}\noindent सिचि च परस्मैपदेषु ।

\marginnote{\scriptsize ७.२.४१}\noindent इट् सनि वा ।

\marginnote{\scriptsize ७.२.४२}\noindent लिङ्सिचोरात्मनेपदेषु ।

\marginnote{\scriptsize ७.२.४३}\noindent ऋतश्च संयोगादेः ।

\marginnote{\scriptsize ७.२.४४}\noindent स्वरतिसूतिसूयतिधूञूदितो वा ।

\marginnote{\scriptsize ७.२.४५}\noindent रधादिभ्यश्च ।

\marginnote{\scriptsize ७.२.४६}\noindent निरः कुषः ।

\marginnote{\scriptsize ७.२.४७}\noindent इण्निष्ठायाम् ।

\marginnote{\scriptsize ७.२.४८}\noindent तीषसहलुभरुषरिषः ।

\marginnote{\scriptsize ७.२.४९}\noindent सनीवन्तर्धभ्रस्जदम्भुश्रिस्वृयूर्णुभरज्ञपिसनाम् ।

\marginnote{\scriptsize ७.२.५०}\noindent क्लिशः क्त्वानिष्ठयोः ।

\marginnote{\scriptsize ७.२.५१}\noindent पूङश्च ।

\marginnote{\scriptsize ७.२.५२}\noindent वसतिक्षुधोरिट् ।

\marginnote{\scriptsize ७.२.५३}\noindent अञ्चेः पूजायाम् ।

\marginnote{\scriptsize ७.२.५४}\noindent लुभो विमोचने ।

\marginnote{\scriptsize ७.२.५५}\noindent जॄव्रश्च्योः क्त्वि ।

\marginnote{\scriptsize ७.२.५६}\noindent उदितो वा ।

\marginnote{\scriptsize ७.२.५७}\noindent सेऽसिचि कृतचृतच्छृदतृदनृतः ।

\marginnote{\scriptsize ७.२.५८}\noindent गमेरिट् परस्मैपदेषु ।

\marginnote{\scriptsize ७.२.५९}\noindent न वृद्भ्यश्चतुर्भ्यः ।

\marginnote{\scriptsize ७.२.६०}\noindent तासि च कॢपः ।

\marginnote{\scriptsize ७.२.६१}\noindent अचस्तास्वत् थल्यनिटो नित्यम् ।

\marginnote{\scriptsize ७.२.६२}\noindent उपदेशेऽत्वतः ।

\marginnote{\scriptsize ७.२.६३}\noindent ऋतो भारद्वाजस्य ।

\marginnote{\scriptsize ७.२.६४}\noindent बभूथाततन्थजगृम्भववर्थेति निगमे ।

\marginnote{\scriptsize ७.२.६५}\noindent विभाषा सृजिदृषोः ।

\marginnote{\scriptsize ७.२.६६}\noindent इडत्त्यर्तिव्ययतीनाम् ।

\marginnote{\scriptsize ७.२.६७}\noindent वस्वेकाजाद्घसाम् ।

\marginnote{\scriptsize ७.२.६८}\noindent विभाषा गमहनविदविशाम् ।

\marginnote{\scriptsize ७.२.६९}\noindent सनिंससनिवांसम् ।

\marginnote{\scriptsize ७.२.७०}\noindent ऋद्धनोः स्ये ।

\marginnote{\scriptsize ७.२.७१}\noindent अञ्जेः सिचि ।

\marginnote{\scriptsize ७.२.७२}\noindent स्तुसुधूञ्भ्यः परस्मैपदेषु ।

\marginnote{\scriptsize ७.२.७३}\noindent यमरमनमातां सक् च ।

\marginnote{\scriptsize ७.२.७४}\noindent स्मिपूङ्रञ्ज्वशां सनि ।

\marginnote{\scriptsize ७.२.७५}\noindent किरश्च पञ्चभ्यः ।

\marginnote{\scriptsize ७.२.७६}\noindent रुदादिभ्यः सार्वधतुके ।

\marginnote{\scriptsize ७.२.७७}\noindent ईशः से ।

\marginnote{\scriptsize ७.२.७८}\noindent ईडजनोर्ध्वे च ।

\marginnote{\scriptsize ७.२.७९}\noindent लिङः सलोपोऽनन्त्यस्य ।

\marginnote{\scriptsize ७.२.८०}\noindent अतो येयः ।

\marginnote{\scriptsize ७.२.८१}\noindent आतो ङितः ।

\marginnote{\scriptsize ७.२.८२}\noindent आने मुक् ।

\marginnote{\scriptsize ७.२.८३}\noindent ईदासः ।

\marginnote{\scriptsize ७.२.८४}\noindent अष्टन आ विभक्तौ ।

\marginnote{\scriptsize ७.२.८५}\noindent रायो हलि ।

\marginnote{\scriptsize ७.२.८६}\noindent युष्मदस्मदोरनादेशे ।

\marginnote{\scriptsize ७.२.८७}\noindent द्वितीयायां च ।

\marginnote{\scriptsize ७.२.८८}\noindent प्रथमायाश्च द्विवचने भाषायाम् ।

\marginnote{\scriptsize ७.२.८९}\noindent योऽचि ।

\marginnote{\scriptsize ७.२.९०}\noindent शेषे लोपः ।

\marginnote{\scriptsize ७.२.९१}\noindent मपर्यन्तस्य ।

\marginnote{\scriptsize ७.२.९२}\noindent युवावौ द्विवचने ।

\marginnote{\scriptsize ७.२.९३}\noindent यूयवयौ जसि ।

\marginnote{\scriptsize ७.२.९४}\noindent त्वाहौ सौ ।

\marginnote{\scriptsize ७.२.९५}\noindent तुभ्यमह्यौ ङयि ।

\marginnote{\scriptsize ७.२.९६}\noindent तवममौ ङसि ।

\marginnote{\scriptsize ७.२.९७}\noindent त्वमावेकवचने ।

\marginnote{\scriptsize ७.२.९८}\noindent प्रतयोत्तरपदयोश्च ।

\marginnote{\scriptsize ७.२.९९}\noindent त्रिचतुरोः स्त्रियां तिसृचतसृ ।

\marginnote{\scriptsize ७.२.१००}\noindent अचि र ऋतः ।

\marginnote{\scriptsize ७.२.१०१}\noindent जराया जरसन्यतरस्याम् ।

\marginnote{\scriptsize ७.२.१०२}\noindent त्यदादीनामः ।

\marginnote{\scriptsize ७.२.१०३}\noindent किमः कः ।

\marginnote{\scriptsize ७.२.१०४}\noindent कु तिहोः ।

\marginnote{\scriptsize ७.२.१०५}\noindent क्वाति ।

\marginnote{\scriptsize ७.२.१०६}\noindent तदोः सः सावनन्त्ययोः ।

\marginnote{\scriptsize ७.२.१०७}\noindent अदस औ सुलोपश्च ।

\marginnote{\scriptsize ७.२.१०८}\noindent इदमो मः ।

\marginnote{\scriptsize ७.२.१०९}\noindent दश्च ।

\marginnote{\scriptsize ७.२.११०}\noindent यः सौ ।

\marginnote{\scriptsize ७.२.१११}\noindent इदोऽय् पुंसि ।

\marginnote{\scriptsize ७.२.११२}\noindent अनाप्यकः ।

\marginnote{\scriptsize ७.२.११३}\noindent हलि लोपः ।

\marginnote{\scriptsize ७.२.११४}\noindent मृजेर्वृद्धिः ।

\marginnote{\scriptsize ७.२.११५}\noindent अचो ञ्णिति ।

\marginnote{\scriptsize ७.२.११६}\noindent अत उपधायाः ।

\marginnote{\scriptsize ७.२.११७}\noindent तद्धितेष्वचामादेः ।

\marginnote{\scriptsize ७.२.११८}\noindent किति च ।

\marginnote{\scriptsize ७.३.१}\noindent देविकाशिंशपादित्यवाड्दीर्घसत्रश्रेयसामात् ।

\marginnote{\scriptsize ७.३.२}\noindent केकयमित्त्रयुप्रलयानां यादेरियः ।

\marginnote{\scriptsize ७.३.३}\noindent न य्वाभ्यां पदान्ताभ्याम् पूर्वौ तु ताभ्यामैच् ।

\marginnote{\scriptsize ७.३.४}\noindent द्वारादीनां च ।

\marginnote{\scriptsize ७.३.५}\noindent न्यग्रोधस्य च केवलस्य ।

\marginnote{\scriptsize ७.३.६}\noindent न कर्मव्यतिहारे ।

\marginnote{\scriptsize ७.३.७}\noindent स्वागतादीनां च ।

\marginnote{\scriptsize ७.३.८}\noindent श्वादेरिञि ।

\marginnote{\scriptsize ७.३.९}\noindent पदान्तस्यान्यतरस्याम् ।

\marginnote{\scriptsize ७.३.१०}\noindent उत्तरपदस्य ।

\marginnote{\scriptsize ७.३.११}\noindent अवयवादृतोः ।

\marginnote{\scriptsize ७.३.१२}\noindent सुसर्वार्धाज्जनपदस्य ।

\marginnote{\scriptsize ७.३.१३}\noindent दिशोऽमद्राणाम् ।

\marginnote{\scriptsize ७.३.१४}\noindent प्राचां ग्रामनगराणाम् ।

\marginnote{\scriptsize ७.३.१५}\noindent संख्यायाः संवत्सरसंख्यस्य च ।

\marginnote{\scriptsize ७.३.१६}\noindent वर्षस्याभविष्यति ।

\marginnote{\scriptsize ७.३.१७}\noindent परिमाणान्तस्यासंज्ञाशाणयोः ।

\marginnote{\scriptsize ७.३.१८}\noindent जे प्रोष्ठपदानाम् ।

\marginnote{\scriptsize ७.३.१९}\noindent हृद्भगसिन्ध्वन्ते पूर्वपदस्य च ।

\marginnote{\scriptsize ७.३.२०}\noindent अनुशतिकादीनां च ।

\marginnote{\scriptsize ७.३.२१}\noindent देवताद्वंद्वे च ।

\marginnote{\scriptsize ७.३.२२}\noindent नेन्द्रस्य परस्य ।

\marginnote{\scriptsize ७.३.२३}\noindent दीर्घाच्च वरुणस्य ।

\marginnote{\scriptsize ७.३.२४}\noindent प्राचां नगरान्ते ।

\marginnote{\scriptsize ७.३.२५}\noindent जङ्गलधेनुवलजान्तस्य विभाषितमुत्तरम् ।

\marginnote{\scriptsize ७.३.२६}\noindent अर्धात् परिमाणस्य पूर्वस्य तु वा ।

\marginnote{\scriptsize ७.३.२७}\noindent नातः परस्य ।

\marginnote{\scriptsize ७.३.२८}\noindent प्रवाहणस्य ढे ।

\marginnote{\scriptsize ७.३.२९}\noindent तत्प्रत्ययस्य च ।

\marginnote{\scriptsize ७.३.३०}\noindent नञः शुचीश्वरक्षेत्रज्ञकुशलनिपुणानाम् ।

\marginnote{\scriptsize ७.३.३१}\noindent यथातथयथापुरयोः पर्यायेण ।

\marginnote{\scriptsize ७.३.३२}\noindent हनस्तोऽचिण्णलोः ।

\marginnote{\scriptsize ७.३.३३}\noindent आतो युक् चिङ्कृतोः ।

\marginnote{\scriptsize ७.३.३४}\noindent नोदात्तोपदेशस्य मान्तस्यानाचमेः ।

\marginnote{\scriptsize ७.३.३५}\noindent जनिवध्योश्च ।

\marginnote{\scriptsize ७.३.३६}\noindent अर्त्तिह्रीब्लीरीक्नूयीक्ष्माय्यातां पुङ्णौ ।

\marginnote{\scriptsize ७.३.३७}\noindent शाच्छासाह्वाव्यावेपां युक् ।

\marginnote{\scriptsize ७.३.३८}\noindent वो विधूनने जुक् ।

\marginnote{\scriptsize ७.३.३९}\noindent लीलोर्नुग्लुकावन्यतरस्यां स्नेहविपातने ।

\marginnote{\scriptsize ७.३.४०}\noindent भियो हेतुभये षुक् ।

\marginnote{\scriptsize ७.३.४१}\noindent स्फायो वः ।

\marginnote{\scriptsize ७.३.४२}\noindent शदेरगतौ तः ।

\marginnote{\scriptsize ७.३.४३}\noindent रुहः पोऽन्यतरस्याम् ।

\marginnote{\scriptsize ७.३.४४}\noindent प्रत्ययस्थात् कात् पूर्वस्यात इदाप्यसुपः ।

\marginnote{\scriptsize ७.३.४५}\noindent न यासयोः ।

\marginnote{\scriptsize ७.३.४६}\noindent उदीचामातः स्थाने यकपूर्वायाः ।

\marginnote{\scriptsize ७.३.४७}\noindent भस्त्रैषाऽजाज्ञाद्वास्वानञ्पूर्वाणामपि ।

\marginnote{\scriptsize ७.३.४८}\noindent अभाषितपुंस्काच्च ।

\marginnote{\scriptsize ७.३.४९}\noindent आदाचार्याणाम् ।

\marginnote{\scriptsize ७.३.५०}\noindent ठस्येकः ।

\marginnote{\scriptsize ७.३.५१}\noindent इसुसुक्तान्तात् कः ।

\marginnote{\scriptsize ७.३.५२}\noindent चजोः कु घिन्ण्यतोः ।

\marginnote{\scriptsize ७.३.५३}\noindent न्यङ्क्वादीनां च ।

\marginnote{\scriptsize ७.३.५४}\noindent हो हन्तेर्ञ्णिन्नेषु ।

\marginnote{\scriptsize ७.३.५५}\noindent अभ्यासाच्च ।

\marginnote{\scriptsize ७.३.५६}\noindent हेरचङि ।

\marginnote{\scriptsize ७.३.५७}\noindent सन्लिटोर्जेः ।

\marginnote{\scriptsize ७.३.५८}\noindent विभाषा चेः ।

\marginnote{\scriptsize ७.३.५९}\noindent न क्वादेः ।

\marginnote{\scriptsize ७.३.६०}\noindent अजिवृज्योश्च ।

\marginnote{\scriptsize ७.३.६१}\noindent भुजन्युब्जौ पाण्युपतापयोः ।

\marginnote{\scriptsize ७.३.६२}\noindent प्रयाजानुयाजौ यज्ञाङ्गे ।

\marginnote{\scriptsize ७.३.६३}\noindent वञ्चेर्गतौ ।

\marginnote{\scriptsize ७.३.६४}\noindent ओक उचः के ।

\marginnote{\scriptsize ७.३.६५}\noindent ण्य आवश्यके ।

\marginnote{\scriptsize ७.३.६६}\noindent यजयाचरुचप्रवचर्चश्च ।

\marginnote{\scriptsize ७.३.६७}\noindent वचोऽशब्दसंज्ञायाम् ।

\marginnote{\scriptsize ७.३.६८}\noindent प्रयोज्यनियोज्यौ शक्यार्थे ।

\marginnote{\scriptsize ७.३.६९}\noindent भोज्यं भक्ष्ये ।

\marginnote{\scriptsize ७.३.७०}\noindent घोर्लोपो लेटि वा ।

\marginnote{\scriptsize ७.३.७१}\noindent ओतः श्यनि ।

\marginnote{\scriptsize ७.३.७२}\noindent क्सस्याचि ।

\marginnote{\scriptsize ७.३.७३}\noindent लुग्वा दुहदिहलिहगुहामात्मनेपदे दन्त्ये ।

\marginnote{\scriptsize ७.३.७४}\noindent शमामष्टानां दीर्घः श्यनि ।

\marginnote{\scriptsize ७.३.७५}\noindent ष्ठिवुक्लम्याचमां शिति ।

\marginnote{\scriptsize ७.३.७६}\noindent क्रमः परस्मैपदेषु ।

\marginnote{\scriptsize ७.३.७७}\noindent इषुगमियमां छः ।

\marginnote{\scriptsize ७.३.७८}\noindent पाघ्राध्मास्थाम्नादाण्दृश्यर्त्तिसर्त्तिशदसदां पिबजिघ्रधमतिष्ठमनयच्छपश्यर्च्छधौशीयसीदाः ।

\marginnote{\scriptsize ७.३.७९}\noindent ज्ञाजनोर्जा ।

\marginnote{\scriptsize ७.३.८०}\noindent प्वादीनां ह्रस्वः ।

\marginnote{\scriptsize ७.३.८१}\noindent मीनातेर्निगमे ।

\marginnote{\scriptsize ७.३.८२}\noindent मिदेर्गुणः ।

\marginnote{\scriptsize ७.३.८३}\noindent जुसि च ।

\marginnote{\scriptsize ७.३.८४}\noindent सार्वधातुकार्धधातुकयोः ।

\marginnote{\scriptsize ७.३.८५}\noindent जाग्रोऽविचिण्णल्ङित्सु ।

\marginnote{\scriptsize ७.३.८६}\noindent पुगन्तलघूपधस्य च ।

\marginnote{\scriptsize ७.३.८७}\noindent नाभ्यस्तस्याचि पिति सार्वधातुके ।

\marginnote{\scriptsize ७.३.८८}\noindent भूसुवोस्तिङि ।

\marginnote{\scriptsize ७.३.८९}\noindent उतो वृद्धिर्लुकि हलि ।

\marginnote{\scriptsize ७.३.९०}\noindent ऊर्णोतेर्विभाषा ।

\marginnote{\scriptsize ७.३.९१}\noindent गुणोऽपृक्ते ।

\marginnote{\scriptsize ७.३.९२}\noindent तृणह इम् ।

\marginnote{\scriptsize ७.३.९३}\noindent ब्रुव ईट् ।

\marginnote{\scriptsize ७.३.९४}\noindent यङो वा ।

\marginnote{\scriptsize ७.३.९५}\noindent तुरुस्तुशम्यमः सार्वधातुके ।

\marginnote{\scriptsize ७.३.९६}\noindent अस्तिसिचोऽपृक्ते ।

\marginnote{\scriptsize ७.३.९७}\noindent बहुलं छन्दसि ।

\marginnote{\scriptsize ७.३.९८}\noindent रुदश्च पञ्चभ्यः ।

\marginnote{\scriptsize ७.३.९९}\noindent अड्गार्ग्यगालवयोः ।

\marginnote{\scriptsize ७.३.१००}\noindent अदः सर्वेषाम् ।

\marginnote{\scriptsize ७.३.१०१}\noindent अतो दीर्घो यञि ।

\marginnote{\scriptsize ७.३.१०२}\noindent सुपि च ।

\marginnote{\scriptsize ७.३.१०३}\noindent बहुवचने झल्येत् ।

\marginnote{\scriptsize ७.३.१०४}\noindent ओऽसि च ।

\marginnote{\scriptsize ७.३.१०५}\noindent आङि चापः ।

\marginnote{\scriptsize ७.३.१०६}\noindent सम्बुद्धौ च ।

\marginnote{\scriptsize ७.३.१०७}\noindent अम्बाऽर्थनद्योर्ह्रस्वः ।

\marginnote{\scriptsize ७.३.१०८}\noindent ह्रस्वस्य गुणः ।

\marginnote{\scriptsize ७.३.१०९}\noindent जसि च ।

\marginnote{\scriptsize ७.३.११०}\noindent ऋतो ङिसर्वनामस्थानयोः ।

\marginnote{\scriptsize ७.३.१११}\noindent घेर्ङिति ।

\marginnote{\scriptsize ७.३.११२}\noindent आण्नद्याः ।

\marginnote{\scriptsize ७.३.११३}\noindent याडापः ।

\marginnote{\scriptsize ७.३.११४}\noindent सर्वनाम्नः स्याड्ढ्रस्वश्च ।

\marginnote{\scriptsize ७.३.११५}\noindent विभाषा द्वितीयातृतीयाभ्याम् ।

\marginnote{\scriptsize ७.३.११६}\noindent ङेराम्नद्याम्नीभ्यः ।

\marginnote{\scriptsize ७.३.११७}\noindent इदुद्भ्याम् ।

\marginnote{\scriptsize ७.३.११८}\noindent औत् ।

\marginnote{\scriptsize ७.३.११९}\noindent अच्च घेः ।

\marginnote{\scriptsize ७.३.१२०}\noindent आङो नाऽस्त्रियाम् ।

\marginnote{\scriptsize ७.४.१}\noindent णौ चङ्युपधाया ह्रस्वः ।

\marginnote{\scriptsize ७.४.२}\noindent नाग्लोपिशास्वृदिताम् ।

\marginnote{\scriptsize ७.४.३}\noindent भ्राजभासभाषदीपजीवमीलपीडामन्यतरस्याम् ।

\marginnote{\scriptsize ७.४.४}\noindent लोपः पिबतेरीच्चाभ्यासस्य ।

\marginnote{\scriptsize ७.४.५}\noindent तिष्ठतेरित् ।

\marginnote{\scriptsize ७.४.६}\noindent जिघ्रतेर्वा ।

\marginnote{\scriptsize ७.४.७}\noindent उरृत् ।

\marginnote{\scriptsize ७.४.८}\noindent नित्यं छन्दसि ।

\marginnote{\scriptsize ७.४.९}\noindent दयतेर्दिगि लिटि ।

\marginnote{\scriptsize ७.४.१०}\noindent ऋतश्च संयोगादेर्गुणः ।

\marginnote{\scriptsize ७.४.११}\noindent ऋच्छत्यॄताम् ।

\marginnote{\scriptsize ७.४.१२}\noindent श‍ृदॄप्रां ह्रस्वो वा ।

\marginnote{\scriptsize ७.४.१३}\noindent केऽणः ।

\marginnote{\scriptsize ७.४.१४}\noindent न कपि ।

\marginnote{\scriptsize ७.४.१५}\noindent आपोऽन्यतरस्याम् ।

\marginnote{\scriptsize ७.४.१६}\noindent ऋदृशोऽङि गुणः ।

\marginnote{\scriptsize ७.४.१७}\noindent अस्यतेस्थुक् ।

\marginnote{\scriptsize ७.४.१८}\noindent श्वयतेरः ।

\marginnote{\scriptsize ७.४.१९}\noindent पतः पुम् ।

\marginnote{\scriptsize ७.४.२०}\noindent वच उम् ।

\marginnote{\scriptsize ७.४.२१}\noindent शीङः सार्वधातुके गुणः ।

\marginnote{\scriptsize ७.४.२२}\noindent अयङ् यि क्ङिति ।

\marginnote{\scriptsize ७.४.२३}\noindent उपसर्गाद्ध्रस्व ऊहतेः ।

\marginnote{\scriptsize ७.४.२४}\noindent एतेर्लिङि ।

\marginnote{\scriptsize ७.४.२५}\noindent अकृत्सार्वधातुकयोर्दीर्घः ।

\marginnote{\scriptsize ७.४.२६}\noindent च्वौ च ।

\marginnote{\scriptsize ७.४.२७}\noindent रीङ् ऋतः ।

\marginnote{\scriptsize ७.४.२८}\noindent रिङ् शयग्लिङ्क्षु ।

\marginnote{\scriptsize ७.४.२९}\noindent गुणोऽर्तिसंयोगाद्योः ।

\marginnote{\scriptsize ७.४.३०}\noindent यङि च ।

\marginnote{\scriptsize ७.४.३१}\noindent ई घ्राध्मोः ।

\marginnote{\scriptsize ७.४.३२}\noindent अस्य च्वौ ।

\marginnote{\scriptsize ७.४.३३}\noindent क्यचि च ।

\marginnote{\scriptsize ७.४.३४}\noindent अशनायोदन्यधनाया बुभुक्षापिपासागर्द्धेषु ।

\marginnote{\scriptsize ७.४.३५}\noindent न च्छन्दस्यपुत्रस्य ।

\marginnote{\scriptsize ७.४.३६}\noindent दुरस्युर्द्रविणस्युर्वृषण्यतिरिषण्यति ।

\marginnote{\scriptsize ७.४.३७}\noindent अश्वाघस्यात् ।

\marginnote{\scriptsize ७.४.३८}\noindent देवसुम्नयोर्यजुषि काठके ।

\marginnote{\scriptsize ७.४.३९}\noindent कव्यध्वरपृतनस्यर्चि लोपः ।

\marginnote{\scriptsize ७.४.४०}\noindent द्यतिस्यतिमास्थामित्ति किति ।

\marginnote{\scriptsize ७.४.४१}\noindent शाछोरन्यतरस्याम् ।

\marginnote{\scriptsize ७.४.४२}\noindent दधातेर्हिः ।

\marginnote{\scriptsize ७.४.४३}\noindent जहातेश्च क्त्वि ।

\marginnote{\scriptsize ७.४.४४}\noindent विभाषा छन्दसि ।

\marginnote{\scriptsize ७.४.४५}\noindent सुधितवसुधितनेमधितधिष्वधिषीय च ।

\marginnote{\scriptsize ७.४.४६}\noindent दो दद् घोः ।

\marginnote{\scriptsize ७.४.४७}\noindent अच उपसर्गात्तः ।

\marginnote{\scriptsize ७.४.४८}\noindent अपो भि ।

\marginnote{\scriptsize ७.४.४९}\noindent सः स्यार्द्धधातुके ।

\marginnote{\scriptsize ७.४.५०}\noindent तासस्त्योर्लोपः ।

\marginnote{\scriptsize ७.४.५१}\noindent रि च ।

\marginnote{\scriptsize ७.४.५२}\noindent ह एति ।

\marginnote{\scriptsize ७.४.५३}\noindent यीवर्णयोर्दीधीवेव्योः ।

\marginnote{\scriptsize ७.४.५४}\noindent सनि मीमाघुरभलभशकपतपदामच इस् ।

\marginnote{\scriptsize ७.४.५५}\noindent आप्ज्ञप्यृधामीत् ।

\marginnote{\scriptsize ७.४.५६}\noindent दम्भ इच्च ।

\marginnote{\scriptsize ७.४.५७}\noindent मुचोऽकर्मकस्य गुणो वा ।

\marginnote{\scriptsize ७.४.५८}\noindent अत्र लोपोऽभ्यासस्य ।

\marginnote{\scriptsize ७.४.५९}\noindent ह्रस्वः ।

\marginnote{\scriptsize ७.४.६०}\noindent हलादिः शेषः ।

\marginnote{\scriptsize ७.४.६१}\noindent शर्पूर्वाः खयः ।

\marginnote{\scriptsize ७.४.६२}\noindent कुहोश्चुः ।

\marginnote{\scriptsize ७.४.६३}\noindent न कवतेर्यङि ।

\marginnote{\scriptsize ७.४.६४}\noindent कृषेश्छन्दसि । ७.४.६५ दाधर्तिदर्धर्तिदर्धर्षिबोभूतुतेतिक्तेऽलर्ष्यापनीफणत्- संसनिष्यदत्करिक्रत्कनिक्रदद्भरिभ्रद्दविध्वतोदविद्युतत्- तरित्रतःसरीसृपतंवरीवृजन्मर्मृज्यागनीगन्तीति च ।

\marginnote{\scriptsize ७.४.६६}\noindent उरत् ।

\marginnote{\scriptsize ७.४.६७}\noindent द्युतिस्वाप्योः सम्प्रसारणम् ।

\marginnote{\scriptsize ७.४.६८}\noindent व्यथो लिटि ।

\marginnote{\scriptsize ७.४.६९}\noindent दीर्घ इणः किति ।

\marginnote{\scriptsize ७.४.७०}\noindent अतः आदेः ।

\marginnote{\scriptsize ७.४.७१}\noindent तस्मान्नुड् द्विहलः ।

\marginnote{\scriptsize ७.४.७२}\noindent अश्नोतेश्च ।

\marginnote{\scriptsize ७.४.७३}\noindent भवतेरः ।

\marginnote{\scriptsize ७.४.७४}\noindent ससूवेति निगमे ।

\marginnote{\scriptsize ७.४.७५}\noindent निजां त्रयाणां गुणः श्लौ ।

\marginnote{\scriptsize ७.४.७६}\noindent भृञामित् ।

\marginnote{\scriptsize ७.४.७७}\noindent अर्तिपिपर्त्योश्च ।

\marginnote{\scriptsize ७.४.७८}\noindent बहुलं छन्दसि ।

\marginnote{\scriptsize ७.४.७९}\noindent सन्यतः ।

\marginnote{\scriptsize ७.४.८०}\noindent ओः पुयण्ज्यपरे ।

\marginnote{\scriptsize ७.४.८१}\noindent स्रवतिश‍ृणोतिद्रवतिप्रवतिप्लवतिच्यवतीनां वा ।

\marginnote{\scriptsize ७.४.८२}\noindent गुणो यङ्लुकोः ।

\marginnote{\scriptsize ७.४.८३}\noindent दीर्घोऽकितः ।

\marginnote{\scriptsize ७.४.८४}\noindent नीग्वञ्चुस्रंसुध्वंसुभ्रंसुकसपतपदस्कन्दाम् ।

\marginnote{\scriptsize ७.४.८५}\noindent नुगतोऽनुनासिकान्तस्य ।

\marginnote{\scriptsize ७.४.८६}\noindent जपजभदहदशभञ्जपशां च ।

\marginnote{\scriptsize ७.४.८७}\noindent चरफलोश्च ।

\marginnote{\scriptsize ७.४.८८}\noindent उत् परस्यातः ।

\marginnote{\scriptsize ७.४.८९}\noindent ति च ।

\marginnote{\scriptsize ७.४.९०}\noindent रीगृदुपधस्य च ।

\marginnote{\scriptsize ७.४.९१}\noindent रुग्रिकौ च लुकि ।

\marginnote{\scriptsize ७.४.९२}\noindent ऋतश्च ।

\marginnote{\scriptsize ७.४.९३}\noindent सन्वल्लघुनि चङ्परेऽनग्लोपे ।

\marginnote{\scriptsize ७.४.९४}\noindent दीर्घो लघोः ।

\marginnote{\scriptsize ७.४.९५}\noindent अत् स्मृदृत्वरप्रथम्रदस्तॄस्पशाम् ।

\marginnote{\scriptsize ७.४.९६}\noindent विभाषा वेष्टिचेष्ट्योः ।

\marginnote{\scriptsize ७.४.९७}\noindent ई च गणः ।

\section*{॥ अध्याय ८ ॥}

\marginnote{\scriptsize ८.१.१}\noindent सर्वस्य द्वे ।

\marginnote{\scriptsize ८.१.२}\noindent तस्य परमाम्रेडितम् ।

\marginnote{\scriptsize ८.१.३}\noindent अनुदात्तं च ।

\marginnote{\scriptsize ८.१.४}\noindent नित्यवीप्सयोः ।

\marginnote{\scriptsize ८.१.५}\noindent परेर्वर्जने ।

\marginnote{\scriptsize ८.१.६}\noindent प्रसमुपोदः पादपूरणे ।

\marginnote{\scriptsize ८.१.७}\noindent उपर्यध्यधसः सामीप्ये ।

\marginnote{\scriptsize ८.१.८}\noindent वाक्यादेरामन्त्रितस्यासूयासम्मतिकोपकुत्सनभर्त्सनेषु ।

\marginnote{\scriptsize ८.१.९}\noindent एकं बहुव्रीहिवत् ।

\marginnote{\scriptsize ८.१.१०}\noindent आबाधे च ।

\marginnote{\scriptsize ८.१.११}\noindent कर्मधारयवत् उत्तरेषु ।

\marginnote{\scriptsize ८.१.१२}\noindent प्रकारे गुणवचनस्य ।

\marginnote{\scriptsize ८.१.१३}\noindent अकृच्छ्रे प्रियसुखयोरन्यतरस्याम् ।

\marginnote{\scriptsize ८.१.१४}\noindent यथास्वे यथायथम् ।

\marginnote{\scriptsize ८.१.१५}\noindent द्वन्द्वं रहस्यमर्यादावचनव्युत्क्रमण- यज्ञपात्रप्रयोगाभिव्यक्तिषु ।

\marginnote{\scriptsize ८.१.१६}\noindent पदस्य ।

\marginnote{\scriptsize ८.१.१७}\noindent पदात् ।

\marginnote{\scriptsize ८.१.१८}\noindent अनुदात्तं सर्वमपादादौ ।

\marginnote{\scriptsize ८.१.१९}\noindent आमन्त्रितस्य च ।

\marginnote{\scriptsize ८.१.२०}\noindent युष्मदस्मदोः षष्ठीचतुर्थीद्वितीयास्थयोर्वान्नावौ ।

\marginnote{\scriptsize ८.१.२१}\noindent बहुवचने वस्नसौ ।

\marginnote{\scriptsize ८.१.२२}\noindent तेमयावेकवचनस्य ।

\marginnote{\scriptsize ८.१.२३}\noindent त्वामौ द्वितीयायाः ।

\marginnote{\scriptsize ८.१.२४}\noindent न चवाहाहैवयुक्ते ।

\marginnote{\scriptsize ८.१.२५}\noindent पश्यार्थैश्चानालोचने ।

\marginnote{\scriptsize ८.१.२६}\noindent सपूर्वायाः प्रथमाया विभाषा ।

\marginnote{\scriptsize ८.१.२७}\noindent तिङो गोत्रादीनि कुत्सनाभीक्ष्ण्ययोः ।

\marginnote{\scriptsize ८.१.२८}\noindent तिङ्ङतिङः ।

\marginnote{\scriptsize ८.१.२९}\noindent न लुट् ।

\marginnote{\scriptsize ८.१.३०}\noindent निपातैर्यद्यदिहन्तकुविन्नेच्चेच्चङ्कच्चिद्यत्रयुक्तम् ।

\marginnote{\scriptsize ८.१.३१}\noindent नह प्रत्यारम्भे ।

\marginnote{\scriptsize ८.१.३२}\noindent सत्यं प्रश्ने ।

\marginnote{\scriptsize ८.१.३३}\noindent अङ्गाप्रातिलोम्ये ।

\marginnote{\scriptsize ८.१.३४}\noindent हि च ।

\marginnote{\scriptsize ८.१.३५}\noindent छन्दस्यनेकमपि साकाङ्क्षम् ।

\marginnote{\scriptsize ८.१.३६}\noindent यावद्यथाभ्याम् ।

\marginnote{\scriptsize ८.१.३७}\noindent पूजायां नानन्तरम् ।

\marginnote{\scriptsize ८.१.३८}\noindent उपसर्गव्यपेतं च ।

\marginnote{\scriptsize ८.१.३९}\noindent तुपश्यपश्यताहैः पूजायाम् ।

\marginnote{\scriptsize ८.१.४०}\noindent अहो च ।

\marginnote{\scriptsize ८.१.४१}\noindent शेषे विभाषा ।

\marginnote{\scriptsize ८.१.४२}\noindent पुरा च परीप्सायाम् ।

\marginnote{\scriptsize ८.१.४३}\noindent नन्वित्यनुज्ञैषणायाम् ।

\marginnote{\scriptsize ८.१.४४}\noindent किं क्रियाप्रश्नेऽनुपसर्गमप्रतिषिद्धम् ।

\marginnote{\scriptsize ८.१.४५}\noindent लोपे विभाषा ।

\marginnote{\scriptsize ८.१.४६}\noindent एहिमन्ये प्रहासे लृट् ।

\marginnote{\scriptsize ८.१.४७}\noindent जात्वपूर्वम् ।

\marginnote{\scriptsize ८.१.४८}\noindent किम्वृत्तं च चिदुत्तरम् ।

\marginnote{\scriptsize ८.१.४९}\noindent आहो उताहो चानन्तरम् ।

\marginnote{\scriptsize ८.१.५०}\noindent शेषे विभाषा ।

\marginnote{\scriptsize ८.१.५१}\noindent गत्यर्थलोटा लृण्न चेत् कारकं सर्वान्यत् ।

\marginnote{\scriptsize ८.१.५२}\noindent लोट् च ।

\marginnote{\scriptsize ८.१.५३}\noindent विभाषितं सोपसर्गमनुत्तमम् ।

\marginnote{\scriptsize ८.१.५४}\noindent हन्त च ।

\marginnote{\scriptsize ८.१.५५}\noindent आम एकान्तरमामन्त्रितमनन्तिके ।

\marginnote{\scriptsize ८.१.५६}\noindent यद्धितुपरं छन्दसि ।

\marginnote{\scriptsize ८.१.५७}\noindent चनचिदिवगोत्रादितद्धिताम्रेडितेष्वगतेः ।

\marginnote{\scriptsize ८.१.५८}\noindent चादिषु च ।

\marginnote{\scriptsize ८.१.५९}\noindent चवायोगे प्रथमा ।

\marginnote{\scriptsize ८.१.६०}\noindent हेति क्षियायाम् ।

\marginnote{\scriptsize ८.१.६१}\noindent अहेति विनियोगे च ।

\marginnote{\scriptsize ८.१.६२}\noindent चाहलोप एवेत्यवधारणम् ।

\marginnote{\scriptsize ८.१.६३}\noindent चादिलोपे विभाषा ।

\marginnote{\scriptsize ८.१.६४}\noindent वैवावेति च च्छन्दसि ।

\marginnote{\scriptsize ८.१.६५}\noindent एकान्याभ्यां समर्थाभ्याम् ।

\marginnote{\scriptsize ८.१.६६}\noindent यद्वृत्तान्नित्यं ।

\marginnote{\scriptsize ८.१.६७}\noindent पूजनात् पूजितमनुदात्तम् (काष्ठादिभ्यः) ।

\marginnote{\scriptsize ८.१.६८}\noindent सगतिरपि तिङ् ।

\marginnote{\scriptsize ८.१.६९}\noindent कुत्सने च सुप्यगोत्रादौ ।

\marginnote{\scriptsize ८.१.७०}\noindent गतिर्गतौ ।

\marginnote{\scriptsize ८.१.७१}\noindent तिङि चोदात्तवति ।

\marginnote{\scriptsize ८.१.७२}\noindent आमन्त्रितं पूर्वम् अविद्यमानवत् ।

\marginnote{\scriptsize ८.१.७३}\noindent नामन्त्रिते समानाधिकरणे (सामान्यवचनम्) ।

\marginnote{\scriptsize ८.१.७४}\noindent (सामान्यवचनं) विभाषितं विशेषवचने (बहुवचनम्) ।

\marginnote{\scriptsize ८.२.१}\noindent पूर्वत्रासिद्धम् ।

\marginnote{\scriptsize ८.२.२}\noindent नलोपः सुप्स्वरसंज्ञातुग्विधिषु कृति ।

\marginnote{\scriptsize ८.२.३}\noindent न मु ने ।

\marginnote{\scriptsize ८.२.४}\noindent उदात्तस्वरितयोर्यणः स्वरितोऽनुदात्तस्य ।

\marginnote{\scriptsize ८.२.५}\noindent एकादेश उदात्तेनोदात्तः ।

\marginnote{\scriptsize ८.२.६}\noindent स्वरितो वाऽनुदात्ते पदादौ ।

\marginnote{\scriptsize ८.२.७}\noindent नलोपः प्रातिपदिकान्तस्य ।

\marginnote{\scriptsize ८.२.८}\noindent न ङिसम्बुद्ध्योः ।

\marginnote{\scriptsize ८.२.९}\noindent मादुपधायाश्च मतोर्वोऽयवादिभ्यः ।

\marginnote{\scriptsize ८.२.१०}\noindent झयः ।

\marginnote{\scriptsize ८.२.११}\noindent संज्ञायाम् ।

\marginnote{\scriptsize ८.२.१२}\noindent आसन्दीवदष्ठीवच्चक्रीवत्कक्षीवद्रुमण्वच्चर्मण्वती ।

\marginnote{\scriptsize ८.२.१३}\noindent उदन्वानुदधौ च ।

\marginnote{\scriptsize ८.२.१४}\noindent राजन्वान् सौराज्ये ।

\marginnote{\scriptsize ८.२.१५}\noindent छन्दसीरः ।

\marginnote{\scriptsize ८.२.१६}\noindent अनो नुट् ।

\marginnote{\scriptsize ८.२.१७}\noindent नाद्घस्य ।

\marginnote{\scriptsize ८.२.१८}\noindent कृपो रो लः ।

\marginnote{\scriptsize ८.२.१९}\noindent उपसर्गस्यायतौ ।

\marginnote{\scriptsize ८.२.२०}\noindent ग्रो यङि ।

\marginnote{\scriptsize ८.२.२१}\noindent अचि विभाषा ।

\marginnote{\scriptsize ८.२.२२}\noindent परेश्च घाङ्कयोः ।

\marginnote{\scriptsize ८.२.२३}\noindent संयोगान्तस्य लोपः ।

\marginnote{\scriptsize ८.२.२४}\noindent रात् सस्य ।

\marginnote{\scriptsize ८.२.२५}\noindent धि च ।

\marginnote{\scriptsize ८.२.२६}\noindent झलो झलि ।

\marginnote{\scriptsize ८.२.२७}\noindent ह्रस्वादङ्गात् ।

\marginnote{\scriptsize ८.२.२८}\noindent इट ईटि ।

\marginnote{\scriptsize ८.२.२९}\noindent स्कोः संयोगाद्योरन्ते च ।

\marginnote{\scriptsize ८.२.३०}\noindent चोः कुः ।

\marginnote{\scriptsize ८.२.३१}\noindent हो ढः ।

\marginnote{\scriptsize ८.२.३२}\noindent दादेर्धातोर्घः ।

\marginnote{\scriptsize ८.२.३३}\noindent वा द्रुहमुहष्णुहष्णिहाम् ।

\marginnote{\scriptsize ८.२.३४}\noindent नहो धः ।

\marginnote{\scriptsize ८.२.३५}\noindent आहस्थः ।

\marginnote{\scriptsize ८.२.३६}\noindent व्रश्चभ्रस्जसृजमृजयजराजभ्राजच्छशां षः ।

\marginnote{\scriptsize ८.२.३७}\noindent एकाचो बशो भष् झषन्तस्य स्ध्वोः ।

\marginnote{\scriptsize ८.२.३८}\noindent दधस्तथोश्च ।

\marginnote{\scriptsize ८.२.३९}\noindent झलां जशोऽन्ते ।

\marginnote{\scriptsize ८.२.४०}\noindent झषस्तथोर्धोऽधः ।

\marginnote{\scriptsize ८.२.४१}\noindent षढोः कः सि ।

\marginnote{\scriptsize ८.२.४२}\noindent रदाभ्यां निष्ठातो नः पूर्वस्य च दः ।

\marginnote{\scriptsize ८.२.४३}\noindent संयोगादेरातो धातोर्यण्वतः ।

\marginnote{\scriptsize ८.२.४४}\noindent ल्वादिभ्यः ।

\marginnote{\scriptsize ८.२.४५}\noindent ओदितश्च ।

\marginnote{\scriptsize ८.२.४६}\noindent क्षियो दीर्घात् ।

\marginnote{\scriptsize ८.२.४७}\noindent श्योऽस्पर्शे ।

\marginnote{\scriptsize ८.२.४८}\noindent अञ्चोऽनपादाने ।

\marginnote{\scriptsize ८.२.४९}\noindent दिवोऽविजिगीषायाम् ।

\marginnote{\scriptsize ८.२.५०}\noindent निर्वाणोऽवाते ।

\marginnote{\scriptsize ८.२.५१}\noindent शुषः कः ।

\marginnote{\scriptsize ८.२.५२}\noindent पचो वः ।

\marginnote{\scriptsize ८.२.५३}\noindent क्षायो मः ।

\marginnote{\scriptsize ८.२.५४}\noindent प्रस्त्योऽन्यतरस्याम् ।

\marginnote{\scriptsize ८.२.५५}\noindent अनुपसर्गात् फुल्लक्षीबकृशोल्लाघाः ।

\marginnote{\scriptsize ८.२.५६}\noindent नुदविदोन्दत्राघ्राह्रीभ्योऽन्यतरस्याम् ।

\marginnote{\scriptsize ८.२.५७}\noindent न ध्याख्यापॄमूर्छिमदाम् ।

\marginnote{\scriptsize ८.२.५८}\noindent वित्तो भोगप्रत्यययोः ।

\marginnote{\scriptsize ८.२.५९}\noindent भित्तं शकलम् ।

\marginnote{\scriptsize ८.२.६०}\noindent ऋणमाधमर्ण्ये ।

\marginnote{\scriptsize ८.२.६१}\noindent नसत्तनिषत्तानुत्तप्रतूर्तसूर्तगूर्तानि छन्दसि ।

\marginnote{\scriptsize ८.२.६२}\noindent क्विन्प्रत्ययस्य कुः ।

\marginnote{\scriptsize ८.२.६३}\noindent नशेर्वा ।

\marginnote{\scriptsize ८.२.६४}\noindent मो नो धातोः ।

\marginnote{\scriptsize ८.२.६५}\noindent म्वोश्च ।

\marginnote{\scriptsize ८.२.६६}\noindent ससजुषो रुः ।

\marginnote{\scriptsize ८.२.६७}\noindent अवयाःश्वेतवाःपुरोडाश्च ।

\marginnote{\scriptsize ८.२.६८}\noindent अहन् ।

\marginnote{\scriptsize ८.२.६९}\noindent रोऽसुपि ।

\marginnote{\scriptsize ८.२.७०}\noindent अम्नरूधरवरित्युभयथा छन्दसि ।

\marginnote{\scriptsize ८.२.७१}\noindent भुवश्च महाव्याहृतेः ।

\marginnote{\scriptsize ८.२.७२}\noindent वसुस्रंसुध्वंस्वनडुहां दः ।

\marginnote{\scriptsize ८.२.७३}\noindent तिप्यनस्तेः ।

\marginnote{\scriptsize ८.२.७४}\noindent सिपि धातो रुर्वा ।

\marginnote{\scriptsize ८.२.७५}\noindent दश्च ।

\marginnote{\scriptsize ८.२.७६}\noindent र्वोरुपधाया दीर्घ इकः ।

\marginnote{\scriptsize ८.२.७७}\noindent हलि च ।

\marginnote{\scriptsize ८.२.७८}\noindent उपधायां च ।

\marginnote{\scriptsize ८.२.७९}\noindent न भकुर्छुराम् ।

\marginnote{\scriptsize ८.२.८०}\noindent अदसोऽसेर्दादु दो मः ।

\marginnote{\scriptsize ८.२.८१}\noindent एत ईद्बहुवचने ।

\marginnote{\scriptsize ८.२.८२}\noindent वाक्यस्य टेः प्लुत उदात्तः ।

\marginnote{\scriptsize ८.२.८३}\noindent प्रत्यभिवादेअशूद्रे ।

\marginnote{\scriptsize ८.२.८४}\noindent दूराद्धूते च ।

\marginnote{\scriptsize ८.२.८५}\noindent हैहेप्रयोगे हैहयोः ।

\marginnote{\scriptsize ८.२.८६}\noindent गुरोरनृतोऽनन्त्यस्याप्येकैकस्य प्राचाम् ।

\marginnote{\scriptsize ८.२.८७}\noindent ओमभ्यादाने ।

\marginnote{\scriptsize ८.२.८८}\noindent ये यज्ञकर्मणि ।

\marginnote{\scriptsize ८.२.८९}\noindent प्रणवष्टेः ।

\marginnote{\scriptsize ८.२.९०}\noindent याज्याऽन्तः ।

\marginnote{\scriptsize ८.२.९१}\noindent ब्रूहिप्रेस्यश्रौषड्वौषडावहानामादेः ।

\marginnote{\scriptsize ८.२.९२}\noindent अग्नीत्प्रेषणे परस्य च ।

\marginnote{\scriptsize ८.२.९३}\noindent विभाषा पृष्टप्रतिवचने हेः ।

\marginnote{\scriptsize ८.२.९४}\noindent निगृह्यानुयोगे च ।

\marginnote{\scriptsize ८.२.९५}\noindent आम्रेडितं भर्त्सने ।

\marginnote{\scriptsize ८.२.९६}\noindent अङ्गयुक्तं तिङ् आकाङ्क्षम् ।

\marginnote{\scriptsize ८.२.९७}\noindent विचार्यमाणानाम् ।

\marginnote{\scriptsize ८.२.९८}\noindent पूर्वं तु भाषायाम् ।

\marginnote{\scriptsize ८.२.९९}\noindent प्रतिश्रवणे च ।

\marginnote{\scriptsize ८.२.१००}\noindent अनुदात्तं प्रश्नान्ताभिपूजितयोः ।

\marginnote{\scriptsize ८.२.१०१}\noindent चिदिति चोपमाऽर्थे प्रयुज्यमाने ।

\marginnote{\scriptsize ८.२.१०२}\noindent उपरिस्विदासीदिति च ।

\marginnote{\scriptsize ८.२.१०३}\noindent स्वरितमाम्रेडितेऽसूयासम्मतिकोपकुत्सनेषु ।

\marginnote{\scriptsize ८.२.१०४}\noindent क्षियाऽऽशीःप्रैषेषु तिङ् आकाङ्क्षम् ।

\marginnote{\scriptsize ८.२.१०५}\noindent अनन्त्यस्यापि प्रश्नाख्यानयोः ।

\marginnote{\scriptsize ८.२.१०६}\noindent प्लुतावैच इदुतौ ।

\marginnote{\scriptsize ८.२.१०७}\noindent एचोऽप्रगृह्यस्यादूराद्धूते पूर्वस्यार्धस्यादुत्तरस्येदुतौ ।

\marginnote{\scriptsize ८.२.१०८}\noindent तयोर्य्वावचि संहितायाम् ।

\marginnote{\scriptsize ८.३.१}\noindent मतुवसो रु सम्बुद्धौ छन्दसि ।

\marginnote{\scriptsize ८.३.२}\noindent अत्रानुनासिकः पूर्वस्य तु वा ।

\marginnote{\scriptsize ८.३.३}\noindent आतोऽटि नित्यम् ।

\marginnote{\scriptsize ८.३.४}\noindent अनुनासिकात् परोऽनुस्वारः ।

\marginnote{\scriptsize ८.३.५}\noindent समः सुटि ।

\marginnote{\scriptsize ८.३.६}\noindent पुमः खय्यम्परे ।

\marginnote{\scriptsize ८.३.७}\noindent नश्छव्यप्रशान् ।

\marginnote{\scriptsize ८.३.८}\noindent उभयथर्क्षु ।

\marginnote{\scriptsize ८.३.९}\noindent दीर्घादटि समानपदे ।

\marginnote{\scriptsize ८.३.१०}\noindent नॄन् पे ।

\marginnote{\scriptsize ८.३.११}\noindent स्वतवान् पायौ ।

\marginnote{\scriptsize ८.३.१२}\noindent कानाम्रेडिते ।

\marginnote{\scriptsize ८.३.१३}\noindent ढो ढे लोपः ।

\marginnote{\scriptsize ८.३.१४}\noindent रो रि ।

\marginnote{\scriptsize ८.३.१५}\noindent खरवसानयोर्विसर्जनीयः ।

\marginnote{\scriptsize ८.३.१६}\noindent रोः सुपि ।

\marginnote{\scriptsize ८.३.१७}\noindent भोभगोऽघोऽपूर्वस्य योऽशि ।

\marginnote{\scriptsize ८.३.१८}\noindent व्योर्लघुप्रयत्नतरः शाकटायनस्य ।

\marginnote{\scriptsize ८.३.१९}\noindent लोपः शाकल्यस्य ।

\marginnote{\scriptsize ८.३.२०}\noindent ओतो गार्ग्यस्य ।

\marginnote{\scriptsize ८.३.२१}\noindent उञि च पदे ।

\marginnote{\scriptsize ८.३.२२}\noindent हलि सर्वेषाम् ।

\marginnote{\scriptsize ८.३.२३}\noindent मोऽनुस्वारः ।

\marginnote{\scriptsize ८.३.२४}\noindent नश्चापदान्तस्य झलि ।

\marginnote{\scriptsize ८.३.२५}\noindent मो राजि समः क्वौ ।

\marginnote{\scriptsize ८.३.२६}\noindent हे मपरे वा ।

\marginnote{\scriptsize ८.३.२७}\noindent नपरे नः ।

\marginnote{\scriptsize ८.३.२८}\noindent ङ्णोः कुक्टुक् शरि ।

\marginnote{\scriptsize ८.३.२९}\noindent डः सि धुट् ।

\marginnote{\scriptsize ८.३.३०}\noindent नश्च ।

\marginnote{\scriptsize ८.३.३१}\noindent शि तुक् ।

\marginnote{\scriptsize ८.३.३२}\noindent ङमो ह्रस्वादचि ङमुण्नित्यम् ।

\marginnote{\scriptsize ८.३.३३}\noindent मय उञो वो वा ।

\marginnote{\scriptsize ८.३.३४}\noindent विसर्जनीयस्य सः ।

\marginnote{\scriptsize ८.३.३५}\noindent शर्परे विसर्जनीयः ।

\marginnote{\scriptsize ८.३.३६}\noindent वा शरि ।

\marginnote{\scriptsize ८.३.३७}\noindent कुप्वोः XकXपौ च ।

\marginnote{\scriptsize ८.३.३८}\noindent सोऽपदादौ ।

\marginnote{\scriptsize ८.३.३९}\noindent इणः षः ।

\marginnote{\scriptsize ८.३.४०}\noindent नमस्पुरसोर्गत्योः ।

\marginnote{\scriptsize ८.३.४१}\noindent इदुदुपधस्य चाप्रत्ययस्य ।

\marginnote{\scriptsize ८.३.४२}\noindent तिरसोऽन्यतरस्याम् ।

\marginnote{\scriptsize ८.३.४३}\noindent द्विस्त्रिश्चतुरिति कृत्वोऽर्थे ।

\marginnote{\scriptsize ८.३.४४}\noindent इसुसोः सामर्थ्ये ।

\marginnote{\scriptsize ८.३.४५}\noindent नित्यं समासेऽनुत्तरपदस्थस्य ।

\marginnote{\scriptsize ८.३.४६}\noindent अतः कृकमिकंसकुम्भपात्रकुशाकर्णीष्वनव्ययस्य ।

\marginnote{\scriptsize ८.३.४७}\noindent अधःशिरसी पदे ।

\marginnote{\scriptsize ८.३.४८}\noindent कस्कादिषु च ।

\marginnote{\scriptsize ८.३.४९}\noindent छन्दसि वाऽप्राम्रेडितयोः ।

\marginnote{\scriptsize ८.३.५०}\noindent कःकरत्करतिकृधिकृतेष्वनदितेः ।

\marginnote{\scriptsize ८.३.५१}\noindent पञ्चम्याः परावध्यर्थे ।

\marginnote{\scriptsize ८.३.५२}\noindent पातौ च बहुलम् ।

\marginnote{\scriptsize ८.३.५३}\noindent षष्ठ्याः पतिपुत्रपृष्ठपारपदपयस्पोषेषु ।

\marginnote{\scriptsize ८.३.५४}\noindent इडाया वा ।

\marginnote{\scriptsize ८.३.५५}\noindent अपदान्तस्य मूर्धन्यः ।

\marginnote{\scriptsize ८.३.५६}\noindent सहेः साडः सः ।

\marginnote{\scriptsize ८.३.५७}\noindent इङ्कोः ।

\marginnote{\scriptsize ८.३.५८}\noindent नुम्विसर्जनीयशर्व्यवायेऽपि ।

\marginnote{\scriptsize ८.३.५९}\noindent आदेशप्रत्यययोः ।

\marginnote{\scriptsize ८.३.६०}\noindent शासिवसिघसीनां च ।

\marginnote{\scriptsize ८.३.६१}\noindent स्तौतिण्योरेव षण्यभ्यासात् ।

\marginnote{\scriptsize ८.३.६२}\noindent सः स्विदिस्वदिसहीनां च ।

\marginnote{\scriptsize ८.३.६३}\noindent प्राक्सितादड्व्यवायेऽपि ।

\marginnote{\scriptsize ८.३.६४}\noindent स्थाऽऽदिष्वभ्यासेन चाभ्यासय ।

\marginnote{\scriptsize ८.३.६५}\noindent उपसर्गात् सुनोतिसुवतिस्यतिस्तौतिस्तोभतिस्थासेनय- सेधसिचसञ्जस्वञ्जाम् ।

\marginnote{\scriptsize ८.३.६६}\noindent सदिरप्रतेः ।

\marginnote{\scriptsize ८.३.६७}\noindent स्तम्भेः ।

\marginnote{\scriptsize ८.३.६८}\noindent अवाच्चालम्बनाविदूर्ययोः ।

\marginnote{\scriptsize ८.३.६९}\noindent वेश्च स्वनो भोजने ।

\marginnote{\scriptsize ८.३.७०}\noindent परिनिविभ्यः सेवसितसयसिवुसहसुट्स्तुस्वञ्जाम् ।

\marginnote{\scriptsize ८.३.७१}\noindent सिवादीनां वाऽड्व्यवायेऽपि ।

\marginnote{\scriptsize ८.३.७२}\noindent अनुविपर्यभिनिभ्यः स्यन्दतेरप्राणिषु ।

\marginnote{\scriptsize ८.३.७३}\noindent वेः स्कन्देरनिष्ठायाम् ।

\marginnote{\scriptsize ८.३.७४}\noindent परेश्च ।

\marginnote{\scriptsize ८.३.७५}\noindent परिस्कन्दः प्राच्यभरतेषु ।

\marginnote{\scriptsize ८.३.७६}\noindent स्फुरतिस्फुलत्योर्निर्निविभ्यः ।

\marginnote{\scriptsize ८.३.७७}\noindent वेः स्कभ्नातेर्नित्यम् ।

\marginnote{\scriptsize ८.३.७८}\noindent इणः षीध्वंलुङ्लिटां धोऽङ्गात् ।

\marginnote{\scriptsize ८.३.७९}\noindent विभाषेटः ।

\marginnote{\scriptsize ८.३.८०}\noindent समासेऽङ्गुलेः सङ्गः ।

\marginnote{\scriptsize ८.३.८१}\noindent भीरोः स्थानम् ।

\marginnote{\scriptsize ८.३.८२}\noindent अग्नेः स्तुत्स्तोमसोमाः ।

\marginnote{\scriptsize ८.३.८३}\noindent ज्योतिरायुषः स्तोमः ।

\marginnote{\scriptsize ८.३.८४}\noindent मातृपितृभ्यां स्वसा ।

\marginnote{\scriptsize ८.३.८५}\noindent मातुःपितुर्भ्यामन्यतरस्याम् ।

\marginnote{\scriptsize ८.३.८६}\noindent अभिनिसः स्तनः शब्दसंज्ञायाम् ।

\marginnote{\scriptsize ८.३.८७}\noindent उपसर्गप्रादुर्भ्यामस्तिर्यच्परः ।

\marginnote{\scriptsize ८.३.८८}\noindent सुविनिर्दुर्भ्यः सुपिसूतिसमाः ।

\marginnote{\scriptsize ८.३.८९}\noindent निनदीभ्यां स्नातेः कौशले ।

\marginnote{\scriptsize ८.३.९०}\noindent सूत्रं प्रतिष्णातम् ।

\marginnote{\scriptsize ८.३.९१}\noindent कपिष्ठलो गोत्रे ।

\marginnote{\scriptsize ८.३.९२}\noindent प्रष्ठोऽग्रगामिनि ।

\marginnote{\scriptsize ८.३.९३}\noindent वृक्षासनयोर्विष्टरः ।

\marginnote{\scriptsize ८.३.९४}\noindent छन्दोनाम्नि च ।

\marginnote{\scriptsize ८.३.९५}\noindent गवियुधिभ्यां स्थिरः ।

\marginnote{\scriptsize ८.३.९६}\noindent विकुशमिपरिभ्यः स्थलम् ।

\marginnote{\scriptsize ८.३.९७}\noindent अम्बाम्बगोभूमिसव्यापद्वित्रिकुशेकुशङ्क्वङ्गुमञ्जि- पुञ्जिपरमेबर्हिर्दिव्यग्निभ्यः स्थः ।

\marginnote{\scriptsize ८.३.९८}\noindent सुषामादिषु च ।

\marginnote{\scriptsize ८.३.९९}\noindent एति संज्ञायामगात् ।

\marginnote{\scriptsize ८.३.१००}\noindent नक्षत्राद्वा ।

\marginnote{\scriptsize ८.३.१०१}\noindent ह्रस्वात् तादौ तद्धिते ।

\marginnote{\scriptsize ८.३.१०२}\noindent निसस्तपतावनासेवने ।

\marginnote{\scriptsize ८.३.१०३}\noindent युष्मत्तत्ततक्षुःष्वन्तःपादम् ।

\marginnote{\scriptsize ८.३.१०४}\noindent यजुष्येकेषाम् ।

\marginnote{\scriptsize ८.३.१०५}\noindent स्तुतस्तोमयोश्छन्दसि ।

\marginnote{\scriptsize ८.३.१०६}\noindent पूर्वपदात् ।

\marginnote{\scriptsize ८.३.१०७}\noindent सुञः ।

\marginnote{\scriptsize ८.३.१०८}\noindent सनोतेरनः ।

\marginnote{\scriptsize ८.३.१०९}\noindent सहेः पृतनर्ताभ्यां च ।

\marginnote{\scriptsize ८.३.११०}\noindent न रपरसृपिसृजिस्पृशिस्पृहिसवनादीनाम् ।

\marginnote{\scriptsize ८.३.१११}\noindent सात्पदाद्योः ।

\marginnote{\scriptsize ८.३.११२}\noindent सिचो यङि ।

\marginnote{\scriptsize ८.३.११३}\noindent सेधतेर्गतौ ।

\marginnote{\scriptsize ८.३.११४}\noindent प्रतिस्तब्धनिस्तब्धौ च ।

\marginnote{\scriptsize ८.३.११५}\noindent सोढः ।

\marginnote{\scriptsize ८.३.११६}\noindent स्तम्भुसिवुसहां चङि ।

\marginnote{\scriptsize ८.३.११७}\noindent सुनोतेः स्यसनोः ।

\marginnote{\scriptsize ८.३.११८}\noindent सदिष्वञ्जोः परस्य लिटि ।

\marginnote{\scriptsize ८.३.११९}\noindent निव्यभिभ्योऽड्व्यावये वा छन्दसि ।

\marginnote{\scriptsize ८.४.१}\noindent रषाभ्यां नो णः समानपदे ।

\marginnote{\scriptsize ८.४.२}\noindent अट्कुप्वाङ्नुम्व्यवायेऽपि ।

\marginnote{\scriptsize ८.४.३}\noindent पूर्वपदात् संज्ञायामगः ।

\marginnote{\scriptsize ८.४.४}\noindent वनं पुरगामिश्रकासिध्रकाशारिकाकोटराऽग्रेभ्यः ।

\marginnote{\scriptsize ८.४.५}\noindent प्रनिरन्तःशरेक्षुप्लक्षाम्रकार्ष्यखदिर- पियूक्षाभ्योऽसंज्ञायामपि ।

\marginnote{\scriptsize ८.४.६}\noindent विभाषौषधिवनस्पतिभ्यः ।

\marginnote{\scriptsize ८.४.७}\noindent अह्नोऽदन्तात् ।

\marginnote{\scriptsize ८.४.८}\noindent वाहनमाहितात् ।

\marginnote{\scriptsize ८.४.९}\noindent पानं देशे ।

\marginnote{\scriptsize ८.४.१०}\noindent वा भावकरणयोः ।

\marginnote{\scriptsize ८.४.११}\noindent प्रातिपदिकान्तनुम्विभक्तिषु च ।

\marginnote{\scriptsize ८.४.१२}\noindent एकाजुत्तरपदे णः ।

\marginnote{\scriptsize ८.४.१३}\noindent कुमति च ।

\marginnote{\scriptsize ८.४.१४}\noindent उपसर्गादसमासेऽपि णोपदेशस्य ।

\marginnote{\scriptsize ८.४.१५}\noindent हिनुमीना ।

\marginnote{\scriptsize ८.४.१६}\noindent आनि लोट् ।

\marginnote{\scriptsize ८.४.१७}\noindent नेर्गदनदपतपदघुमास्यतिहन्तियातिवातिद्रातिप्साति- वपतिवहतिशाम्यतिचिनोतिदेग्धिषु च ।

\marginnote{\scriptsize ८.४.१८}\noindent शेषे विभाषाऽकखादावषान्त उपदेशे ।

\marginnote{\scriptsize ८.४.१९}\noindent अनितेः ।

\marginnote{\scriptsize ८.४.२०}\noindent अन्तः ।

\marginnote{\scriptsize ८.४.२१}\noindent उभौ साभ्यासस्य ।

\marginnote{\scriptsize ८.४.२२}\noindent हन्तेरत्पूर्वस्य ।

\marginnote{\scriptsize ८.४.२३}\noindent वमोर्वा ।

\marginnote{\scriptsize ८.४.२४}\noindent अन्तरदेशे ।

\marginnote{\scriptsize ८.४.२५}\noindent अयनं च ।

\marginnote{\scriptsize ८.४.२६}\noindent छन्दस्यृदवग्रहात् ।

\marginnote{\scriptsize ८.४.२७}\noindent नश्च धातुस्थोरुषुभ्यः ।

\marginnote{\scriptsize ८.४.२८}\noindent उपसर्गाद् बहुलम् ।

\marginnote{\scriptsize ८.४.२९}\noindent कृत्यचः ।

\marginnote{\scriptsize ८.४.३०}\noindent णेर्विभाषा ।

\marginnote{\scriptsize ८.४.३१}\noindent हलश्च इजुपधात् ।

\marginnote{\scriptsize ८.४.३२}\noindent इजादेः सनुमः ।

\marginnote{\scriptsize ८.४.३३}\noindent वा निंसनिक्षनिन्दाम् ।

\marginnote{\scriptsize ८.४.३४}\noindent न भाभूपूकमिगमिप्यायीवेपाम् ।

\marginnote{\scriptsize ८.४.३५}\noindent षात् पदान्तात् ।

\marginnote{\scriptsize ८.४.३६}\noindent नशेः षान्तस्य ।

\marginnote{\scriptsize ८.४.३७}\noindent पदान्तस्य ।

\marginnote{\scriptsize ८.४.३८}\noindent पदव्यवायेऽपि ।

\marginnote{\scriptsize ८.४.३९}\noindent क्षुभ्नाऽऽदिषु च ।

\marginnote{\scriptsize ८.४.४०}\noindent स्तोः श्चुना श्चुः ।

\marginnote{\scriptsize ८.४.४१}\noindent ष्टुना ष्टुः ।

\marginnote{\scriptsize ८.४.४२}\noindent न पदान्ताट्टोरनाम् ।

\marginnote{\scriptsize ८.४.४३}\noindent तोः षि ।

\marginnote{\scriptsize ८.४.४४}\noindent शात् ।

\marginnote{\scriptsize ८.४.४५}\noindent यरोऽनुनासिकेऽनुनासिको वा ।

\marginnote{\scriptsize ८.४.४६}\noindent अचो रहाभ्यां द्वे ।

\marginnote{\scriptsize ८.४.४७}\noindent अनचि च ।

\marginnote{\scriptsize ८.४.४८}\noindent नादिन्याक्रोशे पुत्रस्य ।

\marginnote{\scriptsize ८.४.४९}\noindent शरोऽचि ।

\marginnote{\scriptsize ८.४.५०}\noindent त्रिप्रभृतिषु शाकटायनस्य ।

\marginnote{\scriptsize ८.४.५१}\noindent सर्वत्र शाकल्यस्य ।

\marginnote{\scriptsize ८.४.५२}\noindent दीर्घादाचार्याणाम् ।

\marginnote{\scriptsize ८.४.५३}\noindent झलां जश् झशि ।

\marginnote{\scriptsize ८.४.५४}\noindent अभ्यासे चर्च्च ।

\marginnote{\scriptsize ८.४.५५}\noindent खरि च ।

\marginnote{\scriptsize ८.४.५६}\noindent वाऽवसाने ।

\marginnote{\scriptsize ८.४.५७}\noindent अणोऽप्रगृह्यस्यानुनासिकः ।

\marginnote{\scriptsize ८.४.५८}\noindent अनुस्वारस्य ययि परसवर्णः ।

\marginnote{\scriptsize ८.४.५९}\noindent वा पदान्तस्य ।

\marginnote{\scriptsize ८.४.६०}\noindent तोर्लि ।

\marginnote{\scriptsize ८.४.६१}\noindent उदः स्थास्तम्भोः पूर्वस्य ।

\marginnote{\scriptsize ८.४.६२}\noindent झयो होऽन्यतरस्याम् ।

\marginnote{\scriptsize ८.४.६३}\noindent शश्छोऽटि ।

\marginnote{\scriptsize ८.४.६४}\noindent हलो यमां यमि लोपः ।

\marginnote{\scriptsize ८.४.६५}\noindent झरो झरि सवर्णे ।

\marginnote{\scriptsize ८.४.६६}\noindent उदात्तादनुदात्तस्य स्वरितः ।

\marginnote{\scriptsize ८.४.६७}\noindent नोदात्तस्वरितोदयमगार्ग्यकाश्यपगालवानाम् ।

\marginnote{\scriptsize ८.४.६८}\noindent अ अ इति ।

\begin{center}\textbf{इति}\end{center}

\end{document}